\documentclass[12pt, oneside]{book}  % 13pt는 직접 설정함

% --- Language and Font Settings ---
\usepackage[english]{babel}
\usepackage{fontspec}
\usepackage{kotex}  % Korean typesetting

% --- Font Configuration (Main + Korean) ---
\usepackage{libertinus}
\setmainhangulfont[
  Path = ./,
  UprightFont = *Batang Light.ttf,
  BoldFont    = *Batang Medium.ttf
]{KoPubWorld}

\newfontfamily\kopublatin[
  Path = ./,
  UprightFont = *Dotum Medium.ttf
]{KoPubWorld}

\newhangulfontfamily\kopubhangul[
  Path = ./,
  UprightFont = *Dotum Medium.ttf
]{KoPubWorld}

\newcommand{\emphksans}[1]{%
  {\kopubhangul\kopublatin\textsl{#1}}%
}

\newcommand{\notetitle}[1]{%
  {\kopubhangul\kopublatin #1}%
}

\newcommand{\noteentry}{%
  \par\addvspace{0.35\baselineskip}\noindent
}

% --- Page Geometry: A4 with wide margins for print readability ---
\usepackage[b5paper, margin=1in]{geometry}

% --- Line Spacing ---
\usepackage{setspace}
\setstretch{1.45}
\hyphenation{theory}

% --- Section Formatting: Disable numbering ---
\usepackage{titlesec}
\setcounter{secnumdepth}{0}
\setcounter{tocdepth}{2} % 목차 수준

% Footnotes reset only at each chapter input below.
% Do not reset by section/subsection: the Postscript must run 1--75.

\titleformat{\chapter}[display]
  {\normalfont\Huge\bfseries}
  {}{0pt}{\Huge}
\titleformat{\section}
  {\normalfont\Large\bfseries\centering} % 섹션 제목 가운데 정렬
  {}{0pt}{\Large}
\titleformat{\subsection}
  {\normalfont\normalsize\bfseries\centering} % 서브섹션 제목 가운데 정렬
  {}{0pt}{\normalsize}

\usepackage{etoolbox}
\pretocmd{\section}{\clearpage}{}{}

% --- Math and Tables ---
\usepackage{amsmath, amsthm, mathtools}
\usepackage{graphicx}
\usepackage{enumitem}
\providecommand{\tightlist}{%
  \setlength{\itemsep}{0pt}\setlength{\parskip}{0pt}}
\usepackage{tabularx, booktabs, longtable}
\usepackage{footmisc}

% --- Quotes, Hyperlink, Header/Footer ---
\usepackage{csquotes}
\usepackage[hidelinks]{hyperref}
\usepackage{fancyhdr}
\pagestyle{fancy}
\fancyhf{}
\fancyfoot[C]{\thepage}
\setlength{\headheight}{15pt}

% --- Title Info ---
\title{\Huge\textsc{The Concept of Law} \\[2ex] \Large Third Edition}
\author{\Large H. L. A. Hart}
\date{}

\begin{document}

% --- Title Page ---
\begin{titlepage}
  \centering
  \vspace*{3cm}
  {\Huge\bfseries\textsc{법의 개념}}\\[1.5ex]
  {\Large 제3판}\\[4ex]
  \textsc{H. L. A. Hart}\\[6ex]
  {\small 2026 SUMMER 강의용 한국어판\\
  강좌 외의 사용을 불허함}
  \vfill
\end{titlepage}

\small

\frontmatter
\maketitle
\input{sections/01_1961_PREFACE-translation.tex}
\section{1994년 제2판의 Bulloch와 Raz의 편집자 주 (Editors'
Note)}\label{uxb144-uxc81c2uxd310uxc758-bullochuxc640-razuxc758-uxd3b8uxc9d1uxc790-uxc8fc-editors-note}

\textbf{(제2판을 위해 작성됨)}

{\small

\emph{The Concept of Law}는 출간된 지 불과 몇 년 만에 영어권 세계와 그
너머에서 법리학(jurisprudence)이 이해되고 연구되는 방식을 변모시켰다. 이
책의 막대한 영향은 이 책과 그 교설들(doctrines)을 논의하는 수많은
출판물들로 이어졌으며, 그것은 법이론(legal theory)의 맥락에서뿐 아니라
정치철학(political philosophy)과 도덕철학(moral philosophy)에서도
그러했다.

하트(Hart)는 여러 해 동안 \emph{The Concept of Law}에 한 장(chapter)을
추가할 생각을 염두에 두고 있었다. 그는 그토록 큰 영향을 미친 본문을
자잘하게 손대기를 원하지 않았고, 그의 뜻에 따라 본문은 사소한
정정들(minor corrections)을 제외하고는 변경 없이 여기 출판된다. 그러나
그는 이 책에 관한 많은 논의들에 응답하기를 원했다. 그는 자신의 입장을
잘못 이해한 이들에 맞서 그 입장을 방어하고, 근거 없는 비판을 논박하며,
그가 보기에도 동등하게 중요하게는 정당한 비판의 힘(force)을 인정하고,
그러한 지점들에 대응하기 위해 책의 교설들(doctrines)을 조정하는 방식들을
제안하고자 했다. 새로운 장은 처음에는 머리말(preface)로 생각되었으나
최종적으로는 후기(postscript)로 생각되었고, 그가 사망할 당시
미완성이었다. 이는 부분적으로만 그의 꼼꼼한 완벽주의(perfectionism)
때문이었다. 그것은 또한 이 기획의 현명함에 관한 지속적인 의문들, 그리고
원래 구상된 책의 테제들(theses)이 가진 생동력과 통찰을 자신이 충분히
살려낼 수 있을지에 관한 끈질긴 불확실성 때문이기도 했다. 그럼에도 그는
많은 중단에도 불구하고 후기 작업을 계속했고, 사망 당시 의도된 두 절들
가운데 첫째 절은 거의 완성되어 있었다.

제니퍼 하트(Jennifer Hart)가 우리에게 초고들을 살펴보고 그 안에 출판
가능한 것이 있는지 결정해 달라고 요청했을 때, 우리의 가장 우선적인
생각은 하트가 만족하지 않았을 어떤 것도 출판되도록 해서는 안 된다는
것이었다. 그러므로 우리는 후기의 첫째 절 대부분이 그처럼 완성된 상태에
있음을 발견하고 기뻤다. 우리는 둘째 절을 위한 손글씨 주석들만
발견했는데, 그것들은 출판 가능하기에는 너무 단편적이고
미정형(inchoate)이었다. 대조적으로 첫째 절은 여러 판본들(versions)로
존재했고, 타자되고, 수정되고, 다시 타자되고, 다시 수정되어 있었다. 가장
최근의 판본조차도 그가 최종 상태로 생각한 것은 분명 아니었다. 연필과
바이로펜(Biro)으로 된 수많은 변경사항들(alterations)이 있다. 더욱이
하트는 더 이른 판본들을 버리지 않았고, 손 가까이에 있던 어느 판본에든
계속 작업한 것으로 보인다. 이것이 편집 작업을 더 어렵게 만들었지만, 지난
2년 동안 도입된 변경들은 대부분 문체적 뉘앙스의 변경들이었고, 그 사실
자체가 그가 본문을 있는 그대로 본질적으로 만족스러워했다는 점을
나타냈다.

우리의 과제는 대안적 판본들(alternative versions)을 비교하고, 그것들이
일치하지 않을 때에는 그중 하나에만 나타나는 본문 조각들(segments of
text)이 다른 판본들에서 빠져 있는 이유가 그가 그것들을 버렸기 때문인지,
아니면 모든 수정사항들(emendations)을 편입한 하나의 판본을 그가 한 번도
갖지 못했기 때문인지를 확정하는 것이었다. 출판된 텍스트는 하트가 버리지
않았고 그가 계속 수정한 텍스트 판본들에 나타나는 모든 수정사항들을
포함한다. 때때로 텍스트 자체는 비정합적이었다. 종종 이것은 타자 담당자가
원고를 잘못 읽은 결과였음이 틀림없고, 하트는 그 오류들을 언제나
알아차리지는 못했다. 다른 때에는 그것이 틀림없이, 그가 끝내 살아서 하지
못한 최종 원고 작성 단계에서 정리되었을 법한, 작성 과정에서 문장들이
자연스럽게 꼬이는 방식 때문이었다. 이런 경우 우리는 원래 텍스트를
복원하거나 최소한의 개입으로 하트의 생각(thought)을 되살리려 했다.
제6절(재량(discretion)에 관하여)은 하나의 특별한 문제를 제기했다. 우리는
그 첫 단락의 두 판본을 발견했는데, 하나는 그 지점에서 끝나는 사본 안에
있었고, 다른 하나는 그 절의 나머지를 포함하는 사본 안에 있었다. 잘린
판본(truncated version)은 그의 가장 최근 수정들 중 많은 것을 편입한 사본
안에 있었고, 그가 결코 버린 적이 없었으며, 후기의 일반적 논의와도
부합했으므로, 우리는 두 판본 모두 출판되도록 허용하기로 결정했다.
이어지지 않은 판본은 미주(endnote)에 나타난다.

하트는 주석들을, 대부분 참조들(references)인 그 주석들을, 타자하도록 한
적이 결코 없었다. 그는 손글씨 주석 판본을 가지고 있었고, 그
단서들(cues)은 본문 가장 이른 타자 사본에서 가장 쉽게 추적되었다. 나중에
그는 가끔 여백 논평에 참조들을 덧붙였지만, 대체로 이것들은 불완전했고,
때로는 그 참조를 추적할 필요만을 표시했다. 티머시 엔디컷(Timothy
Endicott)은 모든 참조들을 확인했고, 불완전한 것들을 모두 추적했으며,
하트가 원천(source)을 지목하지 않고 드워킨(Dworkin)을 인용하거나 그를
가깝게 풀어쓴 곳에는 참조들을 추가했다. 엔디컷은 또한 인용문들이
부정확했던 곳에서 본문을 바로잡았다. 폭넓은 조사와
수완(resourcefulness)을 수반한 이 작업 과정에서, 그는 위에서 제시된 편집
지침들에 부합하는 본문에 대한 몇 가지 정정사항들도 제안했고, 우리는 이를
감사히 편입하였다.

우리 마음속에는, 기회가 주어졌다면 하트가 출판 전에 텍스트를 더 다듬고
개선했을 것이라는 데 아무 의심도 없다. 그러나 우리는 출판된
후기(postscript)가 드워킨의 많은 논변들(arguments)에 대한 그의 숙고된
응답을 담고 있다고 믿는다.

}

\vspace{1em}

\begin{flushright}
\textbf{페넬로페 A. 불로흐 (Penelope A. Bulloch) 조셉 라즈 (Joseph Raz)}
\\ 1994
\end{flushright}

\input{sections/03_2012_Third_Preface-translation.tex}
\setcounter{footnote}{0}
\input{sections/04_2012_INTRODUCTION-translation.tex}

\tableofcontents
\clearpage

\mainmatter
\cleardoublepage
\markboth{}{}

\setcounter{footnote}{0}
\input{sections/chapter_1-translation.tex}
\setcounter{footnote}{0}
\input{sections/chapter_2-translation.tex}
\setcounter{footnote}{0}
\input{sections/chapter_3-translation.tex}
\setcounter{footnote}{0}
\section{IV. 주권자와 수범자(SOVEREIGN AND
SUBJECT)}\label{iv.-uxc8fcuxad8cuxc790uxc640-uxc218uxbc94uxc790sovereign-and-subject}

강제적 명령으로서의 법이라는 단순한 모델을 비판하면서, 우리는 지금까지
이 이해방식에 따르면 어떤 사회의 법을 구성하는 것으로 여겨지는
`주권자(sovereign)' 인격 또는 인격들의 존재에 관해서는 아무런 문제도
제기하지 않았다. 실제로, 위협에 의해 뒷받침되는 명령이라는 개념이 법의
다양한 유형들을 설명하기에 충분한지 여부를 논의함에 있어, 우리는
잠정적으로 다음과 같은 가정을 두었다. 즉, 법이 존재하는 모든 사회에는
실제로 하나의 주권자가 존재하며, 그는 복종의 습관(habit of obedience)을
참조하여(by reference to) 다음과 같이 긍정적·부정적으로 특징지어질 수
있다: 사회의 대다수가 그의 명령을 습관적으로 복종하는 반면, 그 자신은
다른 어떤 개인이나 집단의 명령에도 습관적으로 복종하지 않는 인격 또는
인격들의 집단.

이제 우리는 모든 법체계의 기초에 관한 이 일반 이론을 다소 상세히
검토해야 한다. 극단적으로 단순함에도 불구하고, 주권 교설(doctrine of
sovereignty)은 결코 그 중요성이 단순하지 않기 때문이다. 이 이론은, 법이
존재하는 모든 인간 사회에서는, 민주정이든 절대군주정이든 정치적 형태의
다양성 아래에 잠재되어 있는 것으로서, 수범자들이 습관적으로 복종을
행하고 누구에게도 습관적으로 복종하지 않는 주권자 사이의 이 단순한
관계가 궁극적으로 잠재되어 발견된다고 단언한다. 이 이론에 따르면,
주권자와 수범자로 구성된 이러한 수직적 구조는, 법을 보유한 사회에 있어
인간에게 척추가 필수적인 것과 마찬가지로 필수적인 요소이다. 이러한
구조가 현존한다면, 우리는 그 사회를 그 주권자와 함께 하나의 독립된
국가라고 부를 수 있으며, \emphksans{그 사회의} 법(\emph{its} law)에 대해서도
말할 수 있다. 반대로 이러한 구조가 존재하지 않는 경우에는, 이러한 표현들
가운데 어느 것도 적용할 수 없는데, 이는 주권자와 수범자의 관계가 이
이론에 따르면 그러한 표현들의 의미 자체의 일부를 구성하기 때문이다.

이 교설에는 특별히 중요한 두 가지 논점이 있으며, 우리는 이 장의 나머지
부분에서 상세하게 추구될 비판의 방향을 제시하기 위해 여기서 이를
일반적인 용어로 강조하고자 한다. 첫째 논점은 주권자의 법(sovereign's
laws)이 적용되는 자들에게 \textbf{(p.~51)} 요구되는 전부라고 여겨지는
복종의 \emphksans{습관}(a \emph{habit} of obedience)이라는 관념에 관한
것이다. 여기서 우리는 그러한 습관이 대부분의 법체계의 두 가지 두드러진
특징 --- 즉 서로 다른 입법자들이 연속적으로 보유하는 입법 권위(authority
to make law)의 \emphksans{연속성(continuity)}과, 그 법을 만든 자와 그에게
습관적 복종(habitual obedience)을 부여했던 자들이 사망한 이후에도 법이
존속하는 \emphksans{지속성(persistence)} --- 을 설명하기에 충분한지 여부를
검토할 것이다. 둘째 논점은 법 위에 존재하는 주권자의 지위(position)에
관한 것이다. 그는 타인을 위해 법을 창설함으로써 그들에게 법적 책무 또는
`제한(limitations)'을 부과하는 반면, 그 자신은 법적으로 무제한적이며
제한 불가능한(legally unlimited and illimitable) 존재로 간주된다. 여기서
우리는 최고 입법자(supreme lawgiver)의 이러한 법적으로 제한 불가능한
지위가 법의 존재(existence of law)에 필수적인지 여부, 그리고 입법적
권한(legislative power)에 대한 법적 한계(legal limits)의 존재 여부가, 이
이론이 이러한 개념들을 습관과 복종이라는 단순한 용어로 분석하는 방식
속에서 이해될 수 있는지를 검토할 것이다.

\subsection{1. 복종의 습관과 법의
연속성}\label{uxbcf5uxc885uxc758-uxc2b5uxad00uxacfc-uxbc95uxc758-uxc5f0uxc18duxc131}

복종(obedience)이라는 관념은, 엄밀한 검토없이 사용되는 다른 많은
겉보기에는 단순한 관념들과 마찬가지로, 복잡성으로부터 자유롭지 않다.
우리는 이미 지적된 복잡성\footnote{See p.~19 above.}, 즉 `복종'이라는
단어가 종종 위협에 의해 뒷받침된 명령에 대한 단순한 준수(compliance)만이
아니라 권위에 대한 존중을 시사한다는 점은 여기서는 논외로 하겠다.
그럼에도 불구하고, 한 사람이 다른 사람에게 대면하여 단일한 명령(single
order)을 내리는 경우조차도, 특정된 행위(specified act)의
수행(performance)이 복종을 구성하기 위해서는 그 명령의 제시와 그 행위의
수행 사이에 어떠한 정확한 연결이 존재해야 하는지를 명확히 진술하기는
쉽지 않다. 예컨대, 명령을 받은 사람이 어떠한 명령도 없었더라도 바로 그
동일한 일을 확실히 했을 것이라는 사실이, 사실로서 그러한 경우, 어떠한
관련성을 가지는가? 이러한 어려움은 특히 법의 경우에 더욱 첨예하게
나타나는데, 그중 일부는 많은 사람들이 애초에 할 생각조차 하지 않았을
행위를 금지하기 때문이다. 이러한 어려움들이 해결되기 전까지는, 한 나라의
법에 대한 `일반적 복종의 습관(general habit of obedience)'이라는 전체적
관념은 다소 불분명한 상태로 남을 수밖에 없다. 그러나 우리의 현재 목적을
위해서는, `습관'과 `복종'이라는 단어들이 비교적 분명하게 적용될 수
있다고 인정될 법한 매우 단순한 사례를 상정해 볼 수 있을 것이다.

\textbf{(p.~52)} 우리는 다음과 같은 상황을 가정해 보자. 어떤 영토에
인구가 거주하고 있고, 그곳에서는 절대군주인 군주(Rex)가 매우 오랜 기간
동안 통치하고 있다. 그는 일반적 명령을 위협으로 뒷받침하여 사람들을
통제하며, 그들이 그렇지 않았다면 하지 않았을 여러 행위를 하도록
요구하고, 또 그렇지 않았다면 했을 행위를 하지 않도록 요구한다. 통치
초기에는 혼란이 있었으나, 오래전에 사태는 안정되었고, 일반적으로
사람들은 그에게 복종할 것이라고 기대할 수 있다. Rex가 요구하는 내용은
종종 부담스러운 것이며, 복종하지 않고 처벌을 감수하고자 하는 유혹도
상당하기 때문에, 일반적으로 복종이 이루어진다 하더라도 그것이 이 단어의
온전한 의미나 가장 통상적인 의미에서의 `습관(habit)' 또는
`습관적인(habitual)' 복종이라고 보기는 어렵다. 사람들은 실제로 어떤 법에
대한 준수하기 습관을 문자 그대로의 의미에서 획득할 수 있다. 예컨대
영국인에게 도로의 좌측으로 운전하는 것은 그러한 획득된 습관의 전형적인
사례일 것이다. 그러나 세금 납부를 요구하는 법처럼 강한 성향에 반하는
경우에는, 비록 우리가 그것을, 심지어 정기적으로(regular), 결과적으로
준수(compliance)하더라도 그 준수는 습관이 지니는 무반성적이고, 노력이
들지 않으며, 몸에 배어 있는 성격을 갖지 않는다. 그럼에도 불구하고
Rex에게 부여되는 복종은 흔히 이러한 습관적 요소를 결여하고 있을지라도,
다른 중요한 요소들은 지니게 된다. 어떤 사람이 예컨대 아침 식사 때 신문을
읽는 습관을 가지고 있다고 말하는 것은, 그가 상당한 기간 동안 그렇게
해왔고, 또한 그러한 행위를 반복할 개연성이 있다는(likely) 것을 함축한다.
그렇다면 상정된 이 공동체의 대부분의 사람들에 대하여, 초기의 혼란기를
지난 어느 시점 이후에는, 그들이 일반적으로 Rex의 명령에 복종해 왔고
앞으로도 계속 그렇게 할 개연성이 있다는(likely) 점은 참일 것이다.

여기서 주목할 점은, Rex의 통치 하에서의 이 사회적 상황에 대한 설명에
따르면, 복종의 습관(habit of obedience)은 Rex와 각 수범자 간의 개별적
관계라는 점이다. 각 개인은, Rex가 여러 사람에게 내린 명령 중 자신에게
해당하는 것을 정기적으로(regularly) 수행한다. 우리가 ``이 \emphksans{인구
전체(population)}가 그런 습관을 갖고 있다''고 말할 경우, 이는 ``사람들이
토요일 밤마다 주점에 가는 습관이 있다''고 말할 때와 마찬가지로, 단지
대다수 사람들의 습관이 수렴한다(convergent)는 의미일 뿐이다. 즉, 그들
각각이 Rex의 명령에 복종하듯, 각각이 토요일 밤마다 주점을 찾는 것이다.

이 매우 단순한 상황에서는, 공동체가 Rex를 주권자로 성립시키기 위해
요구되는 것은 오직 인구 구성원들의 개인적 복종 행위(personal acts of
obedience)뿐이라는 점에 유의할 필요가 있다. 각 개인은 자기 몫으로 단지
복종하기만 하면 되고, \textbf{(p.~53)} 복종이 정기적으로(regularly)
이루어지기만 한다면, 공동체의 누구도 자기 자신의 복종이나 타인의 Rex에
대한 복종이 어떤 의미에서 옳은지(right), 적절한지(proper), 또는 정당하게
요구되는 것인지(legitimately demanded)에 관하여 어떠한 견해를 가질
필요도, 그것을 표현할 필요도 없다. 분명히 우리가 묘사한 이 사회는 복종의
습관(habit of obedience)이라는 개념에 가능한 한 문자 그대로 적용하기
위해 설정된, 매우 단순한 사회이다. 이는 아마도 실제로는 어디에서도
존재한 적이 없을 만큼 지나치게 단순한 사회일 가능성이 크며, 또한 결코
원시 사회도 아니다. 왜냐하면 원시 사회는 Rex와 같은 절대적
지배자(absolute rulers)를 거의 알지 못하고, 그 구성원들 역시 단지
복종하는 데에만 관심을 두기보다는, 관련된 모든 사람들의 복종이 옳은지에
관하여 뚜렷한 견해를 가지는 것이 보통이기 때문이다. 그럼에도 불구하고,
Rex의 지배하에 있는 이 공동체는 적어도 Rex의 생존 기간 동안에는 법에
의해 지배되는 사회가 지니는 몇 가지 중요한 표지를 분명히 갖추고 있다. 이
공동체는 심지어 일정한 통일성(unity)도 지니고 있어서 `국가(state)'라고
불릴 수도 있다. 이러한 통일성은 구성원들이 복종이 옳은지에 대한 견해를
전혀 공유하지 않더라도, 모두가 동일한 한 사람에게 복종한다는 사실에 의해
구성된다.

이제 성공적인 통치 이후 렉스(Rex)가 사망하고, 그의 아들 Rex 2세가 남아
일반적 명령(general orders)을 발하기 시작했다고 가정해 보자. Rex 1세의
생전 동안 그에게 일반적인 복종의 습관(general habit of obedience)이
존재했다는 사실만으로는, 그것만으로 Rex 2세 역시 습관적으로 복종을 받게
될 가망조차(probable) 말할 수 없다. 따라서 우리가 Rex 1세에 대한 복종의
사실과, 그가 계속 복종을 받을 개연성(likelihood)만을 근거로 삼는다면,
Rex 1세의 마지막 명령에 대해서는 말할 수 있었던 것---즉 그것이
주권자(sovereign)에 의해 내려졌고 따라서 법(law)이었다는 점---을 Rex
2세의 첫 번째 명령에 대해서는 말할 수 없게 된다. 아직 Rex 2세에 대한
확립된 복종의 습관이 존재하지 않기 때문이다. 우리는 이 이론에 따르면,
Rex 2세가 이제 주권자이며 그의 명령이 법이라고 말할 수 있으려면, 그의
부친에게 그랬던 것처럼 Rex 2세에게도 복종이 실제로 부여될 것인지를
기다려 보아야 한다. 출발 시점에서 그를 주권자로 만들어 주는 것은
아무것도 없다. 그의 명령들이 일정 기간 동안 복종을 받아 왔다는 사실을
알게 된 이후에야 비로소, 복종의 습관이 확립되었다고 말할 수 있을 것이다.
그때가 되어서야, 그리고 그때에 이르러서야, 우리는 이후에 내려지는 어떤
명령에 대해서도 그것이 발해지는 즉시, 그리고 실제로 복종되기 이전에도
이미 법이라고 말할 수 있게 된다. 이 단계에 도달하기 전까지는, 어떠한
법도 제정될 수 없는 공백기(interregnum)가 존재하게 된다.

이러한 사태는 물론 가능하며 혼란의 시기에는 때때로 현실화되기도 했다.
그러나 단절(discontinuity)의 위험은 분명하며, 보통은 이를 감수하려 하지
않는다. 그 대신, \textbf{(p.~54)} 절대군주제하에서도 법체계(legal
system)의 특징은 한 입법자에서 다른 입법자로의 이행을 연결하는 규칙들에
의해 입법권(law-making power)의 중단 없는 연속성(uninterrupted
continuity)을 확보하는 데 있다. 이러한 규칙들은 입법자의 자격과 그를
결정하는 방식(mode)을 \emphksans{사전에(in advance)} 규율하며, 이를
지명하거나 일반적 용어로 특정한다. 현대 민주국가에서는 이러한 자격이
매우 복잡하고 구성원이 빈번히 교체되는 입법부의 구성과 관련되지만,
연속성을 위해 요구되는 규칙의 본질은 우리가 상정한 가상의 군주제에
적합한 보다 단순한 형태에서도 확인될 수 있다. 만약 규칙이 장자의 계승을
규정한다면, Rex 2세는 부친의 뒤를 이을 \emphksans{권원(title)}을 가진다. 그는
부친의 사망 시 법을 제정할 \emphksans{권리(right)}를 가지며, 그의 첫 명령이
발해질 때에는 그 개인과 수범자들 사이에 습관적 복종(habitual
obedience)의 관계가 성립할 시간이 아직 없었더라도, 그것들이 이미
법이라고 말할 좋은 이유(good reason)가 있을 수 있다. 실제로 그러한
관계는 끝내 성립하지 않을 수도 있다. 그럼에도 그의 말은 법이 될 수 있다.
왜냐하면 Rex 2세는 첫 명령을 발한 직후 곧바로 사망할 수도 있으며, 그는
복종을 받으며 살 기회를 갖지 못했더라도 법을 제정할 권리를 가졌고, 그의
명령은 법일 수 있기 때문이다.

개별 입법자(individual legislators)의 교체되는 승계를 통해
입법권(law-making power)의 연속성을 설명할 때, `승계 규칙(rule of
succession)', `권원(title)', `승계할 권리(right to succeed)', `법을
제정할 권리(right to make law)'와 같은 표현을 사용하는 것은 자연스럽다.
그러나 분명한 것은, 이러한 표현들과 함께 우리는 복종의 습관(habits of
obedience)에 기초한 일반적 명령(general orders)만으로는 설명될 수 없는
새로운 요소들의 집합을 도입하게 된다는 점이다. 주권 이론(theory of
sovereignty)의 처방을 따라 Rex 1세의 단순한 법적 세계를 구성할 때
사용했던 그러한 요소들로는, 이 새로운 요소들을 설명할 수 없다. 왜냐하면
그 세계에는 규칙이 존재하지 않았고, 따라서 권리나 권원도 존재하지
않았으며, 그러므로 \emphksans{더더군다나(a fortiori)} 승계할 권리나 권원은
존재하지 않았기 때문이다. 그곳에는 단지 Rex 1세가 명령을 발했고, 그의
명령이 습관적으로 복종되었다는 사실만이 있었다. 그의 생존 기간 동안
Rex를 주권자로 구성하고 그의 명령을 법으로 만드는 데에는 그 이상이
필요하지 않았지만, 이것만으로는 그의 후계자의 \emphksans{권리들(rights)}을
설명하기에는 충분하지 않다. 실제로, 복종의 습관이라는 관념은 한 입법자가
다른 입법자를 계승할 때 모든 정상적인 법체계(normal legal system)에서
관찰되는 연속성(continuity)을 두 가지 서로 연관되지만 구별되는 방식으로
설명하지 못한다. 첫째, 한 입법자가 발한 명령에 대한 단순한 복종의
습관만으로는 \textbf{(p.~55)} 새로운 입법자에게 구(舊) 입법자를 계승하여
그의 자리에 서고 명령을 발할 어떤 \emphksans{권리(right)}도 부여할 수 없다.
둘째, 구(舊) 입법자에 대한 습관적 복종만으로는 새로운 입법자의 명령이
복종될 것이라는 가망(probability)을 그 자체로 확보하거나, 그러한
추정(presumption)을 기초지울 수 없다. 만약 승계의 순간에 이러한 권리와
이러한 추정이 존재하려면, 이전 입법자의 통치 기간 동안 사회 어딘가에는
복종의 습관이라는 개념으로는 설명될 수 없는, 그보다 더 복잡한 일반적
사회적 실천(general social practice)이 존재했어야 한다. 즉, 새로운
입법자가 승계할 자격을 가지게 되는 규칙의 수용(acceptance of a rule)이
있었어야 한다.

그렇다면 이러한 보다 더 복합적인 실천(more complex practice)이란
무엇인가? `규칙의 수용(acceptance of a rule)'이란 무엇인가? 여기서
우리는 제1장에서 이미 개괄해 두었던 탐구를 다시 이어가야 한다. 이 질문에
답하기 위해 우리는 한동안 법 규칙의 특수한 경우에서 벗어나야 한다.
습관(habit)과 규칙(rule)은 어떻게 다른가? 사람들이 토요일 밤마다
영화관에 가는 습관이 있다고 말하는 것과, 교회에 들어설 때 남성은 반드시
모자를 벗어야 한다는 규칙이 있다고 말하는 것 사이에는 어떤 차이가
있는가? 우리는 이미 제1장에서 이러한 유형의 규칙을 분석하기 위해 반드시
도입되어야 할 몇 가지 요소들을 언급한 바 있으며, 이제는 그 분석을 한
걸음 더 진전시켜야 할 것이다.

사회적 규칙과 습관 사이에는 확실히 하나의 유사점이 존재한다. 두 경우
모두에서 문제 되는 행위(예: 교회에서 모자를 벗는 행위)는 반드시
전반적으로 일반적이어야 하며, 반드시 항상 불변할 필요는 없다. 즉, 기회가
생기면 대다수 집단 구성원이 그 행위를 반복한다는 점에서 ``그들은 그것을
\emphksans{규칙으로써(as a rule)} 한다''고 말하는 표현은 바로 이 의미를 담고
있다. 하지만 이러한 유사성이 있음에도 불구하고, 양자 사이에는 세 가지
중요한 차이점이 있다.

첫째, 집단이 어떤 \emphksans{습관(habit)}을 갖고 있다고 하기 위해서는, 단지
그들의 행동이 실제로 수렴하면(converge) 충분하다. 즉, 정기적인(regular)
코스에서 벗어나는 일이 있어도 그것이 어떤 비판의 대상이 될 필요는 없다.
그러나 행동의 일반적 수렴 또는 일치만으로는, 그 행동을 요구하는 규칙이
존재한다고 말하기에 충분하지 않다. 규칙이 있는 경우에는, 그로부터의
이탈은 일반적으로 실수(lapse) 또는 과실(fault)로 간주되어 비판의 대상이
되고, 이탈이 예상될 경우 동조를 요구하는 압력이 발생한다. 이러한 비판과
압력의 양상은 규칙의 유형에 따라 다양하다.

둘째로, 그러한 규칙들이 존재하는 경우에는, 실제로 그러한 비판이 이루어질
뿐만 아니라, 그 기준에서의 이탈 자체가 일반적으로 그러한 비판을 제기하기
위한 하나의 \emphksans{좋은 이유(good reason)}로 받아들여진다. 기준에서의
이탈에 대한 비판은 \textbf{(p.~56)} 이와 같은 의미에서
정당하거나(legitimate) 정당화된(justified) 것으로 간주되며, 이탈이
예상될 때 그 기준에의 준수를 요구하는 요구(demands) 또한 마찬가지로
정당한 것으로 간주된다. 더 나아가, 완고한 소수의 상습적 위반자(hardened
offenders)를 제외하면, 이러한 비판과 요구는 그것을 제기하는 사람들뿐만
아니라 그것을 받는 사람들에 의해서도 일반적으로 정당한 것으로, 또는 좋은
이유에 근거한 것으로 받아들여진다. 이러한 여러 방식으로 집단 구성원들 중
몇 명이, 그리고 얼마나 자주 또 얼마나 오랫동안, 행태의 정기적인
방식(regular mode of behaviour)을 비판의 기준(standard of criticism)으로
취급해야 비로소 그 집단이 하나의 규칙(rule)을 가지고 있다고 말할 수
있는지는 명확히 정해진 문제가 아니다. 이는 사람이 대머리라고 불릴 수
있으면서도 머리카락을 몇 가닥이나 가질 수 있는지에 관한 문제와
마찬가지로, 우리를 지나치게 괴롭힐 필요는 없다. 우리가 기억해야 할 것은,
어떤 집단이 특정한 규칙을 가지고 있다는 진술은, 그 규칙을 위반할 뿐만
아니라 그 규칙을 자기 자신이나 타인에 대해 기준으로 보기를 거부하는
소수가 존재하는 것과도 양립 가능하다는 점이다.

사회적 규칙(social rules)을 습관(habits)과 구별하는 세 번째 특징은 이미
앞에서 말한 내용 속에 암묵적으로 포함되어 있지만, 법리학에서 매우
중요함에도 불구하고 자주 간과되거나 오인되어 온 것이므로 여기서 이를 좀
더 상세히 설명하고자 한다. 이 특징을 우리는 이 책 전반에 걸쳐 규칙의
\emphksans{내적 측면(internal aspect)}이라고 부를 것이다. 어떤 습관이 사회
집단 내에서 일반화되어 있을 때, 이러한 일반성은 그 집단 구성원 대다수의
관찰 가능한 행동에 관한 하나의 사실에 불과하다. 그러한 습관이 존재하기
위해서는, 집단의 구성원들 중 누구도 그 일반적인 행동을 어떤 방식으로든
의식할 필요가 없고, 나아가 문제의 행동이 일반적이라는 사실을 알 필요조차
없다. 더더욱 그 행동을 가르치거나 유지하려고 노력하거나 의도할 필요도
없다. 각자가 다른 사람들도 실제로 그렇게 행동하고 있다는 점에서 자신
또한 그렇게 행동하기만 하면 충분하다. 이에 반해, 사회적 규칙이
존재하려면 적어도 일부 구성원들은 문제의 행동을 집단 전체가 따라야 할
일반적 기준(general standard)으로 간주해야 한다. 사회적 규칙은 사회적
습관과 공유하는 외적 측면(external aspect), 즉 관찰자가 기록할 수 있는
정기적이고 획일적인 행동(regular uniform behaviour)이라는 측면에 더하여,
이러한 `내적' 측면을 함께 가진다.

이러한 규칙의 내적 측면은 어떤 게임의 규칙에서 쉽게 예시할 수 있다. 체스
플레이어들은 단지 외부 관찰자가 ``퀸은 저렇게 움직이는구나''라고 기록할
수 있을 정도로 비슷한 방식으로 퀸을 움직이는 습관만 갖고 있는 것이
아니다. 이에 더해, \textbf{(p.~57)} 그들은 이 행동 양식에 대해
성찰적이고 비판적인 태도를 취한다. 즉, 그것을 체스를 두는 사람 모두가
따라야 할 기준으로 본다. 각자는 단지 퀸을 일정한 방식으로 움직이는 것에
그치지 않고, 모두가 그렇게 해야 한다고 생각한다. 이러한 생각은, 다른
사람이 규칙을 어기거나 어길 가능성이 있을 때, 그를 비판하고 동조를
요구하는 방식으로 표현되며, 다른 사람이 자신에게 비판이나 요구를 할 때
이를 정당한 것으로 받아들이는 방식에서도 드러난다. 이러한 비판, 요구,
수용의 표현에는 다양한 규범적 언어(normative language)가 사용된다.
예컨대 ``퀸을 그렇게 움직이면 안 됐어(I/You ought not to have moved the
Queen like that)'', ``그렇게 해야 해(I/You must do that)'', ``그건
맞아(That is right)'', ``그건 틀렸어(That is wrong)'' 등이 이에
해당한다.

규칙의 내적 측면(internal aspect)은 종종 외적으로 관찰 가능한 물리적
행동(physical behaviour)에 대비되는 단순한 `감정(feelings)'의 문제로
오도되어 이해된다. 물론 규칙이 어떤 사회 집단에 의해 일반적으로
수용되고, 사회적 비판과 순응을 요구하는 압력에 의해 일반적으로 지지되는
경우, 개인들은 제약(restriction)이나 강제(compulsion)의 경험과 유사한
심리적 경험을 흔히 가질 수 있다. 사람들이 자신이 특정한 방식으로
행동하도록 `구속되어 있다고 느낀다(feel bound)'고 말할 때, 실제로는
이러한 경험들을 가리키고 있을 수도 있다. 그러나 이러한 감정들은 `구속력
있는(binding)' 규칙의 존재를 위해 필요조건도 아니고 충분조건도 아니다.
사람들이 특정한 규칙을 수용하면서도 그러한 강제의 감정을 전혀 경험하지
않는다고 말하는 데에는 아무런 모순도 없다. 필요한 것은, 일정한 행동
패턴(patterns of behaviour)을 공동의 기준(common standard)으로 간주하는
비판적 성찰적 태도가 존재하는 것이며, 이러한 태도가 비판(자기비판을
포함하여), 순응 요구(demands for conformity), 그리고 그러한 비판과
요구가 정당화될 수 있다는 인정(acknowledgements) 속에서 드러나는 것이다.
이 모든 것은 `\textasciitilde 해야 한다(ought)', `\textasciitilde 해야만
한다(must)', `\textasciitilde 하는 것이 당연하다(should)',
`옳다(right)'와 `그르다(wrong)'라는 규범적 용어(normative
terminology)에서 그 특유한 표현을 발견한다.

이것들이 바로 사회적 규칙(social rules)을 단순한 집단적 습관(mere group
habits)과 구별하는 핵심적 특징들이며, 이러한 점들을 염두에 두고 우리는
다시 법으로 돌아갈 수 있다. 우리는 우리의 사회 집단이 교회에서 모자를
벗는 것에 관한 규칙처럼 특정한 유형의 행위를 기준(standard)으로 만드는
규칙들뿐만 아니라, 일정한 사람의 말---구두로 표현되었든 문서로
표현되었든---을 참조함으로써(by reference to) 보다 간접적인 방식으로
행위의 기준을 식별하도록 하는 규칙을 가지고 있다고 가정할 수 있다. 가장
단순한 형태에서 이러한 \textbf{(p.~58)} 규칙은, Rex가 (아마도 일정한
형식적 방식으로) 지정하는 어떠한 행위들이 행해져야 한다는 내용일 것이다.
이는 우리가 처음에 Rex에 대한 단순한 복종의 습관이라는 관점에서 묘사했던
상황을 변형시킨다. 왜냐하면 이러한 규칙이 수용되는 경우, Rex는 단지
사실상 무엇을 해야 하는지를 지정하는 데 그치지 않고, 그렇게 할
\emphksans{권리}(the \emph{right})를 가지게 되며, 그의 명령에 대한 일반적인
복종이 존재할 뿐 아니라, 그에게 복종하는 것이 \emphksans{옳다(right)}(it is
\emph{right} to obey him)는 점 또한 일반적으로 수용되기 때문이다. Rex는
사실상 입법할 \emphksans{권위(authority)}를 가진 입법자(legislator)가 되며,
즉 집단의 삶 속에 새로운 행위 기준(standards of behaviour)을 도입할
권한을 갖게 된다. 그리고 이제 우리가 `명령'이 아니라 기준(standards)을
다루고 있는 이상, 그가 자신의 입법에 의해 구속되어서는 안 될 이유는 전혀
없다.

이러한 입법적 권위(legislative authority)의 기초를 이루는 사회적
실천들(social practices)은 교회에서 머리를 벗는 행위에 관한 규칙과 같은
단순하고 직접적인 행위 규칙(simple direct rules of conduct)의 기초를
이루는 실천들과, 모든 본질적인 점들에 있어, 동일하며, 우리는 이제 이를
단순한 관습 규칙(mere customary rules)으로 구별할 수 있고, 또한 일반적
습관(general habits)과는 동일한 방식으로 구별된다. 이제 Rex의
\emphksans{말(word)}은 행위의 기준이 되어, 그가 지정한 행위로부터의 이탈은
비판의 대상이 될 수 있으며; 그의 말은 일반적으로 비판과 준수
요구(demands for compliance)를 정당화하는 근거로 참조되고 수용될 것이다.

이러한 규칙들이 입법적 권위의 연속성(continuity of legislative
authority)을 어떻게 설명하는지를 이해하기 위해서는, 어떤 경우에는 새로운
입법자가 실제로 입법을 시작하기도 전에, 그가 하나의 \emphksans{계층(class)}
또는 계보(line)에 속한 사람으로서 차례가 되었을 때 입법할 권리(right)를
가진다는 점을 부여하는 확고히 확립된 규칙이 존재함이 이미 분명할 수
있다는 점에 주목하기만 하면 된다. 예컨대 Rex 1세의 생존 기간 동안, 집단
전체에 의해 일반적으로 수용되는 바에 따르면, 그 말(word)이 복종의 대상이
되는 사람은 개별 인물 Rex 1세에 한정되는 것이 아니라, 일정한 방식으로
자격을 갖춘 사람, 예컨대 특정 조상의 직계(line) 중 생존한 최연장자로서
그 시점에 자격을 갖춘 바로 그 사람이라는 것이다. Rex 1세는 단지 특정한
시점에 그러한 자격을 갖춘 특정 인물에 불과하다. 이러한 규칙은 Rex
1세에게 복종하는 습관과는 달리, 현재의 실제 입법자뿐만 아니라 장차
가능하게 등장할 미래의 입법자들까지 지시한다는 점에서 미래를 향해 있다.

이러한 규칙의 수용(acceptance), 그리고 따라서 그 존재(existence)는 Rex
1세의 생존 기간 동안 부분적으로는 그에 대한 복종을 통해 드러나지만, 또한
일반 규칙(general rule)에 따른 그의 자격(qualification)에 의해 복종이
그가 권리를 가진 어떤 것이라는 점에 대한 인정(acknowledgements)을
통해서도 드러난다. \textbf{(p.~59)} 어떤 집단에 의해 특정 시점에 수용된
규칙의 범위(scope)가 이와 같은 방식으로 일반적으로 입법자 공직(office of
legislator)의 후임자를 향해 열려 있을 수 있다는 바로 그 이유 때문에, 그
규칙의 수용은 후임자가 실제로 입법을 시작하기 이전부터 입법할 권리(right
to legislate)를 가진다는 법적 진술(statement of law)에 대해서도, 그리고
그가 전임자와 동일한 복종을 받을 개연성이 있다는(likely) 사실
진술(statement of fact)에 대해서도 \emphksans{둘 모두(both)}에 대한
근거(grounds)를 제공한다.

물론, 어떤 사회가 어느 한 시점에서 규칙을 수용한다고 해서 그 규칙의
지속적인 존재가 \emphksans{보장되는(guarantee)} 것은 아니다. 혁명이 일어날 수
있으며, 그 사회가 해당 규칙을 더 이상 수용하지 않게 될 수도 있다. 이는
한 입법자, Rex 1세의 생존 기간 중에 발생할 수도 있고, 새로운 입법자 Rex
2세로의 이행 시점에서 발생할 수도 있으며, 만약 그러한 일이 발생한다면
Rex 1세는 입법할 권리(right to legislate)를 상실하게 되거나, Rex 2세는
그 권리를 취득하지 못하게 된다. 상황이 불분명해질 수도 있다는 점은
사실이다. 즉, 단순한 봉기(insurrection)나 기존 규칙의 일시적
중단(temporary interruption)에 직면해 있는 것인지, 아니면 그것의
전면적이고 실효적인 포기(full-scale effective abandonment)에 직면해 있는
것인지가 명확하지 않은 중간적이고 혼란스러운 단계가 존재할 수 있다.
그러나 원칙적으로 이 문제는 분명하다. 새로운 입법자가 입법할 권리를
가진다는 진술(statement)은, 그가 그 권리를 가지게 하는 규칙이 사회 집단
내에 존재함을 전제한다. 만약 현재 그에게 자격을 부여하는 규칙이,
동일하게 그의 전임자에게도 자격을 부여했던 것으로서, 전임자의 생존 기간
동안 수용되었음이 분명하다면, 반대 증거가 없는 한, 그 규칙은 포기되지
않았으며 여전히 존재하고 있다고 가정되어야 한다. 유사한 연속성은
경기(game)에서도 관찰된다. 즉, 득점 기록자는 지난 이닝(last innings)
이후 경기 규칙이 변경되었다는 증거가 없는 한, 새로운 타자가 통상적인
방식으로 평가되어 만들어낸 출루를 그에게 귀속시킨다.

Rex 1세와 Rex 2세의 단순한 법적 세계를 고찰하는 것만으로도, 대부분의
법체계를 특징짓는 입법 권위(legislative authority)의 연속성이 한 규칙의
수용을 구성하는 사회적 실천의 한 형태에 의존하며, 우리가 이미 지적한
방식들에서 단순한 습관적 복종의 사실들과는 구별된다는 점을 보여주기에
충분할 것이다. 논증은 다음과 같이 요약될 수 있다. Rex와 같이 그의 일반적
명령이 습관적으로 복종되는 사람이 입법자(legislator)라 불릴 수 있고 그의
명령들이 법이라 불릴 수 있다는 점을 우리가 인정하더라도(concede), 그러한
입법자들의 연속(succession)에 속한 각자에 대한 복종의 습관만으로는
후임자(successor)가 계승할 \emphksans{권리(right)}와 그에 따른 입법적 권한의
연속성을 설명하기에 충분하지 않다. 첫째, \textbf{(p.~60)} 습관은
`규범적이지(normative)' 않다. 즉, 그것들은 누구에게도 권리(rights)나
권위(authority)를 부여할 수 없다. 둘째, 한 개인에 대한 복종의 습관은,
수용된 규칙(accepted rules)과는 달리, 현재의 입법자뿐만 아니라 장차
연속적으로 등장할 입법자들의 계보(class or line of future successive
legislators)을 참조하거나(refer to), 그들에 대한 복종을
개연적으로(likely) 만들 수 없다. 따라서 한 입법자에 대한 습관적 복종이
존재한다는 사실은, 그의 후임자가 입법할 권리를 가진다는
진술(statement)에 대해서도, 그가 실제로 복종을 받을 개연성이
있다는(likely) 사실 진술(factual statement)에 대해서도, 아무런
근거(grounds)를 제공하지 않는다.

그러나 이 지점에서, 우리가 뒤의 장에서 충분히 전개할 중요한 논점 하나를
주목하지 않을 수 없다. 이는 오스틴의 이론이 지니는 강점 가운데 하나를
이룬다. 수용된 규칙(accepted rules)과 습관(habits) 사이의 본질적 차이를
드러내기 위해 우리는 매우 단순한 형태의 사회를 가정해 왔다.
주권(sovereignty)의 이 측면을 떠나기 전에, 입법할 권위를 부여하는
규칙(rule conferring authority to legislate)의 수용에 대한 우리의 설명이
어느 정도까지 현대적 국가(modern state)에 얼마나 적용될 수 있을지
검토해야 한다. 우리의 단순한 사회를 언급하면서 우리는 마치 대부분의 일반
시민들이 법에 단지 복종할 뿐만 아니라, 법 제정자의 계승을 정당화하는
규칙을 이해하고 수용하고 있는 것처럼 말했다. 단순한 사회에서는 이것이
사실일 수 있다. 그러나 현대 국가에서는, 아무리 법을 잘
지키는(law-abiding) 사람들이라 하더라도, 끊임없이 변화하는 입법권한이
부여된 사람들의 집단의 자격(qualifications)을 구체화하는 규칙들에 대해
인구 대다수가 명확한 인식을 가지고 있다고 생각하는 것은 터무니없을
것이다. 어떤 소규모 부족의 구성원들이 그들의 연속적인 추장에게 권위를
부여하는 규칙을 수용하는 것과 동일한 방식으로, 대중이 이러한 규칙들을
`수용한다(accept)'고 말하는 것은, 보통의 시민들이 실제로는 갖고 있지
않을지도 모르는 헌법적 사안에 대한 이해를 그들의 머릿속에 집어넣는 것을
함의하게 된다. 우리는 이러한 이해를 체계의 공무담당자들 또는
전문가들에게만 요구하면 된다. 즉, 법이 무엇인지를 결정할
책임(responsibility)을 부여받은 법원(courts)과, 일반 시민이 법이
무엇인지 알고자 할 때 자문하는 변호사들(lawyers)이 그것이다.

단순한 부족 사회와 현대 국가 사이의 이러한 차이들은 주목할 가치가 있다.
그렇다면, 연속적으로 교체되는 입법자들의 변화 속에서도 유지되는 의회의
여왕(Queen in Parliament)의 입법적 권위의 연속성을, 일반적으로 수용된
어떤 근본 규칙 또는 규칙들에 기초한 것으로 우리는 어떤 의미에서 이해해야
하는가? 분명히 여기서의 일반적 수용(general acceptance)은
\textbf{(p.~61)} 복합적인 현상으로서, 공무담당자와 일반 시민 사이에 어떤
의미에서 분할되어 있으며, 이들은 서로 다른 방식으로 이에 기여함으로써
법체계의 \emphksans{존재(existence)}에 기여한다. 체계의 공무담당자들은 입법
권위를 부여하는 이러한 근본 규칙들을 명시적으로 인정(acknowledge
explicitly)한다고 말할 수 있다. 입법자들은 자신들에게 권한을
부여하는(empower) 규칙들에 따라 법을 제정할 때 그렇게 하고, 법원은
그러한 자격을 갖춘 자들에 의해 제정된 법을 자신들이 적용할 법으로
식별할(identify) 때 그렇게 하며, 전문가들은 그렇게 제정된 법을 참조하여
일반 시민들의 행위를 향도할(guide) 때 그렇게 한다. 일반 시민은 이러한
공적 작용(official operations)의 결과에 대한 묵인(acquiescence)을 통해
자신의 수용을 주로 표출한다. 그는 이러한 방식으로 제정되고 식별된 법을
지키고(keeps the law), 그 법에 의해 부여된 청구(claims)를 제기하며
권한(powers)을 행사한다. 그러나 그는 그 법의 기원(origin)이나
제정자(makers)에 대해서는 거의 알지 못할 수도 있다. 어떤 이들은 법에
대해 그것이 단지 `법(the law)'이라는 점 이상을 전혀 알지 못할 수도 있다.
법은 일반 시민들이 하고 싶어 하는 일들을 금지하며(forbids), 그들은 이를
위반할 경우 경찰 공무원에 의해 체포되고 판사에 의해 징역형에 처해질 수
있다는 사실을 알고 있다. 위협에 의해 뒷받침된 명령에 대한 습관적 복종이
법체계의 기초라는 점을 고집하는 이 교설의 강점은, 우리가 법체계의
존재라고 부르는 이 복합적 현상의 상대적으로 수동적인 측면을 현실적인
용어로 사유하도록 강제한다는 데 있다. 이 학설의 약점은, 체계의 공직자나
전문가들의 법 제정(law-making), 법 식별(law-identifying), 법
적용(law-applying) 작용에서 주로---비록 배타적으로는
아니지만---드러나는, 다른 상대적으로 능동적인 측면을 흐리게 하거나
왜곡한다는 데 있다. 우리가 이 복합적인 사회적 현상을 실제로 그러한
것으로 이해하려면, 이 두 측면을 모두 염두에 두어야 한다.

\subsection{2. 법의 지속성}\label{uxbc95uxc758-uxc9c0uxc18duxc131}

1944년, 한 여성은 영국에서 점을 치는 행위로 1735년 제정된
「주술법(Witchcraft Act)」을 위반한 혐의로 기소되어 유죄 판결을
받았다.\footnote{\emph{R.} v. \emph{Duncan}{[}1944{]} 1 KB 713.} 이는
매우 익숙한 법적 현상의 하나를 보여주는 다소 인상적인 사례에 불과하다.
즉, 수세기 전에 제정된 법률이 오늘날에도 여전히 법일 수 있다는 것이다.
그러나 이 현상이 아무리 익숙하다고 하더라도, 이러한 방식의 법의
지속성(persistence)은, \textbf{(p.~62)} 법을 습관적으로 복종받는 한
인물(person)이 내린 명령으로 파악하는 단순한 도식으로는 이해 가능하게
만들 수 없다. 우리는 여기서 사실상 방금 살펴본 입법 권위의
연속성(continuity) 문제의 역(converse)에 직면해 있다. 앞의 경우에는,
복종의 습관이라는 단순한 도식에 기초하여, 입법자 지위(office)의
후임자(successor)가 제정한 최초의 법이, 그가 개인적으로 아직 습관적
복종을 받기 이전임에도 \emphksans{이미(already)} 법이라고 말할 수 있는지가
문제였다. 여기서의 문제는 이와 반대로, 오래전에 사망한 선행 입법자가
제정한 법이, 그에게 습관적으로 복종한다고 말할 수 없는 사회에서 어떻게
\emphksans{여전히(still)} 법일 수 있는가 하는 것이다. 첫 번째 경우와
마찬가지로, 우리가 입법자의 생존 기간(lifetime of the legislator)으로
시야를 한정한다면, 이 단순한 도식에는 아무런 어려움도 발생하지 않는다.
실제로 그것은 왜 「주술법」이 영국에서는 법이었지만 프랑스에서는 법이
아니었는지를 매우 훌륭하게 설명하는 것처럼 보인다. 설령 그 법의 용어들이
프랑스에서 점을 치는 프랑스 시민에게까지 확장되었다 하더라도
마찬가지이다. 물론 불운하게도 영국 법원에 회부된 프랑스인들에게는 그
법이 적용될 수 있었을 것이다. 이 단순한 설명에 따르면, 영국에서는 이
법을 제정한 자들에 대한 복종의 습관이 존재했으나, 프랑스에서는 그러한
습관이 존재하지 않았다는 것이다. 따라서 이 법은 영국에서는 법이었지만
프랑스에서는 법이 아니었던 것이다.

그러나 우리는 법을 그 제정자들의 생존 기간으로만 한정하여 볼 수는 없다.
우리가 설명해야 할 특징은 바로 법이 그 제정자들과 그들에게 습관적으로
복종하던 사람들까지 사라진 이후에도 끈질기게 존속할 수 있는
능력(capacity)에 있기 때문이다. 「주술법」이 동시대의 프랑스인들에게는
법이 아니었다면, 왜 그것이 오늘날 우리에게는 여전히 법인 것인가? 분명히,
어떻게 언어를 늘여 써도, 20세기의 영국인들인 우리가 지금 George 2세와
그의 의회(Parliament)에 습관적으로 복종한다고 말할 수는 없다. 이 점에서
오늘날의 영국인들과 당시의 프랑스인들은 동일하다. 즉, 어느 쪽도 이 법의
제정자에게 습관적으로 복종하지 않았거나 복종하고 있지 않다. 「주술법」이
이 재위 기간 동안 제정된 유일한 법률이라고 하더라도, 그것은 오늘날에도
여전히 영국에서 법일 것이다. 이 `왜 여전히 법인가?'라는 문제에 대한
해답은, 원칙적으로, 우리가 앞서 제기한 `왜 이미 법인가?'라는 첫 번째
문제에 대한 해답과 동일하다. 그것은 주권적 인물(sovereign person)에 대한
복종의 습관이라는 지나치게 단순한 개념을, 사회에서 행동의 기준(standard
of behaviour)을 구성할 말(word)을 발할 자들의 계급(class) 또는 계보(line
of persons), 다시 말해 입법할 \emphksans{권리(right)}를 가진 자들을 규정하는,
현재 수용된 근본 규칙들(currently accepted fundamental rules)이라는
개념으로 대체하는 것을 포함한다. 그러한 규칙은 비록 현재(now) 존재해야
하지만, 어떤 의미에서는 시간 초월적(timeless)인 참조(reference)를 가질
수 있다. 즉, 그것은 \textbf{(p.~63)} 장차의 입법자(future legislator)의
입법 작용(legislative operation)을 지시하며 앞으로 나아갈 뿐만 아니라,
과거의 입법자(past one)의 작용을 지시하며 뒤를 돌아볼 수도 있다.

Rex 왕조라는 단순한 틀로 제시하면, 상황은 다음과 같다. Rex 1세, 2세,
3세로 이어지는 입법자들의 계보에 속한 각자는, 직계(direct line)에 있는
생존한 최연장 후손에게 입법할 권리를 부여하는 동일한 일반 규칙(general
rule)에 따라 자격을 갖출 수 있다. 개별 통치자(individual ruler)가
사망하면 그의 입법 작업(legislative work)은 계속 존속한다. 이는 그것이,
각 입법자가 생존해 있던 때마다 사회의 연속적인 세대(successive
generations)가 존중해 온 일반 규칙이라는 기초 위에 놓여 있기 때문이다.
단순한 사례에서 Rex 1세, 2세, 3세는 모두 동일한 일반 규칙에 따라, 입법을
통해 행동 기준(standards of behaviour)을 도입할 자격을 가진다. 대부분의
법체계에서는 사정이 이처럼 단순하지 않다. 과거의 입법을 법으로 인정하는
현재 수용된 규칙(presently accepted rule)이, 동시대의 입법에 관한 규칙과
서로 다를 수도 있기 때문이다. 그러나 그 기초가 되는 규칙(underlying
rule)이 현재 수용되고 있다는 점을 전제로 한다면, 법의 지속성(persistence
of laws)은, 구성원이 변경된 팀들 사이의 토너먼트 첫 라운드에서 내려진
심판의 판정(decision)이, 세 번째 라운드에서 그 자리를 대신한 심판의
판정과 마찬가지로 최종 결과(final result)에 동일한 관련성(relevance)을
갖는다는 사실만큼이나 신비로운 것이 아니다. 그럼에도 불구하고,
신비롭지는 않다 하더라도, 과거와 미래, 그리고 현재의 입법자들의
명령에까지 권위를 부여하는 수용된 규칙의 관념은, 현재의 입법자에 대한
복종의 습관이라는 관념보다 분명히 더 복잡하고 정교하다. 그렇다면 이러한
복잡성을 제거하고, 위협에 의해 뒷받침된 명령이라는 단순한 개념을 어떤
기지 있는 확장(ingenious extension)을 통해 발전시켜, 법의 지속성이
결국에는 현재의 주권자에 대한 습관적 복종이라는 더 단순한 사실들 위에
기초하고 있음을 보여줄 수는 없는가?

이를 달성하려는 한 가지 기지 있는 시도가 제시된 바 있다. 홉스(Hobbes)는,
여기서 벤담(Bentham)과 오스틴에 의해 반복되듯이, ``입법자(legislator)란
법이 처음 제정되었을 때 그 법에 권위(authority)를 부여한 자가 아니라, 그
법이 지금도 법으로 존속하도록 권위를 부여하는 자이다''라고
말하였다.\footnote{\emph{Leviathan}, chap.~xxvi.} 규칙의 개념을 배제하고
보다 단순한 습관의 관념을 택한다면, 입법자의 `권위(authority)'가 그의
`권한(power)'과 구별되어 무엇을 의미할 수 있는지는 즉각적으로 명확하지
않다. 그러나 이 인용문에 표현된 일반적 \textbf{(p.~64)} 논증은 분명하다.
그것은 「주술법」과 같은 법의 원천(source)이나 기원(origin)이 역사적
사실로서 과거의 주권자(past sovereign)의 입법 작용에 있었음에도
불구하고, 20세기 잉글랜드에서 그 법이 현재 법으로서의 지위(status as
law)를 가지는 것은 현재의 주권자에 의해 법으로 승인(recognition)되기
때문이라는 주장이다. 이러한 인정은, 현재 살아있는 입법자들에 의해 제정된
법률의 경우와 같은 \emphksans{명시적(explicit)} 명령의 형태를 취하지 않고,
주권자의 의지의 \emphksans{묵시적(tacit)} 표현의 형태를 취한다. 이는,
주권자가 그럴 수 있음에도 불구하고, 오래전에 제정된 그 법률을 자신의
대리인들(agents)-법원과 경우에 따라서는 행정부-에 의해 집행되는 것을
방해하지 않는다는 사실로 구성된다.

이는 물론, 어느 누구에 의해서도 어느 시점에서도 명령된 것처럼 보이지
않는 일정한 관습적 규칙(customary rules)의 법적 지위를 설명하기 위해
원용되었던, 이미 검토된 묵시적 명령에 관한 동일한 이론이다. 제III장에서
우리가 이 이론에 대해 제기한 비판들은, 과거의 입법(past legislation)이
법으로서 계속 승인되는 것을 설명하는 데 이 이론이 사용될 때 더욱
분명하게 적용된다. 법원이 합당하지 않은(unreasonable) 관습 규칙을 배제할
수 있는 광범위한 재량(wide discretion)을 부여받고 있기 때문에, 특정
사건에서 법원이 실제로 관습 규칙을 적용하기 전까지는 그것이 법으로서의
지위를 갖지 않는다는 견해에는 어느 정도의 그럴듯함(plausibility)이 있을
수 있다. 그러나 과거의 `주권자'에 의해 제정된 법률이, 특정 사건에서
법원에 의해 실제로 적용되고, 현재의 주권자의 묵인 아래 집행되기 전까지는
법이 아니라는 견해에는 거의 그럴듯함이 없다. 만약 이 이론이 옳다면,
법원은 그것이 이미 법이기 때문에 집행하는 것이 아니라는 결론이 따라야
한다. 그러나 이러한 결론은, 현재의 입법자가 과거의 제정(enactments)을
폐지(repeal)할 수 있었음에도 그 권한(power)을 행사하지 않았다는
사실로부터 도출하기에는 터무니없이 부당한(absurd) 추론이다. 빅토리아
시대 법률들과 오늘날 의회의 여왕(Queen in Parliament)에 의해 제정된
법률들은 분명히 현재의 영국에서 정확히 동일한 법적 지위(legal status)를
가진다. 양자는 모두, 그것들이 적용될 사건들이 법원에 제기되기 이전부터
이미 법이다. 그리고 그러한 사건들이 발생했을 때, 법원은 빅토리아 시대
법률과 현대 법률을 모두, 그것들이 이미 법이기 때문에 적용한다. 어느
경우에서도 이들 법률은 법원에 의해 적용된 이후에야 비로소 법이 되는 것이
아니며, 양 경우 모두에서 그 법으로서의 지위는, 이 법률들이 현재 수용된
규칙들에 따라 권위를 가지는(authoritative) 자들에 의해 제정되었다는
사실에 기인하며, 그 사람들이 생존해 있는지 사망했는지의 여부와는
무관하다.

\textbf{(p.~65)} 과거의 제정법이 현재 법으로서의 지위를 법원에 의한
적용에 대해 현 입법부가 묵인하고 있기 때문에 갖게 된다는 이론의
비일관성(incoherence)은, 오늘날의 법원들이 폐지되지(repealed) 않은
빅토리아 시대 법률을 여전히 법으로 취급하는 반면, Edward 7세 치하에서
폐지된 법률을 더 이상 법이 아닌 것으로 구별해야 하는 이유를 설명하지
못한다는 점에서 가장 분명하게 드러난다. 분명히 이러한 구별을 할 때,
법원들은(그리고 그와 함께 이 체계를 이해하는 변호사나 일반 시민들도)
과거와 현재의 입법 작용을 모두 포괄하는, 무엇이 법으로 간주되어야
하는지를 규정하는 근본 규칙(fundamental rule) 또는 규칙들(rules)을
기준(criterion)으로 사용한다: 그들은 두 법률 사이의 구별을, 현재의
주권자가 한쪽은 묵시적으로 사령(즉, 집행을 허용)하고 다른 쪽은 그렇게
하지 않았다는 사실에 대한 앎(knowledge)에 근거시키지 않는다.

다시 말해, 우리가 배척한 이 이론에서 유일하게 미덕이라고 할 만한 것은
현실주의적 요소를 흐릿한 형태로 상기시켜 준다(a realistic reminder)는
점뿐인 듯하다. 이 경우 그 상기(reminder)란, 해당 체계의 공무담당자들,
특히 무엇보다도 법원이, 특정한 입법 작용이 \emphksans{과거이든 현재이든}
권위를 가진다는 규칙을 수용하지 않는다면, 법으로서의 지위에 필수적인
어떤 것이 결여될 것이라는 점을 상기시키는 것이다. 그러나 이러한 단조로운
유형의 현실주의는, 그 주요 특징들이 뒤에서 상세히 논의되는\footnote{See
  pp.~136--47 below.} 이른바 법적 현실주의(Legal Realism)라는 이론으로
과도하게 확대되어서는 안 된다. 이 이론은 일부 버전에서, 어떤 제정법도
법원에 의해 실제로 적용되기 전까지는 \emphksans{어떠한(no)}
제정법(statutes)도 법이 \emphksans{아니라}고 주장한다. 어떤 법률이 법이 되기
위해서는 법원이 특정한 입법 작용이 법을 만든다는 규칙을 수용해야 한다는
진실(truth)과 법원이 특정 사건에서 그것을 적용하기 전까지는 아무것도
법이 아니라는 오도된 이론(misleading theory) 사이에는, 법을 이해하는 데
있어 결정적으로 중요한 차이가 있다. 물론 법적 현실주의 이론의 일부
버전들은, 우리가 비판해 온 법의 지속성에 대한 거짓 설명보다도 한참
넘어선다. 왜냐하면 그것들은 과거의 주권자에 의해 제정되었든
\emphksans{현재의} 주권자에 의해 제정되었든, 법원이 실제로 적용하기 전에는
어떠한 제정법도 법의 지위(status of law)를 가질 수 없다고 전면적으로
부정하기 때문이다. 그러나 완전한 현실주의 이론(full Realist
theory)에까지 이르지는 않으면서도, 과거의 주권자와 구별하여 현재의
주권자가 제정한 제정법은 법원에 의해 적용되기 이전에도 법이라고 인정하는
한편, 법의 지속성을 설명하려는 이론은 \textbf{(p.~66)} 양쪽 세계의
최악을 취하는 것이며, 분명히 터무니없이 부조리하다. 이러한 중간적 입장은
옹호될 수 없다. 왜냐하면 현재의 주권자가 제정한 법률의 법적 지위와, 이전
주권자에 의해 제정되었으나 폐지되지 않은 법률의 법적 지위를 구별할
아무런 근거도 없기 때문이다. 둘 다(보통의 변호사들이 인정하듯이) 법원이
오늘날의 특정 사건에 적용하기 이전부터 법이거나, 아니면 완전한 현실주의
이론이 주장하듯이 둘 다 법이 아니어야 한다.

\subsection{3. 입법권에 대한 법적
한계}\label{uxc785uxbc95uxad8cuxc5d0-uxb300uxd55c-uxbc95uxc801-uxd55cuxacc4}

주권 교설에서 수범자의 일반적 복종의 습관은 주권자에게서의 그러한 습관의
부재(absence)를, 그것의 상보적 짝(complement)으로, 가지게 된다. 그는
수범자들을 위해 법을 만들며, 어떠한 법에도 속하지 않는 지위(position
outside any law)에서 법을 만든다. 그의 법-생성 권한(law-creating
power)에는 법적 한계(legal limits)이 존재하지 않으며, 또 존재할 수도
없다. 주권자의 법적으로 무제한적인 권한은 정의상(by definition) 그에게
귀속된다는 점을 이해하는 것이 중요하다: 이 이론은 단순히 다음을 단언할
뿐인데, 즉 입법자에게 법적 제한이 존재할 수 있는 경우는 오직 그 입법자가
자신이 습관적으로 복종하는 다른 입법자의 명령 아래에 있을 때뿐이라는
것과, 그 경우 그는 더 이상 주권자가 아니라는 것이다. 그가 주권자라면
그는 어떤 다른 입법자에게도 복종하지 않으며, 따라서 그의 입법 권력에는
법적 제한이 존재할 수 없다. 물론 이 이론의 중요성은, 사실에 대해
아무것도 말해 주지 않는 이러한 정의(definitions)와 그로부터 도출되는
단순하고 필연적인 귀결들(simple necessary consequences)에 있는 것은
아니다. 그 중요성은, 법이 존재하는 모든 사회에는 이러한
속성(attributes)을 지닌 주권자가 존재한다는 다음과 같은 주장(claim)에
있다: 모든 법적 권력이 제한되어 있으며, 어떠한 개인이나 집단도
주권자에게 귀속되는 어떠한 법에도 속하지 않는 지위를 차지하고 있지 않은
것처럼 보이게 하는 법적 또는 정치적 형식 배후를 우리는 들여다볼 필요가
있을지도 모른다. 그러나 우리가 집요하게 추적한다면, 이 이론이 주장하듯이
그러한 형식들 배후에 실제로 존재하는 현실(reality)을 발견하게 될 것이다.

우리는 이 이론을 실제로 그것이 주장하는 바보다 더 약하거나 더 강한
주장을 하는 것으로 오해해서는 안 된다. 이 이론은 법적 한계을 받지 않는
주권자가 존재하는 \emphksans{어떤(some)} 사회가 존재한다는 점만을 진술하는
것이 아니라, 법의 존재는 어디에서나 그러한 주권자의 존재를 함의한다는
점을 주장한다. 한편 이 이론은 주권자의 권한에 아무런 한계도 없다고
주장하는 것이 아니라, 오직 그 권한에 \emphksans{법적} 한계(\emph{legal}
limits)가 없다고 주장할 뿐이다. 따라서 주권자는 입법 권한을 행사함에
있어 실제로는 대중적 입장에 따라 자제할 수도 있다. \textbf{(p.~67)} 이는
이를 멸시할 경우의 결과에 대한 두려움 때문일 수도 있고, 혹은 자신이
도덕적으로 이를 존중하도록 도덕적으로 구속된다고 스스로 생각하고 있기
때문일 수도 있다. 매우 다양한 요인들이 이러한 점에서 그에게 영향을 미칠
수 있으며, 만약 대중의 반란에 대한 두려움이나 도덕적 확신(moral
conviction) 때문에 그가 그렇지 않았다면 했을 방식으로 입법하지 않게
된다면, 그는 실제로 이러한 요인들을 자신의 권한에 대한 `한계'라고
생각하고 말할 수도 있다. 그러나 그것들은 법적 한계(legal limits)는
아니다. 그는 그러한 입법을 삼가야 할 법적 책무(legal duty)를 지고 있지
않으며, 법원(law courts)은 자신들 앞에 주권자의 법이 존재하는지를
판단함에 있어, 대중적 입장이나 도덕성의 요구로부터의 분기가 그 법을
법으로서 인정받지 못하게 만든다는 주장을, 주권자가 그렇게 하라고 명령한
경우가 아니라면, 받아들이지 않을 것이다.

이 이론이 법에 대한 일반적 설명(general account of law)으로서 지니는
매력은 분명하다. 이 이론은 두 가지 주요한 질문에 대해 만족스러울 만큼
단순한 형태로 답을 제시해 주는 것처럼 보인다. 습관적 복종을 받지만
누구에게도 복종을 돌려주지 않는 주권자를 발견하게 되면, 우리는 두 가지
일을 할 수 있다. 첫째, 우리는 그의 일반적 명령 속에서 특정 사회의 법을
식별할 수 있고, 그 사회 구성원들의 삶을 또한 규율하는 다른 많은
규칙(rules), 원칙(principles), 또는 기준(standards)---도덕적이거나
단순히 관습적인(moral or merely customary) 것들과---을 그것으로부터
구별할 수 있다. 둘째, 법의 영역(area of law) 안에서, 우리가 독립적인
법체계(independent legal system)에 직면해 있는지, 아니면 어떤 더
광범위한 체계(wider system)의 종속적인 부분(subordinate part)에 불과한
것에 직면해 있는지를 판단할 수 있다.

일반적으로, 단일하고 연속적인 입법적 존재자(single continuing
legislative entity)로서 이해되는 의회의 여왕(Queen in Parliament)은 이
이론의 요건(requirements)을 충족시키며, 의회의 주권성은 바로 그러한 충족
사실에 존재한다고 주장된다. 이 믿음의 정확성 여부가 어떠하든(그 일부
측면은 제VI장에서 나중에 검토한다), 우리는 이 이론이 요구하는 바를 Rex
1세의 상상속의 단순화된 세계 속에서 충분히 일관되게 재현할 수 있다. 보다
복잡한 현대 국가의 사례를 검토하기에 앞서 이를 수행하는 것은 교육적으로
유익할 것인데, 이 방식으로 이 이론의 전체 함의(full implications)가 가장
잘 드러나기 때문이다. 제1절에서 복종의 습관 개념에 대해 제기된 비판들을
수용하기 위해, 우리는 상황을 습관이 아니라 규칙의 관점에서 구상할 수
있다. 이러한 발판 아래에서, 우리는 법원(courts), 공무담당자, 시민에 의해
일반적으로 수용된 규칙이 존재하는 사회를 상정할 것이다. 그 규칙에
따르면, Rex가 무엇이든 하라고 명령할 때마다 그의 말(word)은 집단을 위한
행동의 기준(standard of behaviour)을 구성한다. 이러한 명령들 가운데에서,
Rex가 `공적' 지위(`official' status)를 부여하고자 하지 않는
\textbf{(p.~68)} `사적' 바람(`private' wishes)의 표현과, 그러한 지위를
부여하고자 하는 표현을 구별하기 위해, 보조적 규칙(ancillary rules)들이
채택될 수도 있다. 이 규칙들은 군주(monarch)가 `군주로서의 성격으로(in
the character of a monarch)' 입법할 때 사용해야 할 특별한 방식(style)을
규정하되, 아내나 정부(情婦)에게 사적 명령을 내릴 때에는 이를 사용하지
않도록 명시할 것이다. 입법의 방식과 형식(manner and form)에 관한 이러한
규칙들은 그 목적을 달성하기 위해서는 진지하게 받아들여져야 하며, 때로는
Rex에게 불편(inconvenience)을 초래할 수도 있다. 그럼에도 불구하고,
우리가 그것들을 법적 규칙(legal rules)으로 분류할 수는 있겠지만, 그의
입법 권한(legislative power)에 대한 `한계(limits)'로 간주할 필요는 없다.
요구된 형식(required form)을 따르기만 한다면, 그의 바람(wishes)을
실현하기 위해 입법할 수 없는 주제는 존재하지 않기 때문이다. 그의 입법
권력의 `형식(form)'은 법에 의해 한계지어질 수 있을지 모르나, 그
`범위(area)'는 법에 의해 한계지어지지 않는다.

법에 대한 일반 이론(general theory of law)으로서 이 이론에 제기될 수
있는 반론은, 이 상상된 사회에서의 Rex와 같은 어떠한 법적 제한에도
종속되지 않는 주권자의 존재가 법의 존재를 위한 필요조건(necessary
condition)이나 전제(presupposition)가 아니라는 점이다. 이를 입증하기
위해 우리는 논쟁의 여지가 있거나 이의를 제기할 만한 법의 유형들을 원용할
필요가 없다. 따라서 우리의 논증은, 입법부가 결여되어 있다는 이유만으로
법이라는 지위(title of law)를 부정하고자 하는 이들도 있는 관습법 체계나
국제법에서 도출되는 것이 아니다. 이러한 사례들에 호소할 필요는 전혀
없다. 왜냐하면 법적으로 무제한적인 주권자라는 개념은, 법이 존재한다는
점에 대해 누구도 의문을 제기하지 않을 많은 현대 국가들에서의 법의
성격(character of law)을 잘못 대변하기 때문이다. 이러한 국가들에는
입법부들이 존재하지만, 체계 내에서의 최고 입법 권한이 무제한과는 거리가
먼 경우도 종종 있다. 성문헌법(written constitution)은, 입법의 형식과
방식을 규정함으로써---이를 우리는 제한으로 보지 않아도 된다고 인정할 수
있지만---입법부의 권능(competence)을 제한할 뿐만 아니라, 입법 권능의
범위에서 특정 사안들을 전면적으로 배제함으로써, 실체적 제한을 부과할
수도 있다.

다시 한 번, 현대 국가의 복잡한 사례를 검토하기에 앞서, Rex가 최고
입법자인 단순화된 세계에서 그의 입법 권한에 대한 `법적 제한'이 실제로
무엇을 의미하는지, 그리고 왜 그것이 완전히 일관된 생각인지를 살펴보는
것이 유용하다.

Rex의 단순화된 사회에서는, (성문헌법에 구현되어 있든 그렇지 않든 간에)
Rex의 어떠한 법도 토착 주민을 영토(territory)에서 배제하거나 재판 없는
구금을 규정한다면 유효하지 않으며, 이러한 규정에 반하는 어떠한 제정도
\textbf{(p.~69)} 무효(void)로서 간주되고 모두에 의해 그렇게 취급된다는
규칙이 수용되어 있을 수도 있다. 이러한 경우 Rex의 입법할 권한은 분명히
법적(legal)이라고 할 수밖에 없는 제한을 받게 될 것이다. 설령 우리가
이러한 근본적인 헌법 규칙(fundamental constitutional rule)을
`\emphksans{하나의} 법(\emph{a} law)'이라고 부르기를 꺼린다 하더라도 말이다.
그가 종종 자신의 성향(inclinations)에 반하여서까지도 자제하게 할 수 있는
대중적 입장이나 대중의 도덕적 확신(popular moral convictions)을 무시하는
경우와는 달리, 이러한 특정한 제약들(specific restrictions)을 무시하는
것은 그의 입법을 무효로 만들 것이다. 따라서 법원(courts)은, 입법자의
권한 행사에 대한 다른 단지 도덕적 또는 \emphksans{사실상의(de facto)}
한계들과는 달리, 이러한 한계들에 관해서는 직접적인 관련을 갖게 될
것이다. 그럼에도, 이러한 법적 제한들에도 불구하고, 그 범위(scope) 내에서
이루어진 Rex의 제정(enactments)은 분명히 법(laws)이며, 그의 사회에는
독립적인 법체계가 존재한다.

우리는 이러한 상상의 단순한 사례를 좀 더 자세히 들여다볼 필요가 있다.
그래야만 이러한 유형의 법적 제한이 정확히 무엇인지 분명해지기 때문이다.
우리는 종종 Rex의 입장을 다음과 같이 표현할 수 있다: ``그는 재판 없이
구금을 부여하는 법을 제정\emphksans{할 수 없다(cannot)}.'' 이때의
``\textasciitilde 할 수 없다(cannot)''는 표현은, 어떤 사람이 어떤 것을
하지 않아야 할 법적 책무(legal duty)나 의무(obligation)를 지고 있다는
의미와는 다르다. 후자의 의미에서 ``할 수 없다(cannot)''가 사용되는
경우는 예를 들어 ``자전거를 인도에서 탈 수 없다(You cannot ride a
bicycle on the pavement.)''고 말하는 경우이다. 최고 입법자의 권한을
실질적으로 제한하는 헌법은, 입법자에게 특정 방식의 입법을 시도하지
말라는 법적 책무를 부과하지는 않는다(또는 반드시 부과할 필요는 없다). 그
대신, 그런 입법은 무효가 된다고 규정할 수 있다. 이런 경우 입법자에게
부과되는 것은 법적 책무(duties)가 아니라, 법적 무능력(disabilities)이다.
따라서 여기서 `제한(limits)'이란 것은 \emphksans{책무}의 현존(the presence of
\emph{duty})가 아니라 법적 권한의 부재(the absence of legal power)를
의미한다.

Rex의 입법 권한에 대한 이러한 제약들은 헌법적(constitutional)이라고 불릴
수 있을 것이다. 그러나 그것들은 법원이 관여하지 않는 단순한
관행(conventions)이나 도덕적 사안(moral matters)이 아니다. 그것들은
입법할 권위를 부여하는 규칙(rule conferring authority to legislate)의
일부를 이루며, 법원에게 중대한 관련성을 가진다. 왜냐하면 법원은 그러한
규칙을 자신들 앞에 제시되는 입법으로 의도된 제정들(purported legislative
enactments)의 유효성(validity)을 판단하는 기준(criterion)으로 사용하기
때문이다. 그럼에도 불구하고, 이러한 제한들이 법적이며 단순히
도덕적이거나 관습적인 것이 아님에도 불구하고, 그 존재나 부재는 Rex가
다른 사람들에게 습관적으로 복종하는지 여부라는 용어로는 표현될 수 없다.
Rex는 그러한 제약들의 적용을 받을 수 있으며, 그것들을 회피하려는 시도를
\textbf{(p.~70)} 전혀 하지 않을 수도 있다; 그럼에도 불구하고, 그가
습관적으로 복종하는 대상은 아무도 없을 수 있다. 그는 단지 유효한
법(valid law)을 제정하기 위한 조건들(conditions)을 충족시킬 뿐이다. 혹은
그는 그 제약들과 불일치하는 명령(orders inconsistent with them)을
발함으로써 그 제약들을 회피하려고 시도할 수도 있다. 그러나 그가 그렇게
하더라도 그는 누구에게도 불복종한 것이 아니며, 상위 입법자의 법을
위반하거나 법적 책무를 침해한 것도 아니다. 그는 분명히 유효한 법을
제정하는 데 실패했을 뿐이며(그는 법을 어기는 것이 아니라), 그 반대의
경우도 마찬가지이다. 즉, Rex에게 입법 자격을 부여하는 헌법적 규칙 안에
Rex의 입법할 권위에 대한 법적 제약이 전혀 존재하지 않는다면, 그가 인접한
영토의 왕인 Tyrannus의 명령에 습관적으로 복종한다는 사실은, Rex의
제정(enactments)을 법으로서의 지위(status as law)에서 박탈하지도 않으며,
그것들이 Tyrannus가 최고 권위를 가지는 단일 체계(single system)의 종속적
부분(subordinate parts)임을 보여주지도 않는다.

앞서 논의한 명백한 고찰들은, 비록 단순한 주권 이론에 의해 자주 가려져
있지만, 법체계의 기초를 이해하는 데 있어 결정적으로 중요한 여러 사항들을
확립해 준다. 이를 다음과 같이 요약할 수 있다: 첫째, 입법 권한에 대한
법적 제한은 입법자에게 어떤 상위 입법자의 명령에 복종하라는 책무를
부과하는 것이 아니라, 입법 자격을 부여하는 규칙에 포함된
무능력(disabilities)의 형태로 존재한다.

둘째, 어떤 법률 제정이 유효한 법인지 입증하기 위해, 그것이 어떤 `주권자'
혹은 `무제한적' 입법자에 의해 명시적 또는 묵시적으로 제정된 것임을
추적할 필요는 없다. 여기서 `무제한적'이라는 말은, 그의 입법 권한이
법적으로 제한받지 않는다는 의미이거나, 그가 다른 누구에게도 습관적으로
복종하지 않는다는 의미일 수 있다. 그러나 우리가 실제로 해야 할 일은,
해당 제정 행위가 현존하는 규칙에 따라 입법 자격을 갖춘 입법자에 의해
이루어진 것이며, 그 규칙에 어떤 제한이 없다거나, 적어도 해당
제정행위에는 적용되지 않는 제한만이 존재함을 보여주는 것이다.

셋째, 어떤 체계가 독립된 법체계인지 판별하기 위해, 그 최고 입법자가
법적으로 무제한적이거나, 타인에게 습관적으로 복종하지 않는다는 점을
보여줄 필요는 없다. 필요한 것은 단지, 그 입법자를 자격화하는 규칙들이,
다른 영토에 대해서도 권한을 가진 자들에게 우월한 권위를 부여하지
않는다는 사실을 보여주는 것이다. 반대로, 그가 외부 권위에 복종하지
않는다고 해서, 곧바로 그가 자기 영토 내에서 무제한적 권한을 가진다는
뜻은 아니다.

넷째, 법적으로 무제한적인 입법 권한과, \textbf{(p.~71)} 제한은 있지만 그
체계 내에서 최고 권한을 갖는 입법자는 구분되어야 한다. Rex는 헌법에 의해
일정한 제한을 받는다 하더라도, 그의 입법이 다른 모든 입법을 폐지할 수
있는 의미에서 그 나라 법에 의해 인식되는 최고 입법 권한을 가질 수 있다.

다섯째, 그리고 마지막으로, 입법자의 입법 권한에 대한 제한 규칙의 존재
여부는 핵심적인 반면, 그 입법자가 타인에게 습관적으로 복종하는 습관이
있는지 여부는, 기껏해야 간접적인 증거로서만 일정한 의미를 갖는다. 만약
입법자가 다른 사람에게 습관적으로 복종하지 않는다는 것이 사실이라면,
이는 때때로 그 입법자의 권한이 다른 이의 권위에 헌법적 혹은 법적으로
종속되지 않을 수 있다는 점에 대한 약간의, 그러나 결코 결정적인 것은
아닌, 증거로 기능할 수 있다. 마찬가지로, 그 입법자가 실제로 어떤 다른
사람에게 습관적으로 복종한다는 사실 역시, 그의 입법 권한이 어떤 규칙
아래에서 다른 사람의 권위에 종속되어 있음을 나타내는 일정한 증거가 될 수
있다.

\subsection{4. 입법부 배후의
주권자}\label{uxc785uxbc95uxbd80-uxbc30uxd6c4uxc758-uxc8fcuxad8cuxc790}

현대 세계에는, 통상적으로 그 체계 내에서 최고 입법부로 간주되는 기관이
입법 권한의 행사에 있어 법적 제한을 받는 법체계들이 다수 존재한다.
그럼에도 불구하고, 법률가와 법이론가 모두가 동의하듯이, 그러한 입법부가
자신에게 부여된 제한된 권한의 범위(scope) 내에서의 제정(enactments)은
분명히 법(law)이다. 이러한 경우들에서, 만약 우리가 법이 존재하는 곳마다
법적 제한을 받을 수 없는 주권자(sovereign incapable of legal
limitation)가 존재한다는 이론을 유지하고자 한다면, 우리는 법적으로
제한된 입법부 배후에 있는 그러한 주권자를 찾아야 한다. 그가 실제로
거기에 존재하는지 여부가, 이제 우리가 검토해야 할 문제이다.

우리는 잠시 동안, 모든 법체계가 어떤 형태로든---비록 반드시 성문헌법에
의해서일 필요는 없지만---입법자들의 자격(qualification)과 입법의 `방식과
형식'에 관해 마련해야 하는 규정들을 논외로 둘 수 있다. 이러한 규정들은
입법 권력의 범위에 대한 법적 제한이라기보다는, 입법기관의
동일성(identity)과 그것이 입법하기 위해 무엇을 해야 하는지를
특정(specifications)하는 것으로 생각될 수 있다. 그러나 실제로는,
남아프리카(South Africa)의 경험이 보여주듯이,\footnote{See \emph{Harris}
  v. \emph{Dönges}{[}1952{]} 1 TLR 1245.}, \textbf{(p.~72)} 입법의
`방식과 형식'에 관한 단순한 규정이나 입법기관의 정의(definitions)를
`실질적(substantial)' 제한과 만족스럽게 구별해 주는 일반적
기준(criteria)을 제시하기는 어렵다. 그럼에도 불구하고 실질적 제한의
명확한 사례들은, 미합중국이나 오스트레일리아의 연방헌법에서 찾아볼 수
있다. 이들 헌법에서는 중앙정부와 구성 주(states) 사이의 권한
분할(division of powers), 그리고 또한 일정한 개인적 권리(individual
rights)들이 통상의 입법 절차에 의해 변경될 수 없다. 이러한 경우, 주
입법부나 연방 입법부 중 어느 쪽이든, 연방적 권한 분할에 반하거나 그것을
변경하려는 것으로 주장되는 제정(enactment), 또는 이러한 방식으로
보호되는 개인적 권리들과 양립할 수 없는 제정은 \emphksans{권한을 초과한(ultra
vires)} 것으로 취급될 수 있으며, 헌법 규정과 충돌하는 범위 내에서 법원에
의해 법적으로 유효하지 않다(legally invalid)고 선언된다. 이러한 입법
권력에 대한 법적 제한들 가운데 가장 유명한 것은 미합중국 헌법의 수정
제5조(Fifth Amendment)이다. 이 조항은, 다른 내용들과 함께, 어떠한 사람도
적법절차(due process of law) 없이 `생명, 자유 또는 재산(life liberty or
property)'을 박탈당해서는 안 된다고 규정한다. 그리고 의회(Congress)의
법률(statutes)이 이러한 규정들 또는 헌법이 그 입법 권력에 부과한 다른
제한들과 충돌하는 것으로 판명될 경우, 법원에 의해 유효하지 않다고
선언되어(declared invalid) 왔다.

물론 헌법의 규정들을 입법부의 작용으로부터 보호하기 위한 다양한 장치들이
존재한다. 어떤 경우들에서는, 예컨대 스위스의 경우처럼, 연방 구성 주의
권리와 개인의 권리에 관한 일부 규정들이 형식상으로는 강행적임에도
불구하고, `단지 정치적(merely political)'이거나 권고적인(hortatory)
것으로 취급된다. 이러한 경우 법원에는 연방 입법부의 제정을
`심사하여(review)', 비록 그것이 입법부의 작용의 적절한 권한 범위에 관한
헌법 규정과 명백히 충돌하더라도 이를 유효하지 않다(invalid)고 선언할
관할권(jurisdiction)이 부여되지 않는다.\footnote{See Art. 113 of the
  Constitution of Switzerland.} 미합중국 헌법의 특정 규정들 역시 `정치적
질문(political questions)'을 제기하는 것으로 판시되어 왔으며, 사건이 이
범주에 속하는 경우 법원은 해당 법률이 헌법을 위반하는지 여부를 심리하지
않는다.

최고 입법부의 통상적 작용(normal operations)에 대한 법적 제한이 헌법에
의해 부과되는 경우, 이러한 제한들 자체는 \textbf{(p.~73)} 특정한 형태의
법적 변경(legal change)으로부터 면책(immune)될 수도 있고 그렇지 않을
수도 있다. 이는 헌법이 그 개정(amendment)을 위해 어떠한 규정을 두고
있는지의 성격에 달려 있다. 대부분의 헌법은, 통상의 입법부와 구별되는
기관(body)에 의해 행사되거나, 또는 통상의 입법부 구성원들이 특별한
절차(special procedure)를 사용하여 행사하는, 광범위한 개정 권한(amending
power)을 포함하고 있다. 미합중국 헌법 제5조가, 각 주의 입법부 4분의 3
또는 각 주의 4분의 3으로 소집된 회의(conventions)에 의해
비준되는(ratified) 개정을 규정하고 있는 것은 전자의 유형에 속하는 개정
권한의 한 사례이며, 1909년 남아프리카법(South Africa Act of 1909)
제152조(s. 152)에 규정된 개정 절차는 후자의 유형에 속하는 사례이다.
그러나 모든 헌법이 개정 권한을 포함하고 있는 것은 아니며, 또한 개정
권한이 존재하는 경우에도, 입법부에 한계를 부과하는 헌법의 특정 규정들이
그 개정 권한의 범위 밖에 두어지는 경우가 있다. 이 경우 개정 권한 자체가
제한되는 것이다. 이러한 점은(비록 그중 일부 제한은 더 이상 실질적
중요성을 가지지 않지만) 미합중국 헌법에서도 관찰될 수 있다. 왜냐하면
제5조는, `1808년 이전에 이루어진 어떠한 개정도 제1조 제9절의 제1항과
제4항을 어떠한 방식으로도 변경할 수 없으며, 어떠한 주도 그 동의 없이는
상원(Senate)에서의 동등한 참정권(equal suffrage)을 박탈당하지 않는다'고
규정하고 있기 때문이다.

입법부가 남아프리카의 경우처럼, 입법부 구성원들이 특별한 절차를
작동시킴으로써 제거할 수 있는 제한을 받는 경우에는, 그 입법부를 해당
이론이 요구하는 바와 같은 법적 제한을 받을 수 없는 주권자와 동일시할 수
있다는 주장이 가능하다. 이 이론에 있어 곤란한 사례는, 미합중국의
경우처럼 입법부에 대한 제약들(restrictions)이 특별한 기관에 위임된 개정
권한의 행사를 통해서만 제거될 수 있는 경우이거나, 혹은 그러한 제약들이
모두 개정 권한의 범위 밖에 있는 경우이다.

이러한 사례들을 일관되게 설명할 수 있다는 이 이론의 주장을 검토함에
있어, 종종 간과되지만 반드시 상기해야 할 점은, 오스틴 자신이 이 이론을
전개함에 있어 영국에서조차 주권자를 입법부와 동일시하지는
\emphksans{않았다(not)}는 점이다. 이는 의회의 여왕(Queen in Parliament)이
통상적으로 받아들여지는 교설에 따르면 그 입법 권한에 대해 법적
제한으로부터 자유롭고, 따라서 `경성(rigid)' 헌법에 의해 제한되는
의회(Congress)나 다른 입법부들과 대비되는 의미에서 \textbf{(p.~74)}
`주권적 입법부(sovereign legislature)'가 무엇인지를 보여주는
전형(paradigm)으로 자주 인용됨에도 불구하고 그러하였다. 그럼에도
불구하고 오스틴의 견해에 따르면, 어떤 민주정(democracy)에서도 선출된
대표자(elected representatives)가 주권적 기구(sovereign body)를
구성하거나 그 일부를 이루는 것은 아니며, 주권을 구성하는 것은
선거권자들(electors)이다. 따라서 잉글랜드에서는 ``엄밀하게 말하자면
하원의원(members of the commons house)은 자신들을 선출하고 임명한 집단을
위한 단순한 수탁자(trustees)에 불과하며, 결과적으로 주권은 항상
국왕(King), 귀족원(Peers), 그리고 하원의 선거권자 집단(electoral body of
the commons)에 귀속된다''고 하였다.\footnote{Austin, \emph{Province of
  Jurisprudence Determined}, Lecture VI, pp.~230--1.} 마찬가지로 그는
미합중국에서도 각 주의 주권과, ``연방 결합(Federal Union)으로부터 발생한
더 큰 국가의 주권 역시 주 정부들(states' governments)이 하나의 집합적
기구(aggregate body)를 형성하는 데에 귀속되는데, 여기서 주 정부란 통상의
입법부(ordinary legislature)가 아니라, 통상의 입법부를 임명하는 시민들의
집단(body of citizens)을 의미한다''고 보았다.\footnote{Ibid., p.~251.}

이러한 관점에서 보면, 통상의 입법부(ordinary legislature)가 법적
제한으로부터 자유로운 법체계와, 입법부가 그러한 제한을 받는 법체계
사이의 차이는, 주권적 유권자 집단(sovereign electorate)이 자신의 주권적
권한(sovereign powers)을 행사하는 방식의 차이에 불과한 것으로 보인다. 이
이론에 따르면 영국에서는, 유권자 집단이 주권에 참여하는 유일한 직접적
행사는 의회(Parliament)에 앉을 대표자를 선출하고 그들에게 자신의 주권적
권력을 위임하는 데에 있다. 이러한 위임은, 위임된 권력이 남용되지 않도록
신임(trust)이 부여되어 있다는 점에서는 의미를 가지지만, 그러한 신임은
오직 도덕적 제재(moral sanctions)의 문제일 뿐이며, 입법 권력에 대한 법적
제한과 달리 법원이 관여하는 사안은 아니다. 이에 반해, 통상의 입법부가
법적으로 제한되는 모든 민주정에서와 마찬가지로 미합중국에서는, 선거인
집단(electoral body)이 자신의 주권적 권한 행사를 대표자 선출에만
국한하지 않고, 그 대표자들에게 법적 제약(legal restrictions)을 부과해
왔다. 이 경우 유권자 집단은, 헌법적 제약(constitutional restrictions)을
준수할 법적 의무를 지는 통상의 입법부에 비해 우월한, `통상을 넘어서고
이면에 감추어진 입법부(``extraordinary and ulterior legislature'')'로
간주될 수 있으며, 충돌이 발생하는 경우 법원은 통상의 입법부의
조항들(Acts)을 유효하지 않다(invalid)고 선언할 것이다. 따라서 이 경우,
이 이론이 요구하는 바와 같이 모든 법적 제한으로부터 자유로운 주권자는
바로 유권자 집단(electorate)에 존재하게 된다.

\textbf{(p.~75)} 이 이론의 이러한 추가적 전개 국면에서는, 주권자에 대한
초기의 단순한 개념이 급진적인 변형까지는 아니더라도 일정한 정교화를
거쳤다는 점이 분명하다. `사회 구성원의 대다수가 습관적으로 복종하는
대상인 인격 또는 인격들(the person or persons to whom the bulk of the
society are in the habit of obedience)'로서의 주권자에 대한 기술은, 이
장 제1절에서 우리가 보여주었듯이, Rex가 절대 군주였고 입법자로서의
승계(succession)에 관한 어떠한 규정도 존재하지 않았던 가장 단순한 형태의
사회에 대해서는 거의 문자 그대로 적용될 수 있었다. 그러나 그러한 규정이
마련된 경우에는, 현대 법체계의 매우 두드러진 특징인 입법적 권위의
연속성(continuity)은 복종의 습관이라는 단순한 용어로는 표현될 수 없었고,
실제로 입법을 수행하고 복종을 받기 이전에 이미 후임자(successor)가
입법할 권리를 가진다는 점을 규정하는 수용된 규칙(accepted rule)이라는
개념을 필요로 하였다. 그러나 민주국가의 유권자 집단(electorate)과
주권자를 동일시하는 현재의 설명은, 핵심어인 `복종의 습관'과 `인격 또는
인격들'에, 단순한 사례에 적용되었을 때와는 전혀 다른 의미를 부여하지
않는 한, 전혀 그럴듯함을 갖지 못한다. 그리고 이러한 의미는 수용된
규칙이라는 개념이 은밀하게 도입되지 않는 한 명확해질 수 없다. 복종의
습관과 명령이라는 단순한 도식만으로는 이를 감당할 수 없다.

이 점은 여러 방식으로 입증될 수 있지만, 가장 분명하게 드러나는 예는
다음과 같은 경우이다. 유권자에서 유아와 정신적 결함자를 제외한 모든
구성원이 포함되어 있는 민주주의 사회에서는, 유권자 자체가 인구의
`대다수(bulk)'를 구성한다. 또는, 모든 성인이 투표권을 가진 단순한
사회집단을 상상해 보자. 이러한 경우 유권자를 주권자로 간주하고, 원래
이론의 단순한 정의를 적용하려 하면, 우리는 결국 사회의 `대다수'가
스스로에게 습관적으로 복종한다고 말하게 된다. 따라서 ``법적으로 제한을
받지 않는 주권자가 명령을 내리고, 수범자들은 이에 습관적으로
복종한다''는 명확한 사회 구도는 사라지고, 다수에 의해 내려진 명령에
다수가 복종하거나, 전체 구성원에 의해 내려진 명령에 전체가 복종한다는
모호한 이미지로 대체된다. 그러나 이것은 더 이상 원래 의미의
`명령'(\emphksans{타인(others)}이 특정 방식으로 행동하길 의도한 표현)도
아니고, `복종'도 아니다.

이러한 비판에 대응하기 위해, 사회 구성원들을 개인으로서의 사적
행위능력(private capacity)에서의 개인들과 \textbf{(p.~76)}
선거권자(electors)나 입법자로서의 공적 행위능력에서의 동일한 인격들로
구분할 수 있을 것이다. 이러한 구분은 완전히 이해 가능한 것이며, 실제로
많은 법적·정치적 현상들은 이러한 용어로 가장 자연스럽게 설명된다. 그러나
설령 우리가 한 걸음 더 나아가 공적 행위능력에서의 개인들이 \emphksans{또 다른
인격(another person)}을 구성하며, 그 인격이 습관적으로 복종을 받는다고
말할 준비가 되어 있다 하더라도, 이러한 구분만으로는 주권 이론을 구출할
수 없다. 왜냐하면 어떤 사람들의 집단이 대표자를 선출하거나 명령을 발할
때, 그들이 `개인으로서'가 아니라 `공적 행위능력에서(in their official
capacity)' 행위했다고 말하는 것이 무엇을 의미하는지를 묻는다면, 그에
대한 답변은 오직 특정한 규칙들에 따른 그들의
자격요건들(qualifications)과, 유효한 선거(valid election)나 법을
성립시키기 위해 그들이 따라야 하는 것을 규정하는 다른 규칙들에 대한
준수라는 용어로만 주어질 수 있기 때문이다. 바로 이러한 규칙들에 대한
참조(reference)를 통해서만, 우리는 어떤 행위를 이 사람들의 집단에 의해
이루어진 선거나 법 제정으로 식별할(identify) 수 있다. 이러한 것들은,
개인의 발언이나 서면 명령을 그 개인에게 귀속시키는 데 사용하는 것과
동일한 단순하고 자연적인 기준(simple natural test)에 의해, 그것을
`만든(making)' 집단(body)에게 귀속되는 것이 아니다.

그렇다면 그러한 규칙들이 존재한다는 것은 무엇을 의미하는가? 그것들은
사회 구성원들이 유권자 집단(electorate)으로서 기능하기 위해---그리고
따라서 그 이론의 목적상 주권자로서 기능하기 위해---무엇을 해야 하는지를
규정하는 규칙들이므로, 그 자체가 주권자에 의해 발령된 명령의 지위를 가질
수는 없다. 왜냐하면 그 규칙들이 이미 존재하고 준수되어 있지 않는 한,
어떠한 것도 주권자에 의해 발령된 명령으로 간주될 수 없기 때문이다.

그렇다면 우리는 이러한 규칙들이 인구(population)의 복종의
\emphksans{습관(habits)}에 대한 기술(description)의 일부에 불과하다고 말할 수
있는가? 사회의 대다수가 주권자인 단일한 개인을, 그가 특정한
형식---예컨대 서면으로 작성되고 서명 및 증인이 갖추어진 형식---으로
명령을 내릴 경우에, 그리고 그 경우에만(if, and only if) 복종하는 단순한
사례에서는, (제1절에서 이 맥락에서의 `습관' 개념 사용에 대해 제기된
이의들을 유보한다면) 그가 이러한 방식으로 입법해야 한다는 규칙은 사회의
복종의 습관에 대한 기술의 일부에 불과하다고 말할 수 있을 것이다. 즉,
사회는 그가 이러한 방식으로 명령을 내릴 \emphksans{때(when)} 습관적으로
그에게 복종한다는 것이다. 그러나 주권자인 개인이 규칙들과 독립적으로
식별될 수 없는 경우에는, 이러한 규칙들을 사회가 주권자에게 습관적으로
복종하는 조건이나 항(term)으로서 단지 그렇게 표상할 수는 없다. 이
규칙들은 주권자를 \emphksans{구성하는(constitutive)} 것이지, 주권자에 대한
복종의 습관을 기술할 때 \textbf{(p.~77)} 언급해야 할 사항들에 불과한
것이 아니다. 따라서 현재의 경우, 유권자 집단(electorate)의 절차를
규정하는 규칙들이, 사회가 다수의 개인으로서 유권자 집단으로서의 자기
자신에게 복종하는 조건들을 나타낸다고 말할 수는 없다. 왜냐하면 `유권자
집단으로서의 자기 자신(itself as an electorate)'이라는 표현은 규칙들과
별도로 식별 가능한 어떤 인격(person)에 대한 참조(reference)가 아니기
때문이다. 그것은 선거권자들이 대표자를 선출함에 있어 규칙들을
준수하였다는 사실에 대한 압축된 참조(condensed reference)에 불과하다.
기껏해야 우리는(제1절의 이의들을 유보한다면) 이러한 규칙들이
\emphksans{선출된 인격들(elected persons)}이 습관적으로 복종을 받는 조건들을
제시한다고 말할 수 있을 뿐이다. 그러나 그렇게 되면 우리는
입법부(legislature)가 아니라 유권자 집단이 주권자라는 현재의 설명에서
벗어나, 입법부가 주권자인 이론의 한 형태로 되돌아가게 되며, 그러한
입법부가 그 입법 권력에 대해 법적 제한(legal limitations)을 받을 수
있다는 사실에서 비롯되는 모든 난점들은 여전히 해결되지 않은 채로 남게 될
것이다.

이 장의 앞선 절에서 제시된 논증들과 마찬가지로, 이 이론에 대한 이러한
논증들은, 단지 세부에서 오류가 있다는 주장에 그치지 않고, 명령, 습관,
그리고 복종이라는 단순한 관념만으로는 법에 대한 분석에 충분할 수 없다는
주장을 구성한다는 의미에서 근본적이다. 이에 대신하여 요구되는 것은,
특정한 절차를 준수함으로써 입법할 자격을 갖춘 사람들에게, 제한될 수도
있고 제한되지 않을 수도 있는, 권한을 부여하는 어떤 규칙(a rule
conferring powers)이라는 생각이다.

이 이론이 지니는 일반적 개념적 부적합성과는 별도로, 통상적으로 최고
입법부로 간주되는 기관이 법적으로 제한될 수 있다는 사실을 이 이론 안에
수용하려는 시도에는 많은 보조적 반론들(ancillary objections)이 존재한다.
그러한 경우 주권자를 유권자 집단(electorate)과 동일시해야 한다면, 설령
유권자 집단이 통상 입법부에 대한 모든 제한을 제거할 수 있는 무제한의
개정 권한(unlimited amending power)을 가지고 있는 경우라 하더라도,
그러한 제한들이 법적인 이유가 유권자 집단이 통상 입법부가 습관적으로
복종하는 명령을 부과했기 때문이라고 말할 수 있는지 우리는 물어볼 수
있다. 우리는 입법 권력에 대한 법적 제한들이 명령으로, 그리고 그에 따른
\emphksans{책무들(duties)}로 잘못 표상된다는 우리의 반론을 일단 유보할 수도
있을 것이다. 그렇다 하더라도, 이러한 제한들이 유권자 집단이 입법부에게
이행하라고 심지어 \emphksans{묵시적으로(tacitly)} 명령한 의무라고 가정할 수
있을까? 앞선 장들에서 묵시적 명령(tacit orders)이라는 관념에 대해 제기된
모든 반론들은, 이 맥락에서 그 사용에 대해 더욱 강력한 힘으로 적용된다.
미합중국 헌법에 규정된 것처럼 \textbf{(p.~78)} 그 행사 방식이 극히
복잡한 개정 권한(amending power)을 행사하지 않는다는 사실은, 유권자
집단의 바람들(wishes)을 보여주는 징후로서는 빈약할 수 있으며, 오히려
그들의 무지와 무관심을 보여주는 신뢰할 만한 징후인 경우가 많다. 우리는,
장군이 병사들에게 하사관이 지시하는 바를 따르라고 묵시적으로 명령했다고,
아마도 그럴듯하게, 간주될 수 있는 경우와는 실로 거리가 멀다.

다시 말해, 유권자 집단에 위임된 개정 권한의 범위 바깥에 전적으로 놓여
있는 입법부에 대한 제약이 존재하는 경우, 이 이론의 용어로 우리는
무엇이라고 말해야 하는가? 이는 단지 상상 가능한 경우에 그치지 않고,
실제로 일부 사례들에서는 그러한 상태가 존재한다. 이러한 경우 유권자
집단은 법적 제한의 적용을 받으며, 비록 그것을 통상을 넘어서는
입법부(extraordinary legislature)라고 부를 수는 있을지라도, 법적
제한으로부터 자유롭지 않으므로 주권자는 아니다. 그렇다면 우리는, 사회가
이러한 제한들에 대해 혁명을 일으키지 않았다는 이유로, 전체로서의
사회(the society as a whole)가 주권자이며 이러한 법적 제한들이
\emphksans{그것(it)}에 의해 묵시적으로 명령되어 왔다고 말해야 하는가? 이것이
혁명과 입법 사이의 구별을 유지할 수 없게 만든다는 점만으로도, 아마 이
주장을 배척하기에 충분한 이유가 될 것이다. 마지막으로, 유권자 집단을
주권자로 취급하는 이 이론은, 기껏해야 유권자 집단이 존재하는
민주정에서의 제한된 입법부에 대해서만 설명을 제공한다. 그러나 Rex와 같은
세습 군주(hereditary monarch)가, 체계 내에서 최고(supreme)이면서도
동시에 제한된(limited) 입법 권력을 향유한다는 개념에는 아무런 부조리도
존재하지 않는다.

\subsection{CHAPTER IV 주석}\label{chapter-iv-uxc8fcuxc11d}

50쪽. 오스틴의 주권 이론. 본 장에서 검토된 주권 이론은 오스틴이
『법리학의 범위』(\emph{The Province}) 제5강과 제6강에서 전개한 것이다.
우리는 오스틴이 단지 일정한 형식적 정의들이나 법체계를 논리적으로
배열하기 위한 추상적 도식을 제시한 것이 아니라, 영국이나 미합중국과 같은
법이 존재하는 모든 사회에서는 오스틴이 정의한 속성을 갖춘 주권자가
어딘가에 반드시 존재한다는 \emphksans{사실 주장(factual claim)}을 하고 있다고
해석했다. 다만 이러한 주권자의 존재는 다양한 헌법적·법적 형식들에 의해
가려질 수 있다. 일부 이론가들은 오스틴을 달리 해석하여 그가 그러한 사실
주장을 한 것이 아니라고 본다(예: 스톤, 『법의 영역과 기능』, 제2장 및
제6장, 특히 60, 61, 138, 155쪽. 여기서 오스틴이 여러 공동체에서 주권자를
식별하려 한 시도는 그의 본래 목적에서 벗어난 무의미한 일탈로 다뤄진다).
이러한 해석에 대한 비판은 Morison, `Some Myth about Positivism', loc.
cit., pp.~217--22 참조. 또한 Sidgwick, \emph{The Elements of Politics,}
Appendix (A) `On Austin's Theory of Sovereignty'도 참조.

54쪽. 오스틴에서의 입법적 권위의 연속성. 『\emphksans{법리학의 범위(The
Province)}』에서 `승계의 방식으로 주권을 취득하는 인격들(persons who
``take the sovereignty in the way of succession'')'에 대한 간략한
언급들(제5강(Lecture V), pp.~152--4)은 시사적이지만 불분명하다.
오스틴은, 주권을 취득하는 인물들이 교체되는 일련의 승계를 통해 주권의
연속성을 설명하기 위해서는 그의 핵심 개념들인 `습관적 복종'과
`사령(commands)'에 더하여 무엇인가 추가적인 요소가 필요하다는 점을
인정하는 듯 보이지만, 그 추가 요소를 결코 명확히 식별하지는 않는다. 그는
이와 관련하여 `\emphksans{자격(title)}', 승계에 대한 `\emphksans{청구(claims)}',
그리고 `\emphksans{정당한(legitimate)} 자격'에 관해 말하지만, 이러한 표현들은
통상적으로 사용될 때 모두 승계를 규율하는 \emphksans{규칙}의 존재를 함축하며,
연속적인 주권자들에 대한 단순한 복종의 \emphksans{습관}만을 의미하지는
않는다. 오스틴이 사용하는 이러한 용어들과 `일반적 자격(generic title)'
및 주권을 취득하는 `일반적 방식(the generic mode)'이라는 표현들에 대한
그의 설명은, 주권자의 `확정적(determinate)' 성격에 관한 그의 교설로부터
도출되어야 한다(op. cit., 제5강(Lecture V), pp.~145--55). 여기서 그는
주권자인 개인 또는 개인들이 개별적으로---예컨대 이름으로---식별되는
경우와, `어떤 일반적 기술(generic description)에 부합하는 것으로서'
식별되는 경우를 구별한다. 따라서 (가장 단순한 예를 들면) 세습
군주제에서는 일반적 기술이 특정 조상의 `생존한 장남 후손'일 수 있고,
의회 민주정(parliamentary democracy)에서는 입법부 구성원
자격요건들(qualifications for membership of the legislature)을 재현하는
고도로 복잡한 기술이 될 것이다.

오스틴(Austin)의 견해는, 어떤 사람이 그러한 `일반적' 기술(`generic'
description)을 충족할 때 그는 승계에 대한 `자격(title)' 또는
`권리(right)'를 가진다는 것인 듯하다. 그러나 주권자에 대한 이러한 일반적
\emphksans{기술(description)}에 의존한 설명은, 오스틴이 이 맥락에서
`기술'이라는 말로 승계를 규율하는 수용된 \emphksans{규칙}을 의미하지 않는 한,
그 자체로는 불충분하다. 왜냐하면, 사회의 구성원들이 \emphksans{사실상(as a
matter of fact)} 그때그때 일정한 기술에 부합하는 누구에게든 습관적으로
복종하는 경우와, 그러한 기술에 부합하는 누구든지 복종받을 \emphksans{권리}
또는 \emphksans{자격}을 가진다는 규칙이 수용되어 있는 경우 사이에는 분명한
구별이 존재하기 때문이다. 이는 사람들이 체스 말을 습관적으로 일정한
방식으로 움직이는 경우와, 그렇게 행동할 뿐만 아니라 그것이 말을 움직이는
\emphksans{올바른(right)} 방식이라는 규칙을 수용하고 있는 경우 사이의 차이와
평행(parallel)을 이룬다. 승계에 대한 `권리'나 `자격'이 존재하려면,
승계를 규정하는 규칙(rule providing for the succession)이 반드시 있어야
한다. 오스틴의 일반적 기술(generic descriptions)에 관한 교설은 그러한
규칙을 대체할 수는 없으며, 다만 그 필요성(necessity)을 분명히 드러낼
뿐이다. 입법자로서의 자격을 부여하는 규칙(rule qualifying persons as
legislators)이라는 개념을 인정하지 않은 오스틴에 대한 유사한 비판으로는,
Gray, \emph{The Nature and Sources of the Law}, 제3장(chap.~iii), 특히
제151--157절(ss. 151--7)을 보라. 또한 주권적 기구(sovereign body)의
통일성(unity)과 법인격적(corporate) 또는 `합의체적(collegiate)'
행위능력(capacity)에 관한 오스틴의 제5강(Lecture V)의 설명 역시 동일한
결함을 지니고 있다(이 장 제4절(s. 4) 참조).

55쪽. 규칙과 습관. 여기서 강조된 규칙의 \emphksans{내면적 측면(internal
aspect)}은 제5장 §2(p.88), §3(p.98), 제6장 §1, 제7장 §3에서 더 논의된다.
또한 Hart, ``Theory and Definition in Jurisprudence'', \emph{29 PAS}
Supplementary vol.~(1955), 247--250쪽도 참고. 유사한 관점으로는 Winch,
『사회과학의 개념』(1958), 제2장 57--65쪽, 제3장 84--94쪽, 그리고
Piddington, ``Malinowski의 욕구 이론'', 『인간과 문화』(Firth 편) 참조.

60쪽. 헌법의 기본 규칙에 대한 일반적 수용. 헌법의 수용, 나아가 법체계의
존재와 관련하여, 공무원과 민간 시민들이 법 규칙에 대해 보이는 다양한
태도들은 제5장 §2 (88--91쪽), 제6장 §2 (114--117쪽)에서 더 깊이
논의된다. 또한 Jennings, 『헌법법』(3판) 부록 3: ``법 이론에 대한 노트''
참조.

63쪽. 홉스와 묵시적 사령 이론. 앞선 제3장 §3 및 해당 주석 참조. 또한
Sidgwick, 『정치의 요소들』 부록 A. 현대 입법기관의 제정법조차도
집행되기 전까지는 법이 아니라고 보는 유사한 `현실주의적' 이론으로는
Gray, 『법의 본성과 원천』 제4장, J. Frank, 『법과 현대정신』 제13장
참조.

66쪽. 입법 권한에 대한 법적 제한. 오스틴과 달리, 벤담은 최고 권한이
`명시적 합의(express convention)'에 의해 제한될 수 있으며, 그 합의를
위반하여 제정된 법률은 무효(void)가 될 수 있다고 보았다. 이에 관해서는
『\emphksans{정부에 관한 단편(A Fragment on Government)}』 제4장 제26,
34--38절 참조. 주권자의 권력에 대한 법적 제한의 가능성을 부정하는
오스틴의 논증은, 그러한 제한의 적용을 받는다는 것이 곧
\emphksans{책무(duty)}의 적용을 받는다는 전제에 기초하고 있다. 이에 대해서는
『\emphksans{법리학의 범위(The Province)}』 제6강, 254--268쪽 참조. 그러나
실제로 입법적 권위(legislative authority)에 대한 제한은 책무들(duties)이
아니라 \emphksans{무능력(disabilities)}으로 구성된다(호펠드(Hohfeld),
『\emph{Fundamental Legal Conceptions}』(1923), 제1장 참조).

68쪽. 입법의 방식 및 형식에 관한 규정. 이러한 규정들을 입법 권한에 대한
실질적 제한과 구별하는 것이 어렵다는 점은 제7장 제4절, 149--152쪽에서
추가로 논의된다. 주권 기관의 역량을 `정의하는 것'과 `제한하는 것' 사이의
구분에 대한 포괄적 논의는 마셜, 『의회 주권과
영연방』(\emph{Parliamentary Sovereignty and the Commonwealth}, 1957),
제1--6장 참조.

72쪽. 헌법적 안전장치와 사법 심사. 사법 심사가 허용되지 않는 헌법에
대해서는 휘어, 『현대 헌법』(\emph{Modern Constitutions}) 제7장 참조.
여기에 해당하는 국가는 (주법을 제외한) 스위스, 제3공화국 시기의 프랑스,
네덜란드, 스웨덴 등이 포함된다. 미합중국 대법원이 `정치적 문제(political
questions)'를 야기하는 위헌 주장에 대해 판결을 거부한 사례로는
\emphksans{루터 대 보든}(\emph{Luther v. Borden}), 7 Howard 1, 12 L. Ed. 581
(1849) 참조. 프랑크퍼터의 「대법원」, 『사회과학 백과사전』 제14권,
474--476쪽도 참조.

74쪽. `예외적 입법기관(extraordinary legislature)'으로서의 유권자. 많은
체계에서 통상적 입법기관이 법적 제한을 받는다는 비판을 피하려는 시도로
오스틴이 이 개념을 사용한 것에 대해서는 『법학의 영역』 제6강,
222--233쪽 및 245--251쪽 참조.

76쪽. 사적 행위능력의 입법자와 공적 행위능력의 입법자. 오스틴은 종종
주권 기관의 구성원을 `개별적으로 고려했을 때'와 `구성원으로서 또는
법인적·주권적 행위능력으로 고려했을 때'를 구분한다(『법학의 영역』
제6강, 261--266쪽). 그러나 이 구분은 주권 기관의 입법 활동을 규율하는
규칙 개념을 수반한다. 오스틴은 이러한 공적 또는 합의체적
행위능력(official or collegiate capacity)이라는 생각을 `일반적
기술(generic description)'이라는 불충분한 용어로 설명하려고 암시만 할
뿐이다(54쪽 주석 참조).

78쪽. 개헌 권한의 제한 범위. 미합중국 헌법 제5조의 단서 조항 참조. 독일
연방공화국의 기본법(1949년) 제1조와 제20조는 제79조 제3항이 부여한 개헌
권한의 범위 밖에 놓여 있다. 터키 헌법(1945년)의 제1조 및 제102조도 참조.

\subsection{CHAPTER IV 제3판
주석}\label{chapter-iv-uxc81c3uxd310-uxc8fcuxc11d}

55--57쪽. \emphksans{사회 규칙의 본성에 대하여}. 하트의 `실천 이론(practice
theory)'에 대한 비판은 다음 저작들을 참조할 것: G. J. 워녹(G. J.
Warnock), 『도덕성의 대상(The Object of Morality)』(메튜언, 1971) 제4장;
로널드 드워킨(Ronald Dworkin), 『권리를 진지하게 대하기(Taking Rights
Seriously)』 48--58쪽; 조셉 라즈(Joseph Raz), 『실천적 이성과
규범(Practical Reason and Norms)』 49--58쪽.

하트는 이후의 저술에서 사령에 관한 사고를 발전시켜 규칙에 대한 또 다른
견해를 제안하였다. 그는 사령이 `결정적인(peremptory)' 것으로서, 즉
``청자가 어떤 행위를 할 것인지에 대해 그 장단점을 독자적으로 숙고하는
것을 차단하거나 차단하려는 의도를 가진 것''이며, 또한 `내용
독립적(content-independent)'으로, ``실행될 행위의 성격이나 본질에
관계없이 이유로 기능하도록 의도된 것''이라고 말한다. 그는 이어,
``사회에서 사령자의 발언이 결정적인 행위 이유(peremptory reasons for
action)로 일반적으로 승인되는 것은 사회 규칙의 존재와 동일하다''고
주장한다(H. L. A. Hart, 「사령과 권위 있는 법적 이유(Commands and
Authoritative Legal Reasons)」, 『벤담에 관한 에세이( \emph{Essays on
Bentham})』 제10장, 258쪽 참조). 하트는 `법적·도덕적 의무(Legal and
Moral Obligation)'라는 논문에서 `내용 독립적 이유' 개념을 처음
도입하였으며, 이는 A. I. 멜든(A. I. Melden) 편 『도덕철학 에세이(Essays
in Moral Philosophy)』(워싱턴 대학 출판부, 1958), 82--107쪽에 수록되어
있다. 이 개념은 라즈의 『자유의 도덕성(The Morality of
Freedom)』(옥스퍼드 대학 출판부, 1986) 제2장에서, 그리고 레슬리
그린(Leslie Green)의 『국가의 권위(The Authority of the
State)』(옥스퍼드 대학 출판부, 1990), 36--42쪽에서 더욱 발전되었으며, P.
마크위크(P. Markwick)의 「법과 내용 독립적 이유(Law and
Content-Independent Reasons)」(\emph{Oxford Journal of Legal Studies},
2000년 제20권, 579쪽)에서는 비판받았다.

56--58쪽. \emphksans{법의 `내면적 측면'(internal aspect)}. 하트는 종종
`내면적 관점(internal point of view)'과 `내면적 측면(internal aspect)'을
혼용해서 사용한다(예: 88--89쪽 참조). 이 일반 개념에 대해서는 닐
매코믹(Neil MacCormick), 『H. L. A. Hart』 29--40쪽 참조. 이 개념의
모호성에 대한 논의로는 스티븐 페리(Stephen Perry), 「하트의 사회 규칙과
법의 기초: 내면적 관점의 해방(Hart on Social Rules and the Foundations
of Law: Liberating the Internal Point of View)」(\emph{Fordham Law
Review}, 2006년 제75권, 1171쪽), 스콧 샤피로(Scott J. Shapiro), 「내면적
관점이란 무엇인가?(What is the Internal Point of View?)」(\emph{Fordham
Law Review}, 2006년 제75권, 1157쪽) 참조. 법이론에서 이 개념의 방법론적
중요성은 라즈, 『권위와 해석 사이(Between Authority and
Interpretation)』 제2장에서 논의된다. 이 방법론적 주장에 대한 비판으로는
스티븐 페리, 「법이론에서의 해석과 방법론(Interpretation and Methodology
in Legal Theory)」, 안드레이 마머(Andrei Marmor) 편 『법과 해석(Law and
Interpretation)』(옥스퍼드 대학 출판부, 1997), 존 피니스(John Finnis),
『자연법과 자연권(Natural Law and Natural Rights)』 제1장 등이 있다.
피니스의 비판은 줄리 딕슨(Julie Dickson), 『평가와 법이론(Evaluation and
Legal Theory)』(하트 출판사, 2001), 제2--4장에서 면밀히 분석된다.

66--71쪽. \emphksans{입법권에 대한 법적 제한}. 하트는 『벤담에 관한 에세이(
\emph{Essays on Bentham})』 제9장 「주권과 법적으로 제한된
정부(Sovereignty and Legally Limited Government)」에서 벤담이 입법
제한을 이해하려 한 시도를 분석한다. 오스틴의 무제한성(ilimitability)
개념에 대한 라즈의 비판은 『법체계의 개념(The Concept of a Legal
System)』 제2장에서 다루어진다. 오스틴의 분석이 헌법이론에 미치는 함의에
대해서는 제프리 마셜(Geoffrey Marshall), 『헌법 이론(Constitutional
Theory)』(옥스퍼드 대학 출판부, 1980), 제1, 3장 참조.

\setcounter{footnote}{0}
\input{sections/chapter_5-translation.tex}
\setcounter{footnote}{0}
\section{VI. 법체계의 토대들(THE FOUNDATIONS OF A LEGAL
SYSTEM)}\label{vi.-uxbc95uxccb4uxacc4uxc758-uxd1a0uxb300uxb4e4the-foundations-of-a-legal-system}

\subsection{1. 승인 규칙과 법적 유효성(RULE OF RECOGNITION AND LEGAL
VALIDITY)}\label{uxc2b9uxc778-uxaddcuxce59uxacfc-uxbc95uxc801-uxc720uxd6a8uxc131rule-of-recognition-and-legal-validity}

제4장에서 비판된 이론에 따르면, 한 법체계의 토대들은 사회 집단의 다수가,
스스로는 누구에게도 습관적으로 복종하지 않는 주권자인 인격 또는 인격들의
위협으로 뒷받침된 명령들에 습관적으로 복종하는 상황으로 이루어진다. 이
이론에서 이러한 사회적 상황은 법의 존재에 대한 필요조건이자
충분조건이다. 우리는 이미 이 이론이 현대의 국내법 체계의 몇몇 두드러진
특징들을 설명할 수 없는 무능력(incapacity)을 상당히 상세하게 보여
주었다. 그럼에도 불구하고, 많은 사상가들의 정신을 사로잡아 온 이 이론이
시사하듯이, 그것은 비록 흐릿하고 오도적인 형태이기는 하나, 법의 몇몇
중요한 측면들에 관한 일정한 진리들을 포함하고 있기는 하다. 그러나 이러한
진리들은, 의무의 일차 규칙들을 식별하기(identify) 위해 이차적 승인
규칙(rule of recognition)이 수용되고 사용되는 보다 복잡한 사회적 상황의
관점에서만 비로소 명확하게 제시될 수 있고, 그 중요성 또한 올바르게
평가될 수 있다. 바로 이 상황이야말로, 무엇인가가 한 법체계의
토대들이라고 불릴 만하다면, 그렇게 불릴 자격이 있는 것이다. 이 장에서는
주권 이론 및 다른 곳들에서 부분적으로만 또는 오도된 방식으로 표현되어 온
이 상황의 여러 요소들을 논의할 것이다.

이러한 승인 규칙이 수용되는 곳에서는 어디에서나, 사적 개인들과
공무담당자들 모두에게 의무의 일차 규칙들을 식별하기(identify) 위한 권위
있는 기준들이 제공된다. 앞서 보았듯이, 이렇게 제공되는 기준들은 다음의
다양한 형태들 가운데 하나 또는 그 이상의 형태를 취할 수 있다: 여기에는
권위 있는 문헌에 대한, 입법적 제정에 대한, 관습적 실천에 대한, 특정된
사람들의 일반적 선언에 대한, 또는 개별 사건들에서의 과거 사법적 결정들에
대한 참조(reference)가 포함된다. 제4장에서 묘사된 Rex I의 세계와 같이
매우 단순한 체계에서는, 그가 제정한 것만이 법이며 관습 규칙이나 헌법적
\textbf{(p.~101)} 문헌에 의해 그의 입법 권한에 법적 제한들이 전혀
부과되지 않기 때문에, 법을 식별하는 유일한 기준은 Rex I에 의한
제정이라는 사실에 대한 단순한 참조가 될 것이다. 이러한 단순한 형태의
승인 규칙의 존재는, 공무담당자들이나 사적 개인들이 이 기준에 따라
규칙들을 식별하는 일반적 실천 속에서 분명히 드러난다. 법의 다양한
‘원천들(sources)’이 존재하는 현대의 법체계에서는 승인 규칙도 이에
상응하여 더 복잡해지며, 법을 식별하기 위한 기준들은 복수적이고,
통상적으로 서면 헌법, 입법부에 의한 제정, 그리고 사법적 선례를 포함한다.
대부분의 경우 이러한 기준들 사이에 발생할 수 있는 충돌을 대비하여,
그것들을 상대적 종속성과 우위성(primacy)의 질서 속에 배열하는 규정이
마련되어 있다. 바로 이러한 방식으로, 우리 체계에서는 ‘보통법(common
law)’이 ‘제정법(statute)’에 종속된다.

한 기준이 다른 기준에 대해 갖는 이러한 상대적
\emphksans{종속성(subordination)}은 \emphksans{도출(derivation)}과 구별하는 것이
중요한데, 왜냐하면 이 두 관념을 혼동함으로써, 모든 법은 본질적으로 또는
‘실제로(really)’(설령 단지 ‘묵시적으로(tacitly)’일 뿐이라 하더라도)
입법의 산물이라는 견해에 대해 그럴듯하지만 허위적인 지지가 제공되어 왔기
때문이다. 우리 체계에서 관습과 판례는 입법에 종속되는데, 이는 관습법 및
보통법 규칙들이 제정법(statute)에 의해 그 법적 지위를 박탈당할 수 있기
때문이다. 그러나 이 규칙들은, 그 지위가 불안정하다는 점은 인정하더라도,
법으로서의 지위를 ‘묵시적인’ 입법 권한의 행사에 빚지고 있는 것이 아니라,
그것들에게 이러한 독립적이되 종속된(subordinate) 지위를 부여하는 승인
규칙의 수용에 빚지고 있다. 또, 단순한 경우에서와 마찬가지로, 이러한
상이한 기준들이 위계적으로 배열된 복합적인 승인 규칙의 존재는, 그러한
기준들에 의해 규칙들을 식별하는 일반적 실천 속에서 드러난다.

법체계의 일상적 운영 속에서 승인 규칙(rule of recognition)은 규칙으로서
명시적으로 정식화되는(expressly formulated) 경우가 극히 드물다. 다만
영국의 법원들이 때때로 법의 한 기준이 다른 기준에 대해 차지하는 상대적
지위를 일반적인 표현으로 공표하는(announce) 경우는 있는데, 예컨대 의회
제정법들(Acts of Parliament)이 법의 다른 원천들(sources)이나 법의 후보
원천들(suggested sources of law)에 대해 최고성(supremacy)을 가진다고
단언하는(assert) 때가 그러하다. 대체로 승인 규칙은 진술되지 않지만,
법원이나 다른 공무담당자들, 또는 사적 개인들이나 그들의 조언자들이 개별
규칙들을 식별하는(identify) 방식 속에서 그 존재가
\emphksans{드러난다(shown)}. 물론 승인 규칙이 제공하는 기준들을 법원이
사용하는 방식과, 다른 이들이 사용하는 방식 사이에는 차이가 있다.
왜냐하면 법원이 특정한 규칙이 법으로서 올바르게 식별되었다는 전제하에
어떤 특정한 \textbf{(p.~102)} 결론에 도달할 경우, 그들이 하는 말은 다른
규칙들에 의해 부여된 특별한 권위 있는 지위를 가지기 때문이다. 이러한
점에서, 그리고 다른 많은 점에서, 법체계의 승인 규칙은 게임의 득점
규칙(scoring rule)과 유사하다. 게임의 진행 과정에서 득점을 구성하는
행위들(득점(runs), 골 등)을 규정하는 일반 규칙은 거의 정식화되지 않으며,
대신 경기 임원들과 선수들이 승리에 산입되는 특정 국면들을 식별하는 데
그것을 \emphksans{사용(use)}한다. 여기에서도 경기 임원들(심판이나
득점기록원)의 선언은 다른 규칙들에 의해 부여된 특별한 권위 있는 지위를
가진다. 더 나아가, 두 경우 모두에서 이러한 규칙의 권위 있는 적용들과, 그
규칙의 문언에 비추어 볼 때 명백히 요구되는 바에 대한 일반적 이해 사이에
충돌이 발생할 가능성이 존재한다. 이는 우리가 뒤에서 보게 되겠지만,
이러한 유형의 규칙 체계가 존재한다는 것이 무엇을 의미하는지에 대한
어떠한 설명에서도 반영되어야 할 복잡성이다.

법원과 그 밖의 사람들에 의해, 체계의 개별 규칙들을 식별하는(identify)
데에 명시되지 않은 승인 규칙(rule of recognition)이 사용되는 것은 내부
관점(internal point of view)의 특징이다. 이러한 방식으로 승인 규칙을
사용하는 사람들은, 그것들을 향도규칙들(guiding rules)로 자신들이
수용하고 있음을 드러내며, 이러한 태도와 함께 외부 관점(external point of
view)의 자연스러운 표현들과는 다른, 특징적인 어휘가 수반된다. 아마도
이러한 표현들 가운데 가장 단순한 것은 ‘\ldots 은 법이다(It is the law
that \ldots)’라는 표현일 것인데, 이는 판사들뿐만 아니라 법체계 아래에서
살아가는 보통 사람들의 입에서도, 체계의 어떤 주어진 규칙을 식별할 때
발견될 수 있다. 이는 ‘아웃(Out)’이나 ‘골(Goal)’이라는 표현과 마찬가지로,
다른 사람들과 함께 그 목적에 적절한 것으로 인정하는(acknowledges)
규칙들에 대한 참조(reference)로 한 사람이 상황을 평가하는(assessing a
situation) 언어이다. 이러한 규칙에 대한 공유된 수용(shared acceptance)의
태도는, 사회 집단이 그러한 규칙들을 수용하고 있다는 사실을
\emphksans{외부로부터(ab extra)} 기록할 뿐, 자신은 그 규칙들을 수용하지 않는
관찰자의 태도와 대비된다. 이 외부 관점의 자연스러운 표현은 ‘\ldots 은
법이다(It is the law that \ldots)’가 아니라, ‘영국에서는 의회 내 여왕이
제정하는 것은 무엇이든 법으로 인정한다(In England they recognize as law
\ldots{} whatever the Queen in Parliament enacts\ldots)’와 같은 것이다.
이 두 표현 형식 가운데 첫 번째를 우리는 \emphksans{내부 진술(internal
statement)}이라고 부를 것인데, 이는 내부 관점을 드러내며, 승인 규칙을
수용하고 있으면서도 그것이 수용되고 있다는 사실을 진술하지 않은 채, 그
규칙을 적용하여 체계의 어떤 특정 규칙을 유효한(valid) 것으로 인정하는
사람이 자연스럽게 사용하는 표현이기 때문이다. \textbf{(p.~103)} 두 번째
표현 형식은 \emphksans{외부 진술(external statement)}이라고 부를 것인데, 이는
체계의 승인 규칙을 스스로 수용하지 않으면서, 다른 사람들이 그것을
수용하고 있다는 사실을 진술하는 외부 관찰자의 자연스러운 언어이기
때문이다.

수용된 승인 규칙을 사용하여 내부 진술을 하는 이러한 방식이 이해되고, 그
규칙이 수용되고 있다는 사실에 대한 외부 사실 진술(external statement of
fact)과 신중하게 구별된다면, 법적 ‘유효성(validity)’이라는
관념(notion)을 둘러싼 많은 모호함은 사라진다. 왜냐하면
‘유효한(valid)’이라는 말은 언제나 그런 것은 아니지만, 가장 흔히는 바로
그러한 내부 진술에서 사용되며, 법체계의 어떤 특정 규칙에 대하여,
명시되지는 않았으나 수용된 승인 규칙을 적용하기 때문이다. 어떤 주어진
규칙이 유효하다고 말하는 것은, 그것이 승인 규칙이 제공하는 모든 검사
기준들(tests)을 통과하였음을, 따라서 체계의 규칙임을 인정하는 것이다.
기실 우리는, 어떤 특정 규칙이 유효하다는 진술이란 그것이 승인 규칙이
제공하는 모든 기준들(criteria)을 충족한다는 것을 의미한다고 간단히 말할
수 있다. 이는 그러한 진술들의 내부 성격을 가릴 수 있다는 점에서만
부정확한데, 왜냐하면 이러한 유효성 진술들은 크리켓 선수의
‘아웃(Out)’이라는 표현과 마찬가지로, 화자와 다른 사람들이 수용하고 있는
승인 규칙을 어떤 특정 사례에 적용하는 것이지, 그 규칙이 충족되었다는
사실을 명시적으로 진술하는 것이 아니기 때문이다.

법적 유효성(legal validity)이라는 관념과 관련된 몇몇 난점들은, 법의
유효성과 법의 ‘실효성(efficacy)’ 사이의 관계에 관한 것이라고들 말해진다.
만약 ‘실효성’이란, 특정한 행위를 요구하는 어떤 법규칙이 대체로(more
often than not) 복종된다는 사실을 의미하는 것이라면, 개별 규칙의
유효성과 \emphksans{그 규칙의(its)} 실효성 사이에는 필연적인 연관(necessary
connection)이 없다는 점은 분명하다. 다만 체계의 승인 규칙(rule of
recognition)이 그 기준들(criteria) 가운데에---일부 체계가
그러하듯---어떤 규칙이 오랫동안 실효성을 상실한 경우에는 더 이상 그
체계의 규칙으로 간주되지 않는다는 규정(때때로 실효상실 규칙(rule of
obsolescence)이라 불리는 것)을 포함하는 경우는 예외이다.

어떤 개별 규칙의 실효성 결여(inefficacy)는---그것이 그 규칙의 유효성에
불리하게 작용할 수도 있고 그렇지 않을 수도 있지만---체계 전체의 규칙들에
대한 일반적인 무시(general disregard)와는 구별되어야 한다. 이러한 무시는
그 성격상 완전하고 또한 장기간 지속될 수 있는데, 그 경우 우리는 새로운
체계에 대해서는 그것이 특정 집단의 법체계로서 결코 확립된 적이 없다고
말하게 되거나, 한때 확립되었던 체계에 대해서는 그것이 더 이상 그 집단의
법체계가 아니게 되었다고 말하게 될 것이다. 어느 경우든, 체계의 규칙들에
근거하여 \textbf{(p.~104)} 내부 진술(internal statement)을 하는 데
필요한 정상적인 맥락 또는 배경은 결여되어 있다. 이러한 경우에는, 특정
개인들의 권리들과 책무들을 그 체계의 일차 규칙들을 참조하여 평가하는
것도, 그 체계의 승인 규칙들을 참조하여 그 규칙들 가운데 어느 것의
유효성을 평가하는 것도, 일반적으로 \emphksans{무익한(pointless)} 일이 된다.
실제로 한 번도 실효적이었던 적이 없거나 이미 폐기된 규칙 체계를
적용하겠다고 고집하는 것은, 아래에서 언급할 특수한 사정을 제외하면, 한
번도 수용된 적이 없거나 이미 폐기된 득점 규칙(scoring rule)을 참조하여
경기의 진행 상황을 평가하려는 것만큼이나 무익하다.

체계의 특정 규칙의 유효성에 관하여 내부 진술을 하는 사람은, 그 체계가
일반적으로 실효성을 갖고 있다는 외부 사실 진술의 참임(truth)을
\emphksans{전제(presuppose)}하고 있다고 말할 수 있다. 왜냐하면 내부 진술의
정상적인 사용은 그러한 일반적 실효성의 맥락 속에서 이루어지기 때문이다.
그러나 유효성에 관한 진술이 체계가 일반적으로 실효적이라는 것을
‘의미한다(mean)’고 말하는 것은 잘못일 것이다. 어떤 체계가 한 번도 확립된
적이 없거나 이미 폐기된 경우, 그 체계의 규칙의 유효성에 대해 말하는 것은
대체로 무익하거나 공허한(normally pointless or idle) 말이기는 하지만,
그렇다고 해서 그것이 전혀 의미 없거나(meaningless) 언제나 무익한(always
pointless) 것은 아니다. 로마법(Roman Law)을 가르치는 하나의 생생한
방식은, 그 체계가 여전히 실효적\emphksans{인 것처럼(as if)} 말하며, 개별
규칙들의 유효성을 논하고 그 규칙들에 따라 문제를 해결하는 것이다; 또한
혁명에 의해 파괴된 옛 사회질서의 회복에 대한 희망을 품고 새로운 질서를
거부하는 하나의 방식은, 구체제의 법적 유효성 기준들(criteria of legal
validity)에 집착하는 것이다. 이는 차르 체제 러시아(Tsarist Russia)에서
유효했던 어떤 상속 규칙(rule of descent)에 근거하여 여전히 재산권을
주장하는 백계 러시아인(White Russian)에 의해 암묵적으로 행하여진다.

체계의 특정 규칙이 유효하다는 내부 진술과, 그 체계가 일반적으로
실효적이라는 외부 사실 진술 사이의 정상적인 맥락적 연관을 파악하면,
규칙의 유효성을 단언하는(assert) 것은 그것이 법원에 의해 집행되거나 그
밖의 어떤 공무담당 작용이 취해질 것임을 예측하는 것이라는 통상적 이론을
그에 걸맞은 관점에서 이해하는 데 도움이 된다. 이러한 이론은 우리가 지난
장에서 검토하고 배척한 의무의 예측적 분석과 여러 면에서 유사하다. 두
경우 모두에서 이 예측적 이론을 제시하게 되는 동기는, 오직 이러한
방식으로만 형이상학적 해석을 회피할 수 있다는 다음과 같은 확신이다: 즉,
규칙이 유효하다는 진술은 \textbf{(p.~105)} 경험적 수단으로는 탐지될 수
없는 어떤 신비한 성질을 귀속시키는 것이거나, 아니면 공무담당자들의 장래
행태에 대한 예측이어야 한다는 것이다. 또한 두 경우 모두에서 이 이론의
그럴듯함은 동일한 다음의 중요한 사실에 기인한다: 즉 체계가 일반적으로
실효적이며 앞으로도 그러할 개연성이 높다는, 관찰자가 기록할 수 있는,
외부 사실 진술의 참임(truth)이, 규칙을 수용하고 의무나 유효성에 관한
내부 진술을 하는 사람에 의해 정상적으로(normally) 전제된다는 점이다. 이
둘은 분명히 매우 밀접하게 연관되어 있다. 마지막으로, 두 경우 모두에서 이
이론의 오류는 동일하다. 그것은 내부 진술의 특수한 성격을 간과하고, 이를
공무담당 작용(official action)에 관한 외부 진술로 취급하는 데에 있다.

이러한 오류는 특정 규칙이 유효하다는 판사 자신의 진술이 사법적 결정에서
어떻게 기능하는지를 고려하면 즉각적으로 드러난다. 왜냐하면 이 경우에도
판사는 그러한 진술을 하면서 체계의 일반적 실효성을 전제할 뿐 이를
진술하지는 않지만, 분명히 자기 자신의 또는 타인의 공무담당 작용을
예측하는 데에는 관심이 없기 때문이다. 규칙이 유효하다는 그의 진술은, 그
규칙이 그의 법정에서 무엇이 법으로 간주될 것인지를 식별하기 위한 검사
기준들(tests)을 충족한다는 점을 인정하는 내부 진술이며, 결정에 대한
예언(prophecy)이 아니라 결정의 \emphksans{이유(reason)}의 일부를 구성한다.
기실 이러한 진술이 사적 개인에 의해 이루어지는 경우에는, 규칙이
유효하다는 진술을 예측으로 간주하는 것이 더 그럴듯해 보일 수 있다.
왜냐하면 비공무담당자의 유효성 또는 유효하지 않음에 관한 진술들과 사안을
판단하는 법원의 진술 사이에 충돌이 발생하는 경우, 전자는 철회되어야
한다고 말하는 데에 종종 상당한 일리가 있기 때문이다. 그러나 이 경우에도,
제7장에서 공무담당자들의 선언과 규칙의 명백한 요구 사이의 이러한 충돌의
의미를 검토하게 되면 보게 되듯이, 그것이 법원이 무엇이라고 말할지를
\emphksans{예측하는(predicted)} 데에서 오류를 범했기 때문에 이제
\emphksans{그른(wrong)} 진술로 드러나 철회되는 것이라고 단정하는 것은
독단적일 수 있다. 왜냐하면 진술을 철회하는 이유는 그것이 그르다는
사실만이 아니며, 또한 그른 방식 역시 이러한 설명이 허용하는 것보다 더
다양하기 때문이다.

다른 규칙들의 유효성을 평가하기 위한 기준들을 제공하는 승인 규칙(rule of
recognition)은 중요한 의미에서 \emphksans{궁극적(ultimate)} 규칙이다. 이 점은
우리가 곧 명확히 해설하려 한다. 그리고 통상 그렇듯이 여러 기준들이
상대적 종속성과 우위성(primacy)의 서열로 정렬되어 있는 경우, 그중 하나는
\emphksans{최고(supreme)} 기준이 된다. \textbf{(p.~106)} 승인 규칙의
궁극성(ultimacy)과 그 기준 중 하나의 최고성(supremacy)에 관한 이러한
관념들은 주의 깊게 다룰 만한 가치가 있다. 우리는 이 관념들을, 앞서
우리가 기각했던 이론---즉, 모든 법체계에는, 설령 그것이 법적 형식 배후에
숨어 있다 하더라도, 반드시 법적으로 무제한인 주권적 입법 권력(sovereign
legislative power)이 존재해야 한다는 이론---으로부터 분리해내는 것이
중요하다.

이 두 관념, 즉 최고 기준(supreme criterion)과 궁극적 규칙(ultimate rule)
가운데 전자는 정의하기가 가장 쉽다. 법적 유효성의 기준 또는 법의 원천이
최고(supreme)라고 말할 수 있는 경우는, 그 기준을 참조(reference)하여
식별된 규칙들이 다른 기준들을 참조하여 식별된 규칙들과 충돌하더라도
여전히 체계의 규칙으로 인정되는 반면, 후자의 기준들을 참조하여 식별된
규칙들은 최고 기준을 참조하여 식별된 규칙들과 충돌하는 경우 그러한
인정이 이루어지지 않을 때이다. 우리가 이미 사용해 온 ‘우위의(superior)’
기준과 ‘종속된(subordinate)’ 기준이라는 관념들도 이와 유사한 비교적
설명에 의해 해명될 수 있다. 우위의 기준과 최고 기준이라는 관념들은 단지
척도상에서의 \emphksans{상대적인(relative)} 위치를 가리킬 뿐이며, 법적으로
\emphksans{무제한인(unlimited)} 입법 권한에 대한 어떠한 관념도 함의하지
않는다는 점은 명백하다. 그럼에도 불구하고 ‘최고(supreme)’라는 것과
‘무제한적(unlimited)’이라는 것은---적어도 법이론에서는---쉽게 혼동된다.
그 이유 중 하나는, 보다 단순한 형태의 법체계에서는 궁극적 승인
규칙(ultimate rule of recognition), 최고 기준, 그리고 법적으로
무제한적인 입법부라는 관념들이 서로 수렴하는 것처럼 보이기 때문이다.
헌법적 제한을 받지 않는 입법부가 존재하고, 그 입법부가 자신의
제정(enactment)을 통해 다른 원천들로부터 나온 다른 모든 법규칙들의 법적
지위를 박탈할 권능을 가진(competent) 경우, 그러한 체계의 승인 규칙의
일부로서 그 입법부의 제정이 유효성의 최고 기준이 된다. 헌법이론에
따르면, 이것이 영국(United Kingdom)의 경우이다. 그러나 미합중국과 같이
그러한 법적으로 무제한적인 입법부가 존재하지 않는 체계들조차도, 유효성의
기준들로 이루어진 하나의 집합을 제공하는 궁극적 승인 규칙을 충분히
포함할 수 있으며, 그 기준들 가운데 하나는 최고일 수 있다. 이는 일반
입법부의 입법 권능(legislative competence)이 개정 권한을 포함하지
않거나, 또는 그 개정 권한의 범위 밖에 놓인 조항들을 포함하는 헌법에 의해
제한되는 경우에 해당한다. 이 경우, ‘입법부’라는 개념을 가장 넓게
해석하더라도 법적으로 무제한적인 입법부는 존재하지 않지만, 그 체계는
물론 궁극적 승인 규칙을 포함하고 있으며, 헌법 조항들 속에 유효성의 최고
기준을 포함하고 있다.

\textbf{(p.~107)} 승인 규칙(rule of recognition)이 한 체계의
\emphksans{궁극적(ultimate)} 규칙이라는 의미는, 매우 익숙한 하나의 법적 추론
연쇄를 따라가 보면 가장 잘 이해된다. 어떤 후보 규칙(suggested rule)이
법적으로 유효한지 여부가 문제로 제기되는 경우, 그 질문에 답하기 위해서는
다른 어떤 규칙이 제공하는 유효성의 기준을 사용해야 한다. 옥스퍼드셔
카운티 의회(Oxfordshire County Council)의 이 이른바 조례(purported
by-law)는 유효한가? 그렇다: 왜냐하면 그것은 보건부 장관(Minister of
Health)이 제정한 법정 명령(statutory order)에 의해 부여된 권한을
행사하여, 그 명령에 명시된 절차에 따라 제정되었기 때문이다. 이 첫
단계에서 법정 명령은 조례의 유효성이 평가되는 기준을 제공한다.
실천적으로는 더 나아갈 필요가 없을 수도 있지만, 그렇게 할 수 있는
상존하는 가능성은 있다. 우리는 그 법정 명령의 유효성을 문제 삼고, 그러한
명령을 제정할 권한을 장관에게 부여한 제정법(statute)을 참조하여 그
유효성을 평가할 수 있다. 마지막으로, 그 제정법의 유효성마저 문제 삼아
의회 내 여왕이 제정하는 것은 법이라는 규칙을 참조하여 평가하게 되면,
유효성에 관한 탐구는 여기서 멈추게 된다: 왜냐하면 이는 우리가 다른
규칙들의 유효성을 평가하기 위한 기준을 제공하는 규칙에 도달했기
때문이며, 동시에 그 규칙은 중간 단계의 법정 명령이나 제정법과는 달리, 그
자신의 법적 유효성을 평가하기 위한 기준을 제공하는 어떠한 규칙도 더 이상
존재하지 않는다는 점에서 그들과 구별된다.

기실 이 궁극적 규칙에 대하여 제기할 수 있는 질문들은 많다. 우리는
영국에서 법원들, 입법부들, 공무담당자들, 또는 사적 시민들이 실제로 이
규칙을 궁극적 승인 규칙으로 사용하고 있는 실천이 존재하는지 물을 수
있다. 아니면 우리의 법적 추론 과정은 이제 폐기된 체계의 유효성 기준을
가지고 벌이는 공허한 놀이에 불과한 것인가? 우리는 또한 이러한 규칙을
근간으로 하는 법체계가 만족스러운 형태의 법체계인지 물을 수 있다. 그것은
악보다 선을 더 많이 산출하는가? 그것을 지지할 타산적 이유들(prudential
reasons)은 존재하는가? 그렇게 할 도덕적 의무(moral obligation)는 있는가?
이러한 질문들은 분명히 매우 중요하다. 그러나 이와 동일하게 분명한 것은,
우리가 승인 규칙에 관하여 이러한 질문들을 제기할 때, 우리는 더 이상 그
규칙을 도구로 삼아 다른 규칙들에 관해 답했던 것과 같은 종류의 질문에
답하려 하고 있는 것이 아니라는 점이다. 우리가 어떤 특정한
제정(enactment)이 의회 내 여왕이 제정하는 것은 법이다라는 규칙을
충족하기 때문에 유효하다고 말하는 단계에서, 영국에서 바로 이 마지막
규칙이 법원들, 공무담당자들, 그리고 사적 개인들에 의해 궁극적 승인
규칙으로 사용되고 있다고 말하는 단계로 이동할 때, \textbf{(p.~108)}
우리는 체계의 한 규칙의 유효성을 단언하는 내부 법 진술(internal
statement of law)에서, 그 규칙을 수용하지 않더라도 체계의 관찰자가 할 수
있는 사실에 관한 외부 진술(external statement of fact)로 이동한 것이다.
마찬가지로, 어떤 특정한 제정이 유효하다는 진술에서, 그 법체계의 승인
규칙이 훌륭한 것이며 그것에 기초한 체계가 지지할 가치가 있는 체계라는
진술로 이동할 때, 우리는 법적 유효성에 관한 진술에서 가치에 관한 진술로
이동한 것이다.

일부 저자들은 승인 규칙(rule of recognition)의 법적 궁극성(legal
ultimacy)을 강조하면서, 체계의 다른 규칙들의 법적 유효성은 이를 참조하여
입증될(demonstrated) 수 있는 반면, 승인 규칙 자체의 유효성은 입증될 수
없고 ‘가정된다(assumed)’거나 ‘상정된다(postulated)’거나 하나의
‘가설(hypothesis)’이라고 표현해 왔다. 그러나 이러한 표현은 심각하게
오해를 불러일으킬 수 있다. 법체계의 일상적 삶 속에서 판사, 법률가, 또는
일반 시민에 의해 특정 규칙들에 관하여 행해지는 법적 유효성에 대한
진술들은 기실 일정한 전제들을 수반한다. 이러한 진술들은 체계의 승인
규칙을 수용하는 자들의 관점을 표현하는 내부 법 진술들(internal
statements of law)이며, 그러한 것으로서 체계에 관한 외부 사실
진술들(external statements of fact)에서 진술될 수 있는 많은 내용들을
명시하지 않은 채로 남겨 둔다. 이렇게 명시되지 않은 것들은 법적 유효성
진술들의 정상적인 배경 또는 맥락을 이루며, 따라서 그러한 진술들에 의해
‘전제된다(presupposed)’고 말해진다. 그러나 이러한 전제된 사항들이 정확히
무엇인지를 분명히 파악하고, 그 성격을 흐리게 하지 않는 것이 중요하다.
그것들은 두 가지로 구성된다. 첫째, 어떤 주어진 법규칙, 예컨대 특정한
제정법(statute)의 유효성을 진지하게 단언하는 사람은, 법을 식별하는 데
적절한 것으로 자신이 수용하는 승인 규칙을 스스로 사용하고 있다. 둘째,
그가 특정한 제정법의 유효성을 평가하는 기준으로 삼는 이 승인 규칙은 그
자신에게만 수용되는 것이 아니라, 체계의 일반적 작동 속에서 실제로
수용되고 사용되고 있는 승인 규칙이라는 점이다. 만약 이
전제(presupposition)의 참임이 의심된다면, 그것은 실제 실천을
참조함으로써 확립될 수 있다: 즉, 법원들이 무엇을 법으로 간주할지를
식별하는 방식, 그리고 이러한 식별들에 대한 일반적인 수용 또는
묵인(acquiescence)을 참조함으로써이다.

이 두 전제(presuppositions) 중 어느 것도, 입증할 수 없는
‘유효성(validity)’에 대한 ‘가정(assumption)’으로 제대로 설명될 수는
없다. 우리는 ‘유효성’이라는 단어를 오직 \textbf{(p.~109)} 규칙 체계
\emphksans{내부(within)}에서 어떤 규칙이 그 체계의 구성원으로서의 지위를 갖기
위해 승인 규칙(rule of recognition)에 의해 제공된 특정 기준들을 충족해야
할 때, 그러한 질문에 답하기 위해서만 필요로 하며, 대체로 그러한 경우에만
사용한다. 승인 규칙 자체의 유효성에 대해서는 이러한 질문이 제기될 수
없다. 그것은 유효하거나 유효하지 않은 것이 아니라, 단지 그러한 방식으로
사용되기에 적절한 것으로 수용될 뿐이다. 이러한 단순한 사실을, 그 규칙의
유효성은 ‘가정되었지만 입증될 수 없다(assumed but cannot be
demonstrated)’고 모호하게 표현하는 것은, 미터 단위 측정의
정확성(correctness)에 대한 궁극적 검사 기준(ultimate test)인 파리의 표준
미터 자(metre bar) 자체가 정확하다는 것을 ‘우리가 가정하지만 결코 입증할
수 없다’고 말하는 것과 같다.

보다 심각한 반론은, 궁극적 승인 규칙(ultimate rule of recognition)이
유효하다는 ‘가정(assumption)’에 관한 말이, 법률가들의 유효성 진술 배후에
놓인 두 번째 전제의 본질적으로 사실적(factual) 성격을 은폐한다는 점이다.
의심할 여지 없이, 승인 규칙의 실제 존재를 구성하는 판사들, 공무담당자들,
기타 행위자들의 실천(practice)은 복합적인 사안이다. 우리가 뒤에서 보게
되겠지만, 이러한 종류의 규칙의 정확한 내용과 범위, 나아가 그 존재 여부에
관한 질문들이 명확하거나 확정적인(determinate) 답변을 허용하지 않는
상황들도 분명히 존재한다. 그럼에도 불구하고, 이러한 규칙의 ‘유효성을
가정하는 것(assuming the validity)’과 그 ‘존재를 전제하는
것(presupposing the existence)’을 구별하는 것은 중요하다. 최소한 이러한
구별의 실패가, 그러한 규칙이 \emphksans{존재한다(exists)}고 단언하는 것이
무엇을 의미하는지를 흐리게 만들기 때문이다.

지난 장에서 개괄한 의무의 일차 규칙(primary rules of obligation)으로
이루어진 단순한 체계에서는, 어떤 특정 규칙이 존재한다고 단언하는 것은,
그 규칙을 수용하지 않는 관찰자가 제시할 수 있는 외부 사실 진술(external
statement of fact)에 불과하다. 이러한 진술은, 실제로 어떤 행위 양식이
일반적으로 표준으로 수용되고 있는지, 그리고 우리가 보았듯이 단순한
수렴적 습관들(convergent habits)과 사회적 규칙(social rule)을 구별하는
그러한 특징들이 수반되는지를 확인함으로써 검증될 수 있다. 마찬가지로,
현재 우리가 “잉글랜드에서는---비록 법적 규칙은 아니지만---교회에 들어갈
때 모자를 벗어야 한다는 규칙이 존재한다”는 단언을 해석하고 검증해야
하는 방식도 바로 이러한 방식이다. 이러한 유형의 규칙들이 사회 집단의
실제 실천 속에서 존재하는 것으로 확인된다면, 그 유효성에 관하여 별도로
논의할 문제는 없다. 물론 그 가치나 바람직성은 문제 삼을 수 있다. 그러나
일단 그 존재가 사실로 확립되었다면, 그것들이 유효하다고 말하거나
유효하지 않다고 부정하거나, \textbf{(p.~110)} 혹은 “우리는 그것들의
유효성을 가정했을 뿐 이를 입증할 수는 없다”고 말하는 것은 오히려 사안을
혼란스럽게 만들 뿐이다. 반면에, 성숙한 법체계와 같이 승인 규칙을
포함하는 규칙 체계를 갖는 경우에는, 어떤 규칙이 체계의 구성원으로서의
지위를 갖는지는 이제 승인 규칙이 제공하는 일정한 기준들을 충족하는지
여부에 달려 있다. 이는 ‘존재한다(exist)’라는 단어에 새로운 용법을
가져온다. 이제 어떤 규칙이 존재한다고 말하는 진술은, 관습 규칙의 단순한
경우에서와 같이, 특정한 행위 양식이 실제 실천에서 일반적으로 표준으로
수용되고 있다는 \emphksans{사실(fact)}에 관한 외부 사실 진술이 더 이상 아닐
수 있다. 그것은 이제 수용되었으나 진술되지 않은 승인 규칙을 적용하는
내부 진술일 수 있으며, 대략적으로 말해 ‘그 체계의 유효성 기준들에 비추어
유효하다(valid given the system’s criteria of validity)'는 의미에 불과할
수 있다. 그러나 이러한 점에서, 그리고 다른 점들에서도, 승인 규칙은
체계의 다른 규칙들과는 다르다. 승인 규칙이 존재한다고 단언하는 것은 오직
외부 사실 진술일 수밖에 없다. 종속된 규칙(subordinate rule)은 일반적으로
무시되고 있더라도 유효할 수 있으며, 그 의미에서 ‘존재’할 수 있는 반면,
승인 규칙은 일정한 기준들을 참조하여 법을 식별하는 법원, 공무담당자,
사적 개인들(private persons)의 복합적이지만 통상적으로 일치하는
실천으로서만 존재한다. 그 존재는 하나의 사실 문제이다.

\subsection{2. 새로운 질문들(New
Questions)}\label{uxc0c8uxb85cuxc6b4-uxc9c8uxbb38uxb4e4new-questions}

법체계의 토대들이 법적으로 무제한인 주권자에 대한 복종 습관에 있다고
보는 관점을 버리고, 대신 규칙 체계에 유효성의 기준들을 제공하는 궁극적
승인 규칙이라는 구상(conception)으로 대체하게 되면, 우리는 흥미롭고
중요한 일련의 질문들과 마주하게 된다. 이러한 질문들은 비교적 새로운
것들인데, 그 이유는 오랫동안 법리학과 정치이론이 구시대적인 사유방식에
매여 있었던 동안에는 그러한 질문들이 가려져 있었기 때문이다. 이들은 또한
어려운 질문들이며, 충분한 답변을 위해서는 한편으로는 헌법(constitutional
law)의 몇몇 근본 문제들에 대한 이해가, 다른 한편으로는 법적 형식들(legal
forms)이 조용히 전환되고 변화하는 특유의 방식을 파악할 수 있는 통찰이
요구된다. 따라서 우리는 이 질문들을, 법 개념의 해명에서 일차 규칙들과
이차 규칙들의 결합에 중심적인 위치를 부여해야 한다고 우리가 해왔듯이
고수하는 것이 현명한지 또는 현명하지 않은지와 관련되는 범위 내에서만
탐구할 것이다.

\textbf{(p.~111)} 첫 번째 어려움은 분류(classification)의 문제이다.
왜냐하면 궁극적으로 법을 식별하기(identify) 위해 사용되는 그 규칙은,
흔히 전부를 포괄하는(exhaustive) 것으로 간주되는 법체계를 기술하는
관행적(conventional) 범주들로부터 벗어나 있기 때문이다. 따라서
다이시(Dicey) 이후의 영국 헌법 이론가들은 대체로, 영국의 헌정적 배치들이
한편으로는 엄밀히 그렇게 불리는 법들(laws strictly so called)---즉
제정법들(statutes), 추밀원 명령들(orders in council), 그리고
선례(precedents)에 구현된 규칙들---로, 다른 한편으로는 단지
용례(usages), 양해(understandings), 또는 관습(customs)에 불과한
관행들(conventions)로 구성되어 있다고 반복하여 진술해 왔다. 후자에는,
예컨대 여왕은 상원과 하원에서 적법하게 통과된 법안에 대해 동의를 거부할
수 없다는 것과 같은 중요한 규칙들이 포함된다. 그러나 여왕에게 그러한
동의를 할 법적 책무는 존재하지 않으며, 이러한 규칙들이 관행이라고 불리는
이유는 법원이 그것들을 법적 책무를 부과하는 것으로 인정하지 않기
때문이다. 분명히, 의회 내 여왕이 제정하는 것은 법이다라는 규칙은 이들
어느 범주에도 속하지 않는다. 그것은 법원이 그 규칙과 가장 밀접하게
관련되어 있고 이를 사용하여 법을 식별하는 점에서 관행이 아니며, 동시에
그것이 식별의 기준으로 사용되는 ‘엄밀히 그렇게 불리는 법들’과 동일한
수준의 규칙도 아니다. 설령 이 규칙이 제정법으로 제정되었다고 하더라도,
그것이 단순한 제정법의 수준으로 낮아지는 것은 아니다. 그러한 제정 행위의
법적 지위는 필연적으로, 그 규칙이 해당 제정에 선행하고 그 제정과
독립적으로 존재했다는 사실에 의존하기 때문이다. 더 나아가, 앞 절에서
보였듯이, 그 존재는 제정법의 존재와는 달리 실제 실천(actual
practice)으로 구성되어야 한다.

이러한 측면은 일부에게서 절망의 외침을 끌어낸다: 즉, 분명히 법인 헌법의
근본 규정들이 정말로 법이라는 것을 우리는 어떻게 입증할 수 있는가? 이에
대해 다른 이들은, 법체계의 토대에는 ‘법이 아닌 것’, 즉
‘전법적인(pre-legal) 것’, ‘메타-법적인(meta-legal) 것’, 또는 단순히
‘정치적 사실(political fact)’이 존재한다고 고집한다. 이러한 불안감은
모든 법체계에서 가장 중요한 이 특징을 기술하기 위해 사용되어 온 범주들이
지나치게 거칠다는 확실한 징후이다. 승인 규칙을 ‘법’이라고 부를 수 있다는
근거는, 체계 내의 다른 규칙들을 식별하기 위한 기준을 제공하는
규칙이야말로 법체계를 정의하는 특징(defining feature)으로 충분히 간주될
수 있으며, 따라서 그 자체가 ‘법’이라고 불릴 만하다는 점에 있다. 반대로
그것을 ‘사실’이라고 부를 수 있다는 근거는, 그러한 규칙이 존재한다고
단언하는 것이 기실 \textbf{(p.~112)} ‘실효성 있는(efficacious)’ 체계에서
규칙들이 어떠한 방식으로 식별되는지에 관한 실제 사실(actual fact)에 대한
외부 진술이기 때문이다. 이 두 측면은 모두 주목을 요구하지만, 우리는 ‘법’
또는 ‘사실’이라는 어느 한 명칭(label)을 선택하는 것만으로는 이 둘 모두를
정당하게 다룰 수 없다(do justice). 대신, 우리는 궁극적 승인 규칙을 두
가지 관점에서 바라볼 수 있음을 기억해야 한다. 하나는 그 규칙이 체계의
실제 실천(actual practice) 속에 존재한다는 외부 사실 진술로 표현되는
관점이며, 다른 하나는 그 규칙을 사용하여 법을 식별하는 사람들이 행하는
유효성(validity)에 관한 내부 진술로 표현되는 관점이다.

두 번째 질문군은 특정 국가나 사회 집단 안에 법체계가
\emphksans{존재한다(exists)}고 하는 단언(assertion)의 숨겨진 복잡성과
모호성(vagueness)에서 비롯된다. 우리가 이러한 단언을 할 때, 실제로는
대개 함께 수반되는 서로 이질적인 여러 사회적 사실들(heterogeneous social
facts)을 압축적 포괄 형식(portmanteau form)으로 가리키고 있는 것이다.
오해를 불러일으키는 이론의 그림자 아래 발전된 법적 사고와 정치적 사고의
표준 용어법(standard terminology)은 사실들을 지나치게 단순화하고 가리기
쉽다. 그러나 이러한 용어법이 구성하는 안경(spectacles)을 벗고 사실들을
바라보면, 법체계(legal system)는 인간과 마찬가지로, 어떤 단계에서는 아직
태어나지 않은 상태(unborn)에 있고, 또 어떤 단계에서는 아직 모체로부터
완전히 독립하지 않은 상태(not yet wholly independent)에 있으며, 이후
건강하고 독립적인 존재(healthy independent existence)를 누리다가,
쇠퇴(decay)하고 마침내 사망(die)하기도 한다는 것이 분명해진다. 이러한
‘출생과 정상적 독립 존재 사이’, 또는 ‘그 독립된 존재와 죽음 사이’의 중간
단계들은 우리가 익숙하게 사용하는 법 현상에 대한 설명 방식을 어긋나게
만든다(put out of joint). 그러나 이 중간 단계들은 그만큼 우리가 정상
사례에서 “특정 국가에는 법체계가 존재한다”라는 확신에 찬 참된 단언을
할 때 당연하게 여기는 것의 온전한 복잡성을 부각시키기 때문에,
혼란스러울지라도 연구할 가치가 있다.

이러한 복합성을 인식하는 한 가지 방식은, 명령들에 대한 일반적 복종의
습관이라는 단순한 오스틴식 공식(Austinian formula)이, 사회가 법체계를
갖기 위해 충족해야 하는 최소 조건들을 구성하는 복합적 사실들을 재현하는
데서 정확히 어디에서 실패하거나 왜곡하는지를 살펴보는 데 있다. 우리는 이
공식이 하나의 필요조건을 지목한다(designate)는 점은 인정할 수 있다. 즉,
법들이 의무(obligation)나 책무(duty)를 부과하는 경우, 그 법들에는
일반적으로 복종이 이루어져야 하거나 적어도 일반적으로 불복종이
이루어지지 않아야 한다는 점이다. 그러나 이는 필수적이긴 하나, 법체계가
사적 시민에게 영향을 미치는 지점에서 나타나는 우리가 ‘최종 산물(end
product)’이라 부를 수 있는 것만을 다룰 뿐이며, 법체계의 일상적 존재는
\textbf{(p.~113)} 또한 공무담당자들에 의한 법의 창출(official creation),
공무담당자들에 의한 식별, 그리고 공무담당자들에 의한 사용 및 적용으로도
이루어져 있다. 여기에서 관련되는 법과의 관계를 ‘복종(obedience)’이라
부를 수 있으려면, 그 단어를 통상적 용법을 훨씬 넘어 확장하여, 결국에는
이러한 작용들을 정보를 주는 방식으로 특징짓는 기능을 상실하게 해야 한다.
입법자들이 법들을 제정할 때 자신들의 입법 권한들을 부여하는 규칙들에
부합하는 경우, 그러한 권한들을 부여하는 규칙들이 그것들을 따를 책무를
부과하는 규칙들에 의해 강화되지 않는 한, 어떤 통상적 의미에서도
입법자들이 규칙에 ‘복종’한다고 말할 수는 없다. 또한 이러한 규칙들에
부합하지 않았다고 해서 그들이 법에 ‘불복종하는(disobey)’ 것도 아니며,
다만 법을 만들어내는 데 실패할 뿐이다. 마찬가지로, 판사가 체계의 승인
규칙(rule of recognition)을 적용하여 어떤 제정법(statute)을 유효한
법으로 인정하고 분쟁의 판정(determination)에 이를 사용하는 경우에,
‘복종’이라는 말은 그가 하는 일을 적절히 묘사하지 못한다. 물론 우리는
사실들 앞에서 ‘복종’이라는 단순한 용어를 유지하고자 한다면, 여러 장치를
통해 그렇게 할 수는 있다. 그 한 예로, 판사가 제정법을 인정하는 데서
사용하는 일반적 유효성 기준을, ‘헌법 제정자들(Founders of the
Constitution)’이 내린 명령들에 대한 복종의 사례로 표현하거나, (그러한
‘제정자들’이 없는 경우에는) 지휘관 없는 사령, 즉 ‘탈심리화된
사령(depsychologized command)’에 대한 복종으로 표현할 수 있을 것이다.
그러나 후자의 관념은, 삼촌 없는 조카라는 관념만큼이나, 우리의 주의를
진지하게 요구할 이유가 없을지도 모른다. 다른 한편으로는, 법의 공무담당
측면 전체를 시야에서 밀어내고 입법과 재판에서 이루어지는 규칙의 사용에
대한 서술을 포기한 채, 공무담당 세계 전체를 하나의 인격(‘주권자’)이 여러
대리인이나 대변자를 통해 명령들을 발하는 것으로 생각하고, 시민이 그
명령들에 습관적으로 복종한다고 간주할 수도 있다. 그러나 이는 여전히
서술을 기다리는 복합적 사실들에 대한 편의적 축약에 불과하거나, 아니면
재앙적으로 혼란을 초래하는 신화적 구성물일 뿐이다.

법체계가 존재한다는 것을, 기실 일반 시민과 법의 관계를 특징짓기는 하지만
그것을 언제나 남김없이 서술하지는 않는 습관적 복종(habitual
obedience)이라는 기분 좋을 만큼 단순한 표현으로 설명하려는 시도가 실패한
데 대한 반작용으로, 그 반대의 오류를 범하는 일은 자연스러운 일이다. 그
오류는 공무담당 활동들(official activities)---특히 사법적 태도(judicial
attitude)나 법에 대한 관계---의 특징이지만 역시 남김없는 설명은 아닌
것을 취하여, 그것을 법체계를 가진 사회집단에서 반드시 존재해야 하는 것에
대한 충분한 설명으로 취급하는 데 있다. 이러한 오류는 사회 구성원의
대다수가 법에 대해 단순히 습관적으로 복종한다는 단순한 구상을, 그들이
법들의 유효성을 최종적으로 평가하는 기준을 특정하는 궁극적 승인
규칙(rule of recognition)을 일반적으로 공유하거나, 수용하거나, 구속력
있는 것으로 간주해야 한다는 구상으로 대체하는 데 이른다. 물론 우리는
3장에서 살펴본 것처럼, 법의 원천들에 관한 지식과 이해가 사회 전반에 널리
퍼진 단순한 사회를 상상할 수 있다. 그곳에서는 ‘헌법(constitution)’이
너무 단순해서, 일반 시민뿐 아니라 공무담당자들과 법률가들 역시 그 헌법을
알고 있으며 수용하고 있다고 말해도 아무런 허위가 없을 것이다. 렉스 1세의
단순한 세계에서는, 인구 다수가 그의 말에 단순히 습관적으로 복종한 것 그
이상이 있었다고 말할 수도 있다. 그 사회에서는 수범자들과 체계의
공무담당자들 모두가 렉스의 말(Rex's word)을 사회 전체를 위한 유효한 법의
기준으로 규정하는 하나의 승인 규칙을 명시적이고 의식적인 방식으로
‘수용하고(accepted)’ 있었을 수 있다. 물론 수범자들과 공무담당자들은
수행해야 할 역할과 이 기준에 따라 식별되는 법규칙들에 대한 관계가 서로
달랐겠지만 말이다. 그러나 이렇게 단순한 사회에서 가능한 상태를, 복잡한
현대 국가에서도 항상 또는 일반적으로 존재한다고 주장하는 것은 허구를
고집하는 일일 것이다. 여기서 상황의 현실은 분명 이렇다: 상당수의 일반
시민, 어쩌면 과반수는 법적 구조나 그 유효성 기준들에 대한 일반적 구상을
가지고 있지 않다. 그가 복종하는 당해 법(the law)은 그에게 단지
‘법’으로만 알려져 있는 어떤 것일 뿐이다. 그는 여러 서로 다른 이유들로 그
법에 복종할 수 있으며, 그 이유들 중에는 그러는 것이 자신에게 가장
이롭다는 지식이 자주, 비록 항상은 아니지만, 포함될 수 있다. 그는 불복종
시 발생할 수 있는 일반적으로 개연성 있는 결과들---예컨대 자신을 체포할
수 있는 공무담당자들이 있고, 자신을 재판하여 법을 어겼다는 이유로 감옥에
보낼 수 있는 이들이 있다는 것---을 알고 있을 것이다. 그 체계의 유효성
검사 기준들에 따라 유효한 법들에 인구 대다수가 복종하는 한, 이는 특정한
법체계가 존재한다는 것을 확립하기 위해 필요한 증거 전부이다.

그러나 법체계가 일차 규칙들과 이차 규칙들의 복합적 결합이라는 바로 그
이유 때문에, 이러한 증거만으로는 법체계의 존재에 수반되는 법과의
관계들을 모두 서술하기에 충분하지 않다. 이는 체계의 공무담당자들이
공무담당자로서 관련되는 이차 규칙들에 대해 갖는 관계에 관한
\textbf{(p.~115)} 서술로 보충되어야 한다. 여기서 핵심적인 것은 체계의
유효성 기준들을 포함하는 승인 규칙에 대한 공무담당자들의 통일적이거나
공유된 수용(a unified or shared official acceptance)이 존재해야 한다는
점이다. 그러나 바로 이 지점에서, 일반 시민의 경우에 불가결한 최소를
특징짓기에는 충분했던 일반적 복종이라는 단순한 관념은 부적절해진다.
요점은 ‘복종(obedience)’이라는 표현이 법원과 다른 공무담당자들이 이러한
이차 규칙들을 규칙으로서 존중하는 방식을 가리키는 데 자연스럽게 사용되지
않는다는 ‘언어적(linguistic)’ 문제가 아니며, 적어도 단지 그 문제만은
아니다. 필요하다면, 우리는 ‘따르다(follow)’, ‘준수하다(comply)’, 또는
‘부합하다(conform to)’와 같은 보다 포괄적인 표현을 찾아낼 수 있을
것이다. 이러한 표현들은 시민이 병역에 응하기 위해 법과 관련하여 수행하는
행위와, 판사가 의회 내 여왕이 제정하는 것은 법이라는 전제하에 자신의
법정에서 특정 제정법(statute)을 법으로 식별하는 행위를 모두 특징지을 수
있을 것이다. 그러나 이러한 포괄 용어들(blanket terms)은, 우리가
법체계라고 부르는 이 복합적 사회 현상의 존재에 포함된 최소 조건을
이해하기 위해 반드시 파악되어야 할 중대한 차이들을 단지 가려버릴 뿐이다.

입법자들이 자신들의 권한들을 부여하는 규칙들에 부합하여 입법 활동을
하거나, 법원이 수용된 궁극적 승인 규칙(ultimate rule of recognition)을
적용할 때, 이를 ‘복종(obedience)’이라고 묘사하는 것이 오해를
불러일으키는 이유는 다음과 같다. 어떤 규칙(또는 명령)에 복종한다고 해서,
그 복종하는 사람이 자신이 하는 일이 자기 자신뿐 아니라 다른 사람들에게도
옳은 일이라고 생각해야 할 \emphksans{필요는(need)} 없다. 그는 자신이 하는
행동을 사회집단의 다른 구성원들에게도 적용되는 행동 기준을 실현하는
것으로 보지 않아도 된다. 그는 자신이 보이는 부합 행위를 ‘옳다(right)’,
‘올바르다(correct)’, 또는 ‘의무적이다(obligatory)’라고 생각할 필요가
없다. 즉, 그의 태도에는 사회 규칙이 수용되거나 특정 행위 유형이 일반적
기준으로 간주될 때 수반되는 비판적 성격(critical character)이 전혀 없을
수 있다. 그는 그러한 규칙들을 그 규칙들이 적용되는 모든 사람을 위한
기준으로 수용하는 내부 관점(internal point of view)을 공유하지 않을 수도
있다. 대신 그는 그 규칙을 단지 \emphksans{그에게(him)} 요구되는 것, 즉
제재적-불이익(penalty)의 위협 아래 행동을 요구하는 것으로만 여길 수
있다. 그는 그저 결과에 대한 두려움 때문에, 혹은 관성에 따라 복종할 수
있으며, 자신이나 타인이 그것을 따를 의무가 있다고 여기지 않고, 규칙을
어기는 자신이나 타인을 비판할 의향도 없을 수 있다. 하지만,
\textbf{(p.~116)} 모든 일반 시민이 규칙에 복종할 때 가질 \emphksans{수도
있는(may)} 이러한 개인적인 차원의 관심만으로는, 법원이 자신이 법원으로서
다루는 규칙들에 대해 갖는 태도를 설명할 수 없다. 이는 특히 다른 규칙들의
유효성을 평가하는 데 쓰이는 궁극적 승인 규칙(ultimate rule of
recognition)의 경우에 분명하다. 이 규칙은, 존재하기 위해서 반드시 하나의
공적인(public), 공통의(common) 올바른 사법적 결정의 기준(standard of
correct judicial decision)으로서 내부 관점(internal point of view)에서
수용되어야 하며, 단순히 각 판사가 자기 몫으로만 ‘복종’하는 어떤 것으로
여겨져서는 안 된다. 그 법체계의 법원들 중 개별적인 법원이 때때로 이
규칙들로부터 이탈하는 일이 있을 수는 있으나, 일반적으로 그러한 이탈을
공통적이거나 공적인 기준(common or public standards)에서의
일탈(lapses)로 비판적으로 다루어야 한다. 이는 단지 법체계의 효율성이나
건전성과 관련된 문제가 아니라, 하나의 법체계(single legal system)가
존재한다고 말할 수 있기 위한 논리적으로 필수적인 조건(logically
necessary condition)이다. 가령 어떤 판사들만이 ‘자기 몫으로만’ 행동하며
“의회 내 여왕이 제정하는 것은 법이다”라는 전제에 서고, 이 승인 규칙을
존중하지 않는 이들을 아무도 비판하지 않는다면, 법체계의 특유의 통일성 및
계속성(characteristic unity and continuity)은 사라지고 말 것이다. 이러한
통일성과 계속성은 바로 이 핵심적 지점에서 법적 유효성에 관한 공통된
기준(common standards of legal validity)이 수용되는 데 의존하기
때문이다. 이러한 사법적 행태의 변칙들과, 평범한 사람이 서로 상반되는
사법 명령들에 직면하게 될 때 결국 초래될 혼란 사이의 그 사이 구간에서,
우리는 그 상황을 어떻게 묘사해야 할지 난처할 것이다. 우리는 오직 그것이
평소에는 너무 자명해서 인식되지 않는 것들을 날카롭게 의식하게 해주기
때문에만 생각해볼 만한 가치가 있는 하나의 \emphksans{자연의 기이함(lusus
naturae)}과 마주하게 될 것이다.

따라서 하나의 법체계가 존재하기 위해 필요하고 충분한 최소 조건은 두
가지이다. 첫째, 그 법체계의 궁극적 유효성 기준들(ultimate criteria of
validity)에 따라 유효한 것으로 간주되는 행위 규칙들(rules of
behaviour)에는 일반적으로 복종이 이루어져야 한다. 둘째, 법적 유효성
기준들을 특정하는 승인 규칙들(rules of recognition)과 변경 규칙들(rules
of change) 및 재판 규칙들(rules of adjudication)은, 해당 체계의
공무담당자들 사이에서 공무담당 행위의 공통된 공적 기준들(common public
standards of official behaviour)로 실효적으로 수용되어야 한다. 첫 번째
조건은 일반 시민들이 만족시킬 \emphksans{필요가 있는(need)} 유일한 조건이다:
그들은 각자 ‘자기 몫으로만’ 복종할 수 있으며, 그 동기는 어떠하든
상관없다; 물론 건강한 사회에서는 실제로 많은 시민들이 이 규칙들을 공통의
행동 기준들로 수용하고, 그것들에 복종할 의무가 있음을
인정하거나(acknowledge), 심지어 이러한 의무를 \textbf{(p.~117)} 헌법을
존중할 보다 일반적인 의무로 거슬러 추적하기도 할 것이다. 두 번째 조건은
법체계의 공무담당자들이 반드시 만족시켜야 한다. 그들은 해당 규칙들을
공무담당 행위의 공통된 기준들(common standards of official behaviour)로
간주하고, 자기 자신과 동료들의 일탈(deviations)을 기준에서의
이탈(lapses)로 비판적으로 평가해야 한다. 물론 여기에 덧붙여,
공무담당자가 단지 개인적 행위능력(personal capacity)에서 복종하기만 하면
되는 일차 규칙들도 많다는 점은 사실이다.

따라서 “법체계가 존재한다”는 단언은 이중적인 얼굴을 가진(Janus-faced)
진술이다. 한편으로는 일반 시민들의 복종을, 다른 한편으로는
공무담당자들이 이차 규칙들을 공무담당 행위의 비판적 공통 기준들(critical
common standards of official behaviour)으로 수용하는 것을 함께 가리킨다.
이러한 이중성에 놀랄 필요는 없다. 이것은 단지 법체계의 복합적
성격(composite character)---즉 단순하고 분권적인 전법적 사회 구조와
대비되는---을 반영하는 것일 뿐이다. 단지 일차 규칙들만으로 구성된 더
단순한 사회 구조에서는, 공무담당자가 존재하지 않기에, 해당 규칙들은
집단의 행위에 대한 비판적 기준으로서 널리 수용되어야 한다. 만약 그러한
경우에 내부 관점(internal point of view)이 널리 퍼져 있지 않다면,
논리적으로 규칙이 존재할 수 없다. 하지만, 우리가 주장해온 바와 같이,
법체계를 이해하는 가장 유익한 방식은 일차 규칙들과 이차 규칙들의
결합(union of primary and secondary rules)으로 보는 것이며, 이 경우 집단
전체의 공통 기준으로서 규칙을 수용하는 일은, 일반 개인이 자신 몫으로
복종함으로써 그 규칙에 묵인하는 비교적 수동적인 태도와 분리될 수 있다.
극단적인 경우에는, 법적 언어의 특유의 규범적 사용(예: “이것은 유효한
규칙이다”)을 포함한 내부 관점이 오직 공무담당 세계 내부에만 존재할 수도
있다. 이처럼 보다 복잡한 체계에서는, 오직 공무담당자들만이 법적 유효성에
대한 체계의 기준들을 수용하고 사용할 수 있다. 그러한 사회는 비참할 만큼
양처럼 순응적인(sheeplike) 사회일 수 있고, 결국 도살장으로 향하게 될지도
모른다. 하지만 그 사회가 실제로 존재할 수 없다고 보거나, 그것이
법체계라는 명칭을 가질 수 없다고 단정할 이유는 별로 없다.

\subsection{3. 법체계의 병리학(THE PATHOLOGY OF A LEGAL
SYSTEM)}\label{uxbc95uxccb4uxacc4uxc758-uxbcd1uxb9acuxd559the-pathology-of-a-legal-system}

따라서 법체계의 존재에 대한 증거는 사회생활의 두 가지 서로 다른
영역으로부터 도출되어야 한다. 우리가 법체계가 존재한다고 확신을 가지고
말할 수 있는 정상적이고 문제없는 경우란, 바로 이 두 영역이
\textbf{(p.~118)} 법에 대한 각자의 전형적인 관심사에 있어서 서로
일치(congruent)하고 있음이 분명한 경우이다. 거칠게 말하면, 사실관계는
다음과 같다. 즉, 공무담당자 수준에서 유효한 것으로 인정되는 규칙들에
일반적으로 복종이 이루어진다. 그러나 때로는 공무담당 영역이 사적
영역으로부터 분리될(detached) 수 있는데, 이는 법원에서 사용되고 있는
유효성 기준들(criteria of validity)에 따라 유효한 규칙들에 더 이상
일반적 복종이 존재하지 않는다는 의미에서 그러하다. 이러한 상황이
발생하는 다양한 방식들은 법체계의 병리(pathology of legal systems)에
속한다. 왜냐하면 그것들은, 우리가 법체계가 존재한다는 외부 사실
진술(external statement of fact)을 할 때 지칭되는(referred to)
복합적이고 상호 일치적인 실천의 붕괴를 나타내기 때문이다. 여기에는,
특정한 체계 내부로부터 법에 관한 내부 진술을 할 때마다
전제(presupposed)되는 것의 부분적 실패가 존재한다. 이러한 붕괴는 서로
다른 교란 요인들의 산물일 수 있다. 집단 내부에서 통치에 대한 경쟁적
주장들이 제기되는 ‘혁명(revolution)’은 그 한 사례에 불과하며, 이는
언제나 기존 체계의 일부 법들의 위반을 수반하겠지만, 새로운 헌법이나
새로운 법체계가 아니라 단지 공무담당자들을 새로운 개인 집합으로 법적
권한 없이 대체하는 것만을 수반할 수도 있다. 기존 체계 하에서의 권한 없이
외부로부터 통치에 대한 경쟁적 주장이 제기되는 적군 점령(enemy
occupation)도 또 하나의 사례이며, 통치에 대한 정치적 주장조차 없이
무정부 상태나 도적 행위에 직면하여 질서 있는 법적 통제가 단순히 붕괴되는
경우 역시 또 다른 사례이다.

이러한 각각의 경우에서, 법원이 해당 영토에서든 망명 상태에서든
기능하면서, 한때 안정적으로 확립된 옛 체계의 유효성 기준을 여전히
사용하는 중간 단계들(half-way stages)이 있을 수 있다. 그러나 이러한
법질서들은 해당 영토에서 실효적이지 않다. 이러한 경우에 “법체계가
완전히 소멸되었다”고 말할 수 있는 시점은, 어떤 정확한 확정(exact
determination)도 허용하지 않는 사안이다. 명백하게, 복구될 상당한
가망(chance)이 있거나, 기존 체계의 혼란이 결과가 아직 확정되지 않은
전면적 전쟁의 일부 사건일 경우에는, 그 체계가 더 이상 존재하지 않는다는
무조건적 단언(assertion)은 정당화될 수 없다. 이는 바로 “법체계가
존재한다”는 진술(statement)이 중단을 허용할 만큼 충분히 폭넓고 일반적인
유형이기 때문이다. 이 진술은 짧은 시간 동안 일어나는 일들에 의해
검증되거나 반증되는 것이 아니다.

물론 이러한 중단 이후에 법원과 인구 간의 정상적인 \textbf{(p.~119)}
관계가 회복된 뒤에는 어려운 질문들이 제기될 수 있다. 점령군의 축출이나
반란 정부의 패배 이후 망명 정부가 귀환하면, 그 중단의 시기에 해당
영토에서 무엇이 ‘법’이었는지, 혹은 무엇이 아니었는지에 대한 질문들이
생긴다. 여기서 가장 중요한 점은, 이 질문이 사실(fact)의 문제가 아닐 수도
있다는 점이다. 만약 그것이 사실의 문제라면, 해당 중단이 너무나
장기적이고 완전했기 때문에 기존 법체계가 소멸되었고 망명 귀환 시 유사한
새로운 체계가 수립되었다고 기술해야 하는지를 묻는 방식으로
확정되어야(settled) 할 것이다. 그러나 실제로는 이 질문은 국제법상의
문제로 제기되거나, 다소 역설적으로, 복귀 이후 존재하게 된 바로 그 법체계
내부의 법 문제로 제기될 수 있다. 후자의 경우, 복귀한 체계가, 그 체계가
계속해서 해당 영토의 법이었다고 소급적으로 선언하는 법을 포함할 수도
있다---혹은 좀 더 솔직히 말해, 그 체계가 해당 시기 동안 법이었던 것으로
간주된다(deemed)고 선언할 수 있다. 이러한 선언은, 그 중단 기간이 사실
문제로 취급되었을 경우 도달했을 결론과는 상당히 불일치해 보일 수
있음에도 불구하고, 그렇게 이루어질 수 있다. 그런 경우에, 해당 선언은
중단 기간 중 발생한 사건들과 거래들에 대해 법원이 적용해야 할 법을
정하는 복귀 체계의 하나의 규칙으로서, 얼마든지 존립할(stand) 수 있다.

여기에 역설이 있는 것은 오직, 법체계가 스스로의 과거, 현재, 또는 미래의
존재 단계로 간주되어야 하는 것에 관하여 내리는 법 진술들(statements of
law)을, 외부 관점에서 제시되는 그 존재에 대한 사실 진술과
경쟁하는(rivals) 것으로 생각할 때뿐이다. 자기-참조(self-reference)의
겉보기 수수께끼를 제외하면, 어떤 현행 체계 안에서 그 체계가 존재한
것으로 간주되어야 하는 기간에 관한 규정의 법적 지위는, 한 체계의 법이
다른 나라에 어떤 체계가 아직 존재한다고 선언하는 경우와 다르지 않다.
물론 후자의 경우 실천적 귀결이 많을 개연성은 높지 않다. 예를 들어,
우리는 소련 영토에서 현재 존재하는 법체계가 사실상 차르 체제의 법체계가
아니라는 점을 분명히 알고 있다. 그러나 만약 영국 의회의
제정법(statute)이 차르 시대 러시아의 법이 여전히 러시아 영토의 법이라고
선언한다면, 이 제정법은 소련(USSR)에 관한 영국법의 일부로서 기실 의미와
법적 효과를 가질 것이다. \textbf{(p.~120)} 그러나 이는 직전 문장에 담긴
사실 진술의 참임에 전혀 영향을 미치지 않는다. 그 제정법의 효력(force)과
의미는, 영국 법원이 적용할 법을 규정하고, 그 결과로 영국 내에서 러시아와
관련된 사안에 적용될 법을 정하는 것에 있을 뿐이다.

방금 서술한 상황의 역상황(converse)은, 새로운 법체계가 옛 법체계의
자궁으로부터---때로는 제왕절개(Caesarian operation)를 거친
뒤에야---출현하는 매혹적인 전환기의 순간들에서 볼 수 있다.
영연방(Commonwealth)의 최근 역사는 법체계 발생론(embryology of legal
systems)의 이러한 측면을 연구하기에 훌륭한 장(場)을 제공한다. 이 발전
과정을 도식적이고 단순화하여 개관하면 다음과 같다. 한 시기의 초기에
우리는 지방 입법부(local legislature), 사법부(judiciary),
행정부(executive)를 갖춘 식민지를 발견할 수 있다. 이러한 헌정 구조는
영국 의회(United Kingdom Parliament)의 제정법에 의해 설정된 것으로, 영국
의회는 식민지에 대해 입법할 완전한 법적 권능(legal competence)을
보유하고 있다. 여기에는 지방 법들(local laws)뿐 아니라, 식민지 헌법에
관한(referring to) 제정법들을 포함하여, 자기 자신의 제정법들(statutes)을
개정하거나 폐지할 권한도 포함된다. 이 단계에서 식민지의 법체계는, 의회
내 여왕(Queen in Parliament)이 제정하는(enact) 것은 \emphksans{(다른 것들
중에서도(inter alia))} 식민지에 대하여 법이다라는 궁극적 승인 규칙(rule
of recognition)에 의해 특징지어지는 더 넓은 체계의 명백한
종속된(subordinate) 부분이다. 발전 단계의 말미에 이르면, 궁극적 승인
규칙이 전환되었음(shift)을 확인하게 되는데, 이는 웨스트민스터
의회(Westminster Parliament)가 옛 식민지에 대해 입법할 법적 권능(legal
competence)이 더 이상 그곳의 법원들에서 인정되지 않기 때문이다. 옛
식민지의 헌정 구조 상당 부분이 여전히 웨스트민스터 의회의 원래
제정법(statute)에서 발견된다는 점은 여전히 사실이다. 그러나 이는 이제
단지 역사적 사실에 불과하며, 그것은 더 이상 해당 영토에서의 현재적 법적
지위를 웨스트민스터 의회의 권위에 빚지고 있지 않다. 옛 식민지의 법체계는
이제 ‘지역적 뿌리(local root)’를 가지게 되는데, 이는 법적 유효성의
궁극적 기준들을 특정하는 승인 규칙이 더 이상 다른 영토의 입법부가 제정한
것들(enactments)을 참조(refer)하지 않기 때문이다. 새로운 규칙은, 그것이
해당 지역 체계의 사법적 및 기타 공무담당 작용에서 그러한 규칙으로
수용되고 사용된다는 사실에 단순히 기초해 있다. 이 지역 체계의 규칙들에는
일반적으로 복종이 이루어진다. 따라서 비록 지방 입법부의 구성, 제정 방식,
구조가 여전히 원래 헌법에서 규정된 바와 동일할 수 있다 하더라도, 그
제정물들(enactments)이 현재 유효한 것은 \textbf{(p.~121)} 더 이상
웨스트민스터 의회의 유효한 제정법(statute)이 부여한 권한들의 행사이기
때문이 아니다. 그것들은, 지역적으로 수용된 승인 규칙에 따라, 지방
입법부의 제정이 유효성의 궁극적 기준이기 때문에 유효한 것이다.

이러한 발전은 여러 방식으로 달성될 수 있다. 모국 입법부는 일정 기간 동안
사실상 식민지의 동의 없이는 식민지에 대한 자신의 형식적 입법 권위(formal
legislative authority)를 행사하지 않다가, 마침내 해당 옛 식민지에 대한
입법 권한을 포기함으로써 무대에서 퇴장할 수 있다. 이때 주목할 점은,
웨스트민스터 의회가 이와 같이 스스로의 권한들을 불가역적으로 줄일(cut
down) 법적 권능(legal competence)을 가지고 있는지를 영국 법원이
인정할지에 대해 이론적인 의문이 존재한다는 것이다. 반대로, 이러한 단절은
폭력을 통해서만 이루어질 수도 있다. 그러나 어느 경우이든, 이 발전의
말미에는 두 개의 독립된 법체계가 존재하게 된다. 이것은 하나의 사실
진술(factual statement)이며, 비록 그것이 법체계의 존재에 관한 것이라
할지라도, 그렇다고 해서 덜 사실적인 것은 아니다. 이에 대한 주된 증거는,
이전 식민지에서는 현재 수용되고 사용되는 궁극적 승인 규칙이 더 이상 타
지역의 입법부 작용에 대한 어떠한 참조도 유효성 기준들 가운데 포함하지
않는다는 사실이다.

그런데 다시, 영연방(Commonwealth)의 역사에서 흥미로운 사례들이
보여주듯이, 사실상 식민지의 법체계가 이제 모국으로부터 독립되었음에도
불구하고, 모국의 법체계는 이 사실을 인정하지 않을 수도 있다.
웨스트민스터 의회가 식민지에 대해 입법할 권한을 보유하고 있거나 법적으로
되찾을 수 있다는 것이 여전히 영국법의 일부일 수 있으며, 실제로 그러한
경우에, 만약 웨스트민스터 제정법(Westminster statute)과 현지 입법부의 법
하나가 충돌하는 사건이 영국 국내 법원에 회부된다면, 영국 법원은 이
사태에 대한 이러한 관점에 효력을 부여할 수도 있다(give effect). 이 경우,
영국법명제들(propositions of English law)은 사실과 충돌하는 듯 보인다.
식민지의 법은 사실상 독립된 법체계이며, 자체의 지역적 궁극적 승인 규칙을
가진 체계인데, 영국 법원은 이를 사실상 그러한 바로 그것으로 인정하지
\emphksans{않는다(not)}. 사실상은 두 개의 법체계가 존재하지만, 영국법은 오직
하나만 존재한다고 고집한다(insist). 그러나, 한쪽 단언은 사실
진술(statement of fact)이고 다른 쪽은 (영국)법명제(proposition of
law)이기 때문에, 이 둘은 논리적으로 충돌하지 않는다. 이 입장을 명확히
하기 위해, 우리는 사실 진술이 참이며(true), 영국법명제는 “영국법상
올바른(correct in English law)” 것이라고 말할 수 있다.
\textbf{(p.~122)} 이러한 독립된 법체계의 존재 여부에 대한 사실적
단언(또는 부정)과 법체계의 존재에 관한 법명제들 간의 구별은, 국제공법과
국내법 사이의 관계를 논할 때에도 반드시 유념해야 한다. 어떤 매우 기이한
이론들이 그럴듯해 보이는 유일한 이유는 이 구별을 무시하기 때문이다.

이러한 법체계의 병리학(pathology)과 발생학(embryology)에 대한 거친
개관을 마무리하려면, 이제 우리는, “하나의 법체계가 존재한다”는
무조건적 단언(assertion)에 의해 그 일치(congruence)가 단언되는 정상
조건들이 부분적으로 실패하는 다른 형태들도 주목해야 한다. 법체계
내부에서 법에 대한 내부 진술이 행해질 때 일반적으로 전제되는,
공무담당자들 사이의 통일성(unity)은 부분적으로 붕괴될 수 있다. 헌법적
문제들---그리고 오직 그 문제들에 한해서---에 대해, 공무담당 세계
내부에서 분열이 발생하고 궁극적으로 사법부 내부의 분열로 이어질 수 있다.
법을 식별하는 데 사용되는 궁극적 기준들에 대한 이러한 분열의 시작은
1954년 남아프리카 공화국에서 발생한 헌법적 위기에서 볼 수 있었으며, 이는
\emph{Harris} 대 \emph{Dönges} 사건으로 법원에 회부되었다. 여기서
입법부는 자신이 보유한 법적 권능(legal competence)과 권한들(powers)에
대해 법원과는 다른 견해를 취하며 입법 조치들을 제정하였고, 법원은 그
조치들을 유효하지 않다고(invalid) 선언하였다. 이에 대한 대응으로
입법부는, 입법부의 제정물들(enactments)을 유효하지 않게 한 통상 법원들의
판결들에 대해 항소가 제기될 수 있도록 특별 항소 ‘법원’을 창설하였다. 이
법원은 실제로 그러한 항소를 심리하고 통상 법원들의 판결들을 번복했으며,
다시 통상 법원들은 그 특별 법원들을 창설한 입법을 유효하지 않다고
선언하고 해당 판결들을 법적으로 무효(nullity)라고 판시하였다. 이 절차가
중단되지 않았더라면(정부가 자신의 뜻을 관철하는 이 수단을 계속 추구하는
것이 지혜롭지 못하다고 판단했기 때문에 중단되었다), 우리는 입법부의
권능(competence)과 그에 따라 유효한 법의 기준들에 대한 두 견해 사이에서
무한 반복되는 진동 현상을 보았을 것이다. 법체계의 승인 규칙을 식별하는
것이 오직 그 조건 아래에서만 가능한 공무담당자들 사이의---특히
사법적---조화라는 정상 조건은 일시적으로 정지되었을 것이다. 그럼에도
불구하고, 이 헌법적 문제와 직접적으로 관련되지 않은 대부분의 법적
작용들은 이전과 마찬가지로 계속되었을 것이다. 인구가 분열되고, ‘법과
질서’가 \textbf{(p.~123)} 붕괴되기 전까지는, 원래의 법체계가 소멸했다고
말하는 것은 오해를 불러올 수 있다. 왜냐하면 ‘동일한 법체계(same legal
system)’라는 표현은 너무 넓고 유연하기 때문에, 유효한 법의
\emphksans{모든(all)} 원래 기준들에 대한 공무담당자들의 통일적 합의가
유지되는 것을 그 법체계가 ‘동일한 것’으로 남아 있기 위한 필요조건으로
삼을 수는 없기 때문이다. 우리가 할 수 있는 일은 지금까지 기술해온
방식대로 해당 상황을 설명하고, 그것을 하나의 표준 이하(substandard),
비정상적인 사례로 기록해 두는 것이다. 이 사례는 그 속에 법체계가
해체될(dissolve) 수 있는 잠재적 위협을 품고 있다.

이 마지막 사례는, 우리가 다음 장에서 논의하게 될 보다 넓은 주제의 경계로
우리를 이끈다. 이는 법체계의 궁극적 유효성 기준이라는 중대한 헌법적
사안(high constitutional matter)과, 그 법체계의 ‘일상적(ordinary)’ 법
양자 모두와 관련된다. 모든 규칙에는 특정 사례들을 일반적 용어의 사례로
식별하거나 분류하는 과정이 포함되며, 우리가 어떤 것을 규칙(rule)이라고
부를 준비가 되어 있는 모든 경우에 있어서, 해당 규칙이 분명히 적용되는
명확한 중심 사례들과, 그 적용 여부를 단언할 이유와 부정할 이유가 모두
존재하는 경계 사례들을 구별하는 것이 가능하다. 우리가 구체적인 상황들을
일반 규칙 아래로 포섭할 때, 이와 같은 확실성의 핵심(core of certainty)과
의문의 반음영(penumbra of doubt)이라는 이중성은 결코 제거될 수 없다.
이로 인해 모든 규칙에는 모호함의 가장자리(fringe), 곧 ‘개방적 직조(open
texture)’가 부여되며, 이는 법을 식별하는 데 사용되는 궁극적 기준들을
명시하는 승인 규칙(rule of recognition)에도, 하나의 특정한
제정법(statute)에도 마찬가지로 영향을 미칠 수 있다. 법의 이러한 측면은
흔히, 규칙이라는 틀로 법 개념을 해명하려는 시도가 오해를 낳을 수 있다는
점을 보여주는 것으로 여겨진다. 상황의 현실에도 불구하고 이를 고수하는
것은 흔히 ‘개념주의(conceptualism)’나 ‘형식주의(formalism)’라는 낙인이
찍히곤 한다. 우리는 이제 이와 같은 비판을 평가하는 논의로 넘어가게 될
것이다.

\newpage

\subsection{제6장 주석}\label{uxc81c6uxc7a5-uxc8fcuxc11d}


{\footnotesize
\noteentry 100쪽. \notetitle{승인 규칙(rule of recognition)과 켈젠(Kelsen)의 ‘기초규범(basic norm)’.} 본서의 중심 논제 중 하나는, 법체계의 토대들은
법적으로 무제한인 주권자에 대한 일반적 복종 습관(habit of obedience)에
있는 것이 아니라, 그 체계의 유효한 규칙들을 식별하는 권위 있는 기준들을
제공하는 \emphksans{궁극적 승인 규칙(ultimate rule of recognition)}에 있다는
것이다. 이 논제는 켈젠의 \emphksans{기초규범(basic norm)} 구상과 일정 부분
유사하며, 더 나아가서는 셀몬드(Salmond)의 충분히 전개되지 않은 ‘궁극적
법 원칙들(ultimate legal principles)’ 구상과도 가깝다(켈젠,
\emph{General Theory}, pp.~110--124, 131--134, 369--373, 395--396 참조;
셀몬드, \emph{Jurisprudence}, 11판, p.~137 및 부록 I 참조). 그러나
본서는 켈젠과는 다른 용어 체계를 채택하였으며, 그 이유는 아래의 주요한
점들에서 본서의 입장이 켈젠과 다르기 때문이다.

\begin{enumerate}
\def\labelenumi{\arabic{enumi}.}
\item
  승인 규칙이 존재하는지, 또 그것의 내용---즉 주어진 법체계에서의 유효성
  기준들(criteria of validity)---이 무엇인지는, 본서 전반에서
  경험적이면서도 복합적인 \emphksans{사실 문제(question of fact)}로 간주된다.
  이는 다음의 점이 사실이더라도 마찬가지이다: 즉, 법체계 안에서 활동하는
  법률가가 어떤 특정 규칙이 유효하다고 단언할(assert) 때, 그는 그 특정
  규칙의 유효성을 검사하기 위해 참조한 승인 규칙이 그 체계의 수용된 승인
  규칙으로서 존재한다는 사실을 \emphksans{명시적으로 진술하지는
  않지만(explicitly state)} \emphksans{암묵적으로 전제한다(tacitly
  presupposes)}. 만약 이 점에 대해 이의가 제기된다면, 이처럼
  전제되었지만 진술되지 않은 것은 사실들에 호소함으로써, 즉 적용해야 할
  법을 식별할 때 그 체계의 법원들과 공무담당자들이 보이는 실제 실천에
  호소함으로써 확립될 수 있다. 켈젠이 기초규범을 ‘법리학적 가설(juristic
  hypothesis)’(\emph{ibid.} xv), ‘가설적(hypothetical)’(\emph{ibid.}
  396), ‘상정된 궁극 규칙(postulated ultimate rule)’(\emph{ibid.} 113),
  ‘법리학적 의식 내에 존재하는 규칙(rule existing in the juristic
  consciousness)’(\emph{ibid.} 116), ‘가정(an assumption)’(\emph{ibid.}
  396) 등으로 분류하는 용어는, 이 책에서 강조되는 요점, 곧 “어떤
  법체계에서 법적 유효성의 기준들이 무엇인가”라는 물음이 사실 문제라는
  요점을 흐리게 하며, 실제로 그 요점과 양립하지 않을 수도 있다. 그것은
  사실 문제이지만, 하나의 규칙의 존재와 내용에 \emphksans{관한(about)}
  문제이다. Cf. 아고(Ago), “Positive Law and International Law,”
  \emph{American Journal of International Law} 제51권 (1957),
  pp.~703--707.
\item
  켈젠은 기초규범의 \emphksans{유효성(validity)}을 전제한다고 말한다. 그러나
  본문(pp.~108--110)에서 제시된 이유에 따르면, 일반적으로 수용된 승인
  규칙에 관해 그 존재 여부와는 별도로 그 유효성 또는
  비유효성(invalidity)을 묻는 것은 부적절하며, 그런 질문 자체가 성립할
  수 없다.
\item
  켈젠의 기초규범은 어떤 의미에서는 언제나 동일한 내용을 갖는다.
  왜냐하면 그것은 모든 법체계에서 단지 “헌법 또는 ‘최초의 헌법을 정한
  자들’에게 복종해야 한다”는 규칙이기 때문이다(\emph{General Theory},
  pp.~115--116). 그러나 이러한 일양성과 단순성의 외관은 오해를 낳을 수
  있다. 만약 하나의 헌법이 법의 다양한 원천들(sources of law)을
  명시하며, 그 법체계 내의 법원들과 공무담당자들이 실제로 그 헌법이
  제공하는 기준들에 따라 법을 식별한다는 의미에서 살아 있는 현실(living
  reality)이라면, 그러한 헌법은 수용된 것이며 실제로 존재하는 것이다.
  이런 경우, 헌법 또는 그것을 ‘정한’ 자들에게 복종해야 한다는 별도의
  규칙이 존재한다고 말하는 것은 불필요한 중복처럼 보인다. 특히, 영국처럼
  성문 헌법이 존재하지 않는 경우에는 더욱 그렇다. 이 경우, 법을 식별할
  때 특정한 유효성 기준들, 예컨대 의회 내 여왕에 의한 제정이 사용되어야
  한다는 규칙에 덧붙여 ‘헌법에 복종해야 한다’는 규칙이 자리할 곳은 없어
  보인다. 이 규칙 자체가 수용된 규칙이며, 이 규칙에 복종해야 한다는
  규칙을 말하는 것은 신비화할 뿐이다.
\item
  켈젠의 견해에 따르면(\emph{General Theory}, pp.~373--375, 408--410),
  특정한 법규칙을 유효한 것으로 간주하면서 동시에 그 법규칙이 요구하는
  행위를 금지하는 도덕규칙을 도덕적으로 구속력 있는 것으로 수용하는 것은
  논리적으로 불가능하다. 그러나 본서에서 제시한 법적 유효성에 대한
  설명에서는 그러한 결과가 전혀 도출되지 않는다. ‘\emphksans{기초(basic)}
  규범’이라는 표현 대신 ‘승인 규칙(rule of recognition)’이라는 표현을
  사용하는 이유 중 하나는, 법과 도덕 사이의 충돌에 관한 켈젠의 견해를
  따를 어떤 입장도 전제하지 않기 위함이다.
\end{enumerate}

\noteentry 101쪽. \notetitle{법의 원천들(sources of law).} 일부 저자들은 법들의
‘형식적(formal)’ 또는 ‘법적(legal)’ 원천들과 ‘역사적(historical)’ 또는
‘실질적(material)’ 원천들을 구분한다(Salmond, \emph{Jurisprudence},
11판, 제5장). Allen은 이 구분을 비판하지만(\emph{Law in the Making},
6판, p.~260), 이를 ‘원천(source)’이라는 단어의 두 가지 의미의 구분으로
해석하면 이 구분은 중요한 것이다(Kelsen, \emph{General Theory},
pp.~131--132, 152--153 참조). 한 의미에서(즉 ‘실질적(material)’,
‘역사적(historical)’ 의미에서) 원천이란, 특정한 시간과 장소에서 주어진
법규칙이 존재하게 된 원인적 또는 역사적 영향력들을 의미한다. 이 의미에서
현대 영국법의 특정 규칙들의 원천은 로마법 규칙일 수도 있고, 교회법
규칙이나 심지어 대중 도덕 규칙일 수도 있다. 그러나 ‘제정법(statute)’이
법의 원천이라고 말할 때, ‘원천(source)’이라는 단어는 단순한
역사적·인과적 영향력이 아니라, 문제의 법체계에서 수용된 법적 유효성의
기준들(criteria of legal validity) 중 하나를 가리킨다. 권능을 가진
입법부(competent legislature)에 의해 제정법으로 제정되었다는 점은,
주어진 제정법 규칙(statutory rule)이 단순히 존재하게 된
\emphksans{원인(cause)}이 아니라 그것이 유효한 법(valid law)인
\emphksans{이유(reason)}이다. 주어진 법규칙의 역사적 원인과 그 유효성의 이유
사이의 이러한 구분은, 입법부에 의한 제정, 관습적 실천, 선례 등을 유효한
법의 식별 표지들(identifying marks)로 수용하는 승인 규칙(rule of
recognition)이 체계 내에 존재할 때만 가능하다.

그러나 실제 실천에서는 역사적 또는 인과적 원천들과 법적 또는 형식적
원천들 간의 이 분명한 구분이 흐려질 수 있으며, Allen과 같은 저자들이 이
구분을 비판하게 된 것도 바로 이 점 때문이다. 어떤 체계에서
제정법(statute)이 법의 형식적 또는 법적 원천인 경우, 법원은 사건을
결정하면서 관련 제정법에 주의를 기울여야 할 구속을 받지만, 물론 제정법
문언의 의미를 해석하는 데 상당한 자유가 남겨져 있기는 하다(제7장 제1절
참조). 그러나 때때로 판사에게 남겨지는 것은 해석의 자유보다 훨씬 더
크다. 즉, 자기 앞의 사건을 결정하는 제정법이나 법의 다른 형식적 원천이
없다고 판단할 경우, 그는 예컨대 학설휘찬(Digest)의 한 문헌이나 프랑스
법학자의 저술에 자기 결정을 근거지을 수도 있다(예: Allen, \emph{op.
cit.}, p.~260 이하 참조). 해당 법체계가 그러한 원천들을 사용하라고
그에게 \emphksans{요구하는(require)} 것은 아니지만, 그가 그렇게 하는 것은
완전히 적절한 것으로 수용된다. 그러므로 이러한 저술들은 단순한 역사적
또는 인과적 영향력 이상이다. 그러한 저술들은 결정들을 위한 ‘좋은
이유들(good reasons)’로 인정되기 때문이다. 아마도 우리는 이러한
원천들을, 제정법과 같은 ‘구속적(mandatory)’ 법적 또는 형식적 원천들과
역사적 또는 실질적 원천들 양자와 구별하기 위해, ‘허용적(permissive) 법적
원천들’이라고 부를 수 있을 것이다.

\noteentry 103쪽. \notetitle{법적 유효성과 실효성(Legal validity and efficacy).} 켈젠은
전체적으로 실효적인 법질서의 실효성과 개별 규범의 실효성을
구분한다(\emph{General Theory}, pp.~41--2, 118--22). 그에게 있어 하나의
규범이 유효하다는 것은, 만약 그리고 오직 그 경우에만 (if, and only if),
그것이 \emphksans{전반적으로(on the whole)} 실효적인 체계에 속할 때이다. 그는
이 견해를 좀 더 난해하게 다음과 같이 표현하기도 한다. 즉, 체계 전체의
실효성은 그 체계의 규칙 유효성에 대한 \emphksans{필요조건(conditio sine qua
non)}이지, \emphksans{충분조건(conditio per quam)}은 아니라고
말한다(\emph{sed quaere}). 이 구분의 요점은 본서의 용어로 다음과 같이
정리될 수 있다. 법체계의 일반적인 실효성은 승인 규칙(rule of
recognition)이 제공하는 유효성 기준은 아니지만, 체계 내의 어떤 규칙이 그
유효성 기준들을 참조하여 그 체계의 유효한 규칙으로 식별될 때마다 언제나
암묵적으로 전제되며 명시적으로 진술되지는 않는다. 그리고 체계가
일반적으로 실효성을 갖추고 있지 않다면, 유효성에 관한 의미 있는 진술
자체가 불가능해진다. 본문의 입장은 이 점에서 켈젠과 다르다. 여기서는
체계의 실효성이 유효성 진술이 이루어지는 \emphksans{정상적 맥락(normal
context)}이긴 하지만, 특별한 상황에서는 그 체계가 더 이상 실효적이지
않더라도 유효성 진술이 여전히 의미를 가질 수 있다고 논변되기
때문이다(\emph{ante}, p.~104 참조).

켈젠은 또한 \emphksans{불사용에 의한 폐지(desuetudo)} 항목에서, 하나의
법체계가 특정 규칙의 유효성을 그 규칙의 지속적인 실효성에 의존하도록 할
가능성에 대해서도 논의한다. 이러한 경우, (개별 규칙의) 실효성은 체계의
유효성 기준 중 하나가 되며, 단순한 ‘전제(presupposition)’는
아니다(\emph{op. cit.}, pp.~119--22).

\noteentry 104쪽. \notetitle{유효성과 예측(Validity and prediction).} 법이 유효하다는
진술은 미래의 사법적 행위와 그에 따른 특유의 동기 부여 감정에 대한
예측이라는 견해에 대해서는 Ross, \emph{On Law and Justice}, 제1장 및
제2장 참조. 이에 대한 비판은 Hart, `Scandinavian Realism'
(\emph{Cambridge Law Journal}, 1959년)을 보라.

\noteentry 106쪽. \notetitle{개정권이 제한된 헌법(Constitutions with limited amending powers).} 서독과 터키의 사례는 제4장 주석(\emph{ante}, p.~290)을
참조하라.

\noteentry 111쪽. \notetitle{관행적 범주와 헌정 구조(Conventional categories and constitutional structures).} ‘법(law)’과 ‘관행(convention)’으로의 전부를
포괄한다고 주장되는 구분에 대해서는 Dicey, \emph{Law of the
Constitution}, 10판, pp.~23 이하; Wheare, \emph{Modern Constitutions},
제1장 참조.

\noteentry 111쪽. \notetitle{승인 규칙(rule of recognition): 법인가 사실인가?} 정치적
사실로 분류할 것인지에 대한 찬반 논증은 Wade, `The Basis of Legal
Sovereignty', \emph{Cambridge Law Journal} (1955), 특히 p.~189, 그리고
Marshall, \emph{Parliamentary Sovereignty and the Commonwealth},
pp.~43--46을 참조하라.

\noteentry 112쪽. \notetitle{법체계의 존재, 습관적 복종, 승인 규칙의 수용(The existence of a legal system, habitual obedience, and the acceptance of the rule of recognition).} 일반 시민의 복종과 공무담당자들 쪽에서의 헌정 규칙 수용을
모두 포함하는 복합적 사회 현상을 지나치게 단순화할 위험에 대해서는 제4장
제1절, pp.~60--61 및 Hughes, `The Existence of a Legal System', \emph{35
New York University LR} (1960), p.~1010 참조. Hughes는 이 점에서 Hart가
`Legal and Moral Obligation' (\emph{Essays in Moral Philosophy}, Melden
편, 1958)에서 사용한 용어를 정당하게 비판하고 있다.

\noteentry 118쪽. \notetitle{법질서의 부분적 붕괴(Partial breakdown of legal order).} 법체계가 완전히 정상적으로 존재하는 상태와 전혀 존재하지 않는 상태
사이의 많은 가능한 중간 상태들 중 극히 일부만이 본문에서 언급되었다.
혁명(revolution)에 대한 법적 관점에서의 논의는 Kelsen, \emph{General
Theory}, pp.~117 ff., 219 ff.를 참조하고, Cattaneo, \emph{Il Concetto di
Revoluzione nella Scienza del Diritto} (1960)에서는 상세히 다루어진다.
적군 점령(enemy occupation)에 의한 법체계의 중단은 다양한 형태를 띨 수
있으며, 이 중 일부는 국제법상으로 범주화되어 있다. 이에 대해서는 McNair,
`Municipal Effects of Belligerent Occupation', \emph{57 LQR} (1941),
그리고 Goodhart, `An Apology for Jurisprudence' (\emph{Interpretations
of Modern Legal Philosophies}, pp.~288 ff.)의 이론적 논의를 참조하라.

\noteentry 120쪽. \notetitle{법체계의 발생학(The embryology of a legal system).} Wheare,
\emph{The Statute of Westminster and Dominion Status,} 5판에서 추적되는
식민지에서 도미니언으로의 발전 과정은 법이론 연구에 보람 있는 분야이다.
또한 Latham, \emph{The Law and the Commonwealth} (1949)을 참조하라.
Latham은 영연방(Commonwealth)의 헌정 발전을 ‘지역적 뿌리(local root)’를
가진 새로운 기초규범의 성장이라는 관점에서 최초로 해석하였다. 또한
Marshall, 앞서 인용한 저서, 특히 캐나다에 대한 제7장을 참고하고, Wheare,
\emph{The Constitutional Structure of the Commonwealth} (1960), 제4장
‘자생성(Autochthony)’도 참고하라.

\noteentry 121쪽. \notetitle{입법 권한의 포기(Renunciation of legislative power).} 웨스트민스터 제정법(Statute of Westminster) 제4조의 법적 효력에 대한
논의는 Wheare, \emph{The Statute of Westminster and Dominion Status},
5th edn., pp.~297--8을 참조하라. 또한 \emph{British Coal Corporation} v.
\emph{The King} (1935), AC 500; Dixon, `The Law and the Constitution',
\emph{51 LQR} (1935); Marshall, 앞서 인용한 저서, pp.~146 ff.; 그리고
제7장 제4절도 함께 보라.

\noteentry 121쪽. \notetitle{모국 체계에 의해 인정되지 않은 독립(Independence not recognized by the parent system).} 아일랜드 자유국(Irish Free State)에
대한 논의는 Wheare, 앞서 인용한 저서 참조; \emph{Moore} v. \emph{AG for
the Irish Free State} (1935), AC 484; \emph{Ryan} v. \emph{Lennon}
(1935), IRR 170 참조.

\noteentry 121쪽. \notetitle{법체계의 존재에 관한 사실적 단언들과 법 진술들(Factual assertions and statements of law concerning the existence of a legal system).} Kelsen은 국내법과 국제법 사이의 가능한 관계들(‘국가법의 우위성
또는 국제법의 우위성’)에 대한 자신의 설명(op. cit., pp.~373--83)에서,
법체계가 존재한다는 진술은 반드시 한 법체계의 관점에서 다른 법체계에
대해, 그 다른 체계를 ‘유효한(valid)’ 것으로 그리고 자기 자신과 하나의
단일 체계를 형성하는 것으로 수용하면서 행해지는 법 진술(statement of
law)이어야 한다고 전제한다. 이에 반해 상식적 견해에서는 국내법과
국제법이 별개의 법체계들을 구성한다고 보며, 어떤 법체계(국가적이든
국제적이든)가 존재한다는 진술을 사실 진술(statement of fact)로 취급한다.
Kelsen에게는 이러한 견해가 용납될 수 없는
‘다원주의(pluralism)’이다(Kelsen, loc. cit.; Jones, `The “Pure” Theory
of International Law', \emph{16 BYBIL} 1935 참조). 또한 Hart, `Kelsen's
Doctrine of the Unity of Law', \emph{Ethics and Social Justice},
\emph{Contemporary Philosophical Thought}, vol.~4 (New York, 1970) 참조.

\noteentry 122쪽. \notetitle{남아프리카(South Africa).} 남아프리카 헌정 위기에서 배울 수
있는 중요한 법리학적 교훈을 보다 심층적으로 검토하려면, Marshall, 앞서
인용한 저서, 제11장을 참조하라.

\subsection{제6장 제3판
주석}\label{uxc81c6uxc7a5-uxc81c3uxd310-uxc8fcuxc11d}

\noteentry 100--1쪽. \notetitle{법의 원천들(Sources of law).} 법이
\emphksans{원천들(sources)}을 가진다는 관념의 이론적 중요성에 대해서는 Joseph
Raz, \emph{The Authority of Law}, chap.~3을 보라. 보통법(common law)과
관습(custom)을 법의 원천들로 다룬 논의는 Gerald J. Postema,
\emph{Bentham and the Common Law Tradition} (Oxford University Press,
1986) 및 John Gardner, `Some Types of Law', 그의 저서 \emph{Law as a
Leap of Faith} 제3장을 참조하라.

Hart의 승인 규칙(rule of recognition) 구상(conception)에 대한 비판은
Joseph Raz, \emph{The Concept of a Legal System}, chap.~8, 특히
pp.~197--200; Joseph Raz, \emph{The Authority of Law}, chap.~5, 특히
pp.~90--97; Ronald Dworkin, \emph{Taking Rights Seriously}, pp.~39--45,
62--68을 참조하라. 승인 규칙을 의미론적 교설로 해석한 그의 견해는 Ronald
Dworkin, \emph{Law's Empire}, pp.~31--35에서 볼 수 있다.

\noteentry 101--3쪽. \notetitle{승인 규칙과 법원들의 실천(The rule of recognition and the practice of the courts).} Hart에게서 궁극적 승인 규칙들(ultimate
rules of recognition)은 공무담당자들의 관습적 실천들 속에 내재한다.
그런데 공무담당자란 법에 의해 그렇게 식별되는 사람이라면, 순환성이
다가오는 것처럼 보일 수 있다. 즉, 어떤 것이 법인 것은 오직 그것이 승인
규칙에 의해 식별되는 경우뿐이고, 어떤 것이 승인 규칙인 것은 오직 그것이
공무담당자들에 의해 실천되는 경우뿐이며, 어떤 사람이 공무담당자인 것은
오직 그가 법에 의해 그렇게 권한을 부여받은 경우뿐이라는 식이다. 이
문제의 해소 방안들(resolutions)에 관해서는 Neil MacCormick, \emph{H. L.
A. Hart}, pp.~108--120; Michael Bayles, \emph{Hart's Legal Philosophy},
pp.~81--83을 참조하라.

\noteentry 103--4쪽. \notetitle{법적 유효성(Legal validity).} Hart는
‘유효성(validity)’을 주로 체계 소속(system membership)의 문제로 다룬다.
이에 대비되는 견해로는 Kelsen, \emph{Pure Theory of Law}, pp.~10--15,
193--195가 있다. 그는 “‘유효한(valid)’은 그것이 구속력을 가진다는 것,
곧 한 개인이 그 규범이 정한 방식으로 행위해야 한다는 것을
의미한다”(p.~193)고 말한다. 법적 맥락들에서 ‘유효성(validity)’의 다양한
의미들(senses)에 대해서는 J. W. Harris, \emph{Law and Legal Science: An
Inquiry into the Concepts Legal Rule and Legal System} (Oxford
University Press, 1979), chap.~4를 참조하라. Hart의 유효성 설명에 대한
평가는 Joseph Raz, \emph{The Authority of Law}, chap.~8에 실려 있다.
유효성의 정도들(degrees of validity)이라는 가능성에 대해서는 John
Finnis, \emph{Natural Law and Natural Rights}, pp.~276--281을 보라.

\noteentry 103--5쪽. \notetitle{법의 실효성(The efficacy of law).} Hart의 실효성 정의는
의무부과 규칙들(mandatory rules)에만 적용된다(“특정한 행위를 요구하는
법규칙이 대체로 복종된다는 사실”: p.~103). 그렇다면 허용적
규칙들(permissive rules) 혹은 권한-부여 규칙들(power-conferring rules)의
실효성은 어떻게 되는가? 이에 대해서는 Joseph Raz, \emph{The Concept of a
Legal System}, chap.~9; \emph{The Authority of Law}, chap.~5, 특히
pp.~85--90을 참조하라. 또한 Gerald Postema, `Conformity, Custom, and
Congruence: Rethinking the Efficacy of Law', Matthew H. Kramer \emph{et
al}. 편, \emph{The Legacy of H. L. A. Hart: Legal, Political, and Moral
Philosophy} (Oxford University Press, 2008)도 함께 보라.

\noteentry 105--10쪽. \notetitle{궁극적이고 최고인 승인 규칙(The rule of recognition as ultimate and supreme).} Hart는 ‘승인 규칙(rule of recognition)’이라는
용어를 한 법체계 안의 다른 모든 규칙들의 유효성을 평가할 기준들을
제공하지만, 그 자체는 다른 어떤 규칙에 의해서도 유효화되지(validated)
않는 하나의 관습적 사회 규칙을 지칭하는 데 사용한다. 다음 사항에
유의하라.

\begin{enumerate}
\def\labelenumi{(\arabic{enumi})}
\item
  승인 규칙은 판사들과 기타 공무담당자들의 관습적 규칙이지, 법으로
  제정된 것이 아니며 실정법(positive law)이 아니다. 특히, 그것은 형식적
  헌법(formal constitution)이나 그 일부가 아니다. (그러한 헌법은
  \emphksans{그 자체로} 승인 규칙에 의해 유효화되어야(validated) 한다.) 이
  점에 대해 혼란이 매우 많다. 이에 대한 논의로는 Kent Greenawalt, `The
  Rule of Recognition and the Constitution' (1987) 85 \emph{Michigan Law
  Review} 621; Matthew Adler and Kenneth Himma 편, \emph{The Rule of
  Recognition and the U.S. Constitution} (Oxford University Press, 2009)
  수록 논문들; John Gardner, `Can there be a Written Constitution',
  Leslie Green and Brian Leiter 편, \emph{Oxford Studies in Philosophy
  of Law,} vol.~I (Oxford University Press, 2011) 및 그의 저서 \emph{Law
  as a Leap of Faith} (Oxford University Press, 2012) 제4장을 참조하라.
\item
  Hart는 승인 규칙이 “그것이 그 집단의 규칙임을 보여 주는 확정적 인정
  표지(conclusive affirmative indication)를 제공하는 특징들”을
  명시한다고 말하고(94쪽), 그것이 “일차 규칙들을 확정적으로 식별하기
  위한 규칙”(95쪽)이라고 말하지만, 이를 승인 규칙의 기준들이 완전하고
  확정적이라는 의미로 이해해서는 안 된다. 그는 그것을 명시적으로
  부정한다(147--154쪽, 257--259쪽 참조).
\item
  Hart는 각 법체계에 정확히 하나의 승인 규칙이 존재한다고 아무런 논증
  없이 전제한다. 그러나 이는 의문의 여지가 있다. Joseph Raz, \emph{The
  Authority of Law}, 95--96쪽 참조.
\end{enumerate}

\noteentry 114--16쪽. \notetitle{공무담당자들과 이차 규칙들(Officials and secondary rules).} 흔히 두 가지 서로 다른 질문이 융합된다: (a) 누구의 행위가 승인
규칙을 \emphksans{구성하는가}? (b) 누구의 행위가 승인 규칙에 의해
\emphksans{규율되는가}? (a)에 답할 때 Hart는 때로
‘공무담당자들(officials)’(115, 116쪽)을, 때로는 ‘판사들(judges)’(108,
116쪽, 256쪽 참조)을 언급한다. Joseph Raz는 판사와 같은 법 적용
공무담당자들이 특히 중요하다고 논변한다: \emph{The Authority of Law},
chap.~6 참조. (b)에 관해서는 의문의 여지가 없다. 승인 규칙은 모든
공무담당자를 구속한다. 공무담당자들의 우위성, 그리고 법 적용
공무담당자들의 우위성은 정치적 도덕의 교설이 아니라는 점에 유의하라.
법원들이 실제로 그러한 영향력을 갖는 것이 유감스러울 수 있다는 점도
Hart가 논변하는 모든 것과 양립한다. 이에 대한 정치적 문제는 예컨대 Mark
Tushnet, \emph{Taking the Constitution Away from the Courts} (Princeton
University Press, 1999)를 보라.

\noteentry 116--17쪽. \notetitle{복종과 일반 시민(Obedience and the ordinary citizen).} 이 부분들은 Hart가 법을 사회 제도로 평가하는 맥락을 이해하는 데 핵심이
되는 대목들이다. 그는 법이 도덕적 위험들을 가져온다는 점뿐 아니라, 그
위험들이 법의 본성과 긴밀하게 연결되어 있다는 점을 논변한다. 이에 대한
논의로는 Jeremy Waldron, `All We Like Sheep' (1999) 12 \emph{Canadian
Journal of Law and Jurisprudence} 169 및 Leslie Green, `Positivism and
the Inseparability of Law and Morals' (2008) 83 \emph{New York
University Law Review} 1035, 특히 1052--1054쪽을 참조하라.

\noteentry 117--23쪽. \notetitle{법체계의 병리학(Pathology of a legal system).} John
Finnis, `Revolutions and Continuity of Law', 그의 저서 \emph{Philosophy
of Law: Collected Essays Volume IV} (Oxford University Press, 2011)의
제21장도 보라. 법체계들 사이의 상호작용 가능성들은 Neil MacCormick,
\emph{Questioning Sovereignty: Law, State and Nation in the European
Commonwealth} (Oxford University Press, 1999), chap.~7에서 탐구된다.
}

\setcounter{footnote}{0}
\section{VII. 형식주의(Formalism)와
규칙회의주의(Rule-Scepticism)}\label{vii.-uxd615uxc2dduxc8fcuxc758formalismuxc640-uxaddcuxce59uxd68cuxc758uxc8fcuxc758rule-scepticism}

\subsection{1. 법의 개방적 직조(The Open Texture of
Law)}\label{uxbc95uxc758-uxac1cuxbc29uxc801-uxc9c1uxc870the-open-texture-of-law}

어떤 대규모 집단에서든, 일반 규칙들, 기준들, 원칙들은 사회 통제의 주요
수단이어야 하며, 각 개인에게 별도로 내려지는 개별 지시들이 주요
수단이어서는 안 된다. 수많은 개인들이 추가 지시 없이도, 어떤 경우가
발생했을 때 자신들에게 일정한 행위가 요구된다는 것으로 이해할 수 있는
일반적 행위 기준들을 전달할 수 없다면, 우리가 현재 법으로 인정하는 것은
아무것도 존재할 수 없을 것이다. 따라서 법은 주로, 그러나 결코 전적으로는
아니게, 사람의 \emphksans{부류들(classes)}, 그리고 행위, 사물, 상황의
\emphksans{부류들(classes)}을 가리켜야 하며, 사회생활의 광범위한 영역에서
법이 성공적으로 작동하는 것은, 구체적인 행위들, 사물들, 상황들을 법이
만드는 일반적 분류들의 사례들로 식별할 수 있는 능력(capacity)이 널리
퍼져 있는 데 의존한다.

그러한 일반적 행위 기준들을, 그것들이 적용될 이어지는 경우들에 앞서
전달하기 위해 사용되어 온 두 가지 주요 장치가 있다. 처음 보기에는 서로
매우 달라 보이는 이 두 장치 가운데 하나는 일반적 분류어(classifying
words)를 최대한 사용하고, 다른 하나는 최소한 사용한다. 첫 번째 방식은
우리가 입법(legislation)이라고 부르는 것으로 대표되고, 두 번째 방식은
선례(precedent)로 대표된다. 우리는 다음의 단순한 비법적 사례들에서 이
둘의 구별되는 특징들을 볼 수 있다. 한 아버지는 교회에 가기 전에 아들에게
이렇게 말한다. “모든 남성과 소년은 교회에 들어갈 때 모자를 벗어야
한다.” 다른 아버지는 교회에 들어가면서 자기 모자를 벗으며 말한다.
“보아라. 이런 경우에는 이렇게 행동하는 것이 옳은 방식이다.”

예시에 의해 행위 기준들(standards of conduct)을 전달하거나 가르치는
것은, 우리의 단순한 사례보다 훨씬 더 정교한 여러 형태를 취할 수 있다.
만약 특정한 경우에 아이에게 교회에 들어갈 때 아버지가 한 행동을 옳은
행위의 본보기로 간주하라고 말하는 대신, 아버지가 \textbf{(p.~125)}
아이가 자신을 적절한 행위에 관한 권위로 간주하고 어떻게 행동해야
하는지를 배우기 위해 자신을 관찰할 것이라고 가정한다면(assumed), 우리의
사례는 법에서의 선례(precedent) 사용에 더 가까워질 것이다. 법에서의 선례
사용에 더 가까이 접근하려면, 아버지가 자기 자신과 타인들에 의해 전통적인
행위 기준들을 지지하는 사람으로 이해되고 있으며 새로운 기준들을 도입하는
것이 아니라고 상정해야 한다.

예시에 의한 전달은 그 모든 형태에서, ‘내가 하는 대로 하라’와 같은
얼마간의 일반적 언어 지시가 동반되더라도, 의도된 것이 무엇인지에 관하여,
전달하려는 사람이 스스로 명확히 염두에 둔 사항들에 대해서조차,
가능성들의 범위와 따라서 의문의 범위를 열어둘 수 있다. 수행의 어느
정도까지가 모방되어야 하는가? 모자를 벗을 때 오른손 대신 왼손을 쓰는
것이 문제가 되는가? 느리게 하는가 재빠르게 하는가가 문제가 되는가?
모자를 좌석 아래에 두는 것이 문제가 되는가? 교회 안에서 모자를 다시 쓰지
않는 것이 문제가 되는가? 이 모든 것은 아이가 스스로에게 물을 수 있는
보다 일반적인 질문들의 변형들이다. ‘내 행위가 그의 행위와 어떤 방식으로
유사해야 옳은가?’ ‘그의 행위에서 정확히 무엇이 나의 지침이 되어야
하는가?’ 예시를 이해하면서 아이는 그 예시의 어떤 측면들에 주목하고 다른
측면들에는 주목하지 않는다. 그렇게 할 때 아이는 상식과, 어른들이
중요하다고 생각하는 사물들과 목적들의 일반적 종류에 관한 지식, 그리고 그
경우(교회에 가는 것)의 일반적 성격과 그에 적절한 행위 종류에 대한 이해에
의해 인도된다.

예시들의 불확정성들(indeterminacies)과 대조적으로, 명시적인 일반 언어
형식(예: “모든 남성은 교회에 들어갈 때 모자를 벗어야 한다”)에 의한
일반 기준들의 전달은 분명하고, 신뢰할 수 있으며, 확실한 것처럼 보인다.
행위의 일반적 지침들로 받아들여져야 할 특징들은 여기서 말로 식별된다.
그것들은 구체적인 예시 속에 다른 것들과 함께 묻혀 남겨져 있지 않고,
언어적으로 끄집어내어져 있다. 아이는 다른 경우들에서 무엇을 해야 할지를
알기 위해 더 이상 무엇이 의도되었는지, 또는 무엇이 승인될 것인지를
추측할 필요가 없다. 아이는 자기 행위가 옳기 위해서는 그 예시와 어떤
방식으로 유사해야 하는지를 추측해야 하는 처지에 남겨져 있지 않다. 대신
아이에게는 앞으로 무엇을 해야 하는지, 그리고 언제 그것을 해야 하는지를
골라낼 수 있는 언어적 기술(description)이 있다. 그는 명확한 언어
용어들의 사례들을 식별하고, 구체적 사실들을 일반적 분류 표제들 아래에
‘포섭(subsume)’하며, 단순한 삼단논법적 결론을 도출하기만 하면 된다.
\textbf{(p.~126)} 그는 자기 위험 부담 아래 선택하거나 추가적인 권위 있는
지침을 구해야 하는 선택지 앞에 놓여 있지 않다. 그는 자기 혼자서 자기에게
적용할 수 있는 규칙(rule)을 가지고 있다.

20세기의 법리학의 많은 부분은, 권위 있는 예시(선례(precedent))를 통한
전달의 불확실성과 권위 있는 일반 언어(입법(legislation))를 통한 전달의
확실성 사이의 구분이 이 순진한 대비가 시사하는 것만큼 확고하지 않다는
중요한 사실을 점차 인식하거나(때때로는 과장되게) 강조해온 과정이었다.
언어적으로 정식화된 일반 규칙이 사용될 때조차, 그 규칙이 요구하는 행위의
형태에 관한 불확실성은 특정한 구체적 사례에서 발생할 수 있다. 구체적인
사실 상황은 일반 규칙의 적용 여부가 문제되는 사안의 사례들로 미리 서로
구분되고 명명되어 있는 것이 아니며, 규칙 자체가 나서서 스스로의 사례라고
주장할 수도 없다. 규칙의 영역뿐 아니라 모든 경험의 영역에서는, 일반
언어가 제공할 수 있는 지침에는 언어의 본성에 내재한 한계가 존재한다.
기실 유사한 맥락에서 반복적으로 발생하며 일반적 표현이 분명히 적용가능한
명백한 사례들(예: “무엇이 탈 것(a vehicle)이라면, 자동차는 틀림없이 그
중 하나다(If anything is a vehicle a motor-car is one)”)이 존재하지만,
그 표현이 적용되는지 불분명한 사례들(예: “여기서 ‘탈 것(vehicle)’이라는
표현은 자전거(bicycles)나 비행기(airplanes), 롤러스케이트(roller
skates)도 포함하는가?”)도 있다. 후자의 경우는, 자연이나 인간의 발명에
의해 끊임없이 등장하며, 명백한 사례의 일부 특징만을 가지되 다른 특징은
결여한 사실 상황들이다. ‘해석’의 준칙들(canons of 'interpretation')은
이러한 불확실성을 줄일 수는 있어도 없앨 수는 없다. 왜냐하면 이 준칙들
역시 언어 사용에 관한 일반 규칙들이며, 그 자체가 해석을 요하는 일반
용어들을 사용하기 때문이다. 이러한 규칙들 역시 다른 규칙들과 마찬가지로
자기 자신의 해석을 제공할 수 없다. 해석이 필요 없어 보이며 사례들의
식별이 문제없거나 ‘자동적’으로 보이는 명백한 사례들은, 사실 유사한
맥락에서 반복적으로 나타나는 친숙한 것들일 뿐이다. 이러한 경우에는 분류
용어의 적용에 대해 판단이 일치하는 일반적 합의가 존재한다.

일반 용어들(general terms)은 그러한 친숙하고 일반적으로 다투어지지 않는
사례들이 없다면 의사소통의 매체로서 우리에게 무용할 것이다. 그러나
친숙한 것의 변형들(variants) 역시, \textbf{(p.~127)} 어느 특정 시점에
우리의 언어적 자원의 일부를 이루는 일반 용어들 아래에서의
분류(classification)를 요구한다. 여기서 의사소통에서 일종의 위기가
촉발된다. 일반 용어를 사용하는 데에는 찬성하는 이유들과 반대하는
이유들이 모두 있으며, 분류해야 하는 당사자에게 그 용어의 사용을
명하거나, 반대로 그 사용의 거부를 명하는 확고한 관행이나 일반적 합의는
없다. 이러한 경우에 의문들이 해소되어야 한다면, 그것들을 해소해야 할
자는 열린 대안들(open alternatives) 사이에서 일종의 선택을 해야만 한다.

이 지점에서 규칙이 표현되어 있는 권위 있는 일반 언어(authoritative
general language)는 권위 있는 예시가 그러하듯이 불확실한 방식으로만
인도할 수 있다. 이 지점에서는 규칙의 언어가 우리로 하여금 쉽게 식별
가능한 사례들을 단순히 골라낼 수 있게 해 줄 것이라는 느낌이 무너진다.
포섭(subsumption)과 삼단논법적 결론의 도출(drawing of a syllogistic
conclusion)은 무엇이 해야 할 옳은 일인지를 정하는 데 관여하는 추론의
핵심을 더 이상 특징짓지 못한다. 대신 규칙의 언어는 이제 권위 있는 예시
하나, 곧 명백한 사례(plain case)에 의해 구성되는 예시를 지목하는 것처럼
보일 뿐이다. 이것은 선례와 거의 같은 방식으로 사용될 수 있지만, 규칙의
언어는 선례보다 더 지속적이고 더 촘촘하게 주목을 요구하는 특징들을
제한한다. 공원에서 차량 사용을 금지하는 규칙이 그 적용 여부가
불확정적으로 보이는 어떤 상황들의 결합에 적용되는가라는 질문에 직면할
때, 답해야 하는 사람이 할 수 있는 전부는 (선례를 사용하는 사람이
그러하듯이) 현재 사례가 명백한 사례와 ‘관련된’ 측면들에서 ‘충분히’
유사한지를 고려하는 것이다. 언어가 그에게 남겨 두는 이러한
재량(discretion)은 매우 넓을 수 있다. 그래서 그가 그 규칙을 적용한다면,
그 결론은 자의적이거나 비합리적이지 않을 수 있음에도 불구하고 사실상
하나의 선택이다. 그는 법적으로 관련되고 충분히 가까운 것으로 합리적으로
옹호될 수 있는 유사성들 때문에, 하나의 새로운 사례를 사례들의 계열(line
of cases)에 추가하기로 선택한다. 법규칙들의 경우, 관련성(relevance)과
유사성의 가까움(closeness of resemblance)의 기준들은 법체계 전반을
관통하는 많은 복잡한 요소들과 그 규칙에 귀속될 수 있는 목적들 또는
취지(aims or purpose)에 의존한다. 이것들을 특징짓는 것은 법적 추론(legal
reasoning)에 특유하거나 고유한 것이 무엇이든 그것을 특징짓는 일이 될
것이다.

행위 기준들의 전달을 위해 선례(precedent)든 입법(legislation)이든 어느
장치가 선택되더라도, 그것들은 보통 사례들의 대다수에서는 아무리 매끄럽게
작동한다 하더라도, 그 적용이 문제되는 어떤 지점에서는
불확정적(indeterminate)임이 드러날 것이다. 그것들은 이른바 \emphksans{개방적
직조(open texture)}를 지닐 것이다. 지금까지 우리는 입법의 경우에 관하여
이것을 인간 언어의 일반적 특징으로 제시해 왔다. 경계선(borderline)에서의
불확실성은 사실의 문제들에 관한 어떤 형태의 의사소통에서든 일반적 분류
용어들을 사용하는 데 치러야 할 대가이다. 영어와 같은 자연언어는 그렇게
사용될 때 환원 불가능하게 개방적 직조를 지닌다. 그러나 실제 언어가
개방적 직조의 특성들을 가진다는 점에 의존한다는 이 사실과는 별도로, 특정
사례에 적용되는지 여부가 언제나 사전에 확정되어 있고 실제 적용 지점에서
열린 대안들 사이의 새로운 선택을 전혀 수반하지 않는, 그토록 상세한
규칙이라는 구상을---이상(ideal)으로서조차---왜 우리가 품어서는 안
되는지를 이해하는 것이 중요하다. 짧게 말해 그 이유는, 그러한 선택의
필요성이 우리가 신들(gods)이 아니라 인간들(men)이기 때문에 우리에게
떠맡겨지기 때문이다. 특정한 경우마다 추가적인 공무담당 지시(official
direction) 없이 사용될 일반 기준들에 의해, 어떤 행위 영역을 모호성 없이
사전에 규율하려 할 때마다, 우리가 서로 연결된 두 가지 불리한 조건
아래에서 애쓴다는 점은 인간적 처지(the human predicament)의 특징이며,
따라서 입법적 처지의 특징이다. 첫 번째 불리한 조건은 사실에 대한 우리의
상대적 무지(relative ignorance of fact)이고, 두 번째는 목적의 상대적
불확정성(relative indeterminacy of aim)이다. 만약 우리가 사는 세계가
유한한 수의 특징들로만 특징지어지고, 이 특징들과 그것들이 결합할 수 있는
모든 방식이 우리에게 알려져 있다면, 모든 가능성에 대해 사전에 대비할 수
있을 것이다. 우리는 특정 사례들에 대한 적용이 결코 추가적 선택을
요구하지 않는 규칙들을 만들 수 있을 것이다. 모든 것이 알려질 수 있을
것이고, 모든 것에 대해, 그것이 알려질 수 있기 때문에, 무언가가 행해질 수
있으며 규칙에 의해 사전에 특정될 수 있을 것이다. 이것은
‘기계적(mechanical)’ 법리학(jurisprudence)에 알맞은 세계일 것이다.

명백히 이 세계는 우리의 세계가 아니다. 인간 입법자들은 미래가 가져올 수
있는 모든 가능한 상황 조합들에 관해 그러한 지식을 가질 수 없다. 이러한
예견할 수 없음(inability to anticipate)은 목적의 상대적
불확정성(indeterminacy of aim)을 동반한다. 우리가 어떤 일반적 행위
규칙을 과감히 정식화할 때(예: 어떤 차량도 공원에 들어와서는 안 된다는
규칙), \textbf{(p.~129)} 이 맥락에서 사용되는 언어는 어떤 것이 그 범위
안에 들려면 충족해야 하는 필요조건들을 고정하며, 확실히 그 범위 안에
드는 것의 몇몇 명확한 예들이 우리의 마음에 떠오를 수 있다. 그것들은
전범적(paradigm) 명백한 사례들(clear cases), 곧 자동차, 버스,
오토바이이다. 그리고 우리의 입법 목적은 지금까지는
확정적(determinate)인데, 이는 우리가 일정한 선택을 했기 때문이다. 우리는
적어도 이러한 것들을 배제하는 대가를 치르더라도 공원의 평온과 조용함이
유지되어야 한다는 문제를 처음에 확정해 두었다. 다른 한편, 우리가 공원의
평온이라는 일반 목적을 처음에 예견하지 않았거나 어쩌면 예견할 수 없었던
사례들(예컨대 전기로 움직이는 장난감 자동차)과 결합시키기 전까지, 이
방향에서 우리의 목적은 불확정적(indeterminate)이다. 우리는 예견되지 않은
사례가 발생할 때 제기될 문제를, 우리가 예견하지 않았기 때문에, 아직
확정해 두지 않았다. 즉, 이러한 것들을 사용하는 것이 즐거움이나 이익이
되는 아이들을 위해 공원의 어느 정도 평온을 희생할 것인지, 아니면 그
아이들에 맞서 그 평온을 방어할 것인지의 문제이다. 예견되지 않은 사례가
실제로 발생하면, 우리는 걸려 있는 쟁점들에 직면하고, 그때 서로 경쟁하는
이익들 사이에서 우리에게 가장 만족스러운 방식으로 선택함으로써 그 문제를
확정할 수 있다. 그렇게 함으로써 우리는 우리의 최초 목적을 더 확정적으로
만들게 되며, 부수적으로는 이 규칙의 목적상 어떤 일반 단어의 의미에 관한
문제도 확정하게 된다.

서로 다른 법체계들이나 동일한 법체계라도 서로 다른 시기에는, 일반
규칙들을 특정 사례들에 적용할 때 추가적인 선택 행사가 필요하다는 점을
어느 정도로든, 그리고 더 많거나 더 적게 명시적으로 무시하거나 인정할 수
있다. 법이론에서 형식주의(formalism) 또는 개념주의(conceptualism)로
알려진 폐해는, 일반 규칙이 일단 설정되고 나면 그러한 선택의 필요성을
은폐하고 최소화하려는 언어로 정식화된 규칙들에 대한 태도에 있다. 이를
수행하는 한 가지 방식은 규칙의 의미를 동결하여, 그 일반 용어들이 적용이
문제되는 모든 경우에서 동일한 의미를 가져야 한다고 하는 것이다. 이를
확보하기 위해 우리는 명백한 사례(plain case)에 존재하는 특정 특징들을
붙잡고, 어떤 것이 그러한 특징들을 가지면 그것이 다른 어떤 특징들을 더
가지거나 결여하든, 그리고 이 방식으로 규칙을 적용할 때 사회적 결과가
무엇이든, 그것을 규칙의 범위 안에 포함시키기에 그러한 특징들이
필요충분하다고 주장할 수 있다. 이렇게 하는 것은, 우리가 그 구성에 대해
무지한 장래의 사례들의 범위에서 무엇이 행해져야 하는지를 맹목적으로 미리
판단하는 대가로 일정한 확실성 또는 \textbf{(p.~130)}
예측가능성(predictability)을 확보하는 것이다. 따라서 우리는 기실 사전에
쟁점들을 확정하는 데 성공하겠지만, 동시에 어둠 속에서, 그것들이 발생하고
식별될 때에만 합리적으로 확정될 수 있는 쟁점들을 확정하게 된다. 이러한
기법은 합리적인 사회적 목적들에 효력을 부여하기 위해 우리가 배제하고
싶어 할 사례들, 그리고 우리가 그것들을 덜 경직적으로 정의된 채로
두었더라면 우리 언어의 개방적 직조(open texture)를 지닌 용어들이 우리로
하여금 배제할 수 있게 했을 사례들을, 규칙의 범위 안에 포함하도록 우리를
강제할 것이다. 이로써 우리의 분류들의 경직성은 그 규칙을 가지거나
유지하려는 우리의 목적들과 충돌하게 될 것이다.

이 과정의 완성은 법리학자들의 ‘개념의 천국(heaven of concepts)’이다.
이는 하나의 일반 용어가 단일한 규칙의 모든 적용에서 동일한 의미를 갖는
데 그치지 않고, 법체계 내의 어떤 규칙에서 등장하든 언제나 동일한 의미를
부여받을 때에 도달한다. 그러한 경우에는, 그 용어가 여러 차례 재현될
때마다 서로 다른 이해관계나 쟁점을 고려하여 그 의미를 해석하려는 어떠한
노력도 더 이상 요구되거나 시도되지 않는다.

사실 모든 체계들은 각기 다른 방식으로 두 가지 사회적 필요 사이에서
타협한다. 하나는 넓은 행위 영역들에 걸쳐, 새로운 공무담당 지침(official
guidance)이나 사회적 쟁점들의 형량 없이도 사적 개인들이 자기 자신에게
안전하게 적용할 수 있는 확실한(certain) 규칙들에 대한 필요이고, 다른
하나는 구체적 사례에서 발생할 때에만 비로소 적절하게 이해되고 확정될 수
있는 쟁점들을, 나중에 정보를 갖춘 공무담당 선택(official choice)에 의해
확정되도록 열어 두어야 할 필요이다. 어떤 법체계들에서는 어떤 시기에는
확실성(certainty)에 지나치게 많은 것이 희생되어, 제정법들(statutes)이나
선례에 대한 사법적 해석이 지나치게 형식적이고, 그래서 사회적 목적들의
관점에서 고려될 때에만 보이는 사례들 사이의 유사성과 차이에 응답하지
못할 수 있다. 다른 체계들이나 다른 시기에는, 법원이 선례들에서 지나치게
많은 것을 영구히 열려 있거나 수정 가능한 것으로 취급하고, 입법 언어가 그
개방적 직조에도 불구하고 결국 제공하는 한계들에 지나치게 적은 존중을
보이는 것처럼 보일 수 있다. 이 문제에서 법이론은 기묘한 역사를 가진다.
그것은 법규칙들의 불확정성들(indeterminacies)을 무시하거나 과장하는
경향이 있기 때문이다. 이러한 극단들 사이의 진동에서 벗어나려면, 이
\textbf{(p.~131)} 불확정성의 뿌리에 놓인 미래를 예견하지 못하는 인간적
한계가 행위의 서로 다른 영역들에서 정도를 달리하며, 법체계들이 그 한계에
상응하는 다양한 기법들로 대응한다는 점을 상기해야 한다.

때로 법적으로 통제되어야 할 영역은 처음부터, 개별 사례들의 특징들이
사회적으로 중요한 그러나 예측 불가능한 측면들에서 너무나 다양하게 변할
것이어서, 추가적인 공무담당 지시(official direction) 없이 사례에서
사례로 적용될 일양적 규칙들을 입법부가 사전에 유용하게 정식화할 수 없는
영역으로 인정된다. 따라서 그러한 영역을 규율하기 위해 입법부는 매우
일반적인 기준들을 세우고, 다양한 사례 유형들에 정통한 행정적 규칙제정
기관에, 그 특수한 필요들에 맞추어진 규칙들을 만들어 내는 과업을
위임한다. 예컨대 입법부는 어떤 산업에 일정한 기준들을 유지하라고 요구할
수 있다. 즉 \emphksans{공정한 요율(fair rate)}만을 부과하거나 \emphksans{안전한
작업 체계들(safe systems)}을 제공하라고 요구할 수 있다. 서로 다른
기업들이 이 모호한 기준들을 스스로에게 적용하도록 맡겨 두어,
\emphksans{사후적으로(ex post facto)} 그것들을 위반한 것으로 판정될 위험을
부담하게 하기보다는, 행정기관이 규정으로 특정 산업에 대해 무엇이 ‘공정한
요율’ 또는 ‘안전한 체계’로 간주될 것인지를 특정할 때까지 위반에 대한
제재 사용을 유보하는 것이 최선이라고 판명될 수 있다. 이러한 규칙제정
권한은 해당 산업에 관한 사실들에 대한 사법적 조사와 유사한 것 및 특정한
규제 형식에 찬반하는 논변들의 청문을 거친 뒤에만 행사될 수 있을지도
모른다.

물론 매우 일반적인 기준들(very general standards)이 있더라도, 무엇이
그것들을 충족하거나 충족하지 않는지에 관한 명백하고 다툼의 여지가 없는
예들은 있을 것이다. 무엇이 ‘공정한 요율’ 또는 ‘안전한 체계’인지, 또는
아닌지에 관한 몇몇 극단적 사례들은 언제나 \emphksans{처음부터(ab initio)}
식별 가능할 것이다. 따라서 무한히 다양한 사례들의 범위 한쪽 끝에는 필수
서비스에 대해 공중을 사실상 인질로 잡으면서 기업가들에게 막대한 이윤을
안겨주는 너무 높은 요율이 있을 것이고, 다른 끝에는 기업 운영의 유인을
제공하지 못하는 너무 낮은 요율이 있을 것이다. 이 둘은 서로 다른 방식으로
요율 규제에서 우리가 가질 수 있는 어떤 가능한 목적도 좌절시킬 것이다.
그러나 이것들은 서로 다른 요소들의 범위에서 극단들일 뿐이며 실제로
마주칠 개연성이 높지 않다. 그 사이에 주의를 요구하는 어려운 실제
사례들이 놓여 있다. 관련 요소들의 예견 가능한 조합들은 적고, 이는 공정한
요율 또는 안전한 체계라는 우리의 최초 목적에 상대적 불확정성을 수반하며,
추가적인 공무담당 \textbf{(p.~132)} 선택(official choice)의 필요를
수반한다. 이러한 경우에는 규칙제정 기관(rule-making authority)이
재량(discretion)을 행사해야 한다는 점이 명확하며, 여러 사례들이 제기하는
문제를 많은 상충하는 이익들 사이의 합리적 타협인 답과는 구별되는,
발견되어야 할 하나의 유일하게 올바른 답이 있는 것처럼 다룰 가능성은
없다.

두 번째 유사한 기법은, 통제되어야 할 영역이 일률적으로 행해지거나 삼가야
할 특정 행위들의 부류를 식별하여 그것들을 단순 규칙(simple rule)의
대상으로 삼는 것은 불가능하지만, 상황들의 범위가 매우 다양하더라도 공통
경험의 친숙한 특징들을 포괄하는 경우에 사용된다. 여기서는 무엇이
‘합리적(reasonable)’인지에 관한 공통 판단들이 법에 의해 사용될 수 있다.
이 기법은 여러 예견 불가능한 형태로 발생하는 사회적 요구들(social
claims) 사이에서 형량하고 합리적 균형을 맞추는 과업을, 법원의 교정을
받을 수 있는 개인들에게 남겨 둔다. 이 경우 개인들은 가변적 기준(variable
standard)이 공식적으로 규정되기 \emphksans{전에(before)} 그 기준에 부합할
것을 요구받으며, 그들이 그것을 위반했을 때에만 \emphksans{사후적으로(ex post
facto)} 법원으로부터, 특정 행위들이나 부작위들의 관점에서 자신들에게
요구된 기준이 무엇이었는지를 알게 될 수 있다. 그러한 사안들에 대한
법원의 결정들이 선례로 간주되는 곳에서는, 가변적 기준에 대한 법원의
특정화는 행정기관이 위임된 규칙제정 권한(delegated rulemaking power)을
행사하는 것과 매우 유사하지만, 또한 명백한 차이들도 있다.

영미법(Anglo-American law)에서 이 기법의 가장 유명한 예는
과실(negligence) 사건들에서 상당한 주의(due care)의 기준을 사용하는
것이다. 타인에게 신체적 손해를 입히는 것을 피하기 위해 합리적 주의를
기울이지 않은 자들에게는 민사 제재가, 더 드물게는 형사 제재가 적용될 수
있다. 그러나 구체적 상황에서 합리적 주의 또는 상당한 주의란 무엇인가?
물론 우리는 상당한 주의의 전형적 예들을 들 수 있다. 예컨대 통행이
예상되는 곳에서 멈추고, 살피고, 듣는 것과 같은 일들이다. 그러나 우리
모두는 주의가 요구되는 상황들이 엄청나게 다양하며, 이제는 ‘멈추고,
살피고, 듣기’ 외에 또는 그 대신에 많은 다른 행위들이 요구된다는 점을 잘
알고 있다. 기실 보는 것이 위험을 피하는 데 도움이 되지 않는다면 이것들은
충분하지 않을 수 있고 전혀 쓸모없을 수도 있다. 합리적 주의 기준들을
적용하면서 우리가 추구하는 것은 (1) 중대한 해악을 피하게 할
\textbf{(p.~133)} 예방조치들이 취해지게 하면서도, (2) 그 예방조치들이
적절한 예방조치들의 부담이 다른 존중할 만한 이익들의 지나치게 큰 희생을
수반하지 않도록 하는 것이다. 물론 피를 흘리며 죽어가는 사람이 병원으로
이송되는 경우가 아니라면, 멈추고, 살피고, 듣는 것에 의해 그리 많은 것이
희생되지는 않는다. 그러나 주의가 요구되는 가능한 사례들의 엄청난 다양성
때문에, 해악에 대한 예방조치가 취해져야 한다면 어떤 상황들의 조합이
발생할지, 또는 어떤 이익들이 어느 정도로 희생되어야 할지를 우리는
\emphksans{처음부터(ab initio)} 예견할 수 없다. 따라서 우리는 특정 사례들이
발생하기 전에는, 해악의 위험을 줄이기 위해 이익들이나 가치들의 어떤 희생
또는 타협을 정확히 하고자 하는지를 고려할 수 없다. 다시 말해, 사람들을
해악으로부터 보호하려는 우리의 목적은, 경험만이 우리 앞에 가져다줄
가능성들과 결합시키거나 그것들에 비추어 시험하기 전까지는
불확정적(indeterminate)이다. 경험이 그것을 가져오면, 그때 우리는 하나의
결정에 직면해야 하며, 그 결정은 내려질 때 우리의 목적을 \emphksans{그만큼(pro
tanto)} 확정적으로 만들 것이다.

이 두 가지 기법(techniques)을 고려하면, 가변적 기준(variable standard)이
아니라 특정 행위들을 요구하는 규칙(rule)에 의해, 개방적 직조(open
texture)의 가장자리(fringe)만을 남긴 채 \emphksans{처음부터(ab initio)}
성공적으로 통제되는 넓은 행위 영역들의 특징들이 두드러진다. 그것들은
다음 사실에 의해 특징지어진다. 어떤 구별 가능한 행위들, 사건들, 또는
사태들이 우리에게 회피되어야 하거나 생겨나야 할 것으로서 실천적으로 매우
중요하여, 그것들과 함께 나타나는 부수적 상황들이 우리로 하여금 그것들을
달리 보도록 기울게 하는 경우가 거의 없다는 것이다. 이것의 가장 거친 예는
인간 존재의 살해(killing of a human being)이다. 인간 존재들이 다른 인간
존재들을 살해하는 상황들은 매우 다양하지만, 우리는 가변적 기준(‘인간
생명에 대한 상당한 존중’)을 세우는 대신 살해를 금지하는 규칙을 만들 수
있는 위치에 있다. 이는 생명 보호의 중요성에 대한 우리의 평가보다 더 큰
무게를 갖거나 그 평가를 수정하게 하는 요소들이 우리에게 거의 없는 것으로
보이기 때문이다. 거의 언제나 살해는 말하자면 그것에 수반되는 다른
요소들을 \emphksans{지배한다(dominates)}. 따라서 우리가 그것을
‘살해(killing)’로서 사전에 배제할 때, 우리는 서로 형량되어야 하는
쟁점들을 맹목적으로 미리 판단하고 있는 것이 아니다. 물론 예외들, 이
통상적으로 지배적인 요소를 압도하는 요소들이 있다. 정당방위에서의 살해와
기타 정당화 가능한 살해(justifiable homicide)의 형태들이 그것이다.
그러나 이것들은 적고 상대적으로 단순한 용어들로 식별 가능하며, 일반
규칙에 대한 예외들로 인정된다.

\textbf{(p.~134)} 쉽게 식별 가능한 어떤 행위, 사건, 또는 사태의 지배적
지위가 어떤 의미에서는 관행적(conventional)이거나 인위적(artificial)일
수 있으며, 인간 존재로서 우리에게 그것이 가지는 ‘자연적(natural)’ 또는
‘내재적(intrinsic)’ 중요성 때문이 아닐 수 있다는 점을 주목하는 것이
중요하다. 도로 통행 규칙(rule of the road)이 어느 쪽 도로를 지정하는지는
중요하지 않으며, 부동산 양도증서(conveyance)의 작성(execution)을 위해
어떤 형식요건들이 규정되는지도 (일정 한계 안에서는) 중요하지 않다.
그러나 이러한 사항들에 대해 쉽게 식별 가능하고 일양적인 절차가 있어야
하며, 따라서 명확한 옳고 그름(right and wrong)이 있어야 한다는 점은
대단히 중요하다. 이것이 법에 의해 도입되면, 소수의 예외를 제외하고는
그것을 고수하는 것의 중요성이 최우선적(paramount)이 된다. 왜냐하면
상대적으로 적은 수의 부수적 상황들만이 그것보다 더 큰 무게를 가질 수
있고, 그러한 상황들은 쉽게 예외로 식별되어 규칙으로 환원될 수 있기
때문이다. 영국 부동산법(real property law)은 규칙들의 이러한 측면을 매우
명확하게 예증한다.

권위 있는 예시들(authoritative examples)에 의한 일반 규칙들의 전달은,
우리가 보았듯이, 더 복잡한 종류의 불확정성들을 수반한다. 선례를 법적
유효성의 기준으로 인정하는 것은 서로 다른 체계들에서, 그리고 동일한
체계의 서로 다른 시기들에서 서로 다른 것들을 의미한다. 영국의 선례
‘이론(theory)’에 대한 서술들은 어떤 지점들에서는 여전히 매우 논쟁적이다.
기실 그 이론에서 사용되는 핵심 용어들, 즉 ‘\emphksans{판결 이유(ratio
decidendi)}’, ‘중요 사실(material facts)’, ‘해석(interpretation)’조차도
각자 고유한 불확실성의 반음영(penumbra)을 지닌다. 우리는 어떤 새로운
일반적 서술을 제시하지는 않고, 제정법(statute)의 경우에서와 마찬가지로,
개방적 직조(open texture)의 영역과 그 안에서 이루어지는 창조적 사법
활동을 간략하게 특징짓고자 할 뿐이다.

영국법에서 선례(precedent)의 사용에 대한 정직한 서술은 다음과 같은
대조적 사실들의 쌍들을 위한 자리를 반드시 마련해야 한다. \emphksans{첫째},
주어진 권위 있는 선례(authoritative precedent)가 어떤 규칙에 대한
권위(authority)인지를 결정하는 단일한 방법은 없다. 그럼에도 불구하고
판결이 내려진 사건들의 압도적 다수에서는 의문이 거의 없다.
사건요지(head-note)는 보통 충분히 정확하다. \emphksans{둘째}, 사건들로부터
추출될 어떤 규칙에 대해서도 권위 있거나 유일하게 올바른
정식화(authoritative or uniquely correct formulation)는 없다. 다른 한편,
선례가 후속 사건에 대해 가지는 관련성(bearing)이 문제될 때, 주어진
정식화가 적절하다는 데 매우 일반적인 합의가 있는 경우가 흔하다.
\emphksans{셋째}, 선례에서 추출된 규칙이 어떤 권위 있는 지위(authoritative
status)를 가지든, 그것은 그 규칙에 구속되는 법원들이 \textbf{(p.~135)}
다음 두 유형의 창조적 또는 입법적 활동을 행사하는 것과 양립 가능하다.
한편으로 후속 사건을 판단하는 법원들은 선례에서 추출된 규칙을 좁히고,
이전에 고려되지 않았거나 고려되었더라도 열려 있던 어떤 예외를
인정함으로써 선례와 반대되는 결정에 이를 수 있다. 이처럼 앞선 사건을
‘구별(distinguishing)’하는 과정은 그 사건과 현재 사건 사이에서 법적으로
관련된 어떤 차이를 발견하는 것을 포함하며, 그러한 차이들의 부류는 결코
소진적으로 확정될 수 없다. 다른 한편, 법원들은 앞선 선례를 따르면서도,
앞선 사건에서 정식화된 규칙 안에서 발견되는 어떤 제약(restriction)을,
그것이 제정법(statute)이나 더 앞선 선례에 의해 확립된 어떤 규칙에
의해서도 요구되지 않는다는 근거로 폐기할 수 있다. 이렇게 하는 것이
규칙을 넓히는 것이다. 선례의 구속력(binding force)에 의해 열려 있는
이러한 두 형태의 입법적 활동에도 불구하고, 영국 선례 체계의 결과는 그
사용을 통해 중요도가 큰 것과 작은 것을 모두 포함하는 매우 많은 수의
규칙들의 집합을 산출한 것이며, 그 가운데 많은 규칙들은 어떤 제정법
규칙(statutory rule) 못지않게 확정적(determinate)이다. 이제 그것들은
제정법(statute)에 의해서만 변경될 수 있으며, 법원들 자신도 확립된
선례들의 요구와 ‘실질적 타당성(merits)’이 배치되는 것처럼 보이는
사건들에서 흔히 그렇게 선언한다.

법의 개방적 직조(open texture)는, 기실 많은 행위 영역들에서 사안마다
무게가 달라지는 상충하는 이익들 사이의 균형을 상황에 비추어 맞추는
법원들이나 공무담당자들에 의해 많은 것이 발전되도록 남겨져야 함을
의미한다. 그럼에도 불구하고 법의 삶(life of the law)은 매우 큰 정도로,
가변적 기준들(variable standards)의 적용과는 달리 사례마다 새로운 판단을
요구하지 않는 확정적 규칙들(determinate rules)에 의한 공무담당자들과
사적 개인들 모두의 지도로 이루어진다. 이러한 사회생활의 두드러진 사실은,
어떤 규칙이든(문서화된 것이든 선례에 의해 전달된 것이든) 구체적 사례에
대한 적용가능성을 둘러싸고 불확실성들이 터져 나올 수 있음에도 불구하고
여전히 참이다. 여기, 규칙들의 주변부에서 그리고 선례 이론에 의해 열려
있는 영역들에서, 법원들은 행정기관들이 가변적 기준들을 정교화할 때
중심적으로 수행하는 규칙-생산 기능을 수행한다. \emphksans{선례구속 원칙(stare
decisis)}이 확고히 인정된 체계에서, 법원의 이 기능은 행정기관에 의한
위임된 규칙제정 권한(delegated rule-making powers)의 행사와 매우
유사하다. 영국에서는 이 사실이 형식들(forms)에 의해 자주 가려진다.
법원들이 \textbf{(p.~136)} 흔히 그러한 창조적 기능을 부인하고, 제정법
해석과 선례 사용의 적절한 과업은 각각 ‘입법자의 의도(intention of the
legislature)’와 이미 존재하는 법을 탐색하는 것이라고 고수하기 때문이다.

\subsection{2. 규칙회의주의(rule-scepticism)의
다양상들}\label{uxaddcuxce59uxd68cuxc758uxc8fcuxc758rule-scepticismuxc758-uxb2e4uxc591uxc0c1uxb4e4}

우리는 법의 개방적 직조(open texture)에 대해 상당히 길게 논의했는데,
이는 이 특징을 정당한 관점에서 보는 것이 중요하기 때문이다. 그것을
정당하게 다루지 못하면(do justice) 언제나 법의 다른 특징들을 흐리게 할
과장들이 유발될 것이다. 모든 법체계에는 처음에는 모호한 기준들을
확정적으로 만들고, 제정법들(statutes)의 불확실성들을 해소하며, 권위 있는
선례들(authoritative precedents)에 의해 단지 넓게만 전달된 규칙들을
발전시키고 한정하는 데서, 법원들과 기타 공무담당자들이
재량(discretion)을 행사하도록 넓고 중요한 영역이 열려 있다. 그럼에도
불구하고 이러한 활동들은, 중요하고 충분히 연구되지 않았더라도, 그것들이
일어나는 틀(framework)과 그것들의 주된 최종 산물이 모두 일반 규칙들로
이루어져 있다는 사실을 가려서는 안 된다. 이것들은 개인들이 사례마다 그
적용을 스스로 볼 수 있는 규칙들로서, 추가적인 공무담당 지시(official
direction)나 재량에 더 호소할 필요가 없다.

규칙들이 법체계의 구조에서 중심적 자리를 가진다는 논지(contention)가
진지하게 의심받을 수 있었다는 것은 이상해 보일 수 있다. 그러나
‘규칙회의주의(rule-scepticism)’, 곧 규칙에 관한 말은 법이 단지 법원의
결정들과 그것들에 대한 예측들로 이루어진다는 진실을 가리는(cloaking)
하나의 신화라는 주장(claim)은 법률가의 솔직함에 강력하게 호소할 수 있다.
이 주장이 이차 규칙들과 일차 규칙들을 모두 포괄하도록 무제한적인 일반
형식으로 진술된다면, 그것은 기실 매우 비정합적(incoherent)이다. 법원의
결정들이 존재한다는 단언은 아무 규칙도 전혀 존재하지 않는다는 부정과
일관되게 결합될 수 없기 때문이다. 이는 우리가 보았듯이 법원의 존재가
관할권(jurisdiction)을 계속 바뀌는 일련의 개인들에게 부여하고 그래서
그들의 결정들을 권위 있는 것으로 만드는 이차 규칙들의 존재를 수반하기
때문이다. 어떤 사람들의 공동체가 결정과 결정에 대한 예측이라는
관념들(notions)은 이해하지만 규칙이라는 관념(notion)은 이해하지
못한다면, 그들에게는 \emphksans{권위 있는(authoritative)} 결정이라는
관념(idea)이 결여될 것이며 그와 함께 법원이라는 관념(idea)도 결여될
것이다. 사인의 결정과 법원의 결정을 구별할 아무것도 없을 것이다. 우리는
\textbf{(p.~137)} ‘습관적 복종’이라는 관념(notion)으로, 법원에서
요구되는 권위 있는 관할권의 토대로서 결정의 예측가능성이 가지는 결함을
보충하려 할지도 모른다. 그러나 이렇게 하면, 입법 권한들을 부여하는
규칙의 대체물로 제4장에서 습관을 검토했을 때 드러난 모든 부적합성들을,
이 목적에서도, 습관이라는 관념(notion)이 가진다는 것을 발견하게 될
것이다.

이 이론의 더 온건한 몇몇 형태에서는, 법원들이 존재하려면 그것들을
구성하는 법규칙들이 있어야 하며, 따라서 이러한 규칙들 자체는 단순히
법원의 결정들에 대한 예측들일 수 없다는 점이 인정될 수 있다. 그러나 이
양보만으로는 실제로 거의 진전을 이룰 수 없다. 이러한 유형의 이론의
특징적 단언은 제정법들(statutes)이 법원에 의해 적용되기 전에는 법이
아니라 단지 법의 원천들(sources of law)이라는 것이고, 이는 존재하는
유일한 규칙들은 법원들을 구성하는 데 필요한 규칙들이라는 단언과 일관되지
않기 때문이다. 계속 바뀌는 일련의 개인들에게 입법 권한들을 부여하는 이차
규칙들도 있어야 한다. 왜냐하면 이 이론은 제정법들(statutes)이 있다는
것을 부정하지 않기 때문이다. 기실 그것은 그것들을 법의 단순한
‘원천들(sources)’로 인용하고, 제정법들이 법원에 의해 적용되기 전에는
법이라는 것만 부정한다.

이러한 반론들은 중요하며, 이 이론의 경솔한 형태에 대해서는 적절하지만,
그 모든 형태에 적용되지는 않는다. 규칙회의주의는 사법적 또는 입법적
권한을 부여하는 이차 규칙들의 존재를 부정하려는 의도가 전혀 없었고,
그것들이 결정들이나 결정들에 대한 예측들에 불과하다고 보일 수 있다는
주장에 결부된 적도 전혀 없었을 수 있다. 분명히 이러한 유형의 이론이 가장
자주 의존해 온 예들은 사인들에게 책무를 부과하거나 권리들 또는 권한들을
부여하는 규칙들에서 끌어온 것이다. 그러나 규칙들이 있다는 부정과
규칙들이라 불리는 것이 법원의 결정들에 대한 예측들일 뿐이라는 단언을
이러한 방식으로 제한한다고 상정하더라도, 적어도 한 의미에서는 그것이
명백히 거짓이다. 왜냐하면 적어도 근대 국가의 어떤 행위 영역들과
관련해서는, 개인들이 우리가 내부 관점(internal point of view)이라 불렀던
행위와 태도의 전 범위를 실제로 보인다는 점은 의심될 수 없기 때문이다.
법들은 그들의 삶에서 단순히 습관들이나 법원들의 결정들 또는 다른
공무담당자들의 행위들을 예측하는 토대로 기능하는 것이 아니라, 수용된
법적 행위 기준들(accepted legal standards of behaviour)로 기능한다. 즉
그들은 단지 법이 자신들에게 요구하는 것을 \textbf{(p.~138)} 상당한
정도의 반복적 일양성(tolerable regularity)으로 행할 뿐만 아니라, 그것을
법적 행위 기준으로 바라보고, 타인들을 비판하거나 요구들을 정당화할 때
그것을 참조하며, 타인들이 자신들에게 가하는 비판들과 요구들을 받아들일
때에도 그것을 참조한다. 법규칙들을 이러한 규범적 방식으로 사용할 때
그들은 물론 법원들과 기타 공무담당자들이 해당 체계의 규칙들에 따라
반복적으로 일양적인, 따라서 예측 가능한 일정한 방식들로 계속 결정하고
행동할 것이라고 상정한다. 그러나 개인들이 법원들의 결정들이나 제재가
적용될 개연성을 기록하고 예측하는 외부 관점(external point of view)에
자신들을 한정하지 않는다는 것은 분명히 관찰 가능한 사회생활의 사실이다.
그 대신 그들은 행위의 지침으로서 법에 대한 공유된 수용을 규범적 용어들로
계속 표현한다. 우리는 제3장에서 ‘의무(obligation)’와 같은 규범적
용어들이 공무담당 행위에 대한 예측 이상의 아무것도 의미하지 않는다는
주장을 상세히 검토했다. 우리가 논변했듯이 그 주장이 거짓이라면,
법규칙들은 사회생활 속에서 그러한 것으로 기능한다. 그것들은 습관들이나
예측들의 서술이 아니라 규칙들로서 \emphksans{사용된다(used)}. 물론 그것들은
개방적 직조를 지닌 규칙들이며, 그 직조가 열린 지점들에서 개인들은
법원들이 어떻게 결정할지를 예측하고 그에 따라 자기 행위를 조정할 수 있을
뿐이다.

규칙회의주의는 우리의 진지한 주의를 요구하지만, 오직 사법적
결정(judicial decision)에서 규칙들이 수행하는 기능에 관한 이론으로서만
그렇다. 이 형태에서, 앞서 우리가 주의를 환기한 모든 반론들을
인정하면서도, 그것은 법원들에 관한 한 개방적 직조의 영역을 경계 지을
아무것도 없다는 논지(contention)에 이른다. 따라서 판사들을 그들이 하는
방식으로 사건들을 결정하도록 스스로 규칙들에 종속되어 있거나 규칙들에
‘구속되어(bound)’ 있다고 보는 것은, 무의미하지 않더라도 거짓이라는
것이다. 판사들은 충분히 예측 가능한 정도의 반복적 일양성(regularity)과
일양성(uniformity)을 가지고 행동하여, 다른 사람들이 긴 기간에 걸쳐
법원의 결정들에 의해 규칙들처럼 살아갈 수 있게 할 수 있다. 판사들은
심지어 자신들이 그렇게 결정할 때 강제감(feelings of compulsion)을 경험할
수도 있고, 이러한 감정들 역시 예측 가능할 수 있다. 그러나 이 너머에는
그들이 준수하는 규칙으로 특징지을 수 있는 것은 아무것도 없다. 법원들이
올바른 사법적 행위의 기준들(standards of correct judicial behaviour)로
취급하는 것은 아무것도 없으며, 따라서 그 행위 안에는 규칙들의 수용에
특징적인 내부 관점(internal point of view)을 드러내는 것도 아무것도
없다는 것이다.

이 이론의 이러한 형태는 매우 다양한 성격과 중요도를 지닌 근거들로부터
지지를 얻는다. 규칙회의주의자는 \textbf{(p.~139)} 때로는 실망한
절대주의자(absolutist)이다. 그는 규칙들이 형식주의자의
이상세계(formalist's heaven)에서 그러할 모습 그대로도 아니며, 인간이
신들처럼 되어 모든 가능한 사실의 조합을 예견할 수 있고 그래서 개방적
직조가 규칙들의 필수적 특징이 아닌 세계에서 그러할 모습 그대로도
아니라는 점을 발견한 것이다. 규칙이 존재한다는 것이 무엇인지에 관한
회의주의자의 구상은 도달 불가능한 이상(ideal)일 수 있고, 소위 규칙이라
불리는 것에서 그 이상이 달성되지 않는다는 것을 발견할 때, 그는 규칙들이
존재한다거나 존재할 수 있다는 것을 부정함으로써 자신의 실망을 표현한다.
이렇게 해서, 판사들이 사건을 판결할 때 자신들을 구속한다고 주장하는
규칙들이 개방적 직조를 갖고 있다거나, 사전에 소진적으로 특정될 수 없는
예외들을 가진다는 사실, 그리고 규칙으로부터의 일탈이 판사에게 물리적
제재(physical sanction)를 불러오지 않는다는 사실은 회의주의자의 주장을
확립하는 데 자주 이용된다. 이 사실들은 다음과 같은 점을 보이기 위해
강조된다. “규칙들이 중요한 것은 판사들이 무엇을 할지를 예측하는 데
도움을 주는 한에서이다. 예쁜 장난감(pretty playthings)으로서의 중요성을
제외하면 그것이 규칙들의 중요성의 전부이다.” \footnote{Llewellyn,
  \emph{The Bramble Bush} (2nd edn.), p.~9.}

이런 식으로 논변하는 것은 현실 생활의 어떤 영역에서든 규칙들이 실제로
무엇인지를 무시하는 것이다. 그것은 우리가 다음의 딜레마에 직면해 있다고
시사한다. “규칙들은 형식주의자의 천국(formalist's heaven)에서 그러할
모습 그대로이며 족쇄가 구속하듯 구속하거나, 아니면 규칙들은 없고 예측
가능한 결정들이나 행위 양식들만 있다.” 그러나 분명 이것은 거짓
딜레마다. 우리는 다음 날 친구를 방문하겠다고 약속한다. 그 날이 오자 그
약속을 지키는 것이 위험하게 아픈 누군가를 방치하는 것을 수반하게 된다.
이것이 약속을 지키지 않는 적절한 이유로 수용된다는 사실은, 약속을 지킬
것을 요구하는 규칙은 없고 약속을 지키는 데 어떤 반복적 일양성만 있을
뿐이라는 것을 결코 의미하지 않는다. 그러한 규칙들에 소진적으로 진술될 수
없는 예외들이 있다는 사실로부터, 모든 상황에서 우리가 우리의 재량에
맡겨져 있고 결코 약속을 지키도록 구속되지 않는다는 결론은 따르지 않는다.
‘\ldots 하지 않는 한(unless \ldots)’이라는 말로 끝나는 규칙도 여전히
규칙이다.

때때로 법원들을 구속하는 규칙들의 존재는, 어떤 사람이 특정한 방식으로
행위함으로써 자신에게 그렇게 행위할 것을 요구하는 규칙에 대한 수용을
드러냈는지 여부라는 문제가, 그 사람이 행위하기 전이나 행위하면서 거친
사고 과정들에 관한 심리학적 문제들과 \textbf{(p.~140)} 혼동되기 때문에
부정된다. 매우 자주, 어떤 사람이 규칙을 구속력 있는 것으로, 그리고
자신과 타인들이 자유롭게 바꿀 수 없는 것으로 수용할 때, 그는 주어진
상황에서 그것이 무엇을 요구하는지를 아주 직관적으로 보고, 먼저 그 규칙과
그것이 요구하는 것을 생각하지 않고도 그것을 행할 수 있다. 우리가
체스에서 규칙들에 부합하여 말을 움직이거나 신호등이 빨간색일 때 멈출 때,
우리의 규칙-준수(rule-complying) 행위는 종종 규칙들의 관점에서 계산에
의해 매개되지 않은, 상황에 대한 직접 반응이다. 그러한 행위들이 규칙의
진정한 적용이라는 증거는 그것들이 특정한 상황들 속에 놓여 있다는 점이다.
이 상황들 가운데 일부는 특정 행위에 앞서 있고, 다른 일부는 그 뒤를
따른다. 또 그 일부는 일반적이고 가설적인 용어들로만 진술될 수 있다.
행위하면서 우리가 규칙을 적용했다는 것을 보여 주는 이 요소들 가운데 가장
중요한 것은, \emphksans{만약(if)} 우리의 행위가 문제 제기된다면 우리는 그
행위를 그 규칙을 참조하여 정당화하려는 성향이 있다는 점이다. 그리고
규칙에 대한 우리의 수용의 진정성은 그것에 대한 우리의 과거 및 이후의
일반적 인정(acknowledgements)과 그것에의 부합(conformity)뿐 아니라, 우리
자신의 일탈과 타인의 일탈에 대한 비판에서도 드러날 수 있다. 이러한
증거나 유사한 증거에 근거하여 우리는 기실, 규칙에 대한 우리의 ‘생각
없는’ 준수(compliance)에 앞서 무엇이 해야 할 옳은 일이고 왜 그런지를
말하라는 질문을 받았다면, 정직할 경우 우리는 답변에서 그 규칙을 인용했을
것이라고 결론 내릴 수 있다. 진정으로 규칙의 준수(observance)인 행위를,
단지 우연히 그것과 일치하는(coincide with) 행위와 구별하는 데 필요한
것은, 우리의 행위가 그러한 상황들 속에 놓여 있다는 점이지, 규칙에 대한
명시적 사고가 그 행위에 수반되었다는 점이 아니다. 바로 이렇게 우리는,
수용된 규칙의 준수(compliance)로서의 성인 체스 플레이어의 수와, 단지
말을 옳은 자리로 밀어 넣은 아기의 행위를 구별할 것이다.

이것은 가장이나 ‘겉치레(window dressing)’가 불가능하다는 말이 아니며,
그것은 때때로 성공할 수도 있다. 어떤 사람이 자신이 규칙에 근거하여
행위했다고 단지 \emphksans{사후적으로(ex post facto)} 가장했는지를 가리는
검사들은, 모든 경험적 검사들처럼 본질적으로 오류 가능하지만, 고질적으로
그런 것은 아니다. 어떤 사회에서 판사들이 언제나 먼저 직관적으로 또는
‘직감으로(by hunches)’ 결정을 내리고, 그런 다음 법규칙들의 목록에서 당면
사건과 닮았다고 가장한 규칙 하나를 단지 선택할 수도 있다. 그러고 나서
그들은 이것이 자신들의 결정을 요구하는 것으로 자신들이 간주한 규칙이라고
주장할 수 있지만, 그들의 행위나 말의 다른 어떤 것도 \textbf{(p.~141)}
그들이 그것을 자신들을 구속하는 규칙으로 간주했다는 것을 시사하지 않을
수 있다. 일부 사법적 결정들은 이와 같을 수 있다. 그러나 대부분의 경우
결정들은 체스 플레이어의 수처럼, 결정의 인도 기준들로 의식적으로
받아들인 규칙들에 부합하려는 진정한 노력에 의해 도달되거나, 만약
직관적으로 도달되었다면, 판사가 사전에 준수하려는 성향을 가지고 있었고
당면 사건과의 관련성이 일반적으로 인정될 규칙들에 의해 정당화된다는 점은
분명하다.

규칙회의주의(rule-scepticism)의 마지막이자 가장 흥미로운 형태는
법규칙들의 개방적 성격이나 많은 결정들의 직관적 성격에 근거하지 않고,
법원의 결정이 권위 있는 것(something authoritative)으로서 독특한 지위를
가지며, 최고 법원들(supreme tribunals)의 경우에는 최종적(final)이라는
사실에 근거한다. 우리가 다음 절을 할애할 이 형태의 이론은,
그레이(Gray)가 \emph{The Nature and Sources of Law}에서 그렇게 자주
되풀이한 호들리(Hoadly) 주교의 유명한 구절에 암묵적으로 들어 있다.
“글로 쓰였든 말로 발화되었든 어떤 법들을 해석할 절대적 권위(absolute
authority)를 가진 자가 있다면, 사실상 모든 면에서 입법자(lawgiver)는
바로 그 사람이지 그것들을 처음 쓰거나 말한 사람이 아니다.”

\subsection{3. 사법적 결정에서의 최종성(Finality)과
무오류성(Infallibility)}\label{uxc0acuxbc95uxc801-uxacb0uxc815uxc5d0uxc11cuxc758-uxcd5cuxc885uxc131finalityuxacfc-uxbb34uxc624uxb958uxc131infallibility}

최고 법원(supreme tribunal)은 법이 무엇인지를 말하는 데 최후의 말(last
word)을 가지며, 일단 그렇게 말하면 그 법원이 ‘그르다(wrong)’는 진술은 그
체계 안에서는 아무런 귀결도 갖지 않는다. 그로 인해 누구의 권리들이나
책무들도 변경되지 않는다. 그 결정은 물론 입법(legislation)에 의해 법적
효과(legal effect)를 박탈당할 수 있지만, 이에 호소하는 것이 필요하다는
바로 그 사실은, 법에 관한 한 법원의 결정이 그르다는 진술의 공허한 성격을
보여 준다. 이러한 사실들을 고려하면, 최고 법원의 결정들의 경우 그
최종성(finality)과 무오류성(infallibility)을 구분하는 것은 현학적인
것처럼 보인다. 이것은 법원들이 결정할 때 언제나 규칙들에 구속된다는 것을
부정하는 또 다른 형태로 이어진다. “법(또는 헌법)은 법원들이 그것이라고
말하는 것이다.”

이러한 형태의 이론에서 가장 흥미롭고 교훈적인 특징은 ‘법(또는 헌법)은
법원들이 그것이라고 말하는 것이다’와 같은 진술들의 애매성(ambiguity)을
이용한다는 점, 그리고 이 이론이 일관되려면 \textbf{(p.~142)} 법에 대한
비공식 진술들과 법원의 공식 진술들 사이의 관계에 관해 제시해야 하는
설명이다. 이 애매성을 이해하기 위해 우리는 잠시 돌아서 게임의 경우에
있는 유비를 검토할 것이다. 많은 경쟁적 게임들은 공식 득점기록자(official
scorer) 없이 행해진다. 경쟁하는 이익들에도 불구하고, 선수들은 득점
규칙을 특정 사례들에 적용하는 데 상당히 잘 성공한다. 그들은 보통 자기
판단들에서 일치하며, 해소되지 않은 다툼들은 적을 수 있다. 공식
득점기록자 제도가 도입되기 전에, 어떤 선수가 한 득점 진술은, 그가
정직하다면, 그 게임에서 수용된 특정 득점 규칙을 참조하여 게임의 진행을
평가하려는 노력이다. 그러한 득점 진술들은 득점 규칙을 적용하는 내부
진술들(internal statements)이며, 선수들이 일반적으로 규칙들을 지키고 그
위반에 이의를 제기할 것임을 전제하지만, 이러한 사실들의 진술들이나
예측들은 아니다.

관습 체제에서 성숙한 법체계로의 변화와 마찬가지로, 그 판정들이 최종적인
득점기록자 제도를 마련하는 이차 규칙들(secondary rules)이 게임에
추가되면, 체계 안에는 새로운 종류의 내부 진술이 들어온다. 선수들의 득점
진술들과 달리, 득점기록자의 결정들(determinations)은 이차 규칙들에 의해
다툴 수 없는 지위를 부여받기 때문이다. 게임의 목적상 ‘득점은
득점기록자가 말하는 것이다’라는 말이 참인 것은 \emphksans{이런(this)}
의미에서이다. 그러나 득점 \emphksans{규칙}(scoring \emph{rule})은 이전 그대로
남아 있으며, 득점기록자의 책무(duty)는 그것을 자신이 할 수 있는 최선으로
적용하는 것임을 보는 것이 중요하다. ‘득점은 득점기록자가 말하는
것이다’가, 득점기록자가 자기 재량으로 적용하기로 선택한 것 외에는 득점
규칙이 없다는 뜻이라면 거짓일 것이다. 기실 그러한 규칙을 가진 게임도
있을 수 있고, 득점기록자의 재량이 어느 정도 반복적으로 일양적인 행사
양상을 보인다면 그것을 하는 데서 얼마간 즐거움을 찾을 수도 있을 것이다.
그러나 그것은 다른 게임일 것이다. 우리는 그러한 게임을 ‘득점기록자의
재량(scorer’s discretion)' 게임이라고 부를 수 있다.

득점기록자가 가져오는 분쟁의 신속하고 최종적인 처리라는 이점들이 대가를
치르고 사들여진다는 점은 명백하다. 득점기록자 제도는 선수들을
곤경(predicament)에 직면시킬 수 있다. 이전과 같이 게임이 득점 규칙에
의해 규율되기를 바라는 소망과, 그 적용이 의심스러운 곳에서 그 적용에
관하여 최종적이고 권위 있는 결정들을 바라는 소망은 충돌하는 목적들로
드러날 수 있다. 득점기록자는 정직한 \textbf{(p.~143)} 실수를 할 수도
있고, 술에 취했을 수도 있으며, 자기 능력의 최선으로 득점 규칙을 적용해야
한다는 책무(duty)를 함부로 위반할 수도 있다. 그는 이 이유들 중 어느 하나
때문에, 타자가 전혀 움직이지 않았는데도 ‘득점(run)’을 기록할 수 있다.
상급 권위기관에 대한 항소로 그의 판정들을 교정하는 장치가 마련될 수
있다. 그러나 이것은 어딘가에서 최종적이고 권위 있는 판단(judgment)으로
끝나야 하며, 그것은 오류 가능성이 있는 인간 존재들에 의해 내려질
것이므로 정직한 실수, 남용, 또는 위반의 동일한 위험을 수반할 것이다.
모든 규칙의 위반을 교정하기 위한 규칙을 마련하는 것은 불가능하다.

규칙들을 최종적이고 권위 있게 적용할 권위기관(authority)을 세우는 데
내재한 위험은 어느 영역에서든 현실화될 수 있다. 게임이라는 소박한
영역에서 현실화될 수 있는 위험들은 고려할 가치가 있는데, 그것들이 이
형태의 권위(authority)가 사용되는 곳이면 어디서나 그것을 이해하는 데
필요한 몇몇 구별들을 규칙회의주의자가 무시한다는 점을 특히 분명한
방식으로 보여 주기 때문이다. 공식 득점기록자(official scorer)가 세워지고
그의 득점 판정들(determinations)이 최종적인 것으로 만들어지면,
선수들이나 기타 비공식 인물들(non-officials)이 한 득점에 관한 진술들은
게임 안에서 아무런 지위를 갖지 않는다. 그것들은 그 결과에 무관하다.
그것들이 득점기록자의 진술과 우연히 일치하면 문제될 것이 없지만,
충돌한다면 결과를 산정할 때 무시되어야 한다. 그러나 이러한 매우 명백한
사실들은 선수들의 진술들이 득점기록자의 판정들에 대한 예측들로
분류된다면 왜곡될 것이며, 이러한 진술들이 득점기록자의 판정들과 충돌할
때 무시되는 것을 그것들이 거짓으로 드러난 판정들에 대한 예측들이었다고
말함으로써 설명하는 것은 터무니없을 것이다. 공식 득점기록자가 도입된
뒤에도, 선수는 자기 자신의 득점 진술들을 하면서 이전에 하던 일을 하고
있다. 즉 득점 규칙을 참조하여, 자신이 할 수 있는 최선으로 게임의 진행을
평가하고 있다. 득점기록자 자신도 자기 지위의 책무들을 이행하는 한 바로
이것을 하고 있다. 양자의 차이는 한쪽이 다른 쪽이 무엇을 말할지를
예측한다는 데 있지 않고, 선수들의 진술들은 득점 규칙의 비공식
적용들(unofficial applications)이어서 결과를 산정하는 데 의미가 없지만,
득점기록자의 진술들은 권위 있고 최종적이라는 데 있다. 만약 행해지는
게임이 ‘득점기록자의 재량(scorer’s discretion)'이라면, 비공식 진술들과
공식 진술들 사이의 관계가 \textbf{(p.~144)} 필연적으로 달라진다는 점을
관찰하는 것이 중요하다. 선수들의 진술들은 득점기록자의 판정들에 대한
예측일 \emphksans{것(would)}일 뿐 아니라, 그 밖의 아무것도 \emphksans{될 수
없을(could)} 것이다. 왜냐하면 그 경우 ‘득점은 득점기록자가 말하는
것이다’ 자체가 득점 \emphksans{규칙}일 것이기 때문이다. 선수들의 진술들이
득점기록자가 공식적으로 하는 것의 단순한 비공식 판본들일 가능성은 없을
것이다. 그러면 득점기록자의 판정들은 최종적이고 무오류적(infallible)일
것이다. 아니 오히려 그것들이 오류 가능적인지 무오류적인지를 묻는 질문은
무의미할 것이다. 그에게 ‘옳게’ 또는 ‘그르게’ 할 것이 아무것도 없기
때문이다. 그러나 보통의 게임에서 ‘득점은 득점기록자가 말하는 것이다’는
득점 규칙이 아니다. 그것은 특정 사례들에서 득점 규칙에 대한 그의 적용의
권위와 최종성을 마련하는 규칙이다.

권위 있는 결정(authoritative decision)의 이 예에서 배울 두 번째 교훈은
더 근본적인 문제들에 닿아 있다. 우리는 보통의 게임과 ‘득점기록자의 재량’
게임을 구별할 수 있는데, 이는 득점 규칙이 다른 규칙들과 마찬가지로
득점기록자가 선택을 행사해야 하는 개방적 직조(open texture)의 영역을
가지고 있음에도 불구하고, 확정된 의미의 핵심(core of settled meaning)을
가지고 있기 때문이다. 바로 이것이 득점기록자가 자유롭게 이탈할 수 없는
것이며, 그것이 미치는 한에서 선수의 득점에 관한 비공식 진술들과
득점기록자의 공식 판정들 모두에 대해 올바른 득점과 올바르지 않은 득점의
기준을 구성한다. 바로 이것이 득점기록자의 판정들은 비록 최종적이지만
무오류적이지는 않다고 말하는 것을 참으로 만든다. 법에서도 동일하다.

어느 지점까지는 득점기록자가 내린 몇몇 판정들이 명백히 그르다는 사실은
게임이 계속되는 것과 양립하지 않는 것이 아니다. 그것들은 명백히 옳은
판정들만큼이나 동등하게 산입된다(count as much). 그러나 올바르지 않은
결정들(incorrect decisions)에 대한 관용이 동일한 게임의 계속된 존재와
양립할 수 있는 정도에는 한계가 있으며, 이것은 중요한 법적 유비를 가진다.
고립적이거나 예외적인 공식적 일탈들(official aberrations)이 용인된다는
사실은 크리켓이나 야구가 더 이상 행해지고 있지 않다는 것을 의미하지
않는다. 다른 한편, 이러한 일탈들이 빈번하거나 득점기록자가 득점 규칙을
거부한다면, 어느 지점에는 반드시 선수들이 득점기록자의 일탈적 판정들을
더 이상 수용하지 않거나, 만약 그들이 수용한다면 게임이 달라진 지점이
오게 된다. 그것은 더 이상 크리켓이나 야구가 아니라 ‘득점기록자의
재량’이다. 왜냐하면 이러한 다른 게임들의 규정적 특징(defining
feature)은, 그 개방적 직조가 득점기록자에게 어떤 \textbf{(p.~145)}
재량의 폭(latitude)을 남겨두든, 일반적으로 그 결과들이 규칙의 명백한
의미(plain meaning)가 요구하는 방식으로 평가되어야 한다는 점이기
때문이다. 어떤 상상 가능한 조건에서는 실제로 행해지고 있는 게임이
‘득점기록자의 재량’이라고 말해야 할 것이다. 그러나 모든 게임들에서
득점기록자의 판정들이 최종적이라는 사실은 모든 게임들이 그것이라는 것을
의미하지 않는다.

특정 사례에서 법이 무엇인지에 대한 최종적이고 권위 있는 진술로서 법원의
결정이 가지는 고유한 지위에 근거한 규칙회의주의(rule-scepticism)의
형태를 평가할 때에는 이러한 구별들을 마음에 두어야 한다. 법의 개방적
직조는 법원들에게, 그 결정들이 법창출 선례들(law-making precedents)로
사용되지 않는 득점기록자들에게 남겨진 것보다 훨씬 더 넓고 중요한 법창출
권한(law-creating power)을 남긴다. 법원들이 무엇을 결정하든, 그것이
모두에게 명백해 보이는 규칙 부분 안에 놓인 문제이든 논쟁 가능한 경계에
놓인 문제이든, 입법에 의해 변경될 때까지 존립한다. 그리고 그 입법의
해석에 대해서도 법원들은 다시 동일한 최후의 권위 있는 목소리를 가질
것이다. 그럼에도 불구하고, 법원 체계를 세운 뒤 법은 최고 법원이
적합하다고 생각하는 것이면 무엇이든 된다고 규정하는 헌법과, 실제의 미국
헌법---또는 그 문제에 관해서는 어떤 현대 국가의 헌법---사이에는 여전히
구별이 남아 있다. ‘헌법(또는 법)은 판사들이 그것이라고 말하는 것이
무엇이든 그것이다’가 이 구별을 부정하는 것으로 해석된다면, 그것은
거짓이다. 어느 특정 시점에서 판사들, 심지어 최고 법원의 판사들도 그
규칙들이 중심에서는 올바른 사법적 결정의 기준들을 제공할 만큼 충분히
확정적인 체계의 일부이다. 이러한 기준들은 법원들에 의해, 체계 안에서
다툴 수 없는 결정들을 내릴 권위를 행사할 때 자유롭게 무시할 수 없는
것으로 간주된다. 어떤 개별 판사도 자기 직무에 올 때, 어떤 득점기록자가
자기 직무에 올 때와 마찬가지로, 의회 내 여왕의 제정물들(enactments)은
법이라는 규칙과 같은 규칙이 전통으로 확립되어 있고 그 직무 수행의
기준으로 수용되어 있음을 발견한다. 이것은 그 직무 점유자들의 창조적
활동을 허용하면서도 경계 짓는다. 이러한 기준들은 기실 그 시대 판사들
대부분이 그것들을 고수하지 않는다면 계속 존재할 수 없는데, 특정 시점에서
그것들의 존재는 단순히 올바른 재판의 기준들로 그것들이 수용되고 사용되는
데 있기 때문이다. 그러나 이것이 그 기준들을 사용하는 판사를 그 기준들의
창안자(author)로 만들거나, 호들리(Hoadly)의 언어로 말해 자기가 원하는
대로 \textbf{(p.~146)} 결정할 권능을 가진 ‘입법자(lawgiver)’로 만들지는
않는다. 기준들을 유지하기 위해 판사의 고수는 필요하지만, 판사가 그것들을
만드는 것은 아니다.

물론 사법적 결정들을 최종적이고 권위 있는 것으로 만드는 규칙들의
보호막(shield) 배후에서, 판사들이 결합하여 기존 규칙들을 거부하고 심지어
가장 명확한 의회 법률들(Acts of Parliament)조차도 자신들의 결정들에 어떤
한계도 부과하는 것으로 간주하기를 중단할 가능성은 있다. 그들의 판정들
대다수가 이러한 성격을 가지고 수용된다면, 이것은 크리켓에서
‘득점기록자의 재량’으로 게임이 전환되는 것과 평행한 체계의 변형에 해당할
것이다. 그러나 그러한 변형들의 상존 가능성은, 그 변형이 일어났다면
체계가 그러했을 바로 그것이 현재의 체계라는 것을 보여 주지 않는다. 어떤
규칙들도 위반이나 거부(repudiation)에 대해 보장될 수 없다. 인간 존재들이
그것들을 깨거나 거부하는 것은 결코 심리적으로나 물리적으로 불가능하지
않기 때문이다. 그리고 충분히 많은 사람들이 충분히 오래 그렇게 한다면, 그
규칙들은 존재를 중단할 것이다. 그러나 특정 시점에서 규칙들이 존재한다는
것은 파괴에 대한 이러한 불가능한 보장들이 있어야 한다는 것을 요구하지
않는다. 어떤 시점에 판사들이 의회 법률들(Acts of Parliament) 또는
연방의회 법률들(Acts of Congress)을 법으로 수용할 것을 요구하는 규칙이
있다고 말하는 것은, 첫째, 이 요구사항에 대한 일반적 준수(compliance)가
있고 개별 판사들의 일탈이나 거부는 드물다는 것, 둘째, 그것이 발생하거나
발생한다면 우세한 다수(preponderant majority)에 의해 심각한 비판의
대상으로, 그리고 그른(wrong) 것으로 취급되거나 취급될 것이라는 것을
수반한다. 비록 특정 사례에서 그에 따라 내려진 결정의 결과는, 결정들의
최종성에 관한 규칙 때문에, 그 결정의 올바름(correctness)이 아니라
유효성(validity)을 인정하는 입법에 의해서만 상쇄될 수 있다고 하더라도
말이다. 인간 존재들이 자기 약속들을 모두 어기는 것은 논리적으로
가능하다. 처음에는 아마도 그것이 그른 행위라는 감각을 가지고, 그다음에는
그러한 감각 없이 그럴 수 있다. 그러면 약속을 지키는 것을 의무적인 것으로
만드는 규칙은 존재를 중단할 것이다. 그러나 이것은 현재 그러한 규칙이
존재하지 않고 약속들이 실제로 구속력이 없다는 견해를 빈약하게 지지할
뿐이다. 판사들의 경우에서, 그들이 현재 체계의 파괴를 기획해 이끌어낼
가능성에 근거한 평행 논변은 이보다 더 강하지 않다.

규칙회의주의의 주제를 떠나기 전에 우리는, 규칙들이 법원들의 결정들에
대한 예측들이라는 그 적극적 논지(contention)에 관해 마지막 말을 해야
한다. 여기에 어떤 진리가 있든, 그것은 기껏해야 사적 개인들이나 그
조언자들이 제시하는 법 진술들(statements of law)에 적용될 수 있다는 점은
명백하며 자주 \textbf{(p.~147)} 언급되어 왔다. 그것은 법원들 자신의
법규칙 진술들에는 적용될 수 없다. 이것들은 일부 더 극단적인
‘현실주의자들(Realists)’이 주장했듯이 속박 없는 재량 행사를 위한 언어적
덮개(verbal covering)이거나, 아니면 법원들이 내부 관점에서 올바른 결정의
기준으로 진정으로 간주하는 규칙들의 정식화여야 한다. 다른 한편, 사법적
결정들에 대한 예측들은 법 안에서 부정할 수 없이 중요한 자리를 가진다.
개방적 직조(open texture)의 영역에 도달할 때, ‘이 문제에서 법은
무엇인가?’라는 질문에 대한 답으로 우리가 유익하게 제공할 수 있는 전부가
법원들이 무엇을 할지에 대한 조심스러운 예측인 경우가 매우 흔하다. 더
나아가, 규칙들이 요구하는 것이 모두에게 명확한 경우에도 그것에 대한
진술은 법원의 결정에 대한 예측의 형식으로 이루어지는 경우가 많다. 그러나
주로 후자의 경우에, 그리고 전자의 경우에는 다양한 정도로, 그러한 예측의
토대는 법원들이 법규칙들을 예측들로서가 아니라 결정에서 따라야 할
기준들로 간주한다는 지식이라는 점을 주목하는 것이 중요하다. 그 기준들은
그 개방적 직조에도 불구하고 법원의 재량을 배제하지는 않더라도 제한할
만큼 충분히 확정적이다. 따라서 많은 경우 법원이 무엇을 할지에 대한
예측들은, 체스 선수들이 비숍을 대각선으로 움직일 것이라고 우리가 할 수
있는 예측과 같다. 그것들은 궁극적으로 규칙들의 비예측적(non-predictive)
측면과, 그 예측들이 관련되는 자들이 기준들로서 수용한 규칙들의 내부
관점에 대한 이해에 기초한다(rest on). 이것은, 어떤 사회 집단에서
규칙들이 존재한다는 사실이 예측들을 가능하게 하고 자주 신뢰할 수 있게
만들지만 그것과 동일시될 수는 없다는, 제5장에서 이미 강조한 사실의 또
다른 측면일 뿐이다.

\subsection{4. 승인 규칙에서의
불확실성(UNCERTAINTY)}\label{uxc2b9uxc778-uxaddcuxce59uxc5d0uxc11cuxc758-uxbd88uxd655uxc2e4uxc131uncertainty}

형식주의와 규칙회의주의는 법리학적 이론(juristic theory)의
스킬라(Scylla)와 카립디스(Charybdis)이다. 그것들은 서로를 교정하는
곳에서는 유익한 거대한 과장들이며, 진리는 그 사이에 있다. 이 중도적 길을
유익한 세부 내용으로 특징짓고, 제정법(statute)이나 선례에서 법의 개방적
직조가 법원들에게 남겨 두는 창조적 기능을 수행할 때 법원들이 특징적으로
사용하는 다양한 추론 유형들을 보여 주기 위해서는, 기실 여기에서 시도될
수 없는 많은 일이 행해져야 한다. 그러나 우리는 \textbf{(p.~148)} 이
장에서, 제6장 말미에 미뤄 두었던 중요한 주제로 유익하게 돌아갈 수 있을
만큼은 말했다. 이것은 특정 법규칙들의 불확실성이 아니라 승인 규칙의
불확실성, 따라서 법원들이 유효한 법규칙들을 식별할 때 사용하는 궁극적
기준들의 불확실성에 관한 것이었다. 특정 규칙의 불확실성과 그것을 체계의
규칙으로 식별하는 데 사용되는 기준의 불확실성 사이의 구별은 그 자체로
모든 경우에 명확한 것은 아니다. 그러나 규칙들이 권위 있는 텍스트를 가진
제정법 형식의 제정물들(statutory enactments)인 경우에 그 구별은 가장
명확하다. 제정법(statute)의 문언과 그것이 특정 사례에서 요구하는 것은
완전히 명백할 수 있다. 그런데도 입법부가 이런 방식으로 입법할
권한(power)을 가지는지에 관해서는 의문들이 있을 수 있다. 때때로 이러한
의문들의 해소(resolution)는 입법 권한을 부여한 또 다른 법규칙의 해석만을
요구하며, 이 규칙의 유효성은 의심스럽지 않을 수 있다. 예컨대 종속된
권한기관(subordinate authority)이 만든 제정물(enactment)의 유효성이
문제되는 경우가 그럴 것이다. 종속된 권한기관의 입법 권한들을 규정하는
모법(parent Act of Parliament)의 의미에 관해 의문들이 발생하기 때문이다.
이것은 특정 제정법(statute)의 불확실성 또는 개방적 직조의 사례일 뿐이며
어떤 근본적 문제도 제기하지 않는다.

그러한 통상적 문제들과 구별되어야 할 것은 최고 입법부(supreme
legislature) 자체의 법적 권능(legal competence)에 관한 문제들이다.
이것들은 법적 유효성의 궁극적 기준들(ultimate criteria of legal
validity)에 관한 것이며, 최고 입법부의 권능을 특정하는 성문헌법(written
constitution)이 없는 우리 자신의 체계와 같은 법체계에서도 발생할 수
있다. 압도적 다수의 사례들에서 ‘의회 내 여왕이 제정하는 것은 무엇이든
법이다’라는 공식은 의회의 법적 권능에 관한 규칙의 적절한 표현이며,
그렇게 식별된 규칙들이 그 주변부에서는 아무리 열려 있을 수 있더라도,
법의 식별을 위한 궁극적 기준으로 수용된다. 그러나 그 의미나 범위에 관해
의문들이 발생할 수 있다. 우리는 ‘의회에 의해 제정된(enacted by
Parliament)’이라는 말이 무엇을 의미하는지 물을 수 있으며, 의문들이
발생하면 그것들은 법원들에 의해 해소될 수 있다. 한 법체계의 궁극적
규칙이 이처럼 의심스러울 수 있고 법원들이 그 의문을 해소할 수 있다는
사실로부터, 법체계 안에서 법원들의 위치에 관해 어떤 추론이 도출되어야
하는가? 그것은 법체계의 토대가 법적 유효성의 기준들을 특정하는 수용된
승인 규칙이라는 테제에 어떤 한정을 요구하는가?

\textbf{(p.~149)} 이러한 질문들에 답하기 위해 여기서는 영국
의회주권(parliamentary sovereignty) 교설의 몇몇 측면들을 검토할 것이다.
물론 유사한 의문들은 어떤 체계에서든 법적 유효성의 궁극적 기준들과
관련하여 발생할 수 있다. 법은 본질적으로 법적으로 제약받지 않는 의지의
산물이라는 오스틴주의 교설의 영향 아래에서, 오래된 헌법이론가들은 마치
주권적인 입법부가 존재해야 한다는 것이 논리적 필연성인 것처럼 썼다.
여기서 주권적이라는 것은, 그 입법부가 계속적 기구(continuing body)로
존재하는 모든 순간에 \emphksans{외부로부터(ab extra)} 부과된 법적 제한들뿐
아니라 자기 자신의 선행 입법(prior legislation)으로부터도 자유롭다는
의미이다. 이러한 의미에서 의회가 주권적이라는 점은 이제 확립된 것으로
간주될 수 있으며, 이전 의회는 자기 입법을 폐지하지 못하도록 자신의
‘승계자들(successors)’을 막을 수 없다는 원칙은 법원들이 유효한
법규칙들을 식별할 때 사용하는 궁극적 승인 규칙의 일부를 구성한다. 그러나
그러한 의회가 존재해야 한다고 명하는 것은 논리의 필연성도, 하물며 자연의
필연성도 아니라는 점을 보는 것이 중요하다. 그것은 똑같이 상상 가능한
여러 배치들 가운데 하나일 뿐이며, 우리의 체계에서 법적 유효성의 기준으로
수용되게 된 것이다. 그러한 다른 배치들 가운데에는 아마도 더 나을지도
모르게, 똑같이 ‘주권’이라는 이름을 받을 자격이 있는 또 다른 원칙이 있다.
그것은 의회가 자신의 승계자들의 입법 권능(legislative competence)을
불가역적으로 제한할 수 없는 상태여서는 \emphksans{안 되며(not)}, 반대로
이러한 더 넓은 자기제한 권한(self-limiting power)을 가져야 한다는
원칙이다. 그러면 의회는 그 역사에서 적어도 한 번은 수용되고 확립된
교설이 허용하는 것보다 훨씬 더 큰 입법 권능의 영역을 행사할 능력을
가지게 될 것이다. 의회가 그 존재의 모든 순간에, 심지어 자기 자신에 의해
부과된 것들을 포함한 법적 제한들로부터 자유로워야 한다는 요구는, 결국
법적 전능(legal omnipotence)이라는 애매한 관념의 한 해석일 뿐이다.
그것은 사실상, 잇따른 의회들(successive parliaments)의 입법 권능에
영향을 미치지 않는 모든 사항에서의 \emphksans{계속적(continuing)} 전능과, 그
행사가 오직 한 번만 향유될 수 있는 제한 없는
\emphksans{자기포섭적(self-embracing)} 전능 사이에서 선택한다. 이 두 전능
구상들은 전능한 신에 관한 두 구상에서 평행을 가진다. 한편에는 자기
존재의 모든 순간에 동일한 능력들(powers)을 누리고 따라서 그 능력들을
줄일 수 없는 신이 있고, 다른 한편에는 장래를 향해 자기 전능을 파괴할
능력(power)을 포함하는 능력들을 가진 신이 있다. 우리의 의회가 어떤
형태의 전능---계속적 전능인지 자기포섭적 전능인지---을 누리는지는 법을
식별하는 궁극적 기준으로 수용되는 규칙의 형식에 관한 경험적 질문이다.
그것은 법체계의 토대에 놓인 규칙에 관한 질문이지만, 여전히 어느 주어진
시점에는 적어도 어떤 점들에 관해 상당히 확정적인 답이 있을 수 있는
사실의 문제이다. 따라서 현재 수용된 규칙은 계속적 주권의 규칙이며,
그래서 의회는 자기 제정법들(statutes)을 폐지로부터 보호할 수 없다는 점은
명확하다.

그러나 다른 모든 규칙에서와 마찬가지로, 의회주권(parliamentary
sovereignty)의 규칙이 이 지점에서 확정적(determinate)이라는 사실은
그것이 모든 지점에서 그러하다는 것을 의미하지 않는다. 현재 명확히 옳거나
그른 답이 없는 질문들이 그것에 관해 제기될 수 있다. 이것들은 이 문제에서
결국 그 선택들에 권위가 부여되는 누군가가 하는 선택에 의해서만 확정될 수
있다. 의회주권 규칙 안의 그러한 불확정성들(indeterminacies)은 다음
방식으로 나타난다. 현행 규칙 아래에서는, 의회가 제정법(statute)에 의해
어떤 사항을 장래 의회의 입법 범위에서 불가역적으로 철회할 수 없다는 점이
인정된다. 그러나 단순히 그렇게 한다고 표방하는 제정물(enactment)과,
의회가 여전히 어떤 사항에 관해서든 입법할 가능성을 열어 두면서도 입법의
‘방식과 형식(manner and form)’을 변경한다고 표방하는 제정물 사이에는
구별이 그어질 수 있다. 후자는 예컨대 특정 쟁점들에 대해서는 양원이 함께
앉아 다수결로 통과시키거나 국민투표(plebiscite)로 확인되지 않는 한 어떤
입법도 법으로서 효력을 갖지 못한다고 요구할 수 있다. 그것은 그 조항
자체가 동일한 특별 절차에 의해서만 폐지될 수 있다는 규정으로 그러한
조항을 ‘강화 고착(entrench)’할 수 있다. 입법 과정에서의 그러한 부분적
변경은 의회가 자기 승계자들을 불가역적으로 구속할 수 없다는 현행 규칙과
충분히 양립할 수 있다. 그것이 하는 일은 승계자들을
\emphksans{구속(bind)}한다기보다는, 특정 쟁점들에 \emphksans{관한 한(quoad)}
그들을 제거하고 이 쟁점들에 관한 그들의 입법 권한들을 새로운 특별
기구(body)로 이전하는 것이기 때문이다. 따라서 이러한 특별 쟁점들과
관련하여 의회는 의회를 ‘구속’하거나 ‘족쇄를 채우거나(fettered)’ 그
계속적 전능을 감소시킨 것이 아니라, ‘의회’와 입법하기 위해 행해져야 하는
것을 ‘재규정(redefined)’했다고 말할 수 있다.

명백히 이 장치가 유효하다면, 의회는 \textbf{(p.~151)} 그것을 사용하여
의회가 자기 승계자들을 구속할 수 없다는 수용된 교설이 자신의 권한 밖에
두는 것처럼 보이는 결과들과 매우 유사한 결과들을 달성할 수 있을 것이다.
왜냐하면 의회가 입법할 수 있는 영역을 경계 짓는 것과 입법의 방식과
형식만을 변경하는 것 사이의 차이는 기실 몇몇 경우에는 충분히 명확하지만,
실제 효과상 이 범주들은 서로 점차 맞물리며 경계가 흐려지기
때문이다(shade into each other). 엔지니어들의 최저임금을 정한 뒤,
엔지니어 임금에 관한 어떤 법안도 엔지니어 노동조합(Engineers' Union)의
결의로 확인되지 않는 한 법으로서 효력을 가지지 않는다고 규정하고, 이어서
이 조항을 강화 고착한 제정법(statute)은, 임금을 ‘영원히’ 고정하고 그
폐지를 거칠게 전면 금지한 제정법(statute)이 실제로 할 수 있는 모든 것을
기실 확보할 수도 있다. 그러나 법률가들이 어느 정도 설득력을 가진 것으로
인정할 논변(argument)은, 비록 후자는 현재의 계속적 의회주권 규칙
아래에서는 효력이 없을 것이지만 전자는 그렇지 않을 것임을 보이도록
제시될 수 있다. 이 논변의 단계들은 의회가 무엇을 할 수 있는지에 관한
일련의 논지들(contentions)로 이루어지며, 각 논지는 그 앞의 논지와 일정한
유비를 가지면서도 그 앞의 논지보다 더 적은 동의를 얻을 것이다. 그중 어떤
것도 그르다고 배제될 수 없고, 옳다고 자신 있게 수용될 수도 없다. 우리는
체계의 가장 근본적인 규칙의 개방적 직조 영역 안에 있기 때문이다.
여기서는 어느 순간에도 답은 없고 오직 답들만 있는 질문이 제기될 수 있다.

따라서 의회가 상원을 완전히 폐지함으로써 의회의 현재 구성을 불가역적으로
변경할 수도 있고, 그래서 특정 입법에 대한 상원의 동의를 면제했으며 일부
권위 있는 논자들이 단지 여왕과 하원에 대한 의회 권한들 일부의 철회
가능한 위임으로 해석하기를 선호하는 1911년 및 1949년 의회법들(Parliament
Acts)을 넘어설 수 있다는 점이 인정될지도 모른다. 또한 다이시(Dicey)가
단언했듯이, \footnote{\emph{The Law of the Constitution}(10th edn.),
  p.~68 n.} 의회가 자기 권한들이 끝났음을 선언하고 장래 의회들의 선거를
마련하는 입법을 폐지하는 법률(Act)에 의해 자기 자신을 완전히 파괴할 수
있다는 점도 인정될지도 모른다. 그렇다면 의회는 맨체스터 법인(Manchester
Corporation) 같은 어떤 다른 기관에 자신의 모든 권한들을 이전하는 법률을
통해 이러한 입법적 자살(legislative suicide)을 유효하게 수반시킬 수도
있을 것이다. 의회가 이것을 할 수 있다면, 그보다 적은 것을 실효적으로 할
수는 없는가? 특정 사항들에 관해 \textbf{(p.~152)} 입법할 자신의 권한들을
끝내고, 그것들을 자기 자신과 어떤 추가 기관을 포함하는 새로운 복합
실체(composite entity)에 이전할 수는 없는가? 이러한 기초 위에서,
도미니언(Dominion)에 영향을 미치는 어떤 입법에 대해서도 도미니언의
동의를 요구한 웨스트민스터 제정법(Statute of Westminster) 제4조가,
도미니언을 위해 입법할 의회의 권한들과 관련하여 실제로 이 일을 했을 수도
있지 않은가? 이 조항이 도미니언의 동의 없이 실효적으로 폐지될 수 있다는
논지(contention)는, 상키 경(Lord Sankey)이 말했듯이 ‘현실과 아무 관계가
없는’ ‘이론’일 뿐만 아니라, 나쁜 이론일 수도 있다. 또는 적어도 반대
이론보다 나을 것이 없을 수 있다. 마지막으로, 의회가 자기 자신의 행위에
의해 이런 방식들로 재구성될 수 있다면, 왜 그것은 엔지니어
노동조합(Engineers' Union)이 특정 유형의 입법에서 필수적인 동의
요소(consenting element)가 되도록 규정함으로써 자기 자신을 재구성할 수
없는가?

이 논변에서 의심스럽지만 명백히 잘못된 것은 아닌 단계들을 구성하는 몇몇
의문스러운 명제들이, 언젠가 그 문제를 결정하도록 요청받은 법원에 의해
승인되거나 거부될 가능성은 충분히 있다. 그러면 우리는 그것들이 제기하는
질문들에 대한 답을 갖게 될 것이며, 그 답은 체계가 존재하는 한, 주어질 수
있는 답들 가운데 유일한 권위 있는 지위를 가질 것이다. 법원들은 이
지점에서 유효한 법을 식별하는 궁극적 규칙을 확정적인 것으로 만들었을
것이다. 여기서 ‘헌법은 판사들이 말하는 것이다’는 말은 최고
법원들(supreme tribunals)의 특정 결정들이 다투어질 수 없다는 것만을
의미하지 \emphksans{않는다}. 처음 보기에는 그 광경이 역설적으로 보인다.
여기에는 법원들이, 판사로서 자신들에게 관할권(jurisdiction)을 부여하는
바로 그 법들의 유효성 자체가 그 기준들에 의해 검사되어야 하는 궁극적
기준들을 확정하는 창조적 권한들을 행사하고 있기 때문이다. 어떻게 헌법이
헌법이 무엇인지를 말할 권위를 부여할 수 있는가? 그러나 모든 규칙이 어떤
지점들에서는 의심스러울 수 있지만, 모든 규칙이 모든 지점에서 의심에 열려
있지는 않다는 것이 기실 법체계가 존재하기 위한 필요조건이라는 점을
기억하면 이 역설은 사라진다. 어느 주어진 시점에 법원들이 궁극적 유효성
기준들에 관한 이러한 한계 질문들(limiting questions)을 결정할 권위를
가질 가능성은, 단지 그 시점에는 그 권위를 부여하는 규칙들을 포함하여
광대한 법 영역에 대한 그 기준들의 적용이 아무 의문도 제기하지 않는다는
사실에 의존한다. 비록 그 기준들의 정확한 범위(scope)와 적용범위(ambit)에
대해서는 의문이 제기되더라도 말이다.

\textbf{(p.~153)} 그러나 이 답은 어떤 이들에게는 그 문제를 너무 쉽게
처리하는 것처럼 보일 수 있다. 그것은 법적 유효성의 기준들을 특정하는
근본 규칙들(fundamental rules)의 가장자리에서 법원들이 하는 활동을 매우
부적절하게 특징짓는 것처럼 보일 수 있다. 그 이유는 아마도 그것이 이
활동을, 불확정적인 것으로 드러난 특정 제정법(statute)을 해석하면서
법원들이 창조적 선택을 행사하는 통상적 사례들에 지나치게 가깝게
동화시키기 때문일 수 있다. 그러한 통상적 사례들은 어떤 체계에서든
발생해야 한다는 점은 명확하며, 그래서 법원들이 그 사례들을, 제정법이
열어 둔 대안들 사이에서 선택함으로써 확정할 관할권을 가진다는 점은,
그들이 이 선택을 발견(discovery)으로 위장하기를 선호하더라도, 법원들이
근거하여 행위하는 규칙들의 일부, 비록 암묵적 일부일 뿐이더라도, 명백히
일부인 것처럼 보인다. 그러나 적어도 성문헌법(written constitution)이
없는 경우에는, 유효성의 근본 기준들에 관한 질문들은 종종 이렇게 사전에
염두에 둘 수 있는 성질을 가지지 \emphksans{않는(not)} 것처럼 보인다. 그리고
바로 이러한 사전에 염두에 둘 수 있는 성질이 있어야, 그러한 질문들이
발생할 때 법원들이 이미 기존 규칙들 아래에서 이 종류의 질문들을 확정할
명확한 권위를 가지고 있다고 말하는 것이 자연스럽다.

‘형식주의적(formalist)’ 오류의 한 형태는 어쩌면 법원이 취하는 모든
단계가 그것을 취할 권위를 사전에 부여하는 어떤 일반 규칙에 의해 포괄되어
있어서, 법원의 창조적 권한들이 \emphksans{언제나(always)} 위임된 입법
권한(delegated legislative power)의 한 형태라고 생각하는 바로 그것일 수
있다. 진실은, 법원들이 가장 근본적인 헌법 규칙들에 관한 사전에 염두에
두어지지 않았던 질문들을 확정할 때, 그 질문들이 발생하고 결정이 내려진
뒤에야 그 질문들을 결정할 권위가 수용되게 된다는 것일 수 있다. 여기서는
성공하는 것은 오직 성공뿐이다(all that succeeds is success). 문제된
헌법적 질문이 사회를 너무 근본적으로 분열시켜 사법적 결정에 의한 처리를
허용하지 않을 수도 있다는 것은 상상 가능하다. 1909년 남아프리카법(South
Africa Act)의 고착 조항들(entrenched clauses)에 관한 남아프리카의
쟁점들은 한때 법적 처리를 하기에는 너무 분열적일 위험이 있었다. 그러나
덜 중요한 사회적 쟁점들이 관련된 곳에서는, 법의 바로 그 원천들에 관한
매우 놀라운 사법적 법창출(judicial law-making)도 조용히
‘삼켜질(swallowed)’ 수 있다. 그러한 경우에는,
\emphksans{회고적으로(retrospect)} 법원들이 자신들이 한 일을 할 ‘내재적’
권한(inherent power)을 언제나 가지고 있었다고 자주 말해질 것이고,
진정으로 그렇게 보일 수도 있다. 그러나 그 유일한 증거가 행해진 일의
성공뿐이라면, 이것은 경건한 허구(pious fiction)일 수 있다.

영국 법원들이 선례의 구속력(binding force)에 관한 규칙들을 조작한 것은,
\textbf{(p.~154)} 아마도 이 마지막 방식으로, 권한들을 확보하여
행사하려는 성공적인 시도(successful bid)로서 서술되는 것이 가장 정직할
것이다. 여기서 권한(power)은 성공으로부터 \emphksans{사후적으로(ex post
facto)} 권위(authority)를 획득한다. 따라서 형사항소법원(Court of
Criminal Appeal)이 \emph{Rex} 대 \emph{Taylor} 사건 \footnote{{[}1950{]}
  2 KB 368.}에서 신민의 자유에 관한 사항들에 대해 자기 자신의 선례들에
구속되지 않는다고 판정할 권위를 가진다고 판시하기 전에는, 그 법원이
그러한 권위를 가지는지는 전적으로 열려 있는 것으로 보였을 수 있다.
그러나 그 판정은 내려졌고, 지금은 법으로서 따라지고 있다. 법원이 이런
방식으로 판정할 내재적 권한을 언제나 가지고 있었다는 진술은, 분명 상황을
실제보다 더 정돈되어 보이게 하는 방식일 뿐일 것이다. 여기, 바로 이러한
근본적인 것들의 가장자리에서는, 우리는 규칙회의주의자를 환영해야 한다.
다만 그는 자신이 환영받는 곳이 가장자리라는 점을 잊지 않아야 하며,
법원들이 가장 근본적인 규칙들을 이렇게 눈에 띄게 발전시킬 수 있게 하는
것은 상당 부분 법의 광대한 중심 영역들에서 의심할 여지 없이
규칙-지배적인(rule-governed) 작용들을 해 온 데서 법원들이 모은
위신(prestige)이라는 사실을 우리에게 보지 못하게 해서는 안 된다.

\newpage

\subsection{제7장 주석}\label{uxc81c7uxc7a5-uxc8fcuxc11d}


{\footnotesize
\noteentry 125쪽. \notetitle{규칙의 예시적 전달(Communication of rules by examples).} 이러한 용어로 선례(precedent)의 사용을 특징지은 설명으로는 Levi의 “An
Introduction to Legal Reasoning” 제1절, \emph{University of Chicago Law
Review} 제15권(1948년)을 보라. Wittgenstein은 \emph{Philosophical
Investigations} (특히 제1부 제208--238절)에서 규칙을 가르치고 따르는
일에 관한 관념들(notions)에 대해 여러 중요한 통찰을 제시한다.
Wittgenstein에 대한 논의는 Winch, \emph{The Idea of a Social Science}
제24--33쪽, 91--93쪽 참조.

\noteentry 128쪽. \notetitle{언어적으로 정식화된 규칙들의 개방적 직조(Open texture of verbally formulated rules).} 개방적 직조(open texture)라는 관념에
대해서는 Waismann의 “Verifiability”, \emph{Essays on Logic and
Language} 제1권(Flew 편집), 제117--130쪽을 보라. 법적 추론과 관련한
논의로는 Dewey, “Logical Method and Law,” \emph{Cornell Law Quarterly}
제10권(1924); Stone, \emph{The Province and Function of Law}, 제6장;
Hart, “Theory and Definition in Jurisprudence,” \emph{Proceedings of
the Aristotelian Society} 보충권 제29권(1955), 제258--264쪽;
“Positivism and the Separation of Law and Morals,” \emph{Harvard Law
Review} 제71권(1958), 제606--612쪽 참조.

\noteentry 129쪽. \notetitle{형식주의(formalism)와 개념주의(conceptualism).} 법
문헌들에서 이 용어들과 거의 유의어로 사용되는 표현으로는 ‘기계적 또는
자동적 법리학(mechanical or automatic jurisprudence)’, ‘구상들의
법리학(jurisprudence of conceptions)’, ‘논리의 과도한 사용(excessive use
of logic)’ 등이 있다. 이와 관련하여 Pound, “Mechanical Jurisprudence,”
\emph{Columbia Law Review} 제8권(1908); \emph{Interpretations of Legal
History}, 제6장 참조. 이러한 용어들로 정확히 어떤 결함 또는 폐해가
지칭되는지는 항상 명확하지 않다. 이에 관해서는 Jensen, \emph{The Nature
of Legal Argument}, 제1장, 그리고 Honoré의 서평, \emph{Law Quarterly
Review} 제74권(1958), 제296쪽; Hart, 위에 언급한 \emph{Harvard Law
Review} 제71권, 제608--612쪽을 참조.

\noteentry 131쪽. \notetitle{법적 기준들과 구체적 규칙들(Legal standards and specific rules).} 이러한 두 가지 법적 통제 형태의 성격과 상호 관계에 대한 가장
통찰력 있는 일반적 논의는 Dickinson, \emph{Administrative Justice and
the Supremacy of Law}, 제128--140쪽에 있다.

\noteentry 131쪽. \notetitle{행정 규칙제정에 의해 구체화되는 법적 기준들(Legal standards implemented by administrative rule-making).} 미국에서는 Interstate
Commerce Commission, Federal Trade Commission 등 연방 규제기관들이 ‘공정
경쟁(fair competition)’, ‘공정하고 합리적인 요율(just and reasonable
rates)’ 등의 광범위한 법적 기준들을 구체화하는 규칙들을 제정한다.
(Schwartz, \emph{An Introduction to American Administrative Law},
제6--18쪽, 33--37쪽 참조.) 영국에서도 이와 유사한 규칙제정 기능이
행정부에 의해 수행되지만, 일반적으로 미국에서 익숙한 이해관계자들에 대한
형식적 준사법적 청문(quasi-judicial hearing) 절차 없이 이루어진다.
예컨대, 1957년 \emph{Factories Act} 제46조에 근거한 복지규정(Welfare
Regulations), 같은 법 제60조에 의한 건축규정(Building Regulations)을
참조. 또한, 1947년 \emph{Transport Act}에 따라 이의제기자들(objectors)의
의견을 청취한 후 ‘요금 체계(charges scheme)’를 확정할 수 있는
\emph{Transport Tribunal}의 권한은 미국식 모델에 보다 근접해 있다.

\noteentry 132쪽. \notetitle{주의 기준들(Standards of care).} 주의의무(duty of care)의
구성요소들에 대한 통찰력 있는 분석으로는 Learned Hand 판사의 판결문
\emph{US v. Carroll Towing Co.} (1947), 159 F 2nd 169, 173쪽을 보라.
일반적 기준들을 구체적 규칙들로 대체하는 바람직성에 대해서는 Holmes,
\emph{The Common Law}, 제3강, 제111--119쪽 참조. 이 견해에 대한 비판은
Dickinson, 앞의 책, 제146--150쪽에 제시되어 있다.

\noteentry 133쪽. \notetitle{구체적 규칙들에 의한 통제(Control by specific rules).} 유연한 기준들(flexible standards)보다는 엄격하고 고정된 규칙들(hard and
fast rules)이 적절한 통제 형태가 되는 조건에 대해서는 Dickinson, 앞의
책, 제128--132쪽, 제145--150쪽을 보라.

\noteentry 134쪽. \notetitle{선례와 법원의 입법 활동(Precedent and the legislative activity of Courts).} 영국법에서의 선례(precedent in English law)의
현대적 일반 설명으로는 R. Cross, \emph{Precedent in English Law}
(1961)을 보라. 본문에서 언급된 ‘협소화 과정(narrowing process)’의 잘
알려진 예시는 \emph{L. \& S. W. Railway Co.} v. \emph{Gomm} (1880), 20
Ch.D. 562로, 이는 \emph{Tulk} v. \emph{Moxhay} (1848), 2 Ph. 774에서의
법리를 좁혀 적용하였다.

\noteentry 136쪽. \notetitle{규칙회의주의의 다양상들(Varieties of rule-scepticism).} 이
주제에 대한 미국 법학의 논의는 하나의 논쟁(debate)으로 읽을 때 특히
유익하다. 예컨대 \emph{Law and the Modern Mind} (특히 제1장 및 부록 2,
“Notes on Rule Fetishism and Realism”)에서의 Frank, \emph{The Bramble
Bush}에서의 Llewellyn의 논증은 다음과 같은 저작들을 배경으로 고려할 수
있다: Dickinson, “Legal Rules: Their Function in the Process of
Decision,” \emph{University of Pennsylvania Law Review} 제79권(1931);
“법 배후의 법(The Law behind the Law),” \emph{Columbia Law Review}
제29권(1929); “The Problem of the Unprovided Case,” \emph{Recueil
d'Études sur les sources de droit en l'honneur de F. Geny}, 제11권
제5장; Kantorowicz, “Some Rationalism about Realism,” \emph{Yale Law
Review} 제43권(1934).

\noteentry 139쪽. \notetitle{실망한 절대주의자로서의 회의주의자(The sceptic as a disappointed absolutist).} Miller, “Rules and Exceptions,”
\emph{International Journal of Ethics} 제66권(1956)을 보라.

\noteentry 140쪽. \notetitle{규칙의 직관적 적용(Intuitive application of rules).} Hutcheson, “The Judgement Intuitive”; “The Function of the `Hunch' in
Judicial Decision,” \emph{Cornell Law Quarterly} 제14권(1928)을 보라.

\noteentry 141쪽. \notetitle{‘헌법은 판사들이 말하는 바 그대로이다(The constitution is what the judges say it is).’} 이 표현은 미국의 Hughes 대법원장(Chief
Justice Hughes)에게서 유래한 것으로 보이며, Hendel, \emph{Charles Evan
Hughes and the Supreme Court} (1951), 제11--12쪽 참조. 그러나 Hughes
자신의 저작인 \emph{The Defence Court of the United States} (1966년판)
제37, 41쪽에서는 헌법 해석에 있어 개인의 정치적 견해로부터 독립된
판사들의 책무(duty)를 강조하고 있다.

\noteentry 149쪽. \notetitle{의회 주권에 대한 대안적 분석(Alternative analyses of the sovereignty of Parliament).} H. W. R. Wade, “The Basis of Legal
Sovereignty,” \emph{Cambridge Law Journal} (1955)을 보라. 이 견해는
Marshall, \emph{Parliamentary Sovereignty and the Commonwealth}, 제4장과
제5장에서 비판받고 있다.

\noteentry 149쪽. \notetitle{의회 주권과 신의 전능성(Parliamentary sovereignty and divine omnipotence).} Mackie, “Evil and Omnipotence,” \emph{Mind},
1955년, 제211쪽 참조.

\noteentry 150쪽. \notetitle{의회를 구속하는 것인가, 재규정하는 것인가(Binding or redefining Parliament).} 이 구분에 관해서는 Friedmann, “Trethowan's
Case, Parliamentary Sovereignty and the Limits of Legal Change,”
\emph{Australian Law Journal} 제24권(1950); Cowen, “Legislature and
Judiciary,” \emph{Modern Law Review} 제15권(1952) 및 제16권(1953);
Dixon, “The Law and the Constitution,” \emph{Law Quarterly Review}
제51권(1935); Marshall, 앞의 책, 제4장을 보라.

\noteentry 151쪽. \notetitle{1911년 및 1949년 의회법들(Parliament Acts 1911 and 1949).} 이들 법률이 위임입법(delegated legislation)의 한 형식을 수권하는 것으로
해석하는 견해는 H. W. R. Wade, 위 저작 및 Marshall, 위 저작 제44--46쪽
참조.

\noteentry 152쪽. \notetitle{1931년 웨스트민스터 제정법(Statute of Westminster), 제4조.} 권위 있는 견해들의 무게는, 이 조항의 제정이 도미니언의 동의 없이
도미니언을 위해 입법할 권한의 불가역적 종결을 구성할 수는 없었다는
견해를 지지한다. \emph{British Coal Corporation} v. \emph{The King}
(1935), AC 500; Wheare, \emph{The Statute of Westminster and Dominion
Status}, 제5판, 제297--298쪽; Marshall, 위 저작 제146--147쪽 참조. 이에
반하여, “한 번 부여된 자유는 철회될 수 없다(Freedom once conferred
cannot be revoked)”는 견해는 남아프리카 법원에서 \emph{Ndlwana} v.
\emph{Hofmeyr} (1937), AD 229, 제237쪽에서 표현된 바 있다.

\subsection{제7장 제3판
주석}\label{uxc81c7uxc7a5-uxc81c3uxd310-uxc8fcuxc11d}

\noteentry 124--128쪽. \notetitle{법의 개방적 직조(Open texture of law).} 일반적인
논의로는 Avishai Margalit, “Open Texture” (1979), \emph{Meaning and
Use} 제3권 141쪽을 보라. 하트(Hart)가 ‘개방적 직조(open texture)’ 개념을
어떻게 사용하는지에 대해서는 Brian Bix, \emph{Law, Language, and Legal
Determinacy} (Oxford University Press, 1993), 제1장을 참조하라. 법의
불확정성(indeterminacy)의 다른 원천들에 대해서는 Kelsen, \emph{Pure
Theory of Law}, 제348--356쪽을 참조. 불확정성과 그에 따른
재량(discretion)에 대한 드워킨(Dworkin)의 비판은 \emph{Taking Rights
Seriously} 제31--39쪽, 제4장, 제13장과 \emph{A Matter of Principle}
(Harvard University Press, 1986), 제5장에서 볼 수 있다. 드워킨의 이러한
비판은 David Brink, “Legal Theory, Legal Interpretation and Judicial
Review” (1988), \emph{Philosophy \& Public Affairs} 제17권 105쪽 및
Nicos Stavropoulos, “Hart's Semantics” in Jules Coleman 편,
\emph{Hart's Postscript} 중 특히 제88--98쪽에서 지지를 얻는다. 반면에
Andrei Marmor, \emph{Interpretation and Legal Theory}, 제2판 (Hart
Publishing, 2005), 제7장에서는 이 비판을 수용하지 않는다.

\noteentry 128--134쪽. \notetitle{일반 규칙들과 특정 사례들(General rules and particular cases).} 규칙 기반 결정(rule-based decision-making)의 본질에 대해서는
David Lyons, \emph{Forms and Limits of Utilitarianism} (Oxford
University Press, 1965), Frederick Schauer, \emph{Playing By the Rules:
A Philosophical Examination of Rule-Based Decision-Making in Law and in
Life} (Oxford University Press, 1991), 특히 제5장, 그리고 Larry
Alexander와 Emily Sherwin, \emph{The Rule of Rules: Morality, Rules, and
Dilemmas of Law} (Duke University Press, 2001), 제1장과 제2장을
참조하라.

\noteentry 128--129쪽. \notetitle{전범적 사례들(Paradigm cases).} 전범적 사례들이
수행하는 역할에 대한 하트의 이해를 옹호하는 논의로는 Timothy Endicott,
\emph{Vagueness in Law} (Oxford University Press, 2000), 제7장을 보라.
이에 반대하는 드워킨의 견해는 \emph{Law's Empire}, 제3장, 특히
제90--94쪽에 제시되어 있다. 이에 대한 응답으로는 Endicott, “Herbert
Hart and the Semantic Sting” in Jules Coleman 편, \emph{Hart's
Postscript} (Oxford University Press, 2001)을 보라.

\noteentry 129--130쪽. \notetitle{형식주의(Formalism).} ‘형식주의(formalism)’라는 용어는
다른 방식으로도 사용된다. 관련 논의로는 Martin Stone, “Formalism” in
Jules Coleman and Scott Shapiro 편, \emph{The Oxford Handbook of
Jurisprudence and Philosophy of Law} (Oxford University Press, 2002)을
참조하라.

\noteentry 130--132쪽. \notetitle{불확정성과 재량(Indeterminacy and discretion).} 법은,
이해 가능한 법적 질문에 대해 \emphksans{하나의 유일한(unique)} 답을
정당화하지 못할 때, 즉 법이 불완전할 때 불확정적(indeterminate)이다.
이는 어떤 법적 사례가 ‘어렵다(hard)’는 것과 동일하지 않다. 후자는 유일한
답이 존재하지만 그 답이 분명하지 않거나 이견의 대상이 되는 경우에도 참일
수 있다. 또한 이는 법이 실제로(인과적으로) 사법적 결정들을 결정짓는다는
견해와도 다르다. 법은 전적으로 확정적일 수 있지만 판사들이 그것을 알지
못하거나, 알고도 적용하지 않기로 결정할 수 있다. 법의 확정성(legal
determinacy)은 결정의 예측가능성(decisional predictability)을 함의하지
않는다. (비교: Brian Leiter, “Legal Indeterminacy” (1995), \emph{Legal
Theory} 제1권 481쪽) 하트가 옹호하는 것은 법의 정당화적
불확정성(justificatory indeterminacy)이다. 하트의 테제에 대한 거부는
Ronald Dworkin, “No Right Answer?” in P. M. S. Hacker and Joseph Raz
편, \emph{Law, Morality and Society}, 그리고 Ronald Dworkin, “On Gaps
in the Law” in Paul Amselek and Neil MacCormick 편, \emph{Controversies
About Law's Ontology} (Edinburgh University Press, 1991)에서 볼 수 있다.
이에 대한 옹호는 Joseph Raz, \emph{The Authority of Law}, 제4장에서 찾을
수 있다.

\noteentry 131--133쪽. \notetitle{규칙들과 가변적 기준들(Rules and variable standards).} 이 지점에서 하트(Hart)는 규칙(rule)을 두어야 하는지
\emphksans{여부(whether)}를 논의하는 것이 아니라, 다양한 상황에 어떤
\emphksans{종류(kind)}의 규칙이 적절할지를 논의하고 있다. 미국의 법률가들은,
헨리 하트(Henry Hart, H. L. A. Hart와는 무관)와 앨버트 삭스(Albert
Sacks)의 영향을 받아, 때때로 ‘규칙들(rules)’과 ‘기준들(standards)’을
구분하지만, H. L. A. Hart에게 기준들(standards)은 규칙들의 한 종류이다.
이 구분에 대해서는 Henry M. Hart and Albert Sacks, \emph{The Legal
Process: Basic Problems in the Making and Application of Law} (W. N.
Eskridge, Jr.~and P. P. Frickey 편집, Foundation Press, 1994)
139--141쪽, 그리고 Pierre J. Schlag, “Rules and Standards” (1985)
\emph{UCLA Law Review} 제33권 379쪽을 참조하라. 마찬가지로, 흔히
‘원칙들(principles)’이라 불리는 것들도 하트의 규칙들에 관한 관념(idea of
rules)에 포함되며, 그는 드워킨(Dworkin)이 \emph{Taking Rights Seriously}
제22--28쪽에서 설정한 것과 같은 규칙과 원칙 간의 범주적 구분을 거부한다.
이에 대한 하트의 입장은 \emph{Postscript} 제261--263쪽에 언급되어 있다.
관련 논의로는 Joseph Raz, “Legal Principles and the Limits of Law”
(1972) \emph{Yale Law Journal} 제81권 823쪽, Larry Alexander and Ken
Kress, “Against Legal Principles” (1997) \emph{Iowa Law Review} 제82권
739쪽을 참조하라.

\noteentry 134--135쪽. \notetitle{선례(Precedent).} 이 주제에 관한 문헌은 매우 풍부하다.
참고문헌으로는 다음을 보라: A. W. B. Simpson, “The Common Law and Legal
Theory” in A. W. B. Simpson 편, \emph{Oxford Essays in Jurisprudence};
Ronald Dworkin, \emph{Taking Rights Seriously} 제110--123쪽; Neil
MacCormick, \emph{Legal Reasoning and Legal Theory} (수정판, Oxford
University Press, 1994); Laurence Goldstein 편, \emph{Precedent in Law}
(Oxford University Press, 1987); Stephen Perry, “Judicial Obligation,
Precedent and the Common Law” (1987) \emph{Oxford Journal of Legal
Studies} 제7권 215쪽; Frederick Schauer, “Precedent” (1987)
\emph{Stanford Law Review} 제39권 571쪽; Susan Hurley, “Coherence,
Hypothetical Cases, and Precedent” (1990) \emph{Oxford Journal of Legal
Studies} 제10권 221쪽; Larry Alexander, “Precedent” in Dennis
Patterson 편, \emph{A Companion to the Philosophy of Law and Legal
Theory} (Blackwell, 1996); Neil Duxbury, \emph{The Nature and Authority
of Precedent} (Cambridge University Press, 2008).

\noteentry 136--141쪽. \notetitle{규칙회의주의의 다양상들(Varieties of rule scepticism).} 하트의 관련 에세이로는 “American Jurisprudence through
English Eyes: The Nightmare and the Noble Dream” (\emph{Essays in
Jurisprudence and Philosophy}, 제4장)이 있다. 하트의 저작 이후 등장한
가장 저명한 규칙회의주의자(rule-sceptics)는 미국의 ‘비판법학(Critical
Legal Studies, CLS)’ 운동의 법학자들이었다. 철학적으로는 덜 정교하지만
영향력 있는 논의로는 Roberto Unger, \emph{The Critical Legal Studies
Movement} (Harvard University Press, 1986), 특히 제1--14쪽, 그리고
Duncan Kennedy, \emph{A Critique of Adjudication {[}fin de siècle{]}}
(Harvard University Press, 1997)를 들 수 있다. 초기 CLS 문헌에 대한
비판적 논의는 Ken Kress, “Legal Indeterminacy” (1989) \emph{California
Law Review} 제77권 283쪽을 참조하라. 또한 John Finnis, “On the Critical
Legal Studies Movement”, 그의 \emph{Philosophy of Law} 제13장, Andrew
Altman, \emph{Critical Legal Studies: A Liberal Critique} (Princeton
University Press, 1990)도 보라. Brian Leiter는 “Legal Realism and Legal
Positivism Reconsidered”, 그의 \emph{Naturalizing Jurisprudence}
(Oxford University Press, 2007) 제2장에서 하트의 규칙회의주의 다양상들을
다룬다.
}

\setcounter{footnote}{0}
\input{sections/chapter_8-translation.tex}
\setcounter{footnote}{0}
\section{IX. 법과 도덕 (Laws and
Morals)}\label{ix.-uxbc95uxacfc-uxb3c4uxb355-laws-and-morals}

\subsection{1. 자연법과 법실증주의 (Natural Law and Legal
Positivism)}\label{uxc790uxc5f0uxbc95uxacfc-uxbc95uxc2e4uxc99duxc8fcuxc758-natural-law-and-legal-positivism}

법과 도덕 사이에는 다양한 유형의 관계가 존재하며, 그 중 어느 하나를
\emphksans{그들 사이의} 관계로 특정하여 연구 대상으로 삼는 것은 유익하지
않다. 대신, 법과 도덕이 관련되어 있다고 주장하거나 부인할 때 의미될 수
있는 다양한 사태들을 구분하는 것이 중요하다. 때로는 거의 누구도 부인한
적 없는 일종의 연결성을 주장하는 경우도 있으며, 이러한 논박의 여지 없는
사실이 보다 의심스러운 연결성의 징표로 잘못 수용되거나, 심지어 그것과
혼동되기도 한다. 예컨대, 언제 어디서나 법의 발전이 특정 사회 집단의
관행적 도덕(moral convention)과 이상(ideal), 그리고 당시 통용되는 도덕을
초월한 개인들의 계몽된 도덕적 비판(enlightened moral criticism)에 의해
심대한 영향을 받아왔다는 점은 진지하게 반박될 수 없다. 그러나 이 진리를
부당하게 이용하여 전혀 다른 명제를 정당화하는 근거로 삼는 것도 가능하다:
즉, 법체계는 도덕 혹은 정의와의 어떤 특정한 준수/합치(a specific
conformity)를 반드시 드러\emphksans{내어야 하며}(\emph{must} exhibit), ~또는
그에 복종할 도덕적 의무가 있다는 널리 퍼진 확신(widely diffused
conviction)위에 반드시 기초\emphksans{해야 한다}(\emph{must} rest on)는
명제이다. 물론 이 주장은 어떤 의미에서는 참일 수도 있으나, 특정 법체계
내에서 사용되는 개별 법률의 법적 유효성의 기준이 반드시 명시적으로든
암묵적으로든 도덕 혹은 정의에 대한 참조(a reference to)를 포함해야
한다는 점은 이로부터 도출되지 않는다.

법과 도덕의 관계에 관해서는 이 밖에도 많은 다른 문제들이 제기될 수 있다.
이 장에서는 그 가운데 두 가지만을 논의할 것인데, 비록 두 문제 모두 다른
여러 문제들에 대한 일정한 고찰을 수반하겠지만, 여기서는 그 두 문제에
한정한다. 첫 번째는, 오늘날에는 법과 도덕에 관한 서로 다른 여러 논제를
가리키는 명칭으로 사용되게 되었음에도 불구하고, 여전히 자연법(Natural
Law)과 법실증주의(Legal Positivism) 사이의 쟁점으로 묘사하는 것이 유익할
수 있는 문제이다. 여기서 우리는 법실증주의를, \textbf{(p.~186)} 법이
일정한 도덕적 요구를 재현하거나 충족해야 한다는 것이 어떠한 의미에서도
필연적 진리(necessary truth)는 아니라는 단순한 주장으로 이해하겠다. 물론
실제로는 법이 그러한 요구를 충족해 온 경우가 종종 있었지만 말이다.
그러나 이러한 견해를 취해 온 이들이 도덕의 성격에 관해서는 침묵했거나,
또는 서로 크게 상이한 견해를 제시해 왔기 때문에, 법실증주의가 거부되어
온 두 가지 매우 상이한 형태를 검토할 필요가 있다. 그중 하나는 자연법의
고전적 이론들에서 가장 분명하게 표현되어 있는데, 이는 인간 이성이 발견을
기다리고 있는 일정한 인간 행위의 원리들이 존재하며, 인간이 제정한
법(man-made law)이 유효성(validity)을 가지기 위해서는 그러한 원리들에
합치해야 한다는 주장이다. 다른 하나는 도덕에 대해 이와는 다른, 덜
합리주의적인 관점을 취하면서, 법적 유효성이 도덕적 가치와 연결되는
방식에 대해 상이한 설명을 제시한다. 이 절과 다음 절에서는 이 가운데 첫
번째 입장을 검토할 것이다.

플라톤(Plato)에서 오늘날에 이르기까지, 인간이 어떻게 행위해야 하는지가
인간 이성(human reason)에 의해 발견될 수 있다는 명제의 언명(assertion)과
그 부정을 위해 헌신된 방대한 문헌 속에서, 논쟁의 한쪽에 선 이들은 다른
쪽에 대해 ``이것을 보지 못한다면 당신들은 눈이 멀었다''고 말하는 듯
보이며, 이에 대해 상대방은 ``당신들은 꿈을 꾸고 있다''고 응답하는 것처럼
보인다. 이러한 상황이 발생하는 이유는, 올바른 행위의 참된 원리들이
이성적으로 발견 가능하다는 주장이 보통 독립된 하나의 교리로 제시된 것이
아니라, 무생물과 생물을 포괄하는 자연(nature)에 대한 일반적 관념의
일부로서 처음 제시되었고, 오랫동안 그 일부로서 옹호되어 왔기 때문이다.
이러한 관점은 여러 측면에서 현대 세속적 사유(modern secular thought)의
틀을 구성하는 일반적 자연관과 정면으로 대립한다. 그 결과 자연법
이론(Natural Law theory)은 비판자들에게는 현대 사유가 이미 승리적으로
벗어났다고 여겨지는 깊고 오래된 혼란들로부터 비롯된 것으로 보이는 반면,
옹호자들에게는 비판자들이 피상적인 사소함에만 집착하면서 보다 심원한
진리들을 무시하고 있는 것처럼 보이게 된다.

따라서 많은 현대의 비판자들은, 적절한 행위의 법칙(laws of proper
conduct)이 인간 이성(human reason)에 의해 발견될 수 있다는 주장이
`법(law)'이라는 단어의 단순한 애매성(ambiguity)에 근거하고 있으며, 이
애매성이 드러나는 순간 자연법(Natural Law)은 치명타를 입었다고 생각해
왔다. 존 스튜어트 밀(John Stuart Mill)이 『법의 정신(Esprit des Lois)』
제1장에서, 별과 같은 무생물이나 동물들이 `그들의 자연의 법(law of their
nature)'에 복종하는데 왜 인간은 그러하지 않고 죄(sin)에 빠지는가를
순진하게 묻는 몽테스키외(Montesquieu)를 다룬 방식이 바로 그러하다. 밀은
이것이 \textbf{(p.~187)} 자연의 진행이나 정기성(regularities)을
정식화하는 법칙들과, 인간에게 특정한 방식으로 행동할 것을 요구하는
법칙들 사이의 오래된 혼동을 드러낸다고 보았다. 전자는 관찰과 추론에 의해
발견될 수 있으며 `기술적(descriptive)'이라고 불릴 수 있고, 이를 발견하는
것은 과학자의 역할이다. 반면 후자는 그러한 방식으로 확립될 수 없는데,
그것들은 사실에 대한 진술이나 기술이 아니라 인간이 특정한 방식으로
행동해야 한다는 `처방(prescriptions)' 또는 요구이기 때문이다. 따라서
몽테스키외의 질문에 대한 대답은 단순하다. 처방적 법칙(prescriptive
laws)은 위반될 수 있으며, 그럼에도 여전히 법칙으로 남아 있는데, 이는
단지 인간이 지시받은 대로 행동하지 않았다는 것을 의미할 뿐이다. 그러나
과학에 의해 발견된 자연법칙들에 대해 그것들이 위반될 수 있는지 여부를
말하는 것은 무의미하다. 만약 별들이 그들의 규칙적인 운동을 기술한다고
주장하는 과학적 법칙과는 다른 방식으로 행동한다면, 그 법칙들은 `위반된'
것이 아니라, `법칙'이라 불릴 자격을 상실하고 다시 정식화되어야 한다.
이러한 `법(law)'이라는 말의 의미 차이에 대응하여, `must', `bound to',
`ought', 'should'와 같은 관련 어휘들에서도 체계적인 차이가 나타난다.
그러므로 이러한 관점에서 보면, 자연법에 대한 신념은 매우 단순한 오류로
환원된다. 즉, 법적 함축을 지닌 이러한 단어들이 가질 수 있는 매우 상이한
의미들을 인식하지 못한 데서 비롯된 오류라는 것이다. 이는 마치 어떤
사람이 ``당신은 군복무를 위해 신고해야 한다(You are bound to report for
military service)''와 ``바람이 북쪽으로 돌면 얼게 될 것이 확실하다(It is
bound to freeze if the wind goes round to the north)''라는 문장에서 해당
표현들이 지니는 전혀 다른 의미를 인식하지 못한 것과 같은 경우이다.

벤담(Bentham)과 밀(Mill)과 같은 비판자들, 즉 자연법(Natural Law)을 가장
격렬하게 공격한 이들은, 법(law)의 이러한 상이한 의미들 사이에서
상대방들이 보인 혼동을 자연의 관찰된 정기성(regularities)이 우주의 신적
통치자(Divine Governor)에 의해 처방되거나 명령되었다고 믿는 신념의
잔존에 돌리는 경우가 많았다. 이러한 신정적(theocratic) 관점에서는,
중력의 법칙과 십계명---인간을 위한 신의 법(God's law for Man)---사이의
유일한 차이는, 블랙스톤(Blackstone)이 언명했듯이, 창조된 존재들 가운데
오직 인간만이 이성과 자유의지를 부여받았다는 비교적 사소한 점에
불과했다. 그 결과 인간은 사물들과 달리 신적 처방(divine prescriptions)을
인식할 수 있고 또한 그것을 불복종할 수도 있다는 것이다. 그러나 자연법은
항상 우주의 신적 통치자나 입법자(Divine Governor or Lawgiver)에 대한
신념과 결부되어 온 것은 아니며, 설령 그러한 신념과 결부된 경우라
하더라도 그 고유한 신조(tenets)들이 논리적으로 그 신념에 의존해 온 것은
아니다. 자연법에 포함되는 '자연적(natural)'이라는 단어의 관련된 의미와,
\textbf{(p.~188)} 처방적 법(prescriptive laws)과 기술적 법(descriptive
laws) 사이의 차이를---현대인의 눈에는 매우 명백하고
중요한---최소화하려는 그 일반적 관점은, 이러한 목적과 관련해서는
전적으로 세속적이었던 그리스 사상(Greek thought)에 그 뿌리를 두고 있다.
실제로 어떤 형태의 자연법 교설이 지속적으로 재언명되어 온 이유는
부분적으로, 그것의 호소력이 신적 권위와 인간적 권위 모두로부터
독립적이라는 점, 그리고 오늘날에는 거의 받아들여질 수 없는 용어법과 많은
형이상학을 수반하고 있음에도 불구하고, 도덕성과 법을 이해하는 데 중요한
몇 가지 기본적 진실(elementary truths)을 포함하고 있다는 점에 있다.
우리는 이러한 진실들을 그 형이상학적 틀로부터 분리해 내어, 보다 단순한
용어로 여기에서 다시 진술하고자 한다.

현대의 세속적 사유에 따르면, 무생물과 생물, 동물과 인간의 세계는 일정한
유형의 사건들과 변화들이 반복적으로 발생하며, 그 안에서 일정한 정기적
연결관계(regular connections)가 예시되는 장면이다. 이러한 연결관계들
가운데 적어도 일부는 인간에 의해 발견되고 자연법칙(laws of nature)으로
정식화되었다. 이러한 현대적 관점에서 자연을 이해한다는 것은, 자연의 어떤
부분에 대하여 이러한 정기성(regularities)에 관한 지식을 적용하는 것을
의미한다. 물론 위대한 과학 이론들의 구조가 관찰 가능한 사실, 사건, 또는
변화를 어떤 단순한 방식으로 그대로 반영하는 것은 아니다. 실제로 그러한
이론들의 상당 부분은 관찰 가능한 사실과 직접적으로 대응되는 대상이 없는
추상적인 수학적 정식화(abstract mathematical formulations)로 이루어져
있다. 그럼에도 불구하고 이러한 이론들이 관찰 가능한 사건들과 변화들과
연결되는 지점은, 이러한 추상적 정식화들로부터 관찰 가능한 사건들을
참조하며, 그러한 사건들에 의해 확증되거나 반증될 수 있는 일반화들이
연역될 수 있다는 데에 있다. 따라서 과학 이론이 자연에 대한 우리의 이해를
증진시킨다고 언명할 수 있는 근거는, 궁극적으로는 무엇이 발생할 것인지를
예측할 수 있는 그 이론의 능력에 달려 있으며, 이 능력은 정기적으로
발생하는 것들에 대한 일반화에 기초한다. 중력의 법칙과 열역학 제2법칙은,
현대적 사유에서 볼 때, 관찰 가능한 현상들의 정기성에 관하여 제공하는
정보 덕분에, 단순한 수학적 구성물에 그치지 않고 자연법칙(laws of
nature)으로 간주된다.

자연법 교설은, 관찰 가능한 세계를 단지 그러한 정기성(regularities)의
장면으로만 보지 않고, 자연에 대한 앎을 그 정기성에 관한 앎으로만
이해하지 않는, 보다 오래된 자연 개념의 일부이다. 오히려 이러한 오래된
관점에 따르면, 인간적 존재이든, 생물적 존재이든, 무생물적 존재이든, 이름
붙일 수 있는 모든 종류의 존재하는 것은 단지 자기 자신의 존재를
유지하려는 경향을 지닐 뿐만 아니라, \textbf{(p.~189)} 그것에 고유한
선(good)---곧 그것에 적합한 \emphksans{목적(end)}---에 해당하는 일정한 최적
상태(definite optimum state)를 향해 나아가는 것으로 파악된다.

이는 자연이 그 자체 안에 사물들이 실현해 나가는 탁월성의 단계(levels of
excellence)를 포함하고 있다는 목적론적 자연관(teleological conception of
nature)이다. 어떤 특정한 종류에 속하는 사물이 그 고유하거나 적합한
목적(specific or proper end)에 이르기까지 거치는 단계들은
정기적(regular)이며, 사물의 특징적인 변화 방식, 작용 방식, 또는 발전
방식을 서술하는 일반화로 정식화될 수 있다; 이러한 점에서 목적론적
자연관은 현대적 사고와 중첩된다. 차이는, 목적론적 관점에서는 사물들에게
정기적으로 일어나는 사건들이 \emphksans{단지(merely)} 정기적으로 발생하는
것으로만 이해되지 않으며, 그것들이 \emphksans{실제로(do)} 정기적으로
발생하는지 여부, 그것들이 발생\emphksans{해야 하는지(should)} 여부, 또는
그것들이 발생하는 것이 \emphksans{좋은지(good)} 여부라는 질문들이 서로 분리된
질문으로 간주되지 않는다는 데 있다. 오히려('우연(chance)'에 귀속되는
몇몇 드문 기형적 사례들을 제외하면), 일반적으로 발생하는 것은 해당
사물의 적합한 목적이나 목표를 향해 나아가는 하나의 단계로 제시됨으로써,
설명될 수 있을 뿐만 아니라 선한 것 또는 발생해야 할 것으로 평가될 수
있다. 그러므로 사물의 발전에 관한 법칙들은, 그것이 어떻게 마땅히
행동하거나 변화해야 하는지와, 그것이 실제로 어떻게 정기적으로 행동하거나
변화하는지를 모두 보여주어야 한다.

이와 같은 자연에 대한 사고방식은 추상적으로 진술될 경우 낯설게 보인다.
그러나 오늘날에도 적어도 생명체에 관해서 우리가 여전히 사용하는 몇몇
표현 방식을 상기한다면, 그것은 덜 기이하게 보일 수도 있다. 왜냐하면
목적론적 관점(teleological view)은 여전히 그들의 발전을 서술하는
일상적인 방식들 속에 반영되어 있기 때문이다. 예컨대 도토리(acorn)의
경우, 참나무(oak)로 성장하는 것은 도토리들이 정기적으로 달성하는 어떤
상태일 뿐만 아니라, 또한 (마찬가지로 정기적으로 일어나는)
부패(decay)와는 달리, 성숙의 최적 상태(optimum state of maturity)로서
구별된다. 그리고 이 최적 상태에 비추어 중간 단계들은 설명되며, 선한지
나쁜지로 평가되고, 그 다양한 부분들의 `기능(functions)'과 구조적
변화들이 식별된다. 잎의 정상적인 성장은 `완전한(full)' 또는
`적합한(proper)' 발전에 필요한 수분을 얻기 위해 요구되며, 그러한 수분을
공급하는 것이 잎의 '기능'이다. 따라서 우리는 이러한 성장을 '자연적으로
발생해야 할 것(ought naturally to occur)'으로 사고하고 말한다. 반면
무생물의 작용이나 운동의 경우에는, 그것들이 인간에 의해 어떤 목적을 위해
설계된 인공물(artefacts)이 아닌 한, 이러한 방식의 말하기는 훨씬 설득력이
떨어진다. 돌이 땅으로 떨어질 때 어떤 적합한 \textbf{(p.~190)}
'목적(end)'을 실현하고 있거나, 마치 말이 마구간으로 달려 돌아가듯이
자신의 '고유한 자리(proper place)'로 되돌아가고 있다는 관념은,
오늘날에는 다소 희극적으로 느껴진다.

실로, 자연에 대한 목적론적 관점(teleological view of nature)을 이해하는
데에는 하나의 난점이 있는데, 그것은 이 관점이 정기적으로 일어나는 바에
대한 진술과 일어나야 할 바에 대한 진술 사이의 차이를 최소화하듯이, 또한
현대 사상에서 그토록 중요한 차이로 여겨지는, 스스로의 목적을
\emphksans{가지고(with)} 그것을 의식적으로 실현하려고 노력하는 인간과, 다른
생명체나 무생물 사이의 차이 역시 최소화하기 때문이다. 왜냐하면 세계에
대한 목적론적 관점에서는 인간 역시 다른 사물들과 마찬가지로, 그에게
설정된 특정한 최적 상태(optimum state) 또는 목적(end)을 향해 나아가는
존재로 이해되며, 그가 다른 사물들과 달리 이를 의식적으로 수행할 수
있다는 사실은 그를 자연의 나머지 부분과 근본적으로 구별하는 차이로
간주되지 않기 때문이다. 이러한 특정한 인간의 목적 또는 선(good)은
부분적으로는 다른 생명체들과 마찬가지로 생물학적 성숙의 상태와 발달된
신체적 능력의 조건을 포함하지만, 또한 사유(thought)와 행위(conduct)에
나타나는 정신과 인격의 발달과 탁월성이라는, 인간에게 특유한 요소를
포함한다. 다른 사물들과 달리 인간은 추론(reasoning)과 성찰(reflection)을
통해 이러한 정신과 인격의 탁월성을 달성하는 것이 무엇을 의미하는지를
발견하고, 그것을 욕구(desire)할 수 있다. 그럼에도 불구하고, 이 목적론적
관점에 따르면 이러한 최적 상태는 인간이 그것을 욕구하기 때문에 그의
선이나 목적이 되는 것이 아니라, 오히려 그것이 이미 그의 자연적
목적(natural end)이기 때문에 인간이 그것을 욕구하는 것이다.

다시 말해, 이러한 목적론적 관점의 상당 부분은 우리가 인간을 사고하고
말하는 몇몇 방식 속에 여전히 살아 있다. 그것은 우리가 어떤 것들을 인간의
\emphksans{욕구(needs)}로 식별하여 그것을 충족시키는 것이 \emphksans{좋은(good)}
것이라고 여기거나, 인간에게 행해지거나 인간에 \emphksans{의해(by)} 겪어지는
어떤 것들을 \emphksans{해악(harm)} 또는 \emphksans{상해(injury)}로 규정하는 데
잠재되어 있다. 따라서 일부 사람들은 죽기를 원하기 때문에 먹거나 쉬기를
거부할 수도 있다는 점이 사실이기는 하지만, 우리는 먹고 쉬는 것을 단지
인간이 정기적으로 하는 행위이거나 우연히 욕구하는 것 이상으로 생각한다.
음식과 휴식은, 어떤 이들이 그것이 필요한 때에도 이를 거부한다 하더라도,
인간의 욕구이다. 그러므로 우리는 모든 인간이 먹고 잠자는 것이
자연스럽다고 말할 뿐만 아니라, 모든 인간은 때때로 먹고 쉬어야 한다고,
혹은 이러한 행위를 하는 것이 자연적으로 좋다고 말한다. 이러한 인간
행위에 대한 판단에서 `자연적으로(naturally)'라는 말의 힘은, 그것들을
단순한 관행(convention)이나 인간의 규정(prescription)을 반영하는
판단(`모자를 벗어야 한다')---그 내용은 \textbf{(p.~191)} 사유(thought)나
성찰(reflection)에 의해 발견될 수 없는 것---과도 구별하고, 또한 어떤
특정한 목적을 달성하기 위해 요구되는 것이 무엇인지를 지시할 뿐인 판단,
즉 어느 한 시점에는 한 사람이 우연히 가질 수도 있고 다른 사람은 갖지
않을 수도 있는 목적에 관한 판단과도 구별하는 데 있다. 동일한 관점은 신체
기관의 \emphksans{기능(functions)}에 대한 우리의 개념과, 그것들을 단순한
인과적 속성(causal properties)과 구별하여 긋는 경계에도 나타난다. 우리는
심장의 기능이 혈액을 순환시키는 것이라고 말하지만, 암성 종양(cancerous
growth)의 기능이 죽음을 초래하는 것이라고 말하지는 않는다.

인간 행위에 대한 일상적 사고 속에 여전히 살아 있는 목적론적 요소들을
설명하기 위해 제시된 이러한 소박한 예들은, 인간이 다른 동물들과 공유하는
생물학적 사실이라는 낮은 영역에서 취해진 것이다. 이 사고방식과 표현을
이해 가능하게 만드는 것이 전적으로 자명한 어떤 것이라는 점은 정당하게
지적될 것이다. 그것은 곧 인간 활동의 적절한 목적(end)이
생존(survival)이라는 묵시적 가정이며, 이는 대부분의 인간이 대부분의 시간
동안 자신의 존재를 계속 유지하기를 바란다는 단순한 우연적
사실(contingent fact)에 근거한다. 우리가 자연적으로 행하는 것이 좋다고
말하는 행위들은 생존에 요구되는 것들이며, 인간의 욕구(human need),
해악(harm), 그리고 신체 기관이나 변화의 \emphksans{기능(function)}이라는
개념들도 동일한 단순한 사실에 기초한다. 물론 여기에서 멈춘다면 우리는
자연법(Natural Law)의 매우 희석된 버전만을 얻게 될 것이다. 왜냐하면 이
관점의 고전적 옹호자들은 생존(\emph{自己存在保存(perseverare in esse
suo)})을 인간의 목적(end) 또는 인간에게 선(good)으로서의 훨씬 더
복잡하고 훨씬 더 논쟁적인 개념을 이루는 가장 낮은 층위로만 이해했기
때문이다. 아리스토텔레스(Aristotle)는 여기에 인간 지성의 무욕적
계발(disinterested cultivation of the human intellect)을 포함시켰고,
아퀴나스(Aquinas)는 신에 대한 앎(knowledge of God)을 포함시켰는데, 이
둘은 모두 도전될 수 있었고 실제로 도전되어 온 가치들이다. 그러나 다른
사상가들, 그중에서도 홉스(Hobbes)와 흄(Hume)은 시야를 더 낮추는 데
기꺼이 동의해 왔다. 그들은 생존이라는 소박한 목적 속에서 자연법의
용어법에 경험적 건전성(empirical good sense)을 부여하는 중심적이고
논쟁의 여지가 없는 요소를 보았다. ``인간의 본성은 어떤 경우에도 개인들의
결합(association) 없이 존속할 수 없으며, 그러한 결합은 형평(equity)과
정의(justice)의 법칙들에 대한 고려가 전혀 없다면 결코 성립할 수 없었을
것이다.''\footnote{Hume, \emph{Treatise of Human Nature}, m. ii, `Of
  Justice and Injustice'.}

이 단순한 사유는 사실 법과 도덕 둘 모두의 성격과 매우 깊은 관련을
가지며, \textbf{(p.~192)} 인간에게 있어 목적(end) 또는 선(good)이 하나의
특정한 삶의 방식으로 제시되는 일반적인 목적론적 관점 가운데서, 실제로
사람들 사이에 깊은 불일치가 존재할 수 있는 보다 논쟁적인 부분들로부터
분리해낼 수 있다. 더 나아가 우리는 생존(survival)을 언급함에 있어,
그것이 인간의 고유한 목적이나 목표이기 때문에 인간이 필연적으로 욕구하는
것으로 선험적으로 고정되어 있다는 관념을, 현대적 사고에 비추어 지나치게
형이상학적인 것으로서 배제할 수 있다. 대신 우리는, 일반적으로 인간이
살아 있기를 욕구한다는 것이 달리 그러하지 않을 수도 있었던 하나의 우연적
사실(contingent fact)에 불과하며, 생존을 인간의 목적 또는 목표라고 부를
때에도 인간이 실제로 그것을 욕구한다는 것 이상을 의미하지 않을 수도
있다고 볼 수 있다. 그럼에도 불구하고, 우리가 이러한 상식적인 방식으로
생존을 이해하더라도, 생존은 여전히 인간의 행위와 그에 대한 우리의
사유에서 특별한 지위를 차지하며, 이는 정통적 자연법(Natural Law)의
정식화들에서 생존에 부여되는 두드러진 중요성과 필연성에 상응한다. 이는
단지 압도적 다수의 인간이 극심한 고통을 대가로 치르더라도 살아 있기를
원한다는 사실 때문만이 아니라, 이러한 사실이 우리가 세계와 서로를
기술하는 데 사용하는 사유와 언어의 전체 구조 속에 반영되어 있기
때문이다. 우리는 살아 있으려는 일반적인 욕구를 제거한 채로
위험(danger)과 안전(safety), 해악(harm)과 이익(benefit), 필요(need)와
기능(function), 질병(disease)과 치료(cure)와 같은 개념들을 온전히 유지할
수 없다. 왜냐하면 이러한 개념들은, 하나의 목표로 받아들여진 생존에
기여하는 바에 참조하여 사물을 동시에 기술하고 평가하는 방식들이기
때문이다.

그러나 이러한 것들보다 더 단순하고 덜 철학적인 고려들도 존재하는데, 이는
인간의 법과 도덕에 대한 논의와 보다 직접적으로 관련되는 의미에서, 생존을
하나의 목적(aim)으로 받아들이는 것이 필수적임을 보여준다. 우리는 논의의
용어들 자체가 전제하는 어떤 것으로서 이미 이에 관여(commit)되어 있다.
왜냐하면 우리의 관심사는 자살 클럽의 사회적 장치가 아니라, 지속적인
존재를 위한 사회적 장치들이기 때문이다. 우리는 이러한 사회적 장치들
가운데, 이성에 의해 발견 가능한 자연법(natural laws)으로서 조명될 수
있는 것들이 있는지, 그리고 그것들이 인간의 법과 도덕과 어떤 관계에
있는지를 알고자 한다. 인간들이 \emphksans{어떻게} 함께 살아야 하는지에 관한
이 문제 또는 그 밖의 어떤 문제를 제기하기 위해서는, 일반적으로 말해
그들의 목적이 살아가는 데 있다는 점을 가정하지 않으면 안 된다. 이
지점에서부터 논증은 단순하다. 인간의 본성과 인간이 살아가는 세계에 관한
몇 가지 매우 명백한 일반화들---사실상 자명한 진술들(truisms)---에 대한
성찰은, \textbf{(p.~193)} 그것들이 유효한 한, 어떤 사회 조직이 존속
가능하기 위해서는 반드시 포함해야 할 일정한 행위 규칙들이 존재함을
보여준다. 이러한 규칙들은 실제로, 사회적 통제가 서로 다른 형태로 구별될
정도로 발전한 모든 사회들의 법과 관행적 도덕(conventional morality) 속에
공통 요소를 이루고 있다. 이와 함께, 법과 도덕 양자 모두에서, 특정 사회에
고유한 많은 요소들과, 임의적이거나 단순한 선택의 문제로 보일 수 있는
많은 요소들도 발견된다. 인간, 그들의 자연적 환경, 그리고 목적에 관한
기초적 진실들에 근거를 둔 이러한 보편적으로 승인된 행위 원칙들은, 그
이름 아래 흔히 제시되어 온 보다 웅대하고 더 논쟁적인 구성들과 대비하여,
자연법(Natural Law)의 \emphksans{최소 내용(minimum content)}으로 간주될 수
있다. 다음 절에서는, 이 겸손하지만 중요한 최소 내용이 의존하고 있는 인간
본성의 두드러진 특성들을 다섯 가지 자명한 명제(truisms)의 형태로 검토할
것이다.

\subsection{2. 자연법의 최소 내용 (The Minimum Content of Natural
Law)}\label{uxc790uxc5f0uxbc95uxc758-uxcd5cuxc18c-uxb0b4uxc6a9-the-minimum-content-of-natural-law}

여기에서 제시하는 단순한 자명한 이치들(truisms)과 그것들이 법과 도덕과
맺는 연관을 고찰함에 있어, 각 경우마다 언급된 사실들이 생존을 하나의
목적(aim)으로 전제할 때 법과 도덕이 특정한 내용을 포함해야 하는
\emphksans{이유(reason)}를 제공한다는 점을 관찰하는 것이 중요하다. 논증의
일반적 형태는 단순히, 그러한 내용이 없다면 법과 도덕은 인간들이 서로
결합하여 살아가는 데서 갖는 최소한의 목적, 즉 생존이라는 목적을 증진시킬
수 없다는 것이다. 이러한 내용이 결여될 경우, 인간은 있는 그대로의
상태에서 어떠한 규칙도 자발적으로 준수할 이유를 갖지 못할 것이며, 또한
규칙에 복종하고 이를 유지하는 것이 자신의 이익에 부합한다고 판단하는
사람들에 의해 자발적으로 제공되는 최소한의 협력이 없다면, 자발적으로
따르지 않으려는 다른 사람들에 대한 강제(coercion)는 불가능할 것이다.
이러한 접근에서 자연적 사실들과 법적·도덕적 규칙의 내용 사이에 존재하는
독특하게 합리적인 연결(rational connection)을 강조하는 것은 중요하다.
왜냐하면 자연적 사실들과 법적 또는 도덕적 규칙 사이에는 전혀 다른 형태의
연결이 존재하는지에 대해서도 탐구하는 것이 가능하고 또한 중요하기
때문이다. 예컨대 아직 젊은 과학인 심리학과 사회학은, 일정한
물리적·심리적·경제적 조건들이 충족되지 않는 한---예컨대 어린 아이들이
\textbf{(p.~194)} 가족 내에서 특정한 방식으로 먹여지고 양육되지 않는
한---어떠한 법 체계나 도덕 규범도 성립할 수 없거나, 혹은 오직 특정한
유형에 부합하는 법들만이 성공적으로 기능할 수 있다는 사실을 발견할 수도
있고, 이미 발견했을 수도 있다. 이러한 종류의 자연적 조건들과 규칙 체계
사이의 연결은 \emphksans{이유(reasons)}에 의해 매개되지 않는다. 왜냐하면
그것들은 특정 규칙들의 존재를, 그 규칙들의 주체들이 의식적으로 갖는
목적이나 목표와 연관시키지 않기 때문이다. 유아기에 특정한 방식으로
양육되는 것이 어떤 인구 집단이 도덕적 또는 법적 규범을 발전시키거나
유지하는 데 필요한 조건, 심지어 하나의 \emphksans{원인(cause)}임이 입증될
수는 있겠지만, 그것이 그들이 그렇게 하는 \emphksans{이유}는 아니다. 이러한
인과적 연결들은 목적이나 의식적 목표에 기초한 연결들과 물론 충돌하지
않는다. 오히려 그것들은 자연법(Natural Law)이 출발점으로 삼는 인간의
의식적 목적이나 목표가 왜 형성되는지를 실제로 설명해 줄 수 있다는
점에서, 후자보다 더 중요하거나 더 근본적인 것으로 간주될 수도 있다.
이러한 유형의 인과적 설명들은 자명한 이치들(truisms)에 기초하지도
않으며, 의식적 목표나 목적에 의해 매개되지도 않는다. 그것들은 관찰에
기초한 일반화와 이론의 방법, 그리고 가능한 경우에는 실험에 의존하는
방식으로, 사회학이나 심리학과 같은 다른 과학들이 확립해야 할 문제이다.
그러므로 이러한 연결들은, 다음에 제시되는 자명한 명제들에 언급된
사실들과 특정한 법적·도덕적 규칙들의 내용을 연관시키는 연결과는 다른
종류의 것이다.

 (i) \emphksans{인간의 취약성} (Human vulnerability) 법과 도덕의 공통된
요구사항들은 대부분 적극적 봉사가 아니라 자제(forbearance)의 형태로
나타나며, 이는 대개 부정적 명령, 즉 금지(prohibition)로 표현된다.
사회생활에 있어 가장 중요한 것은 타인에게 신체적 위해를 가하거나
살해하는 데에 폭력을 사용하는 것을 제한하는 규칙들이다. 이러한
규칙들의 기본적 성격은 다음과 같은 질문을 통해 드러난다. 만약 이런
규칙이 없다면, 우리 같은 존재들이 \emphksans{다른 종류의} 어떤 규칙이라도
가질 이유가 있을까? 이 수사적 질문의 힘은 인간이 때때로 타인을
신체적으로 공격하는 경향이 있으며, 일반적으로는 타인의 공격에
취약하다는 사실에 기반한다. 이는 상식적 진술(truism)이지만, 필연적
진리(necessary truth)는 아니다. 상황은 달랐을 수도 있고, 언젠가는
달라질 수도 있다. 어떤 동물 종들은 외골격(exoskeleton)이나 딱딱한
등껍질(carapace)과 같은 신체 구조를 지니고 있어, 동종의 다른 개체의
공격에 사실상 면역이거나, 공격 수단이 전혀 없는 경우도 있다. 만약
인간이 \textbf{(p.~195)} 서로에게 신체적으로 취약하지 않게 된다면,
법과 도덕의 가장 대표적인 규정 가운데 하나인 \emphksans{살인하지 말라}(Thou
shalt not kill)는 조항의 명백한 이유도 사라질 것이다.

 (ii) \emphksans{근사적 동등성} (Approximate equality) 인간은 신체적 힘, 민첩성,
그리고 더욱이 지적 능력(intellectual capacity)에서 서로 다르다.
그럼에도 불구하고, 인간 법제와 도덕의 다양한 형태를 이해함에 있어
결정적으로 중요한 사실은, 어떤 개인도 타인들을 장기간에 걸쳐 협력 없이
지배하거나 제압할 수 있을 정도로 압도적으로 강하지 않다는 점이다. 가장
강한 사람조차도 때때로 잠을 자야 하며, 수면 중에는 일시적으로 그의
우월성을 상실한다. 이와 같은 \emphksans{근사적 동등성}(approximate
equality)의 사실은, 상호 자제(mutual forbearance)와 절충(compromise)의
체계가 필요하다는 점---법적·도덕적 의무의 기반이 되는 점---을 무엇보다
명확히 보여준다. 이러한 자제를 요구하는 규칙들과 함께 살아가는
사회생활은 때로는 불편하고 성가시기도 하지만, 이와 같은 대략적으로
평등한 존재들에게 있어 그것은 제한 없는 공격 상태에서의 삶보다는 덜
악랄하고, 덜 야만적이며, 덜 짧다. 물론 이와 양립하는 또 하나의 상식적
진술은, 이러한 자제 체계가 수립되면 언제나 이를 악용하려는 자들이
존재한다는 것이다. 즉, 그들은 동시에 이 체계의 보호 아래 살면서도 그
제약을 깨뜨리려 한다. 실제로 이는, 뒤에서 설명하겠지만, 단순한 도덕적
통제로부터 조직화된 법적 통제로 이행하는 것이 필연적인 이유가 되는
자연적 사실 가운데 하나이다. 다시 말해, 사태는 다르게 전개될 수도
있었을 것이다. 인간이 대략적으로 평등한 대신, 어떤 사람들은 현존 평균
수준보다 훨씬 강하거나, 휴식을 거의 취하지 않아도 되는 존재일 수도
있었고, 또는 대부분의 사람들이 평균보다 훨씬 열등했을 수도 있다.
이러한 예외적 존재들에게는, 타인과 상호 자제하거나 타협하는 것보다
공격을 통해 더 많은 것을 얻을 가능성이 컸을 것이다. 그러나 우리는
거인들이 소인들 사이에서 사는 환상적인 장면을 상상할 필요도 없다.
왜냐하면 \emphksans{근사적 동등성}의 결정적 중요성은 국제 사회의
사실들---국가 간의 힘과 취약성의 극단적 불균형---에서 오히려 더 잘
드러나기 때문이다. 후술하겠지만, 국제법(international law)의 구성
단위들 사이의 이러한 불평등은, 그것이 국내법(municipal law)과 매우
다른 성격을 갖게 만든 요인 중 하나이며, 국제법이 조직화된 강제
체계(coercive system)로 기능할 수 있는 범위를 제한해왔다.
\textbf{(p.~196)}

 (iii) \emphksans{제한된 이타성} (Limited altruism) 인간은 서로를 말살하려는
욕망에 지배당한 악마가 아니며, 단지 생존이라는 소박한 목적만을
가정하더라도, 법과 도덕의 기본 규칙들이 필수적이라는 점을 입증하는
것은, 인간이 본질적으로 이기적이며 타인의 생존과 복지에 대해
무관심하다는 잘못된 견해와 동일시되어서는 안 된다. 그러나 인간이
악마가 아니라면, 동시에 천사도 아니다. 인간이 이 두 극단 사이의
중간자라는 사실은 상호 자제 체계가 \emphksans{필요하며 가능}하다는 점을
보여준다. 타인에게 해를 끼칠 유혹조차 느끼지 않는 천사들 사이에서는,
자제를 요구하는 규칙이 필요하지 않을 것이다. 반면에, 자기 희생을
개의치 않고 파괴를 일삼는 악마들 사이에서는 그러한 규칙이 아예 성립할
수 없다. 실제 세계에서 인간의 이타성은 범위가 제한적이고 지속적이지
않으며, 공격적 성향은 충분히 자주 발현되어 통제되지 않으면 사회생활을
파괴할 수 있다.

 (iv) \emphksans{제한된 자원} (Limited resources) 인간은 음식, 의복, 거처를
필요로 하며, 이들은 무한히 풍부하게 존재하지 않고, 자연으로부터
경작하거나 획득하거나 인간의 노동으로 만들어져야 한다는 사실은 단지
우연적인(contingent) 사실일 뿐이다. 그러나 이 사실들만으로도, 어떤
최소 형태의 재산 제도(property institution)는 필수적이며(비록 그것이
반드시 개인 재산일 필요는 없더라도), 이를 존중할 것을 요구하는 고유한
규칙 유형이 필요하다. 가장 단순한 재산 개념은 다음과 같은 규칙에서
찾아볼 수 있다. 즉, `소유자'(owner)가 아닌 자들의 토지 출입 또는 사용,
혹은 물질적 사물의 취득 및 사용을 배제하는 규칙이다. 작물이 자라기
위해서는 토지가 무차별적 출입으로부터 보호되어야 하고, 식량은 수확
또는 포획과 소비 사이의 기간 동안 타인에 의해 빼앗기지 않아야 한다.
언제 어디서든, 인간의 삶은 이러한 최소한의 자제 행위에 의존한다. 물론
이 점 또한 상황이 달랐을 수도 있다. 인간 유기체가 식물처럼 대기 중에서
음식물을 추출할 수 있었거나, 인간이 필요로 하는 것이 경작 없이 무한히
자라나는 방식으로 구성되어 있었다면 말이다.

지금까지 논의한 규칙들은 \emphksans{정태적} 규칙(\emph{static} rules)이다.
즉, 그것들이 부과하는 의무 및 그 의무의 귀속은 개인에 의해 변형될 수
없는 의미에서 그러하다. 그러나 최소 규모 집단을 제외한 모든 집단은,
충분한 자원을 확보하기 위해 노동의 분업(division of labour)을 발전시켜야
하며, 이는 개별 구성원이 의무를 생성하거나 그 귀속을 변경할 수 있도록
하는 \textbf{(p.~197)} \emphksans{동태적} 규칙(\emph{dynamic} rules)을 필요로
한다. 그 예로는 생산물을 이전, 교환, 또는 판매할 수 있게 하는 규칙들이
있다. 이러한 거래들은, 재산의 가장 단순한 형태를 정의하는 초기 권리와
의무의 귀속을 변경할 수 있는 능력을 전제하기 때문이다. 동일하게 불가피한
노동 분업과 지속적인 협력의 필요성은, 약속이 의무의 원천(source of
obligation)으로 인정되도록 하는 다른 유형의 동태적 규칙들을 필수적인
것으로 만든다. 이러한 장치를 통해, 개인은 말(구두든 문서든)을 통해
자신을 일정한 방식으로 행동할 의무를 지닌 존재로 만들 수 있으며, 그 의무
불이행 시 비난이나 처벌의 대상이 될 수 있다. 이타성이 무제한이 아니라면,
그러한 자기구속적 행위(self-binding act)를 제도적으로 보장하는 절차가
요구된다. 이는 타인의 미래 행위에 대한 최소한의 신뢰를 가능케 하고,
협력에 필수적인 예측 가능성(predictability)을 확보하기 위한 것이다. 이는
특히 교환되거나 공동으로 기획되는 것이 상호 서비스이거나, 재화(goods)가
동시에 또는 즉각적으로 제공되지 않는 경우에 가장 필요하다.

 (v) \emphksans{제한된 이해력과 의지력} (Limited understanding and strength of
will) 인격(persons), 재산(property), 약속(promises)을 존중하는
규칙들이 사회생활에서 필요하다는 사실은 단순하며, 그 상호적 이익도
명백하다. 대부분의 사람들은 이를 인식할 수 있으며, 이러한 규칙을
준수하는 데 요구되는 단기적 이익의 희생을 감수할 수 있다. 실제로
사람들은 다양한 동기에 따라 복종할 수 있다. 어떤 이는 그러한 희생이
그에 상응하는 이익을 가져온다는 타산적 계산(prudential
calculation)에서, 또 어떤 이는 타인의 복지에 대한 사욕없는
관심(disinterested interest)에서, 또 어떤 이는 해당 규칙 자체가 존중할
가치가 있다고 보고, 그것에 헌신하는 데서 이상(ideals)을 찾기 때문이다.
그러한 이해에 기초한 복종 동기의 실효성(efficacy)을 좌우하는 의지의
강도나 선함은 모든 사람에게 동등하게 분포되어 있지 않다. 누구나 때때로
자신의 즉각적 이익을 우선하려는 유혹을 받으며, 이를 탐지하고 처벌하기
위한 특수한 조직이 없다면, 많은 사람들이 그 유혹에 굴복하게 될 것이다.
물론 상호 자제의 이점이 워낙 명백하기 때문에, \textbf{(p.~198)} 강제적
체계(coercive system) 내에서 자발적으로 협력하려는 사람들의 수와 그
힘은, 일반적으로 범법자들의 가능한 조합보다 클 것이다. 그럼에도
불구하고, 아주 소규모의 밀접하게 조직된 사회를 제외하고는, \emphksans{억제
체계(system of restraints)}에 대한 복종은, 그 의무(obligations)에
복종하지 않으면서 그 이점만을 얻고자 하는 자들을 강제할 조직이 없다면,
어리석은 일이 될 것이다. 따라서 `제재'(sanctions)는 통상적인 복종의
동기로 요구되는 것이 아니라, 자발적으로 복종하려는 자들이 복종하지
않으려는 자들에게 희생되지 않도록 \emphksans{보장하는 장치(guarantee)}로서
요구된다. 이러한 장치 없이 복종한다는 것은 도태될 위험을 감수하는
일이다. 이와 같은 항구적인 위험이 존재하는 조건에서, 이성이 요구하는
것은 \emphksans{강제적} 체계 내에서의 \emphksans{자발적} 협력(\emph{voluntary}
co-operation in a \emph{coercive} system)이다.

이 지점에서 주목해야 할 것은, 인간 사이의 근사적 동등성(approximate
equality)이라는 동일한 자연적 사실이, 조직화된 제재가 효과를 발휘하는 데
결정적으로 중요하다는 점이다. 만약 어떤 이들이 다른 이들보다 훨씬 더
강력하여 타인의 자제에 의존하지 않아도 된다면, 범법자들의 힘이 법과
질서를 지지하는 이들의 힘을 능가할 수 있다. 이러한 불평등이 존재할 경우,
제재의 사용은 성공할 수 없으며, 오히려 그것이 억제하려는 위험만큼이나 큰
위험을 초래할 수 있다. 이러한 조건에서는, 사회생활이 소수 범법자에 대해
간헐적으로 힘을 사용하는 상호 자제 체계에 기반하기보다는, 약자들이
가능한 최상의 조건 하에서 강자에게 굴복하고, 그들의 `보호' 하에 살아가는
체계만이 유일하게 존속 가능하다. 이러한 체계는 자원의 희소성(scarcity of
resources)으로 인해 다수의 상호 충돌하는 권력 중심들이 생기게 하며, 각
중심은 자신만의 `강자'(strong man)를 중심으로 결집하게 된다. 이들은
때때로 상호 간 전쟁을 벌일 수도 있지만, 패배라는 자연적 제재(natural
sanction)의 위험이 결코 무시할 수 없는 것이기에 불안정한 평화를 유지할
수도 있다. `권력자들'(powers)이 직접 전투를 벌이기를 꺼리는 쟁점들에
대해서는, 일종의 규칙들이 수용될 수도 있다. 다시 말해, 우리는 소인국과
거인의 환상적 비유를 상상할 필요 없이, 근사적 동등성의 단순한 작동
논리(logistics)와 그것이 법에 있어 갖는 중요성을 이해할 수 있다. 국가들
간의 힘의 격차가 극심했던 국제 사회(international scene)는 그 예시로
충분하다. 수 세기 동안 국가 간의 이러한 불균형은 조직화된 제재의 작동을
불가능하게 만들었고, 법은 `중대 사안'(vital issues)에 영향을 미치지 않는
영역에만 한정되었다. \textbf{(p.~199)} 향후 원자무기(atomic weapons)가
모든 국가에 의해 사용 가능해질 경우, 이러한 불균형이 어느 정도 시정되고,
국제 통제가 국내 형사법(municipal criminal law)에 더 가까운 형태로
진화하게 될지 여부는 아직 확정되지 않았다.

우리가 지금까지 논의한 이러한 단순한 자명한 이치들(truisms)은,
자연법(Natural Law) 이론 속의 핵심적인 이성적 직관을 드러낼 뿐 아니라,
법과 도덕을 이해하는 데 본질적으로 중요하며, 왜 법과 도덕의 기본 형태를
특정한 내용이나 사회적 필요에 대한 언급 없이 순전히 형식적 정의로만
설명하는 것이 부적절했는지를 설명해 준다. 아마도 이러한 관점이
법리학(jurisprudence)에 주는 가장 큰 공헌은, 법의 성격에 대한 논의를
자주 흐리게 해왔던 몇몇 오해를 불러일으키는 이분법(dichotomies)들로부터
벗어날 수 있는 여지를 제공한다는 점이다. 예컨대, 전통적인 질문인 ``모든
법체계는 제재를 반드시 포함해야 하는가?''라는 문제는, 자연법의 이 단순한
판본이 제공하는 관점에서 새롭고 보다 명료한 방식으로 제기될 수 있다.
우리는 더 이상 다음의 두 가지 부적절한 대안 사이에서 선택을 강요받지
않아도 된다: 하나는, 법(law)이나 법체계(legal system)라는 단어의
`정의'에 비추어 제재가 요구된다고 주장하는 입장이고, 다른 하나는,
대부분의 법체계가 실제로 제재를 포함하고 있다는 것이 단지 `사실'(just a
fact)에 불과하다고 주장하는 입장이다. 이 두 입장은 모두 만족스럽지 않다.
중앙집중적 제재가 없는 체계에 대해 `법'이라는 용어를 사용하는 것을
금지하는 확립된 원칙은 존재하지 않으며, 제재를 포함하지 않는 체계를
`국제법'(international law)이라 부르는 데에도 타당한 이유(비록 강제성은
없지만)가 있다. 그러나 반대로, 인간 존재의 조건에 부합하는 최소한의
목적들을 수행하려는 국내 체계(municipal system)의 경우에는, 그 안에서
제재가 반드시 차지해야 할 위치를 구분해야 한다. 자연적 사실들과
목적들이, 국내 체계 안에서 제재를 가능하고 필수적인 것으로 만든다는 점을
고려하면, 우리는 이것이 하나의 \emphksans{자연적 필연성}(natural
necessity)이라고 말할 수 있다; 그리고 이와 같은 표현은, 인격(persons),
재산(property), 약속(promises)에 대한 최소한의 보호 형태가 국내법의 필수
구성 요소로 기능한다는 점을 전달하는 데도 필요하다. 우리는 바로 이러한
방식으로 ``법은 어떤 내용도 가질 수 있다''(law may have any content)는
실증주의 테제에 응답해야 한다. 왜냐하면 법뿐만 아니라 많은 다른 사회적
제도들을 적절히 기술하기 위해서는, 정의(definitions)나 통상적인 사실
진술들에 더하여, 제3의 범주의 진술을 위한 자리를 남겨 두어야 한다는
중요한 진실이 있기 때문이다. \textbf{(p.~200)} 즉, 그 진리값이 인간과
그들이 살아가는 세계가 현재 지니고 있는 두드러진 성격들을 계속
유지하는지 여부에 의존하는(contingent) 그러한 진술들이다.

\subsection{3. 법적 유효성과 도덕적
가치}\label{uxbc95uxc801-uxc720uxd6a8uxc131uxacfc-uxb3c4uxb355uxc801-uxac00uxce58}

상호 자제(mutual forbearance)의 체계가 제공하는 보호와 이익은, 법과
도덕의 기초를 이루면서도, 각기 다른 사회에서 매우 상이한 범위의
사람들에게 확장될 수 있다. 일정한 제한을 기꺼이 수용하는 한 어떤 인간
집단에게도 이러한 기본적 보호를 부정하는 것은, 모든 현대 국가가 최소한
언어적으로는 존중을 표하는 도덕과 정의의 원칙에 반하는 일이라는 점은
사실이다. 오늘날 널리 공언되는 도덕적 관점은, 적어도 이러한 근본적인
영역에서만큼은 인간은 동등하게 대우받을 자격이 있으며, 차별적 대우는
단순히 타인의 이익을 근거로 정당화될 수 없다는 관념을 바탕에 깔고 있다.

그러나 한편으로는, 어떤 사회의 법이나 통용되는 도덕이 그들의 범위 내에
있는 모든 사람에게 이러한 최소한의 보호와 이익을 제공해야 할 필요는
없으며, 실제로 종종 그러지 않았다. 노예제 사회에서 지배 집단은 노예를
단순한 사용의 대상이 아니라 인간이라는 사실 자체를 인식하지 못하게 될 수
있으며, 동시에 서로 간의 권리와 이해관계에는 매우 민감한 도덕의식을
유지할 수도 있다. 『허클베리 핀(Huckleberry Finn)』에서, 증기선의
보일러가 폭발해 다친 사람이 있는지를 묻는 말에 허크는 ``아뇨, 그냥
깜둥이 한 명이 죽었어요(No'm: killed a nigger)''라고 답한다. 이때 앤트
샐리가 ``그래도 다행이네, 가끔은 사람도 다치거든(Well it's lucky because
sometimes people do get hurt)''이라고 덧붙이는 장면은, 실제로 인간
사회에서 오래도록 지배적이었던 하나의 도덕 체계를 응축해서 보여준다.
그러한 도덕 체계가 지배하는 사회에서, 허크가 뼈저리게 경험했듯이, 노예에
대해 지배 집단 구성원들 사이에서 자연스럽게 작동하는 상호 관심을 연장
적용하려는 시도는 중대한 도덕적 위반으로 간주되며, 그에 상응하는 도덕적
죄책(moral guilt)의 결과를 초래한다. 나치 독일과 남아프리카 공화국의
사례는, 그 시간이 우리에게 불편할 정도로 가깝다는 점에서 유사한 교훈을
제공한다.

물론 어떤 사회의 법은 때로는 통용되는 도덕보다 앞서기도 했지만,
일반적으로 법은 도덕을 뒤따르며, 노예 살인조차도 공적 자원의 낭비이거나,
노예를 소유한 주인의 재산에 대한 침해로만 간주되기도 한다. 심지어
노예제가 공식적으로 인정되지 않는 사회에서도, 인종, 피부색,
\textbf{(p.~201)} 또는 신조 등의 이유로 차별이 존재할 수 있으며, 이러한
차별은 모든 인간이 타인으로부터 최소한의 보호를 받을 자격이 있다는 점을
법체계나 사회적 도덕이 인정하지 않게 만든다.

이러한 인간사의 고통스러운 사실들은, 어떤 사회가 존속 가능하려면 반드시
\emphksans{일부} 구성원에게는 상호 자제의 체계를 제공해야 하지만, 그것이
반드시 모든 이에게 제공되어야 할 필요는 없다는 점을 명확히 보여준다.
앞서 제재(sanction)의 필요성과 가능성에 대해 논의하면서 강조한 바와
같이, 어떤 규칙 체계가 강제력(force)에 의해 시행되려면, 일정 수의
사람들이 그것을 자발적으로 수용해야 한다. 이러한 자발적 협력은 곧
\emphksans{권위(authority)}를 형성하며, 이는 법과 정부의 강제력을 정초하는
기반이다. 하지만 이렇게 형성된 강제력은 두 가지 주요한 방식으로 사용될
수 있다. 첫째, 그 강제력은 규칙으로부터 보호를 받으면서도 이기적으로
그것을 위반하는 범법자(malefactors)에 대해서만 행사될 수 있다. 둘째, 그
강제력은 일정한 수의 피지배 집단(subject group)을 억압하고 그들을
영구적인 열등한 지위에 고정시키는 데 사용될 수도 있다. 이때 피지배
집단의 규모는 지배 집단(master group)의 연대력(solidarity),
규율(discipline), 그리고 강제 수단의 가용성에 따라 크거나 작을 수
있으며, 피지배 집단의 무력함이나 조직화 능력 부족은 그 억압을 더욱
가능하게 만든다. 이러한 억압을 받는 자들에게는, 그 체계가 그들의 충성을
요구할 만한 어떠한 것도 제공하지 않으며, 오직 두려움의 대상일 뿐이다.
그들은 법체계의 수혜자(beneficiaries)가 아니라 피해자(victims)이다.

이 책의 앞선 장들에서 우리는 법체계의 존재는 하나의 사회적 현상(social
phenomenon)이라는 점을 강조했으며, 그것이 항상 두 가지 측면을 지닌다는
점에 주의를 기울여야 한다고 말했다. 하나는 규칙의 자발적 수용(voluntary
acceptance of rules)에 내재된 태도와 행위이고, 다른 하나는 단순한
복종(obedience)이나 묵인(acquiescence)에 따른 보다 단순한 형태의 태도와
행위이다.

따라서 법을 갖춘 사회는, 그 규칙들을 단지 불복종할 경우 공무원에게서
어떤 처벌이 닥칠지를 예측하게 해주는 신뢰할 만한 지표로만 여기는 것이
아니라, 행위의 기준으로 수용된 규범(accepted standards of
behaviour)으로, 즉 내적 관점(internal point of view)에서 바라보는
사람들을 포함한다. 그러나 동시에 그러한 사회는,
범죄자(malefactors)이거나 단순히 그 체계의 무력한 희생자인 사람들처럼,
이러한 법적 규범들이 오직 물리적 강제력 또는 그 위협을 통해 부과되어야만
하는 이들도 포함한다. 이들은 그러한 규칙들을 가능한 처벌의 원천으로만
인식한다. 이 두 부류 사이의 균형은 \textbf{(p.~202)} 다양한 요소들에
의해 결정된다. 만약 그 체계가 공정하고, 순응(compliance)을 요구하는 모든
이들의 중대한 이해관계(vital interests)를 진정성 있게 고려한다면, 그것은
대부분의 사람들에게 대부분의 시간 동안 충성을 얻고 유지할 수 있으며,
결과적으로 안정적인 체계가 될 수 있다. 반대로, 법 체계가 지배 집단의
이익에만 복무하는 협소하고 배타적인 것이라면, 그 체계는 점점 더
억압적이며 불안정해질 것이고, 잠재된 격변의 위협 아래 놓이게 될 것이다.
이 두 극단 사이에는 다양한 법에 대한 태도 조합들이 존재하며, 종종 하나의
개인 내부에도 이러한 이중성이 공존한다.

이러한 측면에 대한 성찰은 하나의 냉정한 진실을 드러낸다. 의무의 일차
규칙(primary rules of obligation)만이 사회적 통제의 유일한 수단인 단순한
사회 형태에서, 중앙집중적으로 조직된 입법부, 법원, 공무원, 그리고 제재를
갖춘 법의 세계로 나아가는 단계는, 확실한 이익을 가져오는 동시에 일정한
대가를 수반한다. 그 이익이란 변화에 대한 적응성, 확실성, 그리고
효율성으로서, 그 규모는 막대하다. 그러나 그 대가는, 중앙집중적으로
조직된 권력이 자신을 지지하지 않아도 되는 다수의 사람들을 억압하는 데
사용될 수 있다는 위험이며, 이는 일차 규칙의 단순한 체제에서는 발생할 수
없었던 방식이다. 이러한 위험이 실제로 현실화되었고 또 다시 현실화될 수
있기 때문에, 우리가 자연법(Natural Law)의 최소 내용(minimum content)으로
제시한 바를 넘어서, 법이 도덕에 반드시 순응\emphksans{해야 하는}(\emph{must}
conform) 어떤 추가적인 방식이 존재한다는 주장은 매우 신중한 검토를
필요로 한다. 이러한 언명들 가운데 다수는, 법과 도덕 사이의 연관성이 어떤
의미에서 필연적이라고 주장되는지를 명확히 하지 못하거나, 면밀히 살펴보면
그것이 참이고 중요하기는 하나, 법과 도덕 사이의 필연적 연결로 제시하는
것은 극히 혼란스러운 의미에 불과한 것으로 드러난다. 우리는 이 장을
마무리하면서 이러한 주장들의 여섯 가지 형태를 검토할 것이다.

 (i) \emphksans{권력과 권위} (Power and authority) 법 체계는 단순히 인간이
인간을 지배하는 힘에 근거할 수 없고 또한 그래서는 안 되므로, 반드시
도덕적 의무감(moral obligation)이나 법체계의 도덕적 가치에 대한
확신(conviction)에 근거해야 한다는 주장이 자주 제기된다. 우리는 이
책의 앞선 장들에서, 단순히 위협에 의해 뒷받침되는 명령(orders backed
by threats)과 복종의 습관(habits of obedience)만으로는 법체계의 기초와
법적 유효성(legal validity)이라는 개념을 설명하기에 불충분하다는 점을
강조하였다. 우리는 \textbf{(p.~203)} 제6장에서 논증한 바와 같이,
이러한 개념들을 이해하려면 수용된 승인 규칙(accepted rule of
recognition)의 개념이 반드시 필요하며, 본 장에서 본 것처럼, 강제력의
존재 조건은 적어도 일부 사람들이 자발적으로 그 체계에 협력하고, 그
규칙들을 수용(acceptance)하는 것이다. 이 의미에서, 법의 강제력은
그것의 수용된 권위(accepted authority)를 전제로 한다는 것은 사실이다.
그러나 '단지 권력에 기반한 법'과 '도덕적으로 구속력을 가지는 것으로
받아들여진 법'이라는 이분법(dichotomy)은 완전히 포괄하는(exhaustive)
구분이 아니다. 매우 많은 수의 사람들이, 그들이 도덕적으로 구속력을
가진다고 생각하지 않는 법에 의해 강제될 수 있으며, 자발적으로 법
체계를 수용하는 이들조차도 반드시 스스로를 도덕적으로 구속받는
존재라고 인식해야 하는 것은 아니다. 물론 그들이 그렇게 인식할 때 가장
안정적인 체계가 될 것이다. 사실, 사람들이 그 체계에 대한
충성(allegiance)을 보이는 이유는 다음과 같이 다양할 수 있다. 장기적
이익에 대한 계산, 타인에 대한 사욕없는 관심(disinterested interest),
반성 없는 전통적 태도, 혹은 단지 타인들이 그렇게 하기에 그에 따르려는
바람 등이다. 법체계의 권위를 수용하는 이들이 양심을 살펴본 끝에,
도덕적으로는 이를 받아들여서는 안 된다고 판단하면서도, 다양한 이유로
여전히 수용을 지속할 수도 있는 것이다.

이러한 상식적인 사실들이 종종 모호해지는 이유는, 사람들이 법적
의무(legal obligation)와 도덕적 의무(moral obligation)를 동일한 어휘로
표현하기 때문이다. 법체계의 권위를 수용하는 이들은 법을 내적
관점(internal point of view)에서 바라보며, 법이 요구하는 바를 다음과
같은 규범적 언어로 표현한다. 이 규범적 언어는 법과 도덕 모두에서
공유되는 것이다.: ``나는/너는 \textasciitilde\textasciitilde 해야
한다(I/You ought)'', ``나는/그는 \textasciitilde\textasciitilde 해야만
한다(I/he must)'', ``나는/그들은 의무를 지닌다(I/they have an
obligation).'' 그러나 그들은 이러한 표현을 통해, 법이 요구하는 바를
실천하는 것이 도덕적으로 정당하다는 \emphksans{도덕적} 판단에 전념하고 있는
것은 아니다. 물론 아무런 부가 설명이 없다면, 어떤 사람이 자신 또는
타인의 법적 의무를 이렇게 표현할 경우, 그것을 이행하는 데 반대할 도덕적
또는 기타 이유가 없다고 생각한다는 추정(presumption)이 가능할 것이다.
그러나 이로부터, 어떤 사항이 법적으로 유효한 의무(legally obligatory)로
인정되기 위해서는 반드시 도덕적으로도 의무적(morally obligatory)으로
수용되어야 한다는 결론이 도출되는 것은 아니다. 우리가 말한 이 추정은,
어떤 사람이 법적 의무를 인정하거나 언급하는 것이, 만약 그가 그것을
이행해서는 안 된다고 여기는 확고한 도덕적 혹은 다른 이유를 가지고 있다면
무의미할 것이라는 점에 기초한 것이다.

 (ii) \emphksans{도덕이 법에 미치는 영향} (The influence of morality on law) 모든
현대 국가의 법은 \textbf{(p.~204)} 수많은 지점에서, 통용되는 사회적
도덕(social morality)과 보다 보편적인 도덕적 이상(moral ideals)의
영향을 보여준다. 이러한 영향은 입법을 통해 갑작스럽고 공공연하게 법
안으로 들어오기도 하고, 사법절차를 통해 조용히 점진적으로 스며들기도
한다. 어떤 체계들, 예컨대 미국의 경우에는, 법적 유효성(legal
validity)의 궁극적 기준에 정의(justice)나 실질적 도덕 가치(substantive
moral values)가 명시적으로 포함되어 있다. 반면 영국과 같은 체계에서는,
최고 입법기관의 입법권에 공식적인 제한이 없음에도 불구하고, 그 입법은
정의나 도덕에 대한 신중한 부합(conformity)을 결코 덜 요구하지 않는다.
법이 도덕을 반영하는 추가적인 방식들은 무수히 많으며, 아직 충분히
연구되지 않았다. 예컨대, 어떤 법률 조항은 단지 법적 외피(legal
shell)에 불과하고, 그 조항의 명문 자체가 도덕적 원칙의 도움을 통해
내용을 보완할 것을 요구하기도 한다. 집행 가능한 계약(enforceable
contracts)의 범위는 도덕성과 공정성(fairness)의 개념에 따라 제한될 수
있으며, 민사 및 형사상 불법행위에 대한 책임(liability)은 지배적인
도덕적 책임관(moral responsibility)의 관념에 따라 조정되기도 한다.
어떤 실증주의자(positivist)도 이러한 사실들---즉, 법체계의
안정성(stability)은 이와 같은 도덕과의 대응 유형에 부분적으로
의존한다는 사실---을 부인할 수는 없다. 만약 이것이 법과 도덕 사이의
필연적 연결(necessary connection)을 의미하는 것이라면, 그 존재는
인정되어야 한다.

 (iii) \emphksans{해석} (Interpretation) 법은 구체적 사안에 적용되기 위해서는
해석이 요구되며, 사법 절차(judicial processes)의 본질을 흐리는 신화를
현실적 연구로 제거하면, 법의 개방적 직조(open texture)가 창의적 활동의
광대한 영역을 남긴다는 점이 분명해진다. 이러한 창의적 활동은 일부에서
입법적(legislative)이라 불리기도 한다. 선례(precedent)나
성문법(statute)을 해석함에 있어, 판사들은 맹목적이고 자의적인
선택이나, 사전에 고정된 의미를 갖는 규칙들로부터의 `기계적
연역'(mechanical deduction) 사이에만 갇혀 있는 것이 아니다. 종종
그들의 선택은 해석 대상 규칙의 목적이 합당할 것이라는 가정에 의해
이끌린다. 즉, 해당 규칙들은 부정의를 야기하거나 정착된 도덕
원칙(settled moral principles)을 침해하도록 설계된 것이 아니라는
전제이다. 특히 고도의 헌법적 중요성(high constitutional import)이 있는
사안에서의 사법적 판단(judicial decision)은, 하나의 도덕 원칙의 단순한
적용이 아니라, 상이한 도덕적 가치들 사이의 선택을 수반하는 경우가
많다. 법의 의미가 불확실한 상황에서, 도덕이 항상 명확한 해답을
제공한다고 믿는 것은 어리석은 일이다. 이 지점에서 판사들은 다시금
자의적이지도, 기계적이지도 않은 선택을 할 수 있으며, 바로 여기에서
종종 \textbf{(p.~205)} 판사 고유의 미덕(judicial virtues)이 드러난다.
이러한 미덕들이 법적 판단에 적합한 이유는, 이와 같은 활동을 단순히
`입법'이라 부르기를 주저하게 만드는 이유이기도 하다. 이러한 미덕에는
다음이 포함된다: 대안들을 고찰함에 있어의 공정성과 중립성(impartiality
and neutrality), 영향을 받게 될 모든 이들의 이익에 대한 고려, 그리고
수용 가능한 일반 원칙을 근거 있는 결정의 기반으로 삼고자 하는
노력이다. 물론 항상 다양한 원칙들이 존재 가능하기 때문에, 어떤 결정이
유일하게 옳다는 것을 \emphksans{증명}(demonstrate)할 수는 없지만, 그것은
정보에 기초한 공정한 선택의 추론된 산물(reasoned product)로서 수용될
수 있다. 이 모든 과정에서 우리는 경쟁하는 이해관계들 간의 정의 실현을
위한 `저울질'(weighing)과 `균형 조정'(balancing)의 과정을 보게 된다.

이러한 요소들은, 결정을 수용 가능하게 만드는 데 있어 도덕적(moral)이라
불릴 수 있으며, 그것들의 중요성은 거의 누구도 부인하지 않을 것이다.
대부분의 법체계에서 해석을 지배하는 느슨하고 변화 가능한 전통 혹은 해석
규범(canon of interpretation) 역시 이러한 요소들을 대체로 모호하게나마
포함하고 있다. 그러나 이러한 사실들이 법과 도덕 사이의 \emphksans{필연적}
연결(\emph{necessary} connection)을 입증하는 증거로 제시될 경우, 우리는
다음 사실을 기억해야 한다. 즉, 동일한 원칙들이 지켜지는 경우만큼이나
위반되는 경우에도 거의 동일하게 '존중'되어 왔다는 점이다. 오스틴(Austin)
이후 오늘날에 이르기까지, 이러한 요소들이 판단을 이끌\emphksans{어야
한다(should)}고 주장해 온 이들은, 법관의 법 형성(judicial law-making)이
종종 사회적 가치에 대해 맹목적이거나, '자동적'이거나, 이성적 근거가
부족하다는 점을 지적한 비판자들이었다.


 (iv) \emphksans{법에 대한 비판} (The criticism of law) 때로 법과 도덕 사이의
필연적 연결이 존재한다는 주장은, 단지 ``\emphksans{좋은}(good) 법체계는
정의와 도덕의 요구에 일정한 지점에서 부합(conform)해야 한다''는 주장에
불과하기도 하다. 앞 단락에서 언급한 바와 같은 지점들에서 말이다. 일부
사람들은 이것을 자명한 상식(truism)으로 간주할지도 모르지만, 그것은
동어반복(tautology)은 아니다. 실제로 법에 대한 비판 과정에서는, 어떤
도덕적 기준이 적절한지, 그리고 어떤 지점에서의 부합이 요구되는지를
둘러싸고 의견 차이가 있을 수 있다. 즉, 법이 부합해야 할
도덕(morality)이란 해당 집단 내에서 수용된 통념적 도덕인가? 설령
그것이 미신에 기초하거나, 노예나 피지배 계층에게 이익과 보호를
부여하지 않는 경우일지라도? 아니면, 도덕이란 사실에 대한 이성적 믿음에
기초하고, 모든 인간이 동등한 고려와 존중을 받을 자격이 있다고 인정하는
의미에서 계몽된(enlightened) 기준인가?

\textbf{(p.~206)} 법체계가 그 적용 범위에 속하는 모든 인간을 일정한
기본적 보호와 자유를 향유할 자격이 있는 존재로 취급해야 한다는 주장은,
오늘날 법 비판에서 자명하고도 중요한 이상(ideal)으로 일반적으로
받아들여지고 있다. 실제 실천이 이를 따르지 않더라도, 이 이상에 대한
형식적 동의(lip service)는 통상적으로 제공된다. 심지어 모든 인간이
동등한 고려의 권리를 갖는다는 견해를 수용하지 않는 도덕 이론은,
철학적으로 내적 모순, 독단성, 혹은 비이성성을 내포하고 있음이 논증될
수도 있다. 만약 그렇다면, 이러한 권리를 인정하는 계몽된 도덕(enlightened
morality)은 진정한 도덕(true morality)으로서 특별한
정당성(credentials)을 가지며, 단지 가능한 여러 도덕 중 하나가 아니다.
이러한 주장은 여기서 검토할 수 없지만, 비록 이들이 받아들여지더라도, 그
사실은 다음의 사실을 바꾸거나 흐리게 해서는 안 된다. 즉, 기본
규칙(primary rules)과 이차 규칙(secondary rules)의 구조를 특징으로 하는
국내 법체계(municipal legal systems)는, 이러한 정의의 원칙들을
무시하면서도 오랫동안 지속되어 왔다는 점이다. 부정의한 규칙(iniquitous
rules)을 '법'으로 인정하지 않음으로써 무엇을 얻게 되는지에 대한 논의는,
아래에서 다룬다.

 (v) \emphksans{법다움과 정의의 원칙들} (Principles of legality and justice)
도덕성과 정의에 일정한 지점에서 부합하는(conform) 좋은 법체계와 그렇지
않은 법체계 사이의 구분은, 인간 행위가 일반 규칙에 의해 공공연히
공표되고 사법적으로 적용되는 방식으로 통제되는 한, 정의의 최소한이
필연적으로 실현된다는 이유에서, 오류를 포함한 구분이라고 말할 수도
있다. 실제로 우리는 이미 정의(justice)의 개념을 분석하면서\footnote{p.~160
  above.}, 그 가장 단순한 형태인 `법의 적용에서의 정의'(justice in the
application of the law)가 단지 ``서로 다른 많은 사람들에게 적용되어야
할 것이 동일한 일반 규칙(general rule)''이라는 생각을 편견, 이해관계,
또는 변덕에 휘둘리지 않고 진지하게 받아들이는 것에 지나지 않음을
지적한 바 있다. 이러한 공정성(impartiality)은, 영어 및 미국 법률가들이
`자연적 정의'(Natural Justice)의 원칙들로 알고 있는 절차적
기준(procedural standards)이 보장하고자 하는 바로 그 요소이다. 따라서
가장 혐오스러운 법률조차도 정당하게 적용될 수 있으며, 일반 법규칙을
적용한다는 그 자체의 최소한의 관념 속에는 적어도 정의의 씨앗(germ of
justice)이 존재한다고 할 수 있다.

이러한 정의의 최소 형태, 곧 `자연적'(natural)이라 불릴 수 있는 양상들은,
실제로 \textbf{(p.~207)} 어떤 형태의 사회 통제가 작동하려면 어떤
조건들이 전제되어야 하는지를 살펴볼 때 드러난다. 이는 게임의 규칙이든
법이든, 기본적으로는 일반적인 행위 기준(general standards of conduct)을
사람들의 특정 집단에 전달하고, 그 집단이 추가적인 공식 지시 없이 그것을
이해하고 준수(conformity)할 것을 기대하는 방식의 통제이다. 이러한 유형의
사회 통제가 기능하기 위해서는, 규칙은 다음의 조건을 충족해야 한다: 이해
가능해야 하며(intelligible), 대부분의 사람들이 준수할 수 있어야 하고,
일반적으로는 소급 적용(retrospective)이 되어서는 안 된다. 예외적으로만
가능하다. 이 말은, 규칙 위반으로 처벌받는 사람들은 대부분 그것을 준수할
능력과 기회를 가지고 있었다는 것을 의미한다. 명백히 이러한 규칙에 의한
통제의 특징들은, 법률가들이 법다움의 원칙(principles of legality)이라
부르는 정의의 요구사항들과 밀접히 관련되어 있다. 실제로 어떤 실증주의
비판자는, 이러한 규칙에 의한 통제의 양상들 속에서 법과 도덕 사이의
필연적 연결(necessary connection)의 의미를 찾았고, 그것을 `법의 내적
도덕성'(the inner morality of law)이라 불러야 한다고 제안하기도 했다.
다시 말해, 만약 이것이 법과 도덕의 필연적 연결이라는 것이 의미하는
바라면, 우리는 그것을 수용할 수도 있다. 그러나 유감스럽게도, 그것은
극심한 부정의(inquity)와도 양립할 수 있다.

 (vi) \emphksans{법적 유효성(legal validity)과 법에 대한 저항} 그들의 일반적
관점을 다소 조심성 없이 제시했다고 하더라도, 실증주의자로 분류되는
법이론가들 중 앞의 다섯 항목에서 논의한 법과 도덕 사이의 연결 양태들을
부정하고자 했던 사람은 드물다. 그렇다면, 법실증주의의 유명한 구호들이
지향했던 바는 무엇인가? 예컨대, ``법의 존재는 한 가지 문제고, 그
장점(merit) 여부는 또 다른 문제다.''\footnote{Austin, \emph{The
  Province of Jurisprudence Defined}, Lecture V, pp.~184--5.};
``국가의 법은 이상이 아니라 실제로 존재하는 것이다 \ldots{} 그것은
되어야 할 것이 아니라, 있는 그대로의 것이다.''\footnote{Gray,
  \emph{The Nature and Sources of the Law}, s. 213.}; ``법규범(legal
norms)은 어떠한 내용(content)도 가질 수 있다.''\footnote{Kelsen,
  \emph{General Theory of Law and State}, p.~113.}

이러한 사상가들이 주로 옹호하고자 했던 것은, 도덕적으로 부정의하더라도
적법한 형식으로 제정되었고, 의미가 명확하며, 해당 법체계의 모든 인정된
유효성(validity) 기준을 충족하는 특정 법률의 존재가 제기하는 이론적 및
도덕적 문제들을 명료하고 정직하게 서술하려는 시도였다. 그들의 입장은,
이러한 법률에 대해 고찰할 때, 이론가든 그 법률을 적용하거나 준수해야
하는 공무원 또는 시민이든 간에, \textbf{(p.~208)} 그것을 `법' 또는
'유효한 법'이라는 명칭으로 부르지 말라는 초대는 혼란만을 야기할 것이라는
것이었다. 그들은, 이 문제에 직면할 때, 더 단순하고 더 솔직한 방법들이
있으며, 그것이 관련된 모든 지적·도덕적 고려사항들을 더 명확히 조명해 줄
것이라고 생각했다. 즉 우리는 ``이것은 법이다. 그러나 그것은 너무나
부정의하여 적용하거나 복종할 수 없다''라고 말해야 한다는 것이다.

이에 반대되는 입장은, 혁명이나 대규모 정치적 격변 이후에, 법원이 이전
체제하에서 공무원이나 민간인이 법 형식을 통해 저지른 도덕적 불의에 대해
어떤 태도를 취할지를 고찰해야 할 때, 매력적으로 보이기도 한다. 이들의
처벌은 사회적으로 바람직하게 느껴질 수 있으나, 그것을 달성하기 위해서는,
과거 체제 하에서 허용되거나 심지어 요구되었던 행위를 소급적으로
범죄화하는 입법이 필요하며, 이는 실현이 어렵거나, 스스로 도덕적으로
혐오스럽거나, 실행 불가능할 수도 있다. 이러한 상황에서, 법적 어휘 속에
잠재된 도덕적 함의들을 활용하는 것이 자연스럽게 보일 수도 있다. 특히
\emph{ius, recht, diritto, droit}와 같은 단어들은 자연법(Natural Law)의
이론으로 충만한 용어들이므로 더욱 그러하다. 이 경우, 부정의를 명령하거나
허용했던 법률들은, 비록 그것들이 제정된 법체계가 입법권에 대한 어떠한
제한도 인정하지 않았더라도, 유효(valid)하거나 법의 성격(quality of
law)을 갖는 것으로 인정되어서는 안 된다(recognized)는 주장을 하게 되는
유혹에 빠지기 쉽다. 이러한 형태로, 제이차 세계대전 이후 독일에서는 나치
정권의 부정의와 그 패배로 인해 남겨진 첨예한 사회 문제들에 대응하기 위해
자연법 논증들이 부활했다. 이기적 목적을 위해, 나치 정권 하에서 제정된
괴물같은 법령을 위반한 범죄를 저지른 사람들을 투옥시키기 위해 밀고한
사람들은, 처벌받아야 하는가? 그러한 법령들이 자연법을 위반했기 때문에
무효(void)이며, 그로 인해 그 법령을 위반한 것에 따른 투옥은 실질적으로
불법(unlawful)이었고, 그것을 야기한 행위는 그 자체로 범죄였다고 하여,
전후 독일 법정에서 그들을 유죄로 판단할 수 있었는가?\footnote{See the
  judgment of 27 July 1949, Oberlandesgericht Bamberg, 5
  \emph{Süddeutsche Juristen-Zeitung}, 207: discussed at length in H. L.
  A. Hart, `Legal Positivism and the Separation of Law and Morals', in
  71. \emph{Harvard L. Rev.} (1958), 598, and in L. Fuller, `Positivism
  and Fidelity to Law', ibid., p.~630. But note corrected account of
  this judgment below, pp.~303--4.} 이처럼 \textbf{(p.~209)}
``도덕적으로 부정의한 규칙은 법이 될 수 없다''는 견해를 받아들이는
입장과 거부하는 입장 사이의 논쟁은 표면상 간단해 보이지만, 그 일반적
성격에 대한 이해는 종종 불분명하다. 여기서 문제되는 것은, 도덕적으로
부정의한 규칙을 적용하거나, 복종하거나, 타인이 그것을 항변하는 것을
허용하지 않겠다는 도덕적 결정을 표현하는 \emphksans{대안적 방식들} 간의
차이이다. 그러나 이 문제는 단순히 언어적 표현에 관한 문제로 제시될 수는
없다. 논쟁의 양측 누구도, 단순히 ``당신 말이 맞습니다. 당신이 주장한
바는 영어(또는 독일어)에서 그 종류의 주장을 표현하는 올바른
방식입니다''라는 답변에 만족하지 않을 것이다. 따라서 실증주의자는 다음과
같이 주장할 수도 있다. 즉, 영어 사용 관행에 비추어 보아, ``어떤 법규칙이
너무나 부정의하여 복종할 수 없다''고 주장하는 데에는 모순이 없으며, ``그
규칙에 복종할 수 없다''는 명제가 ``그 규칙은 유효한 법규칙이 아니다''는
결론으로 이어지는 것은 아니라고 주장할 수 있다. 그러나 그 반대편에 선
이들은 이 주장으로 인해 논쟁이 종결되었다고 보지는 않을 것이다.

명백히, 우리는 이 문제를 언어 사용의 적절성(propriety)에 관한 것으로만
본다면, 그것을 충분히 다룰 수 없다. 진정으로 문제가 되는 것은,
사회생활에서 일반적으로 효과적으로 작동하는 하나의 규칙 체계에 속한
규칙들을 분류하는 방식으로서의 넓은 개념과 좁은 개념, 즉 경쟁하는 두
규범 개념들 중 어느 것이 이론적 탐구나 도덕적 숙고---혹은 그 둘
모두---를 더 잘 보조하는지에 따라 선택되어야 한다는 점이다.

이 두 경쟁 개념 중에서 더 넓은 개념은 좁은 개념을 포함한다. 더 넓은
개념을 채택할 경우, 우리는 이론적 탐구에서 다음과 같은 방식을 취하게
된다: 일차적(primary) 및 이차적(secondary) 규칙의 체계가 제시하는 형식적
기준에 따라 유효(valid) 한 모든 규칙들을---비록 그 일부가 특정 사회의
도덕성이나 우리가 진정한 도덕성 또는 계몽된 도덕성이라 여기는 것에
어긋나더라도---`법'(law)으로 묶어 함께 고찰하게 된다. 반면, 좁은 개념을
채택하면 그러한 도덕적으로 불의한 규칙들은 '법'의 범주에서 제외된다.
그러나 이처럼 좁은 개념을 채택하는 것은, 법을 사회적 현상으로 연구하는
이론적·과학적 탐구에 있어서 아무런 이익도 가져다주지 않는 듯하다.
왜냐하면 이는 비록 법의 복잡한 특징들을 모두 갖추었더라도, 특정 규칙들을
법에서 배제하게 되기 때문이다. 이러한 규칙들을 다른 학문 분야로 넘기자는
제안에서 비롯될 수 있는 것은 오직 혼란뿐이며, 실제로 그 어떤 역사적
연구나 법학의 다른 분과도 이러한 시도를 유익하다고 판단하지 않았다. 반면
넓은 개념의 법 개념을 채택하면, \textbf{(p.~210)} 우리는 도덕적으로
불의한 법이 지니는 특수한 성격들과 사회가 그것에 보이는 반응까지
포함하여 포괄적으로 연구할 수 있다. 따라서 좁은 개념을 사용할 경우,
일차·이차 규칙 체계에서 구현되는 사회통제의 특정 방식이 어떻게 발전하고
어떤 가능성을 지니는지를 이해하려는 우리의 노력은 혼란스럽게 분열될
수밖에 없다. 그것의 사용에 대한 연구는 필연적으로 그것의 오용(misuse)에
대한 연구를 포함해야 하기 때문이다.

그렇다면 좁은 개념이 도덕적 숙고에 있어 가지는 실천적 유익성은 무엇인가?
도덕적으로 부정의한 요구에 직면했을 때, '이것은 어떤 의미에서도 법이
아니다'라고 판단하는 것이 '이것은 법이지만 너무 부정의하여 복종하거나
적용할 수 없다'라고 판단하는 것보다 나은가? 이 판단은 사람들을 더
명료하게 만들거나, 도덕이 요구할 때 더 쉽게 불복종하도록 만들 것인가?
혹은 나치 정권이 남긴 문제들을 더 나은 방식으로 처리하게 할 것인가? 물론
사상은 영향력을 가진다. 그러나 도덕적으로 부정의하지만 유효(valid) 한
법이 존재할 수 없다는 좁은 법적 유효성(validity) 개념을 사람들에게
교육하고 훈련시키는 노력이, 조직화된 권력의 위협에 직면하여 악에
저항하는 의지를 강화시키거나, 복종이 요구될 때 도덕적으로 어떤 문제가
제기되는지를 더 명확하게 인식하게 만들 것이라 보기는 어렵다. 인간이
일부의 협조를 통해 타인을 지배할 수 있는 한, 그들은 법의 형식(form)을 그
수단 중 하나로 사용할 것이다. 사악한 자들은 사악한 규칙을 제정하고, 다른
자들은 그것을 집행할 것이다. 이러한 권력의 공식적 남용에 맞서 사람들을
명확하게 인식하게 만들기 위해 무엇보다 필요한 것은, 어떤 규칙이 법적으로
유효하다고 인증(certification)되었다는 사실이 복종의 문제를 최종적으로
해결하는 것은 아니라는 점을 인식하는 것이다. 그리고 공식적 체계가 아무리
위엄이나 권위를 띠더라도, 그 요구는 결국 도덕적 심판(moral scrutiny)에
복종되어야 한다는 점을 자각하는 것이다. 이러한 감각---곧 공식적 체계
밖에 어떤 기준이 존재하며, 개인은 최종적으로 그것에 의거해 복종의 문제를
해결해야 한다는 인식---은 법규칙이 부정의할 수 있다(may be iniquitous)는
사고에 익숙한 사람들 사이에서 훨씬 더 생생히 유지될 가능성이 크다. 반면,
어떤 부정의한 것도 결코 법의 지위를 가질 수 없다고 생각하는 사람들
사이에서는 그렇지 않다.

그러나 우리가 `이것은 법이지만 부정의하다'라고 사고하고 표현할 수 있게
해주는 법의 넓은 개념을 선호해야 할 더 강력한 이유는 따로 있다.
\textbf{(p.~211)} 그것은 부정의한 규칙에 대해 법적 인정을 거부함으로써,
그것이 초래하는 도덕적 문제들의 다양성과 복잡성을 지나치게 단순화할
위험이 있기 때문이다. 벤담(Bentham)과 오스틴(Austin)과 같은 고전
이론가들은 `법이란 무엇인가'와 `법은 어떻게 되어야 하는가' 사이의 구분을
강조했는데, 이는 사람들이 이 구분을 유지하지 않으면, 사회에 어떤 대가가
따르든 무시한 채 법이 유효하지 않다고 성급히 판단하고 복종하지 않으려 할
수 있기 때문이었다. 이러한 무정부 상태의 위험은 그들이 과대평가했을 수도
있으나, 또 다른 형태의 과도한 단순화의 위험도 있다. 만약 우리가 관점을
좁혀 '악한 규칙에 \emphksans{복종(obey)}해야 하는 사람'만을 고려한다면, 그가
그것이 '법'이라는 유효한 규칙인지 아닌지를 어떻게 인식하든, 도덕적
부정의를 인식하고 도덕이 요구하는 대로 행동한다면 그것은 무관하다고
판단할 수도 있다. 그러나 '복종'이라는 도덕적 문제(``내가 이 부정의한
일을 해야 하는가?'') 외에도, 소크라테스가 던진 문제---``불복종에 대한
처벌을 감수해야 하는가, 아니면 탈출해야 하는가?''---가 있으며, 또한 전후
독일 법원이 마주했던 질문---``당시 시행 중이던 부정의한 규칙들이 허용한
악행을 저지른 자들을 우리는 처벌해야 하는가?''---도 존재한다. 이러한
질문들은 서로 매우 다른 도덕적·정의적 문제들을 제기하며, 각각 독립적으로
고려되어야 한다. 그 어떤 목적을 위해서도 부정의한 법률을 유효하다고
인정하지 않겠다(refuse to recognize)는 일회적 선언으로 해결될 수 없다.
이것은 복잡하고 미묘한 도덕적 문제들에 대해 지나치게 거친 접근 방식이다.

법의 유효성(validity)을 그 도덕성(morality)으로부터 구별할 수 있도록
해주는 법 개념은, 이처럼 서로 다른 문제들의 복잡성과 다양성을 인식할 수
있게 한다. 반면, 부정의한 규칙은 결코 법적으로 유효할 수 없다고 주장하는
좁은 법 개념은, 오히려 이러한 문제들을 가리는 역할을 할 수도 있다.
독일의 밀고자들이 이기적인 동기에서 괴물 같은 법 아래에서 타인을
처벌받게 한 것이 도덕적으로 잘못된 일임은 분명하지만, 동시에 도덕은
국가는 오직, 그 당시 법이 금한 악행을 저지른 사람만을 처벌해야 한다고
요구할 수도 있다. 이것이 바로 \emph{`법 없이는 형벌도 없다'(nulla poena
sine lege)라는} 원칙이다. 만약 이 원칙을 희생하지 않고서는 더 큰 악을
피할 수 없다면, 그 희생이 정당화되는 조건들이 명확히 규정되어야 한다.
소급 처벌의 사례가, 마치 당시 이미 불법이었던 행위에 대한 처벌인 것처럼
보이게 해서는 안 된다. \textbf{(p.~212)} 실증주의적 간명한 입장이, 즉
'도덕적으로 부정의한 규칙들도 여전히 법일 수 있다'는 주장이 가지는
최소한의 의의는 바로 이것이다. 그것은 극단적 상황에서 악 중의
선택(choice between evils)이 불가피할 때, 그 선택을 덮거나 위장하지
않는다는 점이다.

\subsection{CHAPTER IX 주석}\label{chapter-ix-uxc8fcuxc11d}

\emph{185쪽. 자연법(Natural Law).} 고전적·스콜라철학적·근대적 자연법
개념과 '실증주의(positivism)'라는 표현의 다의성에 대한 방대한 주석
문헌이 존재하는 탓에, 자연법(Natural Law)이 법실증주의(Legal
Positivism)와 대립할 때 정확히 어떤 쟁점이 걸려 있는지를 파악하기 어려울
때가 많다. 본문에서는 그 중 하나의 쟁점을 식별하려는 시도를 한다. 그러나
이 주제를 논의함에 있어 \emphksans{이차 문헌}만을 읽는다면 얻을 수 있는 것은
극히 제한적이다. \emphksans{자연법의 일차 문헌들}에 나타난 어휘와 철학적
전제들을 직접적으로 이해하는 것이 필수적이다. 다음은 쉽게 접근 가능한
최소한의 자료들이다: 아리스토텔레스(Aristotle), 『자연학(Physics)』
제2권 제8장 (Ross 역, Oxford); 아퀴나스(Aquinas), 『신학대전(Summa
Theologica)』 제1-2부, 문제 90--97 (D'Entrèves, 『아퀴나스: 선별된
정치적 저작』(Aquinas: Selected Political Writings), Oxford, 1948에 번역
수록); 그로티우스(Grotius), 『전쟁과 평화의 법(On the Law of War and
Peace)』, 서론(Prolegomena) (The Classics of International Law, 제3권,
Oxford, 1925 수록 번역본); 블랙스톤(Blackstone), 『주석(Commentaries)』,
서론 제2절.

\emph{185쪽. 법실증주의(Legal Positivism).} 현대 영미권 문헌에서
`실증주의(positivism)'라는 표현은 다음과 같은 하나 혹은 그 이상의 주장을
가리키는 데 사용된다: (1) 법은 인간에 의해 부과된 명령(command)이다; (2)
법과 도덕, 또는 존재하는 법(law as it is)과 마땅히 그래야 할 법(law as
it ought to be) 사이에는 필연적 연결이 없다; (3) 법 개념의 의미를
분석하거나 연구하는 것은 역사적 탐구, 사회학적 탐구, 도덕·사회적
목적·기능에 따른 비판적 평가와는 구별되는 (그러나 적대적인 것이 아닌)
중요한 연구 분야이다; (4) 법체계는 `폐쇄된 논리 체계(closed logical
system)'이며, 올바른 판결은 오직 미리 정해진 법규로부터 논리적
방법만으로 도출될 수 있다; (5) 도덕판단은 사실 명제(statements of
fact)처럼 합리적 논변, 증거, 또는 입증을 통해 정립될 수 없다(윤리적
비인지주의, non-cognitivism in ethics). 벤담(Bentham)과 오스틴(Austin)은
(1), (2), (3)의 견해를 지지했으나 (4), (5)는 지지하지 않았다.
켈젠(Kelsen)은 (2), (3), (5)를 수용했으나 (1), (4)는 수용하지 않았다.
주장 (4)는 흔히 `분석법학자(analytical jurists)'에게 귀속되곤 하지만,
이는 타당한 근거 없이 이루어진 것이다. 대륙 유럽 문헌에서는
`실증주의(positivism)'라는 표현이 흔히 \emphksans{이성에 의해서만 발견 가능한
인간 행위의 어떤 원칙이나 규칙의 존재}를 전면적으로 부정하는 입장을
가리키는 데 사용된다. `실증주의' 개념의 다의성에 대한 아고(Ago)의 유익한
논의는 \emph{American Journal of International Law} 제51권(1957) 참조.

\emph{186쪽. 밀(Mill)의 자연법 비판.} 『자연(Nature), 종교의 효용성,
유신론』에 수록된 ``자연(Nature)'' 에세이 참조.

\emph{187쪽. 블랙스톤과 벤담의 자연법에 대한 견해.} 블랙스톤, 위
인용(loc. cit.); 벤담, 『주석에 대한 논평(Comment on the Commentaries)』
제1--6절.

\emph{193쪽. 자연법의 최소한의 내용.} 자연법의 이 경험주의적 버전은
홉스(Hobbes), 『리바이어던(Leviathan)』 제14·15장과 흄(Hume),
『인성론(Treatise of Human Nature)』 제3권 제2부, 특히 제2절 및
제4--7절에 근거한다.

\emph{200쪽. 『허클베리 핀(Huckleberry Finn)』.} 마크 트웨인(Mark
Twain)의 이 소설은 사회적 도덕이 개인의 공감(sympathy)이나
인도주의(humanitarianism)에 반할 경우 발생하는 도덕적 딜레마에 대한
심오한 탐구다. 모든 도덕성을 인도주의와 동일시하는 경향에 대한 귀중한
반론이 된다.

\emph{200쪽. 노예제(Slavery).} 아리스토텔레스에게 있어 노예는 '살아있는
도구'였다. (\emphksans{정치학(Politics)} 제1권 제2--4장).

\emph{203쪽. 도덕성이 법에 미치는 영향.} 법의 \emphksans{발전(development)}이
도덕성에 의해 어떻게 영향을 받아왔는지를 다룬 유익한 연구들로는
에임스(Ames), ``법과 도덕(Law and Morals)'', \emph{Harvard Law Review}
제22권(1908); 파운드(Pound), 『법과 도덕(Law and Morals)』(1926);
굿하트(Goodhart), 『영국법과 도덕법(English Law and the Moral
Law)』(1953) 등이 있다. 오스틴은 이와 같은 사실적 혹은 인과적 연관을
충분히 인식하고 있었다. 『법의 영역(The Province of Jurisprudence
Determined)』 제5강, 162쪽 참조.

\emph{204쪽. 해석(Interpretation).} 법 해석에서 도덕적 고려가 차지하는
위치에 대해서는 라몽(Lamont), 『가치판단(The Value Judgment)』
296--301쪽; 웩슬러(Wechsler), ``헌법 해석의 중립 원칙을 향하여'',
\emph{Harvard Law Review} 제73권, 960쪽; 하트(Hart), 앞서 인용한
\emph{HLR} 제71권, 606--615쪽 및 풀러(Fuller)의 비판, 동일호 661쪽
\emphksans{끝 부분(ad fin.)} 참조. 오스틴의 `경쟁하는 유추(competing
analogies)' 사이에서 판사가 선택할 수 있는 여지를 인정한 부분과,
공리성의 기준에 그 판결이 부합하지 않는다는 비판은 『강의(The
Lectures)』 제37·38강 참조.

\emph{205쪽. 법에 대한 비판과 모든 인간의 평등한 고려에 대한 권리.}
이러한 권리가 단지 다양한 도덕성 중 하나에 속하는 것이 아니라
\emphksans{진정한 도덕성의 정의적 특징(defining feature)}이라는 견해는
벤(Benn)과 피터스(Peters), 『사회 원칙과 민주국가(Social Principles and
the Democratic State)』 제2·5장, 그리고 바이어(Baier), 『도덕적 관점(The
Moral Point of View)』 제8장을 참조.

\emph{206쪽. 법다움과 정의의 원칙(Principles of legality and justice).}
법다움 원칙에 대해서는 홀(Hall), 『형법의 원칙(Principles of Criminal
Law)』 제1장을, '법의 내재적 도덕성(the internal morality of law)'에
대해서는 풀러(Fuller), \emph{Harvard Law Review} 제71권(1958),
644--648쪽 참조.

\emph{208쪽. 전후 독일에서의 자연법 이론의 부활.} G.
라트브루흐(Radbruch)의 후기 견해, 하트(Hart), 그리고 이에 대한
풀러(Fuller)의 응답에 대한 논의는 모두 \emph{HLR} 제71권(1958) 참조. 이
논의는 1949년 7월 바이에른주 고등법원(Oberlandsgericht Bamberg)의 판결을
중심으로 진행되었는데, 이 사건은 1934년 나치 법률 위반 혐의로 남편을
신고한 아내가 남편의 자유를 불법적으로(unlawfully) 박탈한 혐의로 유죄
판결을 받은 사건이다. 이 사건에 대한 설명은 \emph{HLR} 제64권(1951)
1005쪽에 기술된 것으로, 해당 법원이 1934년 법률을 무효로 판단했다는
전제에 따라 이루어진다. 그러나 이 설명의 정확성은 파페(Pappe)의
비판(``나치 시대의 사법판단의 유효성에 대하여'', \emph{Modern Law
Review} 제23권(1960))에 의해 최근에 도전받았다. 파페 박사의 비판은
정당하며, 하트가 논한 사례는 엄밀히 말하면 가상의 것으로 간주되어야
한다. 파페 박사에 따르면(위 저서, 263쪽), 실제 사건에서 항소법원은
법률이 자연법(Natural Law)을 위반한 경우 불법이 될(be unlawful) 수
있다는 이론적 가능성을 받아들이면서도, 쟁점이 된 나치 법률이 이를
위반했다고 보지는 않았다. 피고는 정보를 제공할 \emphksans{의무(duty)}가
없음에도, 순전히 사적 동기에서 그렇게 했으며, 그러한 행위가 '모든 양심
있는 인간의 건전한 양식과 정의감에 반하는 것임을 인식했음이 분명하다'는
이유로 \emphksans{자유의 불법적 박탈(unlawful deprivation of liberty)} 혐의로
유죄 판결을 받았다. 파페 박사가 유사 사건에 대한 독일 연방대법원의
판결을 세밀하게 분석한 내용(268쪽 \emphksans{끝 부분})도 반드시 참고할 필요가
있다.

\subsection{CHAPTER IX 제3판
주석}\label{chapter-ix-uxc81c3uxd310-uxc8fcuxc11d}

\emph{185--6쪽. 실증주의(Positivism)와 `분리가능성 테제(separability
thesis)'.} 하트는 ``법이 도덕의 특정 요구를 재현하거나 충족시키는 것이
필연적 진리인 것은 결코 아니다''라고 말한다. (홈즈 강연에서는 ``법과
도덕 사이에는 필연적 연결이 없다''고 말한다: H. L. A. Hart, ``Positivism
and the Separation of Law and Morals'' (1957) \emph{Harvard Law Review}
제71권 593쪽, 601쪽 주석 25 참조.) 이 \emphksans{분리가능성 테제(separability
thesis)}는 혼란스러운 개념인데, 그 이유 중 하나는 하트 자신이 \emphksans{법과
정의 사이의 규칙적 연관성}(159--60쪽 주석 참조)과 \emphksans{최소한의 내용
명제(minimum content thesis)}(193쪽 주석 참조)를 통해 두 가지 '필연적
연결'을 옹호하는 것처럼 보이기 때문이다. 이 난제를 정리하려는 다양한
시도로는 다음이 있다: John Gardner, 「Legal Positivism: 5½ Myths」,
\emph{Law as a Leap of Faith} 제2장; Matthew H. Kramer, 「On The
Separability of Law and Morality」 (2004) \emph{Canadian Journal of Law
and Jurisprudence} 제17권 315쪽; James Morauta, 「Three Separation
Theses」 (2004) \emph{Law and Philosophy} 제23권 111쪽; Leslie Green,
「Positivism and the Inseparability of Law and Morals」 (2008) \emph{New
York University Law Review} 제83권 1035쪽.

\emph{185--93쪽. 고전적 자연법 이론(Classical theories of natural law).}
하트가 묘사한 견해의 후손 격이나 \emphksans{신토미즘적(neothomist)} 성향이
가미된 입장으로는 John Finnis의 『Natural Law and Natural Rights』, Mark
C. Murphy의 『Natural Law in Jurisprudence and Politics』(Cambridge
University Press, 2006)이 있다. Nigel Simmonds의 『Law as a Moral
Idea』(Oxford University Press, 2007)와 비교할 것. 자연법과 실증주의의
관계에 대해서는 John Finnis, 「The Truth in Legal Positivism」,
『Philosophy of Law』 제7장 참조. 자연법적 도덕 이론, 법 해석 이론,
자연법 법철학 이론 간의 (결여된) 관계에 대해서는 Philip Soper, 「Some
Confusions about Natural Law」 (1992) \emph{Michigan Law Review} 제90권
2393쪽.

\emph{193--200쪽. `최소한의 내용' 명제(The `minimum content' thesis).}
인간의 본성과 우리가 살아가는 세계를 고려할 때, 보편적으로 가치 있는
것들에 관한 \emphksans{최소한의 내용(minimum content)}을 다루지 않는 규칙
체계는 어떤 경우에도 유효한 법체계(legal system)가 될 수 없다. 하트의
목록을 수정한 사례로는 Neil MacCormick, 『H. L. A. Hart』 92--99쪽, John
Finnis, 『Natural Law and Natural Rights』 제4장 참조. 법의 필수적
내용에 대한 다른 접근은 Joseph Raz, 『Practical Reason and Norms』
162--170쪽 참조.

\emph{203쪽. 법과 도덕의 공유 어휘(The shared vocabulary of law and
morality).} `의무(obligation)', `권리(right)', `자유(liberty)' 등의
개념이 법과 도덕에서 공통적으로 사용되는 이유에 대한 다른 설명은 Joseph
Raz, 『The Authority of Law』 제8장 참조. 이에 대한 하트의 응답은
『Essays on Bentham』 153--161쪽에 실려 있다.

\emph{203--4쪽. 도덕의 법 내포(The incorporation of morality in law).}
이 부분, 그리고 더 강조되는 \emphksans{후기(Postscript)}에서 하트는 다음과
같이 말한다: ``일부 법체계에서는 \ldots{} 법적 유효성(legal validity)의
궁극적 기준이 정의 원칙이나 실질적 도덕 가치(substantive moral values)를
명시적으로 내포하고 있다.'' 하트는 이를 `연성 실증주의(soft
positivism)'라 부르지만, 현재는 일반적으로 `내포 명제(incorporation
thesis)' 또는 '포괄적 법실증주의(inclusive legal positivism)'라 불린다.
관련 문헌은 기술적으로 복잡할 수 있다. 훌륭한 개관으로는 K. E. Himma,
「Inclusive Legal Positivism」, Jules Coleman \& Scott Shapiro 편, 『The
Oxford Handbook of Jurisprudence and Philosophy of Law』 수록. 하트의
입장에 대한 비판은 Joseph Raz, 『The Authority of Law』 제3장, 『Between
Authority and Interpretation』 제7장 참조. 포괄적 실증주의를 발전시킨
논의로는 다음이 있다: Jules Coleman, 「Negative and Positive
Positivism」 (1982) \emph{Journal of Legal Studies} 제11권 139쪽; 『The
Practice of Principle: In Defence of a Pragmatist Approach to Legal
Theory』(Oxford University Press, 2001) 제8장; W. J. Waluchow,
『Inclusive Legal Positivism』(Oxford University Press, 1994).
드워킨(Dworkin)은 하트의 접근(Philip Soper와 David Lyons가 채택한
방식)을 『Taking Rights Seriously』 345--350쪽에서 비판하며, Coleman의
입장은 「Thirty Years On」 (2002) \emph{Harvard Law Review} 제115권
1655쪽에서 비판한다.

\emph{206--7쪽. 법다움과 정의의 원칙(Principles of legality and
justice).} 론 풀러(Lon L. Fuller)의 논의가 영향력 있다: 『법의
도덕성(The Morality of Law)』(개정판, Yale University Press, 1969)
46--91쪽; 하트의 서평 「Lon L. Fuller, `The Morality of Law'」는
『Essays in Jurisprudence and Philosophy』 제16장에 수록. 하트와 풀러 간
논쟁의 중요성에 대해서는 Peter Cane 편, 『The Hart--Fuller Debate: 50
Years On』(Hart Publishing, 2010) 수록 논문들 참조. \emphksans{법의 지배(rule
of law)}와 법의 본성 사이의 관계에 대해서는 Raz, 『The Authority of
Law』 제11장 참조. 드워킨은 「Hart's Postscript and the Character of
Political Philosophy」 (2004) \emph{Oxford Journal of Legal Studies}
제24권 1, 23--37쪽에서 '법다움(legality)'이라는 가치에 기초하여 자신의
이론을 재구성한다.

\emph{207--12쪽. 법적 유효성과 법 저항(Legal validity and resistance to
law).} 정의를 확보하려는 최소한의 시도조차 하지 않는 법은 \emphksans{유효하지
않다(invalid)}는 라드브루흐(Gustav Radbruch)의 논제를 현대적으로 계승한
사례로는 Robert Alexy, 『The Argument from Injustice: A Reply to Legal
Positivism』 (Bonnie L. Paulson \& Stanley L. Paulson 번역, Oxford
University Press, 2002). 법이 우리의 복종(obedience)을 요구할 수 있는
도덕적 정당성에 대해서는 다음 문헌 참조: Joseph Raz, 『The Authority of
Law』 제12--15장; A. J. Simmons, 『Moral Principles and Political
Obligations』 (Princeton University Press, 1979); Leslie Green, 『The
Authority of the State』; W. A. Edmundson, 『Three Anarchical Fallacies:
An Essay on Political Authority』 (Cambridge University Press, 1998).

\emph{208--12쪽. '광의의 법 개념(broad concept of law)'을 옹호하기.}
하트는 \emphksans{실증주의적 법 개념}이 ``도덕적 숙고(moral deliberation)를
촉진하고 명확하게 할 수 있다''고 주장한다 (209쪽). 하트는 이 주장을
자신의 분석의 타당성(validity)을 증명하기 위한 필수 요소로 보지는
않으며, 부차적 고려이거나, 광의의 법 개념이 정치적으로 위험하다고
비판하는 사람들에 대한 응답으로 간주한다. 문헌에서는 \emphksans{법 개념(the
concept of law)}이 도덕적으로 바람직한 결과를 위해 조정되어야 한다는
주장과 \emphksans{법의 내용(the content of the law)}이 도덕적으로 좋은 결과를
위해 조정되어야 한다는 주장 사이에 모호함이 존재한다 (예: 판사의 직무를
오직 \emphksans{근원적 자료(source-based materials)} 적용에 한정하려는 시도
등). 비교 자료: Neil MacCormick, 「도덕주의적 법 비판, 비도덕주의적 법
옹호(A Moralistic Case for A-Moralistic Law)」 (1985) \emph{Valparaiso
Law Review} 제20권 1쪽; Thomas Campbell, 『윤리적 실증주의의 법이론(The
Legal Theory of Ethical Positivism)』 (Dartmouth, 1996); Jeremy Waldron,
「규범적(또는 윤리적) 실증주의」, Jules Coleman 편, 『Hart's
Postscript』; Liam Murphy, 「법 개념에 대한 정치적 물음」, 동일 저서
수록.

\setcounter{footnote}{0}
\section{X. 국제법 (INTERNATIONAL
LAW)}\label{x.-uxad6duxc81cuxbc95-international-law}

\subsection{1. 의문의 원천들 (SOURCES OF
DOUBT)}\label{uxc758uxbb38uxc758-uxc6d0uxcc9cuxb4e4-sources-of-doubt}

본서에서 그토록 중요한 자리를 부여받은 일차 규칙들과 이차 규칙들의
결합이라는 관념(idea)은, 법리학적 극단들 사이의 중간으로 간주될 수 있다.
법이론은 때로는 위협으로 뒷받침된 명령이라는 단순한 관념에서, 때로는
도덕(morality)이라는 복잡한 관념에서 법 이해의 열쇠를 찾으려 해왔다.
법은 이 둘 모두와 분명 많은 친연성과 연결을 가진다. 그러나 우리가
보았듯이, 이것들을 과장하고 법을 다른 사회통제 수단들과 구별하는 특수한
특징들을 흐릴 상존하는 위험이 있다. 우리가 중심적인 것으로 취한 관념의
미덕은, 법, 강제(coercion), 도덕 사이의 복수적 관계들을 그것들이 실제로
그러한 대로 보게 해 주고, 이것들이 있다면 어떤 의미에서 필연적인지를
새롭게 고려하게 해 준다는 데 있다.

일차 규칙들과 이차 규칙들의 결합이라는 관념이 이러한 미덕들을 가지고
있고, 또 이러한 특징적인 규칙 결합의 존재를 ‘법체계(legal system)’라는
표현 적용의 충분조건으로 다루는 것이 용례(usage)에 부합할 것임에도,
우리는 ‘법(law)’이라는 단어가 반드시 그 결합의 관점에서 정의되어야
한다고 주장하지 않았다. 우리가 ‘법’이나 ‘법적’ 같은 단어들의 사용을 이런
방식으로 식별하거나 규율해야 한다는 주장을 하지 않기 때문에, 이 책은
자연스럽게 이 표현들의 사용을 위한 하나의 규칙 또는 여러 규칙들을 제공할
것으로 기대될 수 있는 ‘법’의 정의가 아니라, 법 \emphksans{개념(concept)}의
해명으로 제시된다. 이 목적에 부합하게, 우리는 지난 장에서 독일
사례들에서 제기된 주장(claim), 곧 어떤 규칙들이 현존하는 일차 및 이차
규칙 체계에 속하더라도 그 도덕적 사악함(moral iniquity) 때문에 유효한
법이라는 명칭(title)을 그 규칙들로부터 보류해야 한다는 주장을
검토하였다. 결국 우리는 이 주장을 거부하였다. 그러나 그렇게 한 것은
그러한 체계에 속하는 규칙들은 반드시 ‘법’이라고 불려야 한다는 견해와 그
주장이 충돌하기 때문도 아니고, 용례의 무게와 충돌하기 때문도 아니었다.
\textbf{(p.~214)} 오히려 우리는 도덕적으로 사악한 것을 밀어냄으로써
유효한 법들의 부류를 좁히려는 시도를 비판하였다. 그 근거는 그렇게 하는
것이 이론적 탐구나 도덕적 숙고 어느 쪽도 진전시키거나 명료하게 하지
못한다는 데 있었다. 이러한 목적들을 위해서는, 많은 용례와 부합하고
규칙들이 아무리 도덕적으로 사악하더라도 그것들을 법으로 간주할 수 있게
하는 더 넓은 개념이, 검토 결과 적절한 것으로 드러났다.

국제법(international law)은 우리에게 그 반대의 사례를 제시한다. 지난
150년의 용례와 부합하게 여기서 ‘법’이라는 표현을 사용하더라도, 국제
입법부, 강제관할권을 가진 법원들, 중앙적으로 조직된 제재들의 부재는
적어도 법이론가들의 마음속에 의구심들을 불러일으켜 왔다. 이러한 제도들의
부재는 국가들에 대한 규칙들이, 오직 의무의 일차 규칙들만으로 이루어진
단순한 사회구조의 형식과 닮아 있음을 뜻한다. 개인들의 사회에서 이런
형식을 발견할 때 우리는 그것을 발전된 법체계와 대조하는 데 익숙하다.
우리가 보이겠지만, 국제법은 입법부와 법원들을 마련하는 변경과 재판의
이차 규칙들을 결여할 뿐 아니라, 법의 ‘원천들(sources)’을 특정하고 그
규칙들의 식별을 위한 일반 기준들을 제공하는 통합적 승인 규칙도
결여한다고 기실 논변할 수 있다. 이러한 차이들은 기실 두드러지며,
‘국제법은 정말로 법인가?’라는 질문은 거의 제쳐 놓을 수 없다. 그러나 이
경우에도 우리는 많은 사람들이 느끼는 의문들을 기존 용례를 단순히
환기시키는 것으로 물리치지 않을 것이며, 일차 및 이차 규칙들의 결합이
‘법체계’라는 표현을 올바르게 사용하기 위한 충분조건일 뿐 아니라
필요조건이기도 하다는 근거 위에서 그 의문들을 단순히 확증하지도 않을
것이다. 그 대신 우리는 느껴져 온 의문들의 상세한 성격을 탐구하고, 독일
사례에서처럼, ‘국제법’을 말하는 공통의 더 넓은 용례가 어떤 실천적 또는
이론적 목적을 방해할 개연성이 있는지를 물을 것이다.

우리가 국제법의 성격에 관한 이 문제에 단 하나의 장만을 할애하겠지만,
일부 저자들은 이 질문을 훨씬 더 짧게 처리할 것을 제안해 왔다. 그들에게는
‘국제법은 정말로 법인가?’라는 질문이 오직 단어들의 의미에 관한 사소한
질문을 사물들의 본성에 관한 \textbf{(p.~215)} 중대한 질문으로 오인했기
때문에 생겨났거나 살아남은 것처럼 보였다. 곧 국제법을 국내법(municipal
law)과 구별하는 사실들은 명료하고 잘 알려져 있으므로, 확정되어야 할
유일한 질문은 우리가 기존 관행을 준수할 것인지 아니면 그것에서 벗어날
것인지이며, 이것은 각자가 스스로 확정할 문제라는 것이다. 그러나 이
질문을 이렇게 짧게 처리하는 방식은 확실히 너무 짧다. 이론가들이
‘법’이라는 단어를 국제법으로 확장하는 데 주저하게 된 이유들 가운데, 같은
단어를 많은 서로 다른 것들에 적용하는 일을 무엇이 정당화하는가에 관한
지나치게 단순하고 기실 터무니없는(absurd) 견해가 어느 정도 역할을 한
것은 사실이다. 일반적인 분류 용어들의 확장을 통상적으로 인도하는 원칙
유형들의 다양상(variety)은 법리학에서 너무 자주 무시되어 왔다. 그럼에도
국제법에 관한 의문의 원천들은 단어 사용에 관한 이러한 잘못된 견해들보다
더 깊고 더 흥미롭다. 더욱이 이 질문을 짧게 처리하는 방식이 제시하는 두
대안, 곧 ‘기존 관행을 준수할 것인가, 아니면 그것에서 벗어날 것인가?’는
망라적이지 않다. 그 둘 외에도, 실제로 기존 용례를 인도해 온 원칙들을
명시하고 검토하는 대안이 있기 때문이다.

제안된 이 짧은 처리 방식은 우리가 고유명사를 다루고 있었다면 기실
적절했을 것이다. 누군가 ‘런던’이라 불리는 장소가 \emphksans{정말로}
런던인지를 묻는다면, 우리가 할 수 있는 일은 그에게 관행(convention)을
환기시키고 그가 그것을 따르거나 취향에 맞는 다른 이름을 선택하도록 남겨
두는 것뿐이다. 그런 경우에 어떤 원칙에 따라 런던이 그렇게 불리게
되었는지, 또 이 원칙이 받아들일 만한 것인지를 묻는 것은 부조리할 것이다.
이것이 부조리한 이유는, 고유명사의 배분은 \emphksans{오직} 임시적(\emph{ad
hoc}) 관행에만 기초하는 반면, 어떤 진지한 학문분야의 일반 용어들의
확장은, 그 원칙 또는 이론적 근거(rationale)가 무엇인지 분명하지 않을
수는 있어도, 결코 원칙이나 근거 없이 이루어지지 않기 때문이다. 현재
경우처럼 사실상 ‘우리는 그것이 법이라고 불린다는 것은 알지만, 그것이
정말 법인가?’라고 말하는 이들이 그 확장에 의문을 제기할 때, 요구되는
것은, 분명 흐릿하게나마, 그 원칙을 명시하고 그 자격
근거들(credentials)을 심사하라는 것이다.

우리는 국제법의 법적 성격에 관한 두 가지 주요한 의문의 원천들과,
그것들과 함께, 이론가들이 이러한 의문들에 대응하기 위해 취해 온 조치들을
검토할 것이다. \textbf{(p.~216)} 두 형태의 의문은 모두, 법이 무엇인가에
관한 명료하고 표준적인 사례로 간주되는 국내법과 국제법을 불리하게
비교하는 데서 생겨난다. 첫째 의문은 법을 근본적으로 위협으로 뒷받침된
명령들의 문제로 보는 구상(conception)에 깊이 뿌리박고 있으며, 국제법의
\emphksans{규칙들(rules)}의 성격을 국내법의 규칙들의 성격과 대조한다. 둘째
형태의 의문은, 국가들이 근본적으로 법적 의무의 수범자들(subjects)이 될
수 없다는 모호한 믿음에서 솟아나며, 국제법의 \emphksans{수범자들(subjects)}의
성격을 국내법의 수범자들의 성격과 대조한다.

\subsection{2. 의무들과 제재들 (OBLIGATIONS AND
SANCTIONS)}\label{uxc758uxbb34uxb4e4uxacfc-uxc81cuxc7acuxb4e4-obligations-and-sanctions}

우리가 검토할 의문들은 국제법 책들의 첫 장들에서 흔히 ‘국제법은 어떻게
구속력 있을 수 있는가?’라는 질문 형식으로 표현된다. 그러나 이렇게 즐겨
쓰이는 질문 형식에는 매우 혼란스러운 점이 있다. 그것을 다루기 전에
우리는 답이 결코 분명하지 않은 선행 질문에 직면해야 한다. 이 선행 질문은
이렇다. 법체계 전체에 관하여 그것이 ‘구속력 있다(binding)’고 말할 때
무엇을 뜻하는가? 한 체계의 특정 규칙이 특정 사람에게 구속력 있다고 하는
진술은 법률가들에게 익숙하고 의미상 상당히 명료하다. 우리는 그것을,
문제의 규칙이 유효한 규칙이며 그 규칙 아래에서 문제의 사람에게 어떤
의무(obligation) 또는 책무(duty)가 있다는 단언으로 바꾸어 말할 수 있다.
이 외에도 이 형식의 더 일반적인 진술들이 이루어지는 몇몇 상황들이 있다.
우리는 어떤 상황들에서 어느 법체계가 특정 사람에게 적용되는지 의문을
가질 수 있다. 그러한 의문들은 국제사법(conflict of laws)이나
공국제법(public international law)에서 생길 수 있다. 전자의 경우 우리는
특정 거래에 관하여 프랑스법과 영국법 중 어느 것이 특정 사람에게 구속력
있는지를 물을 수 있고, 후자의 경우 우리는 예컨대 적국이 점령한 벨기에의
주민들이 망명정부가 벨기에법이라고 주장한 것에 구속되었는지, 아니면
점령국의 명령들(ordinances)에 구속되었는지를 물을 수 있다. 그러나 이 두
경우 모두에서 질문들은 어떤 법체계, 곧 국내법이든 국제법이든 어떤 체계
\emphksans{내부에서(within)} 생겨나고 그 체계의 규칙들이나 원칙들을 참조하여
확정되는 법문제들(questions of law)이다. 그것들은 규칙들의 일반적 성격을
문제 삼지 않고, 다만 주어진 상황에서 특정 사람들 또는 거래들에 대한
규칙들의 범위 또는 적용가능성(applicability)을 문제 삼는다.
\textbf{(p.~217)} 분명히 ‘국제법은 구속력 있는가?’라는 질문과 그 동류
질문들인 ‘국제법은 어떻게 구속력 있을 수 있는가?’ 또는 ‘무엇이 국제법을
구속력 있게 만드는가?’는 다른 차원의 질문들이다. 그것들은 적용가능성이
아니라 국제법의 일반적 법적 지위에 관한 의문을 표현한다. 이 의문은
‘이러한 규칙들이 언제라도 의미 있게 그리고 참되게 의무들을 발생시킨다고
말해질 수 있는가?’라는 형식으로 더 솔직하게 표현될 것이다. 책들의
논의들이 보여 주듯이, 이 점에 관한 의문의 한 원천은 단순히 그 체계에
중앙적으로 조직된 제재들이 부재한다는 것이다. 이것은 국내법과의 불리한
비교 지점 하나인데, 국내법의 규칙들은 의심할 여지 없이 ‘구속력 있는’
것으로 그리고 법적 의무의 전형들로 간주된다. 이 단계에서 더 나아가는
논변은 단순하다. 이러한 이유로 국제법의 규칙들이 ‘구속력 있는’ 것이
아니라면, 그것들을 법으로 분류하는 일을 진지하게 받아들이는 것은 확실히
방어될 수 없다는 것이다. 일상 언어의 방식들이 아무리 관대하더라도,
이것은 간과하기에는 너무 큰 차이다. 법의 본성에 관한 모든
사변(speculation)은, 법의 존재가 적어도 일정한 행위를
의무적인(obligatory) 것으로 만든다는 가정에서 출발한다.

이 논변을 검토하면서 우리는 국제체계의 사실들에 관한 모든 의문점을 그
논변에 유리하게 보아 줄 것이다. 우리는 국제연맹 규약 제16조도 유엔 헌장
제7장도 국내법의 제재들과 동일시될 수 있는 어떤 것도 국제법에 도입하지
않았다고 받아들일 것이다. 한국전쟁에도 불구하고, 또 수에즈 사건에서 어떤
교훈(moral)이 도출될 수 있더라도, 우리는 헌장의 법집행
규정들(provisions)은 그 사용이 중요할 때마다 거부권에 의해 마비될
개연성이 있으며, 단지 종이 위에만 존재한다고 말해야 한다고 상정할
것이다.

국제법은 조직화된 제재들이 결여되어 있기 때문에 구속력 있지 않다고
논변하는 것은, 법이 본질적으로 위협으로 뒷받침된 명령들의 문제라는 이론
속에 들어 있는 의무 분석을 묵시적으로 받아들이는 것이다. 이 이론은,
우리가 보았듯이, ‘의무를 가진다’ 또는 ‘구속된다’를 ‘불복종에 대해 위협된
제재나 처벌을 받을 개연성이 있다’와 동일시한다. 그러나 우리가
논변했듯이, 이러한 동일시는 모든 법적 사고와 담론에서 \textbf{(p.~218)}
의무(obligation)와 책무(duty)의 관념들(ideas)이 수행하는 역할을
왜곡한다. 실효적인 조직화된 제재들이 있는 국내법에서조차, 우리는
제III장에서 제시된 다양한 이유들 때문에, ‘나는(너는) 불복종 때문에
고통을 겪을 개연성이 있다’라는 외부 예측 진술의 의미와, ‘나는(너는)
이렇게 행위할 의무를 가진다’라는 내부 규범 진술의 의미를 구별해야 한다.
후자는 특정 사람의 상황을 행위의 지도 기준들로 수용된 규칙들의 관점에서
평가한다. 모든 규칙들이 의무들 또는 책무들을 발생시키는 것은 아니라는
점은 참이다. 또한 그렇게 하는 규칙들이 일반적으로 사적 이익의 어느 정도
희생을 요구하고, 일반적으로 부합(conformity)하라는 진지한 요구들과
일탈들에 대한 집요한 비판에 의해 뒷받침된다는 점도 참이다. 그러나 우리가
예측적 분석과, 법을 본질적으로 위협으로 뒷받침된 명령으로 보는 그 모태
구상(parent conception)으로부터 벗어나고 나면, 의무라는 규범적 관념을
조직화된 제재들에 의해 뒷받침되는 규칙들로 제한할 좋은 이유는 없어
보인다.

그러나 우리는 그 논변의 또 다른 형태를 고려해야 한다. 그것은 위협된
제재들의 개연성으로 의무를 정의하는 데 묶여 있지 않기 때문에 더
그럴듯하다. 회의론자는, 우리가 스스로 강조했듯이, 국내체계 안에는
정당하게 필수적이라고 불리는 일정한 규정들(provisions)이 있다고 지적할
수 있다. 그 가운데에는 폭력의 자유로운 사용을 금지하는 의무의 일차
규칙들과, 이 규칙들 및 다른 규칙들에 대한 제재로서 힘(force)의 공식적
사용을 마련하는 규칙들이 있다. 그러한 규칙들과 그것들을 뒷받침하는
조직화된 제재들이 이 의미에서 국내법에 필수적이라면, 그것들은 국제법에도
똑같이 필수적이지 않은가? 그렇다는 주장은 이것이 ‘구속력
있는(binding)’이나 ‘의무(obligation)’ 같은 단어들의 바로 그 의미로부터
나온다고 고집하지 않고도 유지될 수 있다.

이 형태의 논변에 대한 답은, 국내법의 지속적인 심리적·물리적 배경을
이루는 인간 존재와 그 환경에 관한 그 기초적 진리들(elementary
truths)에서 발견된다. 신체적 힘과 취약성에서 근사적으로 동등한 개인들의
사회에서 물리적 제재들은 필수적이면서도 가능하다. 그것들은 법의 제약들에
자발적으로 복종하려는 이들이, 그러한 제재들이 없다면 스스로는 법을
존중하지 않으면서 타인들의 법 존중에서 오는 이익을 거둘
악행자들(malefactors)의 단순한 희생자가 되지 않게 하기 위해 요구된다.
서로 가까운 곳에서 사는 개인들 사이에서는 \textbf{(p.~219)} 공개적
공격에 의하지 않더라도 속임수로 타인들에게 손해를 가할 기회들이 매우
크고, 빠져나갈 가망(chances)도 상당하므로, 가장 단순한 사회 형태들을
제외하면 단순한 자연적 억제요인들만으로는 법에 복종하기에는 너무
사악하거나, 너무 어리석거나, 너무 약한 이들을 제어하기에 적절할 수 없다.
그러나 같은 근사적 동등성의 사실과 제약 체계에 복종할 때의 분명한 이점들
때문에, 악행자들의 어떤 결합도 그 체계의 유지에 자발적으로 협력하려는
이들의 힘을 넘어설 개연성은 없다. 국내법의 배경을 이루는 이러한
상황들에서는, 제재들이 비교적 작은 위험으로 악행자들에게 성공적으로
사용될 수 있고, 그 제재들의 위협은 존재할 수 있는 자연적 억제요인들에
많은 것을 더할 것이다. 그러나 개인들에게 성립하는 단순한 자명한 이치들이
국가들에게는 성립하지 않고, 국제법의 사실적 배경이 국내법의 사실적
배경과 매우 다르기 바로 때문에, 제재들에 대한 유사한 필수성도
없고(국제법이 그것들에 의해 뒷받침되는 것이 바람직할 수는 있더라도),
그것들을 안전하고 실효적으로 사용할 유사한 전망도 없다.

그 까닭은 국가들 사이의 침략이 개인들 사이의 침략과 매우 다르기
때문이다. 국가들 사이의 폭력 사용은 공적일 수밖에 없고, 국제 경찰력이
없더라도 그것이, 경찰력이 없는 경우의 살인이나 절도처럼, 침략자와 피해자
사이의 문제로 남으리라는 확실성은 거의 있을 수 없다. 전쟁을 개시하는
것은 최강대국에게조차 합리적 확신을 가지고 예측하기 드문 결과를 위해
많은 것을 위험에 거는 일이다. 다른 한편, 국가들의 불평등 때문에 국제질서
편에 선 이들의 결합된 힘이 침략의 유혹을 받는 세력들(powers)보다 우세할
개연성이 있다는 상시적 보장은 있을 수 없다. 따라서 제재들의 조직과
사용은 두려운 위험들을 수반할 수 있고, 그것들의 위협은 자연적
억제요인들에 거의 보태는 것이 없을 수 있다. 이렇게 매우 다른 사실의
배경에 맞서 국제법은 국내법과 다른 형식으로 발전해 왔다. 현대 국가의
인구 안에서 범죄의 조직화된 억압과 처벌이 없다면 폭력과 절도는 매시간
예상될 것이다. 그러나 국가들의 경우에는 참혹한 전쟁들 사이에 오랜 평화의
해들이 끼어들어 왔다. 이 평화의 해들은 전쟁의 위험과 이해관계 및
국가들의 상호 필요들을 고려할 때에만 \textbf{(p.~220)} 합리적으로 기대될
수 있다. 그러나 그것들은, 다른 것들 가운데 어느 중앙 기관에 의한 집행도
마련하지 않는다는 점에서 국내법의 규칙들과 다른 규칙들에 의해 규율할
가치가 있다. 그럼에도 이 규칙들이 요구하는 것은 의무적인(obligatory)
것으로 생각되고 말해진다. 규칙들에 부합(conformity)하라는 일반적 압력이
있고, 청구들(claims)과 시인들(admissions)은 그것들에 근거하며, 그것들의
위반은 보상에 대한 집요한 요구들뿐 아니라 보복조치들과 대응조치들을
정당화하는 것으로 여겨진다. 규칙들이 무시될 때, 그것은 규칙들이 구속력
있지 않다는 근거 위에서 이루어지는 것이 아니다. 오히려 사실들을 감추려는
노력들이 이루어진다. 물론 그러한 규칙들은 국가들이 싸우려 하지 않는
쟁점들에 관한 한에서만 실효적이라고 말할 수 있다. 그럴 수 있으며, 그것은
그 체계의 중요성과 인류에 대한 가치에 불리하게 반영될 수 있다. 그러나
그럼에도 이만큼이라도 확보될 수 있다는 사실은, 조직화된 제재들이
국내법에 필수적이라는 점, 곧 국내법의 물리적·심리적 사실 배경에서의 그
필수성으로부터, 매우 다른 배경의 국제법은 그것들 없이는 어떤 의무도
부과하지 않고 ‘구속력 있지’ 않으며 따라서 ‘법’이라는 명칭을 받을 가치가
없다는 결론으로 가는 단순한 추론은 할 수 없음을 보여 준다.

\subsection{3. 의무와 국가들의 주권 (OBLIGATION AND THE SOVEREIGNTY OF
STATES)}\label{uxc758uxbb34uxc640-uxad6duxac00uxb4e4uxc758-uxc8fcuxad8c-obligation-and-the-sovereignty-of-states}

영국, 벨기에, 그리스, 소비에트 러시아는 국제법 아래에서 권리들과
의무들을 가지며, 따라서 그 수범자들(subjects)에 속한다. 그것들은
일반인은 독립적이라고 생각하고 법률가는 ‘주권적(sovereign)’이라고 인정할
국가들의 임의적 사례들이다. 국제법의 의무적 성격에 관한 가장 지속적인
당혹의 원천들 가운데 하나는, 주권적인 국가가 국제법에 의해
‘구속될(bound)’ 수도 있고 국제법 아래에서 의무를 가질 수도 있다는 사실을
받아들이거나 설명하는 데서 느껴진 어려움이었다. 이러한 형태의 회의론은
어떤 의미에서는 국제법이 제재들을 결여하기 때문에 구속력 있지 않다는
반론보다 더 극단적이다. 왜냐하면 그 반론은 언젠가 국제법이 제재 체계에
의해 강화된다면 대응될 수 있지만, 현재의 반론은 동시에 주권적이고 법의
적용을 받는(subject to law) 국가라는 구상 안에 존재한다고 말해지거나
느껴지는 근본적 비일관성에 기초하기 때문이다.

이 반론의 검토는, 주권이라는 관념(notion)을 \textbf{(p.~221)} 국가
\emphksans{안의(within)} 입법부나 다른 어떤 요소 또는 인격에 적용하는 것이
아니라 국가 자체에 적용하는 경우를 면밀히 살피는 일을 포함한다.
법리학에서 ‘주권적(sovereign)’이라는 단어가 나타날 때마다, 그것을
자기보다 낮은 자들 또는 수범자들(subjects)에 대하여 자신의 말이 법인 법
위의 인격이라는 관념(idea)과 결부시키는 경향이 있다. 우리는 이 책의 앞
장들에서 이 매혹적인 관념(notion)이 국내법체계의 구조에 얼마나 나쁜
안내자인지를 보았다. 그러나 그것은 국제법 이론에서는 훨씬 더 강력한
혼란의 원천이었다. 물론 국가를 이러한 방식으로, 마치 일종의 초인, 곧
본래 법에 매이지 않지만 그 수범자들에게는 법의 원천인 존재인 것처럼
생각하는 것은 \emphksans{가능하다}. 16세기 이래 국가와 군주의 상징적
동일시(`L'état c'est moi')는 이러한 관념을 부추겼을 수 있고, 이 관념은
정치이론뿐 아니라 법이론의 많은 부분에 의심스러운 영감을 제공해 왔다.
그러나 국제법의 이해를 위해서는 이러한 연상들을 떨쳐 내는 것이 중요하다.
‘국가’라는 표현은 본래 또는 ‘본성상’ 법 바깥에 있는 어떤 인격이나 사물의
이름이 아니다. 그것은 두 사실들을 지칭하는 방식이다. 첫째, 한 영토에
거주하는 인구가 입법부, 법원들, 일차 규칙들이라는 특징적 구조를 가진
법체계가 마련한 질서 있는 정부 형식 아래에서 산다는 사실, 둘째, 그
정부가 모호하게 규정된 정도의 독립성을 누린다는 사실이다.

‘국가(state)’라는 단어가 분명 그 자체의 넓은 모호성 영역을 가지지만,
지금까지 말한 것으로 그 중심 의미를 드러내기에는 충분할 것이다. 다시
임의적 사례들을 들면, 영국이나 브라질, 미국이나 이탈리아 같은 국가들은
자기 국경 밖의 어떤 권위들이나 인격들에 의한 법적 통제와 사실적 통제
양자 모두로부터 매우 큰 정도의 독립성을 가지며, 국제법에서
‘주권국가들’로 자리매겨질 것이다. 다른 한편, 미국의 주들처럼 연방연합의
구성원인 개별 주들은 여러 다른 방식으로 연방정부와 연방헌법의 권위와
통제에 종속된다. 그러나 이 연방 주들조차 보유하는 독립성은, 이를테면
‘국가’라는 단어가 전혀 사용되지 않을 잉글랜드 카운티의 지위와 비교하면
크다. 카운티는 자기 구역을 위해 입법부의 기능들 일부를 수행하는
지방의회를 가질 수 있지만, 그 빈약한 권한들은 의회의 권한들에
\textbf{(p.~222)} 종속되며, 어떤 사소한 측면들을 제외하면 그 카운티의
구역은 나라의 나머지 부분과 같은 법들 및 정부의 적용을 받는다.

이 극단들 사이에는 질서 있는 정부를 가진 영토 단위들 사이의 종속성,
따라서 독립성의 많은 서로 다른 유형들과 정도들이 있다. 식민지, 보호령,
종주권관계(suzerainties), 신탁통치지역, 국가연합들은 이 관점에서
흥미로운 분류 문제들을 제시한다. 대부분의 경우 한 단위가 다른 단위에
종속되어 있다는 점은 법적 형식들로 표현되므로, 종속 단위의 영토 안에서
법인 것은 적어도 일정한 쟁점들에 관해서는 궁극적으로 다른 단위의 법제정
작용들(law-making operations)에 의존할 것이다.

그러나 어떤 경우에는 종속된 영토의 법체계가 그 종속성을 반영하지 않을 수
있다. 이는 그 영토가 단지 형식적으로 독립되어 있고 실제로는 꼭두각시들을
통해 외부에서 통치되기 때문일 수도 있고, 또는 종속된 영토가 내부
사무들에 대해서는 실질적 자율성을 가지지만 외부 사무들에 대해서는 그렇지
않으며, 외부 사무들에서 다른 나라에 대한 그 종속성이 그 국내법의 일부로
표현될 필요가 없기 때문일 수도 있다. 그러나 이러한 여러 방식으로 한 영토
단위가 다른 영토 단위에 종속되는 것이 그 독립성이 제한될 수 있는 유일한
형식은 아니다. 제한 요인은 그러한 다른 단위의 권력(power)이나
권위(authority)가 아니라, 서로 똑같이 독립적인 단위들에 영향을 미치는
국제적 권위일 수 있다. 국제적 권위의 많은 서로 다른 형식들과 그에
상응하는 국가 독립성의 많은 서로 다른 제한들을 상상할 수 있다. 그
가능성들에는, 많은 다른 것들 가운데, 모든 단위의 내부/외부 사무들을
규율할 법적으로 무제한적인 권한들(powers)을 가진, 영국 의회를 모델로 한
세계 입법부, 특정 사안들에 대해서만 법적 권능(legal competence)을
가지거나 구성 단위들의 특정 권리 보장에 의해 제한되는, 미국 의회를
모델로 한 연방 입법부, 법적 통제의 유일한 형식이 모두에게 적용가능한
것으로 일반적으로 수용된 규칙들로 이루어지는 체제, 마지막으로 인정되는
의무의 유일한 형식이 계약적 또는 자기부과적인 것이어서 국가의 독립성이
오직 자신의 행위에 의해서만 법적으로 제한되는 체제가 포함된다.

이러한 가능성들의 범위를 고려하는 것은 유익하다. 종속성과 독립성의
가능한 형식들과 \textbf{(p.~223)} 정도들이 많다는 점을 깨닫기만 해도,
국가들이 주권적이기 때문에 국제법의 적용을 받거나 국제법에 의해 구속될
수 ‘\emphksans{없다}’거나 어떤 특정 형식의 국제법에 의해서만 구속될 수
‘\emphksans{있다}’는 주장(claim)에 답하는 한 걸음이 되기 때문이다. 여기서
‘주권적’이라는 단어는 ‘독립적’이라는 말 이상을 뜻하지 않으며, 후자처럼
그 의미상 힘(force)이 부정적이다. 주권국가는 일정 유형들의 통제에
종속되지 \emphksans{않는} 국가이고, 그 주권은 그 국가가 자율적인 행위
영역이다. 우리가 보았듯이, 국가라는 단어의 바로 그 의미가 어느 정도의
자율성을 불러들이지만, 이것이 무제한적이어야 ‘\emphksans{한다}’거나 일정
유형들의 의무에 의해서만 제한될 수 ‘\emphksans{있다}’는 논지(contention)는
기껏해야 국가들이 여타 모든 제약들로부터 자유로워야 한다는 주장(claim)의
단언이고, 최악의 경우에는 이유 없는 독단이다. 실제로 국가들 사이에 어떤
주어진 형식의 국제적 권위가 존재한다는 것을 우리가 발견한다면, 국가들의
주권은 그만큼 제한되며, 그것은 규칙들이 허용하는 바로 그 정도를 가진다.
따라서 어떤 국가들이 주권적인지, 그리고 그 주권의 범위가 무엇인지는
규칙들이 무엇인지 알 때에만 알 수 있다. 이는 영국인이나 미국인이
자유로운지 그리고 그의 자유의 범위가 무엇인지가 영국법이나 미국법이
무엇인지 알 때에만 알 수 있는 것과 같다. 국제법의 규칙들은 많은 점에서
기실 모호하고 충돌하므로, 국가들에게 남겨진 독립성의 영역에 관한 의문은
국내법 아래 시민의 자유 범위에 관한 의문보다 훨씬 더 크다. 그럼에도
이러한 어려움들은, 국제법을 참조하지 않고 국가들에게 속한다고 가정된
절대적 주권으로부터 국제법의 일반적 성격을 연역하려는 \emphksans{선험적(a
priori)} 논변을 타당화하지 않는다.

주권 관념의 무비판적 사용이 국내법 이론과 국제법 이론 양쪽에 유사한
혼란을 퍼뜨려 왔고, 양쪽에서 유사한 교정을 요구한다는 점은 관찰할 가치가
있다. 그 영향 아래에서 우리는 모든 국내법체계에는 법적 제한들에 종속되지
않는 주권 입법자가 반드시 있어야 \emphksans{한다고} 믿게 된다. 마찬가지로
우리는 국가들이 주권적이고 스스로에 의한 제한을 제외하고는 법적 제한을
받을 수 없기 때문에 국제법은 반드시 어떤 성격을 가져야 \emphksans{한다고}
믿게 된다. 두 경우 모두, 법적으로 무제한적인 주권자가 필연적으로
존재한다는 믿음은 실제 \textbf{(p.~224)} 규칙들을 검토할 때에만 답할 수
있는 질문을 선취 판단한다. 국내법에 관한 질문은 이것이다. 이 체계에서
인정되는 최고 입법 권위(supreme legislative authority)의 범위는
무엇인가? 국제법에 관한 질문은 이것이다. 규칙들이 국가들에게 허용하는
최대 자율성 영역은 무엇인가?

따라서 현재의 반론에 대한 가장 단순한 답은, 그것이 질문들이 고려되어야
할 순서를 뒤집는다는 것이다. 국제법의 형식들이 무엇인지, 그리고 그것들이
단지 공허한 형식들인지 아닌지를 알기 전까지는 국가들이 어떤 주권을
가지는지 알 길이 없다. 많은 법리학적 논쟁은 이 원칙이 무시되었기 때문에
혼란스러워졌고, 이 원칙에 비추어 ‘의사주의(voluntarist)’ 이론들 또는
‘자기제한(auto-limitation)’ 이론들로 알려진 국제법 이론들을 고려하는
것은 유익하다. 이러한 이론들은 국가들의 (절대적) 주권을 국제법의 구속력
있는 규칙들의 존재와 조화시키려 하면서, 모든 국제적 의무들을 약속에서
생겨나는 의무처럼 자기부과적인 것으로 다루었다. 그러한 이론들은 사실
국제법에서 정치학의 사회계약론들에 대응하는 것이다. 후자는, ‘자연적으로’
자유롭고 독립적인 개인들이 그럼에도 국내법에 의해 구속된다는 사실들을,
법에 복종할 의무를 구속되는 이들이 서로, 그리고 몇몇 경우에는 자기
통치자들과 맺은 계약에서 생겨나는 의무로 다루어 설명하고자 했다. 우리는
여기서 이 이론을 문자 그대로 취할 때의 잘 알려진 반론들이나, 단지 해명적
유비로 취할 때의 그 가치에 대해서는 고려하지 않을 것이다. 그 대신 우리는
그 역사에서 의사주의 국제법 이론들에 반대하는 세 갈래 논변을 끌어낼
것이다.

첫째, 이러한 이론들은 국가들이 자기부과적 의무들에 의해서만 구속될 수
‘\emphksans{있다}’는 점이 어떻게 알려지는지, 또는 국제법의 실제 성격에 대한
어떤 검토보다 앞서 왜 그들의 주권에 관한 이러한 견해가 수용되어야
하는지를 전혀 설명하지 못한다. 그것을 뒷받침할 것이, 그것이 자주
반복되었다는 사실 외에 더 있는가? 둘째, 국가들이 그 주권 때문에
스스로에게 부과한 규칙들의 적용을 받거나 그 규칙들에 의해 구속될 수 있을
\emphksans{뿐}임을 보이도록 설계된 논변에는 비일관적인 점이 있다. ‘자기제한’
이론의 몇몇 매우 극단적인 형식들에서는 한 국가의 합의나 조약상 약정들이
그 국가가 제안하는 장래 행위에 관한 단순한 선언들로 취급되고, 이행하지
않음은 어떤 의무의 위반으로도 간주되지 않는다. \textbf{(p.~225)} 이것은
사실들과 매우 어긋나지만, 적어도 일관성이라는 이점(merit)을 가진다.
그것은 국가들의 절대적 주권이 어떤 종류의 의무와도 비일관적이어서,
의회처럼 국가도 자신을 구속할 수 없다는 단순한 이론이다. 그러나 국가가
약속, 합의 또는 조약에 의해 스스로에게 의무들을 부과할 수 있다는 덜
극단적인 견해는, 국가들이 이렇게 스스로에게 부과한 규칙들의 적용을
받거나 그 규칙들에 의해 구속될 수 있을 뿐이라는 이론과 비일관적이다.
말해진 것이든 쓰인 것이든 말들이 일정한 상황에서 약속, 합의 또는
조약으로 기능하여 의무들을 발생시키고 타인들이 청구할 수 있는 권리들을
부여하려면, 한 국가가 적절한 말들로 하겠다고 맡은 것은 무엇이든 하도록
구속된다고 규정하는 \emphksans{규칙들}이 이미 존재해야 하기 때문이다.
자기부과적 의무라는 바로 그 관념(notion) 안에 전제되어 있는 그러한
규칙들은, 그것들에 복종해야 한다는 자기부과적 의무에서 \emphksans{그것들의}
의무적 지위(obligatory status)를 도출할 수 없음이 분명하다.

주어진 국가가 해야 하도록 구속된 모든 특정 \emphksans{행위(action)}가 이론상
약속에서 그 의무적 성격을 도출할 수 있다는 것은 참이다. 그럼에도 이는
약속 등이 의무들을 창설한다는 \emphksans{규칙}이 어떤 약속과도 독립적으로 그
국가에 적용되는 경우에만 그럴 수 있다. 개인들로 이루어졌든 국가들로
이루어졌든 어떤 사회에서든, 약속, 합의 또는 조약의 말들이 의무들을
발생시키기 위해 필요하고 충분한 것은, 이를 마련하고 이러한 자기구속
작용들(self-binding operations)을 위한 절차를 특정하는 규칙들이,
보편적으로 그럴 필요는 없지만 일반적으로 인정되는 것이다. 그것들이
인정되는 곳에서는 이러한 절차들을 알고서 사용하는 개인이나 국가는,
구속되기를 선택하든 선택하지 않든 그 절차들에 의해 구속된다. 그러므로
사회적 의무의 이 가장 자발적인 형식조차, 그 규칙들에 의해 구속되는
당사자의 선택과 독립적으로 구속력 있는 몇몇 규칙들을 포함하며, 이는
국가들의 경우 그들의 주권이 그러한 모든 규칙들로부터의 자유를 요구한다는
상정과 비일관적이다.

셋째, 사실들이 있다. 우리는 방금 비판한 \emphksans{선험적} 주장, 곧 국가들은
자기부과적 의무들에 의해서만 구속될 수 ‘\emphksans{있다}’는 주장과, 국가들이
다른 체계 아래에서는 다른 방식들로 구속될 수 있었더라도 사실상 현행
국제법 규칙들 아래에서는 국가들을 위한 다른 어떤 의무 형식도 존재하지
않는다는 주장을 구별해야 한다. 물론 그 \textbf{(p.~226)} 체계가 이
완전히 합의적인(consensual) 형식의 것일 수도 있으며, 그 성격에 관한 이
견해의 단언들과 부인들은 법리학자들의 저술, 판사들의 의견, 심지어
국제법원들의 의견, 그리고 국가들의 선언들에서 발견된다. 오직 국가들의
실제 실천에 대한 냉정한 조사만이 이 견해가 옳은지 아닌지를 보여 줄 수
있다. 현대 국제법이 매우 큰 부분에서 조약법이라는 점은 참이며, 국가들의
사전 동의 없이도 국가들에게 구속력 있는 것으로 보이는 규칙들이 실제로는
동의에 기초해 있다는 것을 보이려는 정교한 시도들이 이루어져 왔다. 다만
이 동의는 오직 ‘묵시적으로(tacitly)’ 주어진 것이거나 ‘추론되어야’ 하는
것일 수 있다. 모두가 허구인 것은 아니지만, 국제적 의무의 형식들을 하나로
환원하려는 이 시도들 중 적어도 일부는, 우리가 보았듯이 국내법의
유사하지만 더 명백히 허위적인 단순화를 수행하도록 설계된 ‘묵시적
명령’이라는 관념(notion)과 같은 의심을 불러일으킨다.

모든 국제적 의무가 구속되는 당사자의 동의에서 생겨난다는 주장에 대한
상세한 검토는 여기서 수행할 수 없지만, 이 교설에 대한 두 가지 명료하고
중요한 예외들은 주목되어야 한다. 첫째는 신생국가의 경우이다. 1932년의
이라크와 1948년의 이스라엘처럼 새롭고 독립적인 국가가 존재로 출현할 때,
그 국가는 그중에서도 조약들에 구속력을 부여하는 규칙들을 포함한 국제법의
일반 의무들에 구속된다는 점은 결코 의문시된 적이 없다. 여기서 신생국가의
국제적 의무들을 ‘묵시적’ 또는 ‘추론된’ 동의에 기초시키려는 시도는 완전히
빈약해 보인다. 둘째는 국가가 영토를 획득하거나 어떤 다른 변화를 겪고, 그
변화가 이전에는 그 국가가 준수하거나 위반할 기회도 없었고 동의를 주거나
보류할 계기도 없었던 규칙들 아래에서 처음으로 의무들의
귀속양상(incidence)을 가져오는 경우이다. 이전에 바다에 접근할 수 없던
국가가 해양 영토를 획득한다면, 이것이 그 국가를 영해와 공해에 관한
국제법의 모든 규칙들의 수범자로 만들기에 충분하다는 점은 분명하다. 이
밖에도 주로 일반 조약이나 다자 조약들이 비당사자들에게 미치는 효과와
관련된 더 논쟁적인 사례들이 있다. 그러나 이 두 중요한 예외들만으로도,
모든 국제적 의무가 자기부과적이라는 일반 이론이 사실들에 대한 존중은
너무 적고 추상적 독단에 지나치게 이끌려 왔다는 의심은 정당화하기에
충분하다.

\subsection{4. 국제법과 도덕 (INTERNATIONAL LAW AND
MORALITY)}\label{uxad6duxc81cuxbc95uxacfc-uxb3c4uxb355-international-law-and-morality}

\textbf{(p.~227)} 제V장에서 우리는 의무의 일차 규칙들만으로 이루어진
단순한 사회구조 형식을 검토하였고, 가장 작고 가장 긴밀하게 결속되어
있으며 고립된 사회들을 제외하면 그것이 중대한 결함들을 겪는다는 점을
보았다. 그러한 체제는 정태적이어야 하고, 그 규칙들은 성장과 쇠퇴의 느린
과정들에 의해서만 변경된다. 규칙들의 식별은 불확실해야 하고, 특정
사례들에서 그 규칙들의 위반 사실을 확인하는 일과 위반자들에게 사회적
압력을 적용하는 일은 우발적이고, 시간을 낭비하며, 약할 수밖에 없다.
우리는 국내법에 특징적인 승인, 변경, 재판의 이차 규칙들을 이러한 서로
다른 결함들에 대한 서로 다르지만 관련된 치유책들(remedies)로 파악하는
것이 해명적임을 발견하였다. 국제법은 형식상 그러한 일차 규칙들의 체제를
닮아 있다. 그럼에도 그 흔히 정교한 규칙들의 내용은 원시 사회의 규칙들과
매우 다르고, 그 개념들, 방법들, 기법들 가운데 많은 것은 현대 국내법의
그것들과 같다. 법리학자들은 매우 자주 국제법과 국내법 사이의 이러한
형식적 차이들이 전자를 ‘도덕(morality)’으로 분류함으로써 가장 잘 표현될
수 있다고 생각해 왔다. 그러나 이런 방식으로 차이를 표시하는 것은 혼란을
불러들이는 것으로 보임이 분명하다.

때로는 국가들 사이의 관계들을 규율하는 규칙들이 오직 도덕 규칙들일
뿐이라는 고집이, 위협으로 뒷받침된 명령들로 환원될 수 없는 모든 사회구조
형식은 오직 ‘도덕’의 한 형식일 수밖에 없다는 오래된 독단주의에서 영감을
받는다. 물론 ‘도덕’이라는 단어를 이처럼 매우 포괄적인 방식으로 사용하는
것은 가능하다. 그렇게 사용되면 그것은 게임, 클럽, 예절의 규칙들, 헌법과
국제법의 기본 규정들(provisions), 그리고 잔혹함, 부정직, 거짓말의 공통된
금지처럼 우리가 통상 도덕적인 것들로 생각하는 규칙들과 원칙들이 함께
들어갈 개념적 폐지통을 제공한다. 이 절차에 대한 반론은, 이렇게
‘도덕’으로 함께 분류된 것들 사이에는 형식과 사회적 기능 양쪽에서 너무
중요한 차이들이 있어서, 그토록 조악한 분류로는 생각할 수 있는 어떤
실천적 또는 이론적 목적도 봉사될 수 없다는 것이다. 이렇게 인위적으로
넓혀진 도덕의 범주 안에서는, 그것이 흐리는 오래된 구별들을 새로 다시
획정해야 할 것이다.

\textbf{(p.~228)} 국제법의 특수한 경우에는 그 규칙들을 ‘도덕’으로
분류하는 데 저항할 여러 서로 다른 이유들이 있다. 첫째는 국가들이 종종
부도덕한 행위에 대해 서로를 비난하거나, 국제 도덕의 기준에 부응했다는
이유로 자신들이나 타국들을 칭찬한다는 점이다. 국가들이 보일 수도 있고
보이지 못할 수도 있는 미덕들 가운데 \emphksans{하나}가 국제법을 준수하는
것임은 의심할 여지가 없지만, 그것이 그 법이 도덕이라는 뜻은 아니다.
실제로 도덕의 관점에서 국가들의 행위를 평가하는 일은, 국제법 규칙들
아래에서 청구들, 요구들, 권리들과 의무들의 인정들을 정식화하는 일과 식별
가능하게 다르다. 제IX장에서 우리는 사회적 도덕의 규정적 특징들(defining
characteristics)로 취해질 수 있는 일정한 특징들을 열거하였다. 그
가운데에는 도덕 규칙들이 일차적으로 뒷받침되는 도덕적 압력의 독특한
형식이 있었다. 이것은 두려움에 대한 호소나 보복의 위협이나 보상 요구로
이루어지지 않고, 말이 건네진 사람이 문제된 도덕 원칙을 일단 상기하게
되면 죄책감이나 수치심에 이끌려 그것을 존중하고 잘못을 바로잡게(make
amends) 될 수 있으리라는 기대 속에서 이루어지는 양심에 대한 호소로
이루어진다.

국제법 아래에서의 청구들은 이러한 용어들로 표현되지 않는다. 물론
국내법에서처럼 그것들은 도덕적 호소와 결합될 수 있다. 국가들이 국제법의
다툼 있는 사안들에 관하여 서로에게 제시하는, 흔히 기술적인 논변들에서
우세한 것은 선례들, 조약들, 법리학 저술들에 대한 참조들이다. 흔히 도덕적
옳음이나 그름, 좋음이나 나쁨에 대한 언급은 전혀 이루어지지 않는다.
따라서 베이징 정부가 포모사에서 국민당 세력들을 추방할 국제법 아래의
권리를 가지고 있다거나 가지고 있지 않다는 주장은, 이것이 공정한지,
정의로운지, 또는 도덕적으로 좋은 일인지 나쁜 일인지의 질문과 매우
다르며, 특징적으로 서로 다른 논변들에 의해 뒷받침된다. 의심할 여지 없이
국가들 사이의 관계들에는, 사적 생활에서 인정되는 예의와 공손의 기준들과
유비적인, 분명히 법인 것과 분명히 도덕인 것 사이의 중간지대들이 있다.
개인적 사용을 위한 물품을 무관세로 받을 수 있도록 외교사절에게 확대된
특권에서 예시되는 국제 예양(international comity)의 영역이 그러하다.

더 중요한 구별 근거는 다음과 같다. 국제법의 규칙들은 국내법의 규칙들처럼
\textbf{(p.~229)} 흔히 도덕적으로 상당히 무관하다. 하나의 규칙은 그것이
다루는 주제들에 관하여 어떤 명료하고 고정된 규칙을 가지는 것이
편리하거나 필수적이기 때문에 존재할 수 있지, 그 특정 규칙에 어떤 도덕적
중요성이 부착되기 때문에 존재하는 것은 아닐 수 있다. 그것은 가능했던
많은 규칙들 중 하나일 수 있고, 그중 어느 하나라도 똑같이 잘 기능했을 수
있다. 따라서 국내적·국제적 법규칙들은 흔히 많은 구체적 세부사항을 담고,
도덕 규칙들이나 원칙들 안의 요소들로서는 이해될 수 없을 자의적 구분들을
긋는다. 사회적 도덕의 가능한 내용에 관하여 우리가 독단적이어서는 안
된다는 점은 참이다. 제IX장에서 보았듯이 한 사회 집단의 도덕은 현대
지식의 빛에서 보면 터무니없거나(absurd) 미신적인 것으로 보일 수 있는
많은 명령을 포함할 수 있다. 그러므로 우리와 매우 다른 일반적 믿음을 가진
사람들이 도로의 오른쪽이 아니라 왼쪽으로 운전하는 데 \emphksans{도덕적}
중요성을 부착하게 되거나, 두 명의 증인이 지켜본 약속을 어기면 도덕적
죄책감을 느끼지만 한 명이 지켜본 경우에는 그러한 죄책감을 느끼지 않게
된다고 상상하는 것은 어렵지만 가능하다. 그러한 낯선 도덕들이
가능하더라도, 도덕은 그것을 따르는 이들이 일반적으로 대안들보다 결코 더
선호할 만하지 않고 아무 내재적 중요성도 없다고 여기는 규칙들을
논리적으로 포함할 수 없다는 점은 여전히 참이다. 그러나 법은 도덕적으로
중요한 많은 것을 포함하기도 하지만, 바로 그러한 규칙들을 포함할 수 있고
실제로 포함한다. 따라서 도덕의 일부로 이해하기는 가장 어려울 자의적
구분들, 형식절차들(formalities), 고도로 구체적인 세부사항은 법의
자연스럽고 쉽게 이해되는 특징들이다. 도덕과 달리 법의 전형적 기능 가운데
하나는 확실성과 예측가능성을 극대화하고 청구들의 증명 또는 평가를
용이하게 하기 위해 바로 이러한 요소들을 도입하는 것이기 때문이다.
형식들과 세부사항에 대한 지나친 고려는 법으로 하여금 ‘형식주의’와
‘법률주의’라는 비난을 얻게 했다. 그러나 이러한 악덕들이 법의 독특한
성질들 일부를 과장한 것임을 기억하는 것이 중요하다.

이 이유 때문에, 유효하게 작성된 유언장이 몇 명의 증인을 가져야 하는지를
도덕이 아니라 국내법체계가 말해 주기를 기대하듯이, 교전국 선박이 급유나
수리를 위해 \textbf{(p.~230)} 중립항에 머물 수 있는 날짜 수, 영해의 폭,
그것을 측정하는 데 사용될 방법들과 같은 것들은 도덕이 아니라 국제법이
말해 주기를 기대한다. 이 모든 것들은 \emphksans{법규칙들}이 정하는 것이
필수적이고 바람직한 규정 사항들(provisions)이다. 그러나 그러한 규칙들이
몇몇 형식들 중 어느 것을 취해도 똑같이 괜찮을 수 있거나, 특정 목적들에
이르는 많은 가능한 수단들 가운데 하나로서만 중요하다는 감각이 유지되는
한, 그것들은 개인적 또는 사회적 삶에서 도덕에 특징적인 지위를 가지는
규칙들과 구별되어 남아 있다. 물론 국제법의 모든 규칙들이 이런
형식적이거나 자의적이거나 도덕적으로 중립적인 종류인 것은 아니다. 요점은
오직 법규칙들은 이런 종류일 \emphksans{수} 있고 도덕 규칙들은 그럴 \emphksans{수
없다}는 것이다.

국제법과 우리가 자연스럽게 도덕으로 생각하는 것 사이의 성격 차이에는 또
다른 측면이 있다. 일정한 실천들을 요구하거나 금지하는 법의 효과가
궁극적으로 한 집단의 도덕에 변화를 일으키는 것일 수는 있지만, 도덕
규칙들을 만들거나 폐지하는 입법부라는 관념(notion)은, 제VII장에서
보았듯이, 부조리한 것이다. 입법부는 자기 \emph{fiat}에 의해 새 규칙을
도입하고 그것에 도덕 규칙의 지위를 부여할 수 없다. 같은 수단으로 한
규칙에 전통의 지위를 부여할 수 없는 것과 마찬가지이다. 다만 왜
그러한지는 두 경우에서 같지 않을 수 있다. 따라서 도덕은 단지 입법부를
결여하거나 우연히 입법부를 가지지 않는 것이 아니다. 인간의 입법적
\emph{fiat}에 의한 변경이라는 바로 그 관념이 도덕이라는 관념에 배치된다.
우리가 도덕을 인간 행위들, 입법적 행위들이든 그 밖의 행위들이든 그것들을
평가하는 궁극적 기준으로 파악하기 때문에 그러하다. 국제법과의 대조는
명료하다. 국제법의 본성이나 기능 안에는 그 규칙들이 입법적 변경의 대상이
될 수 있다는 관념과 유사하게 비일관적인 것이 없다. 입법부의 결여는 많은
이들이 언젠가 수선되어야 할 결함으로 생각하는 바로 그런 결여일 뿐이다.

끝으로 우리는 국제법 이론에서, 제IX장에서 비판된 논변, 곧 국내법의 특정
규칙들이 도덕과 충돌할 수 있더라도 체계 전체는 그 규칙들에 복종할 도덕적
의무가 있다는 일반적으로 확산된 확신에 기초해야 하며, 이것은 특별한
예외적 사례들에서만 압도될 수 있다는 논변과의 병행성을 주목해야 한다.
국제법의 ‘토대들(foundations)’에 관한 논의에서, \textbf{(p.~231)}
최종적으로 국제법의 규칙들은 그것들에 복종할 도덕적 의무가 있다는
국가들의 확신에 기초해야 \emphksans{한다}고 자주 말해져 왔다. 그러나 이것이
국가들이 인정하는 의무들이 공적으로 조직된 제재들에 의해 집행가능하지
않다는 것 이상을 뜻한다면, 그것을 받아들일 이유는 없어 보인다. 물론 한
국가가 국제법에 의해 요구되는 어떤 행위 과정을 도덕적으로 의무적인
것으로 여겼고 그 이유로 행위했다고 우리가 말하는 것을 확실히 정당화할
상황들을 생각할 수 있다. 예컨대 조약들에 대한 신뢰가 심하게 흔들리면
인류에게 명백한 해악이 따를 것이기 때문에, 또는 자신이 그 차례에
타인들에게 부담이 떨어졌을 때 과거에 이익을 얻었던 규약(code)의 성가신
부담들을 짊어지는 것이 공정할 뿐이라는 감각 때문에, 한 국가는 부담스러운
조약의 의무들을 계속 이행할 수 있다. 도덕적 확신에 관한 그러한
사안들에서 정확히 누구의 동기들, 생각들(thoughts), 감정들이 국가에
귀속되어야 하는지는 여기서 우리를 붙들어 둘 필요가 없는 질문이다.

그러나 그러한 도덕적 의무의 감각이 있을 \emphksans{수} 있다 하더라도,
국제법의 존재 조건으로 그것이 왜 또는 어떤 의미에서 반드시 존재해야
\emphksans{하는지} 보기는 어렵다. 국가들의 실천에서 일정한 규칙들이 일정한
희생을 대가로 하면서도 반복적으로 일양적인 존중을 받는다는 점, 그것들을
참조하여 청구들이 정식화된다는 점, 규칙들의 위반은 위반자를 중대한
비판에 노출시키고 보상 또는 보복에 대한 청구들을 정당화하는 것으로
여겨진다는 점은 분명하다. 이것들은 확실히 국가들 사이에 그들에게
의무들을 부과하는 규칙들이 존재한다는 진술을 뒷받침하는 데 요구되는 모든
요소들이다. 어떤 사회에서 ‘구속력 있는’ 규칙들이 존재한다는 입증은,
단순히 그것들이 그러한 것으로 생각되고, 말해지고, 기능한다는 것이다.
‘토대들’의 방식으로 무엇이 더 요구되며, 더 요구된다면 왜 그것은 도덕적
의무의 토대여야 하는가? 물론 압도적 다수가 규칙들을 수용하고 그것들의
유지에 자발적으로 협력하지 않는 한, 규칙들은 국가들 사이의 관계에서
존재하거나 기능할 수 없다는 것은 참이다. 또한 규칙들을 깨뜨리거나
깨뜨리겠다고 위협하는 이들에게 행사되는 압력은 흔히 비교적 약하고,
대체로 탈중심화되어 있거나 조직화되어 있지 않다는 점도 참이다. 그러나
훨씬 더 강하게 강제적인 국내법체계를 자발적으로 수용하는 개인들의 경우와
마찬가지로, 그러한 체계를 자발적으로 지지하는 동기들은 극도로 다양할 수
있다. 어떤 \textbf{(p.~232)} 형식의 법질서든 그것에 부합하는 것이
도덕적으로 의무적이라는 감각이 일반적으로 퍼져 있을 때 가장 건강할 수는
있다. 그럼에도 법에 대한 고수(adherence)는 그것에 의해 동기부여되지 않을
수 있고, 장기적 이익의 계산이나 전통을 계속하려는 바람, 또는 타인들에
대한 사심 없는 관심에 의해 동기부여될 수 있다. 개인들 사이에서든 국가들
사이에서든, 이들 중 어느 하나를 법의 존재에 필요한 조건으로 식별할 좋은
이유는 없어 보인다.

\subsection{5. 형식과 내용의 유비들 (ANALOGIES OF FORM AND
CONTENT)}\label{uxd615uxc2dduxacfc-uxb0b4uxc6a9uxc758-uxc720uxbe44uxb4e4-analogies-of-form-and-content}

순진한 눈에는, 입법부, 강제관할권을 가진 법원들, 공적으로 조직된
제재들을 결여한 국제법의 형식적 구조가 국내법의 그것과 매우 달라 보인다.
그것은 우리가 말했듯이 내용에서는 전혀 아니지만 형식에서는 일차법 또는
관습법의 단순한 체제를 닮아 있다. 그러나 어떤 이론가들은 국제법이
‘법’이라고 불릴 자격(title)을 회의론자에 맞서 방어하려는 초조함 속에서,
이러한 형식적 차이들을 최소화하고 국제법 안에서 입법이나 국내법의 다른
바람직한 형식적 특징들과 관련하여 발견될 수 있는 유비들을 과장하려는
유혹에 굴복하였다. 그래서 패전국이 영토를 양도하거나, 의무들을
부담하거나, 어떤 축소된 독립 형식을 받아들이는 조약으로 끝나는 전쟁은
본질적으로 입법적 행위라고 주장되어 왔다. 입법처럼 그것은 부과된 법적
변경이기 때문이다. 오늘날 이 유비에 감명받거나, 그것이 국제법도 국내법과
동등하게 ‘법’이라고 불릴 자격(title)을 가진다는 점을 보이는 데 도움이
된다고 생각할 사람은 거의 없을 것이다. 국내법과 국제법 사이의 두드러진
차이들 가운데 하나는, 전자는 보통 폭력으로 강취된 합의들의 유효성을
인정하지 않지만 후자는 인정한다는 점이기 때문이다.

‘법’이라고 불릴 자격(title)이 그것들에 달려 있다고 생각하는 이들은 여러
다른, 더 존중할 만한 유비들을 강조해 왔다. 국제사법재판소와 그 전신인
상설국제사법재판소의 판결들이 거의 모든 사례에서 당사자들에 의해 제대로
이행되어 왔다는 사실은, 국내법 법원들과 대조적으로 어떤 국가도 그 사전
동의 없이는 이러한 국제재판소들 앞에 세워질 수 없다는 사실을 그것이
어떻게든 상쇄하는 것처럼 자주 강조되어 왔다. 또한 국내법에서 제재로서
법적으로 규율되고 공적으로 집행되는 \textbf{(p.~233)} 힘의 사용과
‘탈중심화된 제재들’, 곧 국제법 아래의 자기 권리들이 타국에 의해
침해되었다고 주장하는 국가가 전쟁이나 무력 보복에 의존하는 것 사이에서도
유비들이 발견되어 왔다. 어느 정도 유비가 있다는 것은 명백하다. 그러나 그
의의는, 국내법 법원은 ‘자력구제’의 옳음과 그름을 조사하고 그것에 그르게
의존한 일을 처벌할 강제관할권을 가지는 반면, 어떤 국제법원도 유사한
관할권을 가지지 않는다는 똑같이 명백한 사실에 비추어 평가되어야 한다.

이 의심스러운 유비들 중 일부는 국가들이 유엔 헌장 아래에서 부담한
의무들에 의해 훨씬 강화된 것으로 여겨질 수 있다. 그러나 여기서도, 그
강도에 대한 어떤 평가도, 종이 위에서는 훌륭한 헌장의 법집행
규정들(provisions)이 거부권 및 강대국들의 이념적 분열과 동맹들에 의해
마비되어 온 정도를 무시한다면 거의 가치가 없다. 때로 국내법의 법집행
규정들(provisions)도 총파업에 의해 마비될 \emphksans{수도} 있다고 답해지지만,
이것은 거의 설득력 있지 않다. 국내법과 국제법 사이의 비교에서 우리는
사실상 존재하는 것에 관심을 가지며, 여기서 사실들은 부인할 수 없이
다르기 때문이다.

그러나 국제법과 국내법 사이에 제안된 하나의 형식적 유비는 여기서 어느
정도 검토할 가치가 있다. 켈젠과 많은 현대 이론가들은 국내법처럼 국제법도
‘기초규범(basic norm)’, 또는 우리가 승인 규칙이라고 부른 것을 가지고
있고 기실 가져야만 한다고 고집한다. 체계의 다른 규칙들의 유효성은 그것을
참조하여 평가되며, 규칙들은 그것 덕분에 단일한 체계를 구성한다는 것이다.
반대 견해는 이러한 구조의 유비가 거짓이라는 것이다. 국제법은 단지 이러한
방식으로 통합되어 있지 않은, 분리된 의무의 일차 규칙들의
\emphksans{집합(set)}으로 이루어져 있다는 것이다. 국제법은 국제법학자들의
통상 용어로, 조약들에 구속력을 부여하는 규칙이 그중 하나인 관습 규칙들의
집합이다. 그 과제에 착수한 이들이 국제법의 ‘기초규범’을 정식화하는 데
매우 큰 어려움들을 발견했다는 것은 잘 알려져 있다. 이 지위의 후보들에는
\emph{pacta sunt servanda} 원칙이 포함된다. 그러나 이것은, 그 용어를
아무리 넓게 해석하더라도 국제법 아래의 모든 \textbf{(p.~234)} 의무들이
'\emph{pacta}'에서 생겨나는 것은 아니라는 사실과 양립하지 않는 것으로
보이기 때문에, 대부분의 이론가들에 의해 포기되었다. 그래서 그것은 덜
친숙한 어떤 것, 곧 이른바 ‘국가들은 자신들이 관습적으로 행위해 온 대로
행위해야 한다’는 규칙으로 대체되었다.

우리는 국제법의 기초규범에 관한 이러한 공식화들과 다른 경쟁적 공식화들의
이점들(merits)을 논의하지 않을 것이다. 그 대신 우리는 국제법이 그러한
요소를 반드시 포함해야 한다는 가정을 문제 삼을 것이다. 여기서 물어야 할
첫 질문이자 어쩌면 마지막 질문은 이것이다. 왜 우리는 이러한
\emphksans{선험적} 가정, 곧 그것은 바로 선험적 가정인데, 그것을 하고 그리하여
국제법 규칙들의 실제 성격을 선취 판단해야 하는가? 구성원들에게 의무들을
부과하는 규칙들을 ‘구속력 있는’ 것으로 삼아 살아가는 사회가, 그러한
규칙들을 어떤 더 기초적인 규칙에 의해 통합되거나 그 규칙에서 유효성을
도출하는 것이 아니라 단순히 분리된 규칙들의 집합으로 여길 수 있다는 것은
확실히 상상 가능하다. 어쩌면 그런 경우가 자주 있었을 수도 있다. 규칙들의
단순한 존재가 그러한 기초규칙의 존재를 수반하지 않는다는 것은 분명하다.
대부분의 현대 사회에는 예절 규칙들이 있고, 우리는 그것들이 의무들을
부과한다고 생각하지 않지만 그러한 규칙들이 존재한다고 말할 수 있다.
그러나 우리는 분리된 규칙들의 유효성이 도출될 수 있는 예절의 기초규칙을
찾지 않을 것이며, 찾을 수도 없을 것이다. 그러한 규칙들은 체계를 형성하지
않고 단순한 집합을 형성한다. 물론 예절보다 더 중요한 사안들이 걸려 있는
곳에서 이러한 사회통제 형식의 불편들은 상당하다. 그것들은 이미 제V장에서
서술되었다. 그러나 규칙들이 실제로 행위 기준들로 수용되고, 의무적
규칙들에 특징적인 적절한 형식의 사회적 압력에 의해 뒷받침된다면,
그것들이 구속력 있는 규칙들임을 보이기 위해 더 요구되는 것은 없다. 비록
이 단순한 사회구조 형식에서는 국내법에서 우리가 가지는 것, 곧 체계의
어떤 궁극적 규칙을 참조하여 개별 규칙들의 유효성을 논증하는 방식을
가지지 못하더라도 그러하다.

물론 체계가 아니라 단순한 집합을 구성하는 규칙들에 대해서도 우리가 물을
수 있는 여러 질문들이 있다. 예컨대 우리는 그것들의 역사적 기원에 관한
질문들이나 규칙들의 성장을 촉진한 인과적 영향들에 관한 질문들을 물을 수
있다. 또한 우리는 그것들에 따라 살아가는 이들에게 그 규칙들이 가지는
가치, 그리고 그들이 자신들을 그 규칙들에 복종하도록 도덕적으로 구속된
것으로 여기는지 \textbf{(p.~235)} 아니면 어떤 다른 동기에서 복종하는지를
물을 수 있다. 그러나 더 단순한 경우에는, 국내법이 그러하듯 기초규범 또는
이차적 승인 규칙에 의해 풍부해진 체계의 규칙들에 관하여 물을 수 있는 한
종류의 질문을 물을 수 없다. 더 단순한 경우 우리는 이렇게 물을 수 없다.
‘분리된 규칙들은 그 유효성 또는 “구속력(binding force)”을 체계의 어떤
궁극적 규정(provision)으로부터 도출하는가?’ 그러한 규정이 없고, 있을
필요도 없기 때문이다. 그러므로 기초규칙 또는 승인 규칙이 의무의 규칙들
또는 ‘구속력 있는’ 규칙들의 존재에 일반적으로 필요한 조건이라고 상정하는
것은 오류이다. 이것은 필수성이 아니라 사치이며, 그 구성원들이 분리된
규칙들을 하나씩 받아들이게 될 뿐 아니라 일반적 유효성 기준들에 의해
획정된 규칙들의 일반 부류들을 미리 수용하기로 결속되어 있는 발전된
사회체계들에서 발견된다. 더 단순한 사회 형식에서는 한 규칙이 규칙으로
수용되는지 아닌지를 기다려 보아야 한다. 기초적 승인 규칙을 가진
체계에서는 규칙이 실제로 만들어지기 전에, 그것이 승인 규칙의 요건들에
부합한다면 \emphksans{유효할 것}이라고 말할 수 있다.

같은 요점은 다른 형식으로 제시될 수 있다. 그러한 승인 규칙이 분리된
규칙들의 단순한 집합에 추가될 때, 그것은 체계의 이점들과 식별의 용이성을
함께 가져올 뿐 아니라, 처음으로 진술의 새로운 형식을 가능하게 한다.
이것들은 규칙들의 유효성에 관한 내부 진술들이다. 우리는 이제 새로운
의미에서 ‘체계의 어떤 규정(provision)이 이 규칙을 구속력 있게 만드는가?’
또는 켈젠의 언어로 ‘체계 안에서 그것의 유효성의 근거는 무엇인가?’라고
물을 수 있기 때문이다. 이 새로운 질문들에 대한 답들은 기초적 승인 규칙에
의해 제공된다. 그러나 더 단순한 구조에서 규칙들의 유효성이 어떤 더
기초적인 규칙을 참조하여 이렇게 논증될 수 없다고 하더라도, 이것은
규칙들이나 그 구속력 또는 유효성에 관하여 설명되지 않은 채 남겨지는 어떤
질문이 있다는 뜻이 아니다. 그러한 단순한 사회구조에서 규칙들이 왜 구속력
있는지에 관한 어떤 신비가 있고, 우리가 기초규칙을 찾을 수만 있다면
그것이 그 신비를 해소할 것이라는 말이 아니다. 단순한 구조의 규칙들은, 더
발전된 체계들의 기초규칙처럼, 수용되고 그러한 것으로 기능한다면 구속력
있다. 그러나 서로 다른 사회구조 형식들에 관한 이러한 단순한 진리들은,
이러한 바람직한 요소들이 사실상 발견되지 않는 곳에서 통일성과
\textbf{(p.~236)} 체계에 대한 완고한 탐색에 의해 쉽게 가려질 수 있다.

기실 기초규칙 없이 존재하는 가장 단순한 사회구조 형식들을 위해
기초규칙을 만들어 내려고 기울이는 노력에는 우스꽝스러운 점이 있다.
그것은 마치 벌거벗은 야만인이 \emphksans{정말로는} 어떤 보이지 않는 종류의
현대적 의복을 입고 있어야 \emphksans{한다고} 우리가 고집하는 것과 같다.
불행히도 여기에는 상존하는 혼동 가능성도 있다. 우리는 관련 사회,
개인들의 사회이든 국가들의 사회이든, 그것이 일정한 행위 기준들을 의무적
규칙들로 준수한다는 단순한 사실을 텅 빈 반복으로 말하는 어떤 것을
기초규칙으로 다루도록 설득될 수 있다. 국제법에 대해 제안되어 온 기묘한
기초규범, 곧 ‘국가들은 자신들이 관습적으로 행위해 온 대로 행위해야
한다’는 바로 이런 지위를 가지는 것이 분명하다. 그것은 일정한 규칙들을
수용하는 이들이 그 규칙들이 준수되어야 한다는 규칙도 준수해야 한다는 것
이상을 말하지 않기 때문이다. 이것은 규칙들의 집합이 국가들에 의해 구속력
있는 규칙들로 수용된다는 사실의 단순하고 무익한 중복일 뿐이다.

다시, 국제법은 기초규칙을 포함해야 \emphksans{한다}는 가정으로부터 우리가
해방되고 나면, 직면해야 할 질문은 사실의 질문이다. 국가들 사이의
관계들에서 기능하는 규칙들의 실제 성격은 무엇인가? 관찰될 현상들에 관한
서로 다른 해석들은 물론 가능하다. 그러나 여기서 제출되는 견해는,
국제법의 규칙들에 대한 일반적 유효성 기준들을 제공하는 기초규칙은 없고,
실제로 작동하는 규칙들은 체계가 아니라 규칙들의 집합을 구성하며, 그
가운데 조약들의 구속력을 마련하는 규칙들이 있다는 것이다. 많은 중요한
사안들에서 국가들 사이의 관계들이 다자조약들에 의해 규율되고, 이것들이
당사자가 아닌 국가들도 구속할 수 있다고 때때로 논변되는 것은 사실이다.
이것이 일반적으로 인정된다면, 그러한 조약들은 사실상 입법적 제정물들이
될 것이고 국제법은 그 규칙들에 대한 독자적인 유효성 기준들을 가지게 될
것이다. 그때에는 체계의 실제 특징을 나타내고 규칙들의 집합이 사실상
국가들에 의해 준수된다는 사실의 공허한 재진술 이상이 될 기초적 승인
규칙이 정식화될 수 있을 것이다. 아마 국제법은 현재 이것 및 구조상 그것을
국내체계에 더 가까워지게 할 다른 형식들의 수용으로 향하는 이행 단계에
있을지도 모른다. 만일, \textbf{(p.~237)} 그리고 이 이행이 완료될 때에는,
현재로서는 빈약하고 심지어 기만적으로 보이는 형식적 유비들이 실질을
획득할 것이며, 국제법의 법적 ‘성질(quality)’에 관한 회의론자의 마지막
의문들은 그때 잠재워질 수 있다. 그 단계에 도달할 때까지 유비들은 분명
형식의 유비들이 아니라 기능과 내용의 유비들이다. 기능의 유비들은 우리가
국제법이 도덕과 다른 방식들을 성찰할 때 가장 명료하게 나타나며, 그
일부를 우리는 앞 절에서 검토하였다. 내용의 유비들은 국내법과 국제법
모두에 공통되고 법률가들의 기법을 한쪽에서 다른 쪽으로 자유롭게 이전
가능하게 만드는 원칙들, 개념들, 방법들의 범위에 있다. ‘국제법’이라는
표현의 발명자인 벤담은 그것이 국내법과 ‘충분히 유비적’이라고 말하는
것만으로 그 표현을 방어하였다.\footnote{\emph{Principles of Morals and
  Legislation}, XVII. 25, n.~1.} 여기에 두 가지 언급을 덧붙일 가치가
있을 수 있다. 첫째, 그 유비는 형식의 유비가 아니라 내용의 유비라는 것,
둘째, 이 내용의 유비에서 국제법의 규칙들만큼 국내법에 가까운 다른 사회
규칙들은 없다는 것이다.

\newpage

\subsection{제10장 주석}\label{uxc81c10uxc7a5-uxc8fcuxc11d}


{\footnotesize
\noteentry 214쪽. \notetitle{‘국제법은 정말 법인가?’} --- 이것이 단지 사실문제로 오인된
언어적 질문에 불과하다는 견해는 Glanville Williams, 앞서 인용, \emph{22
BYBIL} (1945)을 참조하라.

\noteentry 215쪽. \notetitle{의문의 원천들(Sources of doubt).} --- 건설적인 일반적
개관으로는 A. H. Campbell, `International Law and the Student of
Jurisprudence' in \emph{35 Grotius Society Proceedings} (1950); Gihl,
`The Legal Character and Sources of International Law' in
\emph{Scandinavian Studies in Law} (1957)을 참조하라.

\noteentry 216쪽. \notetitle{‘국제법은 어떻게 구속력 있을 수 있는가?’} --- 이
질문(때로는 국제법의 구속력(binding force)의 문제로 지칭됨)은 Fischer
Williams, \emph{Chapters on Current International Law}, pp.~11--27;
Brierly, \emph{The Law of Nations}, 제5판 (1955), 제2장; \emph{The Basis
of Obligation in International Law} (1958), 제1장에서 제기된다. 또한
Fitzmaurice, `The Foundations of the Authority of International Law and
the Problem of Enforcement' in \emph{19 MLR} (1956)을 참조하라. 이
저자들은 규칙들의 체계가 구속력 있다(또는 구속력 있지 않다)는 단언의
의미를 명시적으로 논의하지 않는다.

\noteentry 217쪽. \notetitle{국제법에서의 제재(Sanctions).} --- 국제연맹 규약(League of
Nations Covenant) 제16조 하에서의 상황은 Fischer Williams, `Sanctions
under the Covenant' in \emph{17 BYBIL} (1936)을 참조하라. 유엔 헌장
제7장 하에서의 제재에 관해서는 Kelsen, `Sanctions in International Law
under the Charter of U.N.', \emph{31 Iowa LR} (1946), 그리고 Tucker,
`The Interpretation of War under present International Law', \emph{4 The
International Law Quarterly} (1951)을 참조하라. 한국전쟁에 대해서는
Stone, \emph{Legal Controls of International Conflict} (1954), 제9장,
담론 14(Discourse 14)를 참조하라. ‘평화를 위한 단결(유엔총회 결의)’이
유엔이 ‘마비되었다(paralyzed)’는 주장을 반박하는 것이라는 해석도
가능하다.

\noteentry 220쪽. \notetitle{의무적인 것으로 생각되고 말해지는 국제법.} --- Jessup,
\emph{A Modern Law of Nations}, 제1장, 그리고 `The Reality of
International Law', \emph{118 Foreign Affairs} (1940)을 참조하라.

\noteentry 220쪽. \notetitle{국가들의 주권(The Sovereignty of States).} --- ‘주권은 단지
국제 영역 중 법에 의해 국가들의 개별 행위에 남겨진 만큼에 부여된 이름일
뿐이다’라는 견해에 대한 명료한 설명은 Fischer Williams, 앞서 인용,
pp.~10--11, 285--99, \emph{Aspects of Modern International Law},
pp.~24--26, 그리고 Van Kleffens, `Sovereignty and International Law',
\emph{Recueil des Cours} (1953), I, pp.~82--83을 참조하라.

\noteentry 221쪽. \notetitle{국가(The State).} --- ‘국가(state)’라는 관념(notion)과
종속국가들(dependent states)의 유형에 대해서는 Brierly, \emph{The Law of
Nations}, 제4장을 참조하라.

\noteentry 224쪽. \notetitle{의사주의(voluntarist) 및 자기제한(auto-limitation) 이론들.} --- 주요 저자들로는 Jellinek, \emph{Die Rechtliche Natur der
Staatsverträge}; Triepel, `Les Rapports entre le droit interne et le
droit internationale', \emph{Recueil des Cours} (1923)이 있다. 극단적인
견해는 Zorn, \emph{Grundzüge des Völkerrechts}에 나타나며, 이에 대한
비판적 논의는 Gihl, 앞서 인용, \emph{Scandinavian Studies in Law}
(1957); Starke, \emph{An Introduction to International Law}, 제1장;
Fischer Williams, \emph{Chapters on Current International Law},
pp.~11--16을 참조하라.

\noteentry 224쪽. \notetitle{의무(obligation)와 동의(consent).} --- 어떤 국제법 규칙도
한 국가의 사전 동의, 명시적이든 묵시적이든 그러한 동의 없이는 그 국가에
구속력 있지 않다는 견해는 영국 법원들에서 표현되었으며(\emph{R.} v.
\emph{Keyn} 1876, 2 Ex. Div. 63, `The Franconia' 사건 참조),
상설국제사법재판소에서도 표현되었다. 이에 관해서는 \emph{The Lotus},
PCIJ Series A, No.~10을 참조하라.

\noteentry 226쪽. \notetitle{신생국 및 해양영역을 획득한 국가들.} --- 이에 대해서는
Kelsen, \emph{Principles of International Law}, pp.~312--13을 참조하라.

\noteentry 226쪽. \notetitle{비당사자들(non-parties)에 대한 일반 국제조약들의 효과.} ---
이에 대해서는 Kelsen, 앞서 인용, pp.~345 ff.; Starke, 앞서 인용, 제1장;
Brierly, 앞서 인용, 제7장, pp.~251--2를 참조하라.

\noteentry 227쪽. \notetitle{‘도덕(morality)’이라는 용어의 포괄적 사용.} --- 이에
대해서는 Austin, \emph{The Province} 제5강(Lecture V), pp.~125--9,
141--2에서 ‘실정 도덕(positive morality)’에 관한 논의를 참조하라.

\noteentry 230쪽. \notetitle{국제법에 복종할 도덕적 의무(moral obligation).} --- 이것이
국제법의 ‘토대(the foundation)’라고 보는 견해는 Lauterpacht, Brierly의
\emph{The Basis of Obligation in International Law} 서문(xviii), 그리고
Brierly, 같은 책 제1장을 참조하라.

\noteentry 232쪽. \notetitle{힘에 의해 부과된 조약을 입법으로 보는 것.} --- 이에
대해서는 Scott, `The Legal Nature of International Law', \emph{American
Journal of International Law} (1907), pp.~837, 862--4를 참조하라. 일반
조약들을 ‘국제 입법(international legislation)’이라고 부르는 통상적
서술에 대한 비판은 Jennings, `The Progressive Development of
International Law and its Codification', \emph{24 BYBIL} (1947),
p.~303에서 확인할 수 있다.

\noteentry 233쪽. \notetitle{탈중심화된 제재들(decentralized sanctions).} --- 이에
대해서는 Kelsen, 앞서 인용, p.~20; Tucker, 앞서 인용, \emph{4
International Law Quarterly} (1951)을 참조하라.

\noteentry 233쪽. \notetitle{국제법의 기초규범(basic norm).} --- 이를 \emph{pacta sunt
servanda}로 공식화한 논의는 Anzilotti, \emph{Corso di diritto
internazionale} (1923), p.~40을 참조하라. 이를 ‘국가들은 자신들이
관습적으로 행위해 온 대로 행위해야 한다’는 식으로 대체하려는 논의는
Kelsen, \emph{General Theory}, p.~369, 그리고 \emph{Principles of
International Law}, p.~418을 참조하라. 이에 대한 중요한 비판적 논의는
Gihl, \emph{International Legislation} (1937) 및 \emph{Scandinavian
Studies in Law} (1957), pp.~62 ff.에 있다. 국제법에 기초규범이 없다는
해석을 보다 충분히 전개한 논의로는 Ago, `Positive Law and International
Law', \emph{51 American Journal of International Law} (1957),
\emph{Scienza giuridica e diritto internazionale} (1958)을 참조하라.
Gihl은 국제사법재판소 규정 제38조에도 불구하고 국제법은 법의 형식적
원천들을 갖지 않는다고 결론짓는다. 국제법에 대해 하나의 ‘초기
가설(initial hypothesis)’을 정식화하려는 시도로서, 본문에서 제기된 것과
유사한 비판들에 열려 있는 것으로 보이는 견해는 Lauterpacht, \emph{The
Future of Law in the International Community}, pp.~420--3을 참조하라.

\noteentry 237쪽. \notetitle{국제법과 국내법 사이의 내용의 유비(analogy of content).} --- 이에 대해서는 Campbell, 앞서 인용, \emph{35 Grotius Society
Proceedings} (1950), p.~121 \emph{ad fin.}, 그리고 조약들, 영토취득,
시효, 임대, 위임통치, 지역권(servitudes) 등에 관한 규칙들을 논의한
Lauterpacht, \emph{Private Law Sources and Analogies of International
Law} (1927)을 참조하라.

\subsection{제10장 제3판
주석}\label{uxc81c10uxc7a5-uxc81c3uxd310-uxc8fcuxc11d}

\noteentry 220--226쪽. \notetitle{주권국가들과 국제법(Sovereign states and international law).} 주권(sovereignty)과 국제법의 제약들(constraints)의 양립가능성에
관해서는 Timothy Endicott, 「The Logic of Freedom and Power」, Samantha
Besson 및 John Tasioulas 편, \emph{The Philosophy of International Law}
(Oxford University Press, 2010)을 참조하라. 국가 주권(state
sovereignty)과 법의 지배(rule of law) 간의 관계에 대해서는 Jeremy
Waldron, 「The Rule of International Law」, (2006) \emph{Harvard Journal
of Law and Public Policy} 제30권, 15쪽을 참조하라. 초국가적
권위(transnational authority)의 다양한 등장 형태들에 비추어 주권을
재해석하려는 시도들에 대해서는 Neil MacCormick, \emph{Questioning
Sovereignty}를 참조하라. 현대 세계에서 ‘주권적’ 국가들의 사실상의
권력(de facto power)에 대한 의문들에 대해서는 Saskia Sassen,
\emph{Losing Control? Sovereignty in an Age of Globalization} (Columbia
University Press, 1996)을 참조하라.

\noteentry 227--232쪽. \notetitle{국제법과 도덕 (International law and morality).} 국제
영역에서 정의의 원칙들(principles of justice)의 적용에 대한 매우 영향력
있는 논의는 John Rawls, \emph{The Law of Peoples} (Harvard University
Press, 1999)을 참조하라. 또한 Allen Buchanan, \emph{Justice, Legitimacy,
and Self-Determination: Moral Foundations for International Law} (Oxford
University Press, 2003)을 참조하라. 국가들이 오직 자국의 국익을 증진하기
위해서만 국제법을 준수하고(comply with), 또 준수해야 한다는
‘현실주의(realist)’ 관점에 대해서는 Jack L. Goldsmith 및 Eric A. Posner,
\emph{The Limits of International Law} (Oxford University Press, 2005)를
참조하라.

\noteentry 232--237쪽. \notetitle{형식과 내용의 유비들 (Analogies of form and content).} Hart가 저술하던 시기에 비해 오늘날의 국제법은 더 체계화되었다는 시사에
대해서는 Samantha Besson 및 John Tasioulas가 편집한 \emph{The Philosophy
of International Law} 서문의 'Introduction'과 거기에서 인용된 원천들을
참조하라. 또한 Allen Buchanan 및 David Golove, 「Philosophy of
International Law」 in Jules L. Coleman 및 Scott Shapiro 편, \emph{The
Oxford Handbook of Jurisprudence and Philosophy of Law}를 참조하라.
}

\setcounter{footnote}{0}
\section{후기 (POSTSCRIPT)}\label{uxd6c4uxae30-postscript}

\subsection{서론 (INTRODUCTORY)}\label{uxc11cuxb860-introductory}

\textbf{(p.~238)} 이 책이 처음 출판된 지 32년이 지났다. 그 이후
법리학(jurisprudence)과 철학은 훨씬 더 가까워졌고, 법이론(legal
theory)이라는 주제는 이 나라와 미국 양쪽에서 크게 발전하였다. 학계의
법률가들과 철학자들 사이에서 이 책의 주요 교설들(doctrines)에 대한
비판자들이 그 교설들에 동조한 이들만큼이나 많았더라도, 나는 이 책이
이러한 발전을 자극하는 데 어느 정도 도움을 주었다고 생각하고 싶다.
어쨌든, 내가 원래 이 책을 영국 학부생 독자들을 염두에 두고 썼지만, 이
책은 훨씬 더 넓게 유통되었고 영어권 세계와 이 책의 번역본들이 나온 여러
나라에서 비판적 논평의 방대한 부차 문헌을 낳았다. 이 비판 문헌의 많은
부분은 법학 및 철학 저널의 논문들로 이루어져 있지만, 그 밖에도 이 책의
여러 교설들이 비판의 표적이자 비판자들 자신의 법이론 전개의 출발점으로
취급된 중요한 책들이 여럿 출판되었다.

나는 내 비판자들 가운데 몇몇, 특히 고 론 풀러 교수\footnote{그의
  \emph{The Morality of Law} (1964)에 대한 나의 서평을 보라,
  \emph{Harvard Law Review} 제78권 (1965) 1281면. 이 글은 \emph{Essays
  in Jurisprudence and Philosophy} (1983)에도 재수록됨, 343면. {[}주:
  대괄호로 표시된 각주는 편집자들이 추가한 것임.{]}}와 R. M. 드워킨
교수\footnote{나의 다음 글을 참조하라: “Law in the Perspective of
  Philosophy: 1776--1976”, \emph{New York University Law Review} 제51권
  (1976) 538면; “American Jurisprudence through English Eyes: The
  Nightmare and the Noble Dream”, \emph{Georgia Law Review} 제11권
  (1977) 969면; “Between Utility and Rights”, \emph{Columbia Law
  Review} 제79권 (1979) 828면. 이들 모두 \emph{Essays in Jurisprudence
  and Philosophy}에 재수록되어 있다. 또한, “Legal Duty and
  Obligation”, \emph{Essays on Bentham} (1982)의 제6장, 그리고 R.
  Gavison 편, \emph{Issues in Contemporary Legal Philosophy} (1987) 수록
  “Comment”(35면)를 참조하라.}를 향해 몇 발의 견제 사격을 가한 적은
있지만, 지금까지 그들 중 누구에게도 일반적이고 포괄적인 답변을 하지는
않았다. 나는 일부 비판자들이 \textbf{(p.~239)} 나와 달랐던 만큼이나
서로도 달랐던, 매우 유익한 계속 논쟁을 지켜보고 거기서 배우는 편을
선호했다. 그러나 이 후기에서는 드워킨이 \emph{Taking Rights Seriously}
(1977)와 \emph{A Matter of Principle} (1985)에 수록된 여러 선구적 논문들
및 그의 책 \emph{Law's Empire} (1986)\footnote{이하에서는 각각
  \emph{TRS}, \emph{AMP}, \emph{LE}로 인용함.}에서 제기한 광범위한
비판들 가운데 일부에 답하려고 한다. 이 후기에서 나는 주로 드워킨의
비판들에 초점을 맞춘다. 그는 이 책의 거의 모든 독특한 테제들이
근본적으로 잘못되었다고 논변했을 뿐 아니라, 이 책에 암묵적으로 들어 있는
법이론의 전체 구상(conception)과 법이론이 무엇을 해야 하는지의 전체
구상을 문제 삼았기 때문이다. 이 책의 주요 주제들에 대한 드워킨의
논변들은 여러 해 동안 대체로 일관되었지만, 몇몇 논변들의 실질과 그것들이
표현되는 용어법 양쪽에는 중요한 변화들이 있었다. 그의 초기 논문들에서
두드러졌던 몇몇 비판들은, 명시적으로 철회되지는 않았지만, 그의 후기
작업에는 나타나지 않는다. 그러나 그러한 초기 비판들은 널리 통용되고 매우
영향력이 있으므로, 나는 그의 후기 비판들뿐 아니라 그것들에도 답하는 것이
적절하다고 생각했다. 이 후기의 첫째이자 더 긴 부분은 드워킨의 논변들을
다룬다. 그러나 둘째 부분에서는, 내 테제들 가운데 일부에 대한 나의 전개에
불명료함과 부정확성이 있을 뿐 아니라 어떤 지점들에서는 실제 비일관성과
모순도 있다는 여러 다른 비판자들의 주장들을 검토한다.\footnote{{[}하트는
  본 절에서 언급된 두 개 절 중 두 번째를 완성하지 못함. 편집자 주
  참조.{]}} 여기서 나는, 생각하고 싶지 않을 만큼 많은 경우에 내
비판자들이 옳았음을 인정해야 하며, 이 후기를 기회로 삼아 불명료한 것을
명료화하고 내가 원래 쓴 것 가운데 비일관적이거나 모순적인 곳을
수정하고자 한다.

\subsection{1. 법이론의 본성 (THE NATURE OF LEGAL
THEORY)}\label{uxbc95uxc774uxb860uxc758-uxbcf8uxc131-the-nature-of-legal-theory}

이 책에서 나의 목표는 일반적이면서도 기술적인(descriptive), 법이
무엇인가에 관한 이론을 제공하는 것이었다. 그것이 \emphksans{일반적}이라는
것은, 그것이 어떤 특정 법체계나 법문화에 묶여 있지 않고,
규칙-지배적이고(rule-governed) 그 의미에서 ‘규범적(normative)’인 측면을
가진 복합적 사회·정치 제도로서의 법에 관하여 설명적이고 명료화하는
설명(account)을 제시하려 한다는 의미이다. 이 \textbf{(p.~240)} 제도는
서로 다른 문화와 서로 다른 시대의 많은 변이에도 불구하고 동일한 일반적
형식과 구조를 취해 왔지만, 명료화를 요구하는 많은 오해들과 흐리게 하는
신화들이 그 주변에 모여 있었다. 이 명료화 작업의 출발점은, 이 책 3쪽에서
내가 어느 교육받은 사람에게나 귀속시킨 근대적 국내법체계의 두드러진
특징들에 관한 널리 퍼진 공통 지식이다. 나의 설명(account)은 도덕적으로
중립적이고 정당화적 목적들을 가지지 않는다는 점에서 \emphksans{기술적}이다.
그것은 법에 관한 나의 일반적 설명에 나타나는 형식들과 구조들을 도덕적
근거 또는 다른 근거 위에서 정당화하거나 권장하려 하지 않는다. 다만 나는
이것들에 대한 명료한 이해가 법에 대한 어떤 유익한 도덕적 비판에도 중요한
예비 단계라고 생각한다.

이 기술적 기획을 수행하는 수단으로서, 내 책은 \emphksans{책무부과
규칙들(duty-imposing rules), 권한부여 규칙들(power-conferring rules),
승인 규칙들(rules of recognition), 변경 규칙들(rules of change),
규칙들의 수용(acceptance of rules), 내부 관점과 외부 관점(internal and
external points of view), 내부 진술들과 외부 진술들(internal and
external statements)}, 그리고 \emphksans{법적 유효성(legal validity)} 같은
여러 개념들을 반복해서 사용한다. 이 개념들은 여러 법제도들과 법적
실천들이 해명적으로 분석될 수 있는 항들이 되는 요소들에 주의를
집중시키고, 이러한 제도들과 실천들에 관한 성찰이 촉발한 법의 일반적
본성에 관한 질문들에 해답들이 제시될 수 있게 한다. 여기에는 다음과 같은
질문들이 포함된다. 규칙들이란 무엇인가? 규칙들은 단순한 습관들이나
행위의 반복적 일양성들과 어떻게 다른가? 법규칙들에는 근본적으로 서로
다른 유형들이 있는가? 규칙들은 어떻게 관련될 수 있는가? 규칙들이 하나의
체계를 형성한다는 것은 무엇인가? 법규칙들과 그것들이 가진 권위는
한편으로 위협들과, 다른 한편으로 도덕적 요구사항들과 어떻게
관련되는가?\footnote{H. L. A. Hart, “Comment”, 위의 Gavison 편, 주석
  2, 35면.}

이런 방식으로 기술적이면서 일반적인 것으로 구상된 법이론은, 드워킨이
법이론을, 또는 그가 자주 부르듯 ‘법리학(jurisprudence)’을, 부분적으로
평가적이고 정당화적이며 ‘특정 법문화에 향한(addressed to a particular
legal culture)’ 것으로 보는 구상과 근본적으로 다른 기획이다.\footnote{\emph{LE}
  102면.} 그 법문화는 대개 이론가 자신의 법문화이고, 드워킨의 경우에는
영미법의 법문화이다. 이렇게 구상된 법이론의 중심 과제는 드워킨에 의해
‘해석적(interpretive)’이라고 불리며\footnote{\emph{LE} 제3장.}
부분적으로 평가적이다. 왜냐하면 그것은 한 법체계의 확립된 법(settled
law)과 법적 실천들에 가장 잘 ‘부합(fit)’하거나 그것들과 정합하는(cohere)
동시에 그것들에 대한 최선의 도덕적 정당화를 제공하는 \textbf{(p.~241)}
원칙들을 식별하는 데, 그리하여 법을 ‘그 최선의 빛 아래’ 보이는 데 있기
때문이다.\footnote{\emph{LE} 90면.} 드워킨에게 이렇게 식별된 원칙들은
법에 관한 이론의 일부일 뿐 아니라 법 자체의 암시적 부분들이기도 하다.
그래서 그에게 “법리학은 재판(adjudication)의 일반적 부분이며, 법에 의한
어떤 결정(any decision at law)에나 붙는 침묵의 서곡”이다.\footnote{\emph{LE}
  90면.} 그의 초기 작업에서 그러한 원칙들은 단순히 “법에 관한 가장
건전한 이론(the soundest theory of law)”으로 지칭되었지만,\footnote{\emph{TRS}
  66면.} 그의 최신 작업인 \emph{Law's Empire}에서 그는 이러한 원칙들과
그것들에서 따라 나오는 특정 법명제들(propositions of law)을 ‘해석적
의미’에서의 법으로 특징짓는다. 그러한 해석이론이 해석해야 할 확립된 법적
실천들 또는 법의 전형들(paradigms)은 드워킨에 의해
‘전해석적(preinterpretive)’이라고 불리며,\footnote{\emph{LE} 65--66면.}
이론가는 그러한 전해석적 자료들을 식별하는 데 어려움도 없고 수행해야 할
이론적 과제도 없는 것으로 여겨진다. 그것들은 특정 법체계들의 법률가들의
일반적 합의의 문제로 확립되어 있기 때문이다.\footnote{그러나 드워킨은
  그러한 전해석적 법(preinterpretive law)을 식별하는 일 자체가 해석을
  수반할 수 있다고 경고한다. \emph{LE} 66면 참조.}

나 자신의 법이론 구상과 드워킨의 법이론 구상처럼 그렇게 서로 다른 기획들
사이에 왜 중대한 충돌이 있어야 하는지, 또는 기실 있을 수 있는지는
분명하지 않다. 따라서 \emph{Law's Empire}를 포함한 드워킨 작업의 많은
부분은, 법, 곧 ‘과거의 정치적 결정들’\footnote{\emph{LE} 93면.}이 강제를
정당화하는 방식에 관한 세 가지 서로 다른 설명들의 비교상 이점들을
정교하게 전개하는 데 바쳐져 있으며, 그로부터 그가
‘관행주의(conventionalism)’, ‘법적 실용주의(legal pragmatism)’,
‘통합성으로서의 법(law as integrity)’이라고 부르는 세 가지 서로 다른
법이론 형식들이 나온다.\footnote{\emph{LE} 94면.} 그가 이 세 유형의
이론에 관해 쓰는 모든 것은 평가적·정당화적 법리학에 대한 기여로서 큰
흥미와 중요성을 가진다. 나는, 이 책에서 제시된 것과 같은 실증주의적
법이론이 그러한 해석이론으로 해명적으로 재진술될 수 있다고 그가 주장하는
한을 제외하고, 이러한 해석적 관념들을 그가 전개한 것에 이의를 제기하려는
것이 아니다.\footnote{그러나 일부 비평가들, 예컨대 Michael Moore는
  “\,`The Interpretive Turn in Modern Theory: A Turn for the
  Worse?'\,”, \emph{Stanford Law Review} 제41권 (1989) 871면에서
  947--948면, 드워킨의 의미에서 법실천이 해석적이라는 점에는
  동의하면서도, 법이론은 해석적일 수 없다고 주장한다.} 이 후자의 주장은
\textbf{(p.~242)} 내 견해로는 잘못되었으며, 아래에서 내 이론의 그러한
해석적 판본에 반대하는 이유들을 제시한다.

그러나 드워킨은 그의 책들에서 일반적이고 기술적인 법이론을 잘못 이끌린
것이거나, 기껏해야 단순히 쓸모없는 것으로 배제하는 듯하다. 그는 “유용한
법이론들”은 “역사적으로 발전하는 실천의 특정 단계에 대한 해석”이라고
말하며,\footnote{\emph{LE} 102면. 또한, “우리에게 법의 일반 이론들은
  우리 자신의 사법 실천에 대한 일반 해석들이다”라는 진술도 비교하라.
  \emph{LE} 410면.} 이전에는 “기술(description)과 평가(evaluation)
사이의 평면적 구별(flat distinction)”이 “법이론을 약화시켰다”고
썼다.\footnote{\emph{AMP} 148면. 또한, “\,‘법이론들은 사회적 실천에
  대한 중립적 설명들로 \ldots{} 합당하게 이해될 수 없다’\,”는 진술도
  비교하라. Ronald Dworkin, “A Reply by Ronald Dworkin”, Marshall
  Cohen 편, \emph{Ronald Dworkin and Contemporary Jurisprudence} (1983),
  이하 \emph{RDCJ}로 인용, 247면 중 254면.}

나는 드워킨이 기술적 법이론, 또는 그가 자주 부르듯 ‘법리학’을 거부하는
정확한 이유들을 따라가기 어렵다. 그의 중심 반론은, 법이론은 하나의
법체계 안의 내부자 또는 참여자의 관점인 법에 대한 내부 관점(internal
perspective)을 고려해야 하며, 참여자의 관점이 아니라 외부 관찰자의
관점을 가진 기술 이론으로는 이 내부 관점에 대한 적절한 설명이 제공될 수
없다는 것처럼 보인다.\footnote{{[}\emph{LE} 13--14면 참조.{]}} 그러나 내
책에서 예시된 기술적 법리학의 기획에는, 비참여자 외부 관찰자가
참여자들이 그러한 내부 관점에서 법을 바라보는 방식들을 기술하는 일을
배제하는 것이 사실상 전혀 없다. 그래서 나는 이 책에서 참여자들이 법을
자기 행위에 대한 지침들과 비판 기준들을 제공하는 것으로 수용함으로써
자신의 내부 관점을 드러낸다고 상당히 길게 설명했다. 물론 기술적
법이론가는 그러한 자격으로는 참여자들이 이러한 방식으로 법을 수용하는
것을 스스로 공유하지 않는다. 그러나 그는 그러한 수용을 기술할 수 있고 또
기술해야 하며, 기실 나는 이 책에서 그렇게 하려고 시도하였다. 이 목적을
위해 기술적 법이론가는 내부 관점을 채택한다는 것이 무엇인지
\emphksans{이해}해야 하고, 그러한 제한된 의미에서 그는 내부자의 자리에 자신을
둘 수 있어야 한다는 것은 참이다. 그러나 이것은 법을 수용하거나, 내부자의
내부 관점을 공유하거나 지지(endorse)하거나, 어떤 다른 방식으로든 자신의
기술적 입장을 포기하는 것이 아니다.

드워킨은 기술적 법리학에 대한 자신의 비판에서, 외부 관찰자가 참여자의
내부 관점을 이런 기술적 방식으로 고려할 수 있다는 \textbf{(p.~243)} 이
명백한 가능성을 배제하는 듯하다. 앞서 말했듯이 그는 법리학을
‘재판(adjudication)의 일반적 부분’과 동일시하고, 이것은 법리학 또는
법이론 자체를 그 사법적 참여자들의 내부 관점에서 보인 한 체계의 법의
일부로 취급하는 것이기 때문이다. 그러나 기술적 법이론가는 법에 관한
내부자의 내부 관점을 채택하거나 공유하지 않고도 그것을 이해하고 기술할
수 있다. 설령 (닐 매코믹\footnote{{[}\emph{Legal Reasoning and Legal
  Theory} (1978), 63--64면 및 139--140면 참조.{]}}과 많은 다른
비판자들이 논변했듯이) 행위 지침들과 비판 기준들을 제공하는 것으로 법을
수용하는 데서 드러나는 참여자의 내부 관점이, 법의 요구사항들에 부합할
\emphksans{도덕적} 이유들이 있고 법의 강제 사용에 대한 \emphksans{도덕적} 정당화가
있다는 믿음도 필연적으로 포함한다 하더라도, 이것 역시 도덕적으로
중립적인 기술적 법리학이 기록할 것이지 지지하거나 공유할 것은 아니다.

그러나 드워킨이 ‘해석적’이라고 부르는 부분적으로 평가적인 쟁점들이
법리학과 법이론의 유일하게 적절한 쟁점들이 아니며 일반적이고 기술적인
법리학에도 중요한 자리가 있다는 나의 주장에 대해, 그는 이것이 그렇다는
점을 인정하였다. 그리고 그는 ‘법리학은 재판의 일반적 부분’이라는 자신의
관찰들과 같은 것들이 한정될 필요가 있다고 설명하였다. 그가 이제 말하듯,
이것은 오직 ‘뜻의 문제(question of sense)에 관한 법리학에 대해서만
참’이기 때문이다.\footnote{R. M. Dworkin, “Legal Theory and the Problem
  of Sense”, R. Gavison 편, \emph{Issues in Contemporary Legal
  Philosophy: The Influence of H. L. A. Hart} (1987), 19면.} 이것은
법이론의 유일하게 적절한 형식은 해석적이고 평가적인 것이라는 과도하고,
기실 드워킨 자신이 표현했듯 ‘제국주의적’인 주장으로 보였던 것에 대한
중요하고 환영할 만한 교정이다.

그러나 드워킨이 이제 자신의 겉보기에 제국주의적인 주장의 철회와 결부시킨
다음의 주의를 주는 말들(cautionary words)의 함의는 나에게 여전히 매우
당혹스럽다. “그러나 하트의 것과 같은 일반 이론들이 주로 논의해 온
쟁점들에서 그 {[}뜻의{]} 문제가 얼마나 널리 퍼져 있는지를 강조할 가치가
있다.”\footnote{앞 인용과 동일.} 이 주의의 관련성은 분명하지 않다. 내가
논의한 쟁점들(위 p.~240의 목록 참조)에는 \textbf{(p.~244)} 한편으로
강제적 위협들과, 다른 한편으로 도덕적 요구사항들과 법의 관계 같은
질문들이 포함된다. 드워킨의 주의의 요점은, 그러한 쟁점들을 논의할 때
기술적 법이론가조차 법명제들(propositions of law)의 뜻(sense) 또는
의미(meaning)에 관한 질문들에 직면해야 하며, 이 질문들은 해석적이고
부분적으로 평가적인 법이론에 의해서만 만족스럽게 답될 수 있다는 것처럼
보인다. 이것이 정말 사실이라면, 어떤 주어진 법명제의 뜻을 결정하기 위해
기술적 법이론가조차 “이 명제가 확립된 법에 가장 잘 부합하고 그것을 가장
잘 정당화하는 원칙들에서 따라 나오려면, 이 명제에는 어떤 의미가
배정되어야 하는가?”라는 해석적이고 평가적인 질문을 묻고 답해야 할
것이다. 그러나 내가 언급한 종류의 질문들에 대한 답을 찾는 일반적이고
기술적인 법이론가가 많은 서로 다른 법체계들에서 법명제들의 의미를
결정해야 한다는 것이 참이라 하더라도, 이것이 드워킨의 해석적이고
평가적인 질문을 묻는 일에 의해 결정되어야 \emphksans{한다}는 견해를 받아들일
이유는 없어 보인다. 더구나 일반적이고 기술적인 법이론가가 고려해야 하는
모든 법체계들의 판사들과 법률가들 자신이 사실상 의미의 질문들을 이러한
해석적이고 부분적으로 평가적인 방식으로 확정한다고 하더라도, 이것은
그러한 법명제들의 의미에 관한 자기의 일반적 기술적 결론들을 기초 지을
하나의 사실로서 일반적 기술 이론가가 기록해야 할 일일 것이다. 이러한
결론들이 그렇게 기초 지어졌다는 이유로, 그 결론들 자체가 해석적이고
평가적이어야 하며, 그것들을 제시하면서 이론가가 기술의 과제에서 해석과
평가의 과제로 옮겨 갔다고 상정하는 것은 물론 중대한 오류일 것이다.
기술(description)은 기술되는 것이 평가(evaluation)일 때에도 여전히
기술일 수 있다.

\subsection{2. 법실증주의의 본성 (THE NATURE OF LEGAL
POSITIVISM)}\label{uxbc95uxc2e4uxc99duxc8fcuxc758uxc758-uxbcf8uxc131-the-nature-of-legal-positivism}

\subsubsection{\texorpdfstring{(i) \emphksans{의미론적 이론으로서의 실증주의}
(Positivism as a Semantic
Theory)}{(i) 의미론적 이론으로서의 실증주의 (Positivism as a Semantic Theory)}}\label{i-uxc758uxbbf8uxb860uxc801-uxc774uxb860uxc73cuxb85cuxc11cuxc758-uxc2e4uxc99duxc8fcuxc758-positivism-as-a-semantic-theory}

내 책은 드워킨에 의해 현대 법실증주의의 대표작으로 취급된다. 이 현대
법실증주의는 주로 벤담과 오스틴의 명령 이론들(imperative theories of
law)과, 모든 법이 법적으로 무제한적인 주권적 입법 \textbf{(p.~245)} 인격
또는 기관으로부터 나온다는 그들의 구상을 거부한다는 점에서 그들 같은
초기 판본들과 구별된다. 드워킨은 나의 법실증주의 판본에서 서로 다르지만
관련된 많은 오류들을 발견한다. 이 오류들 가운데 가장 근본적인 것은, 법적
권리들과 법적 책무들(legal duties)을 서술하는 법명제들(propositions of
law) 같은 것들의 참은 개인적 믿음들과 사회적 태도들에 관한 사실들을
포함한 명백한 역사적 사실에 관한 질문들에만 의존한다는
견해이다.\footnote{\emph{LE} 6면 이하.} 법명제들의 참이 의존하는
사실들은 드워킨이 ‘법의 근거들(grounds of law)’이라고 부르는 것을
구성하며,\footnote{\emph{LE} 4면.} 그에 따르면 실증주의자는 이것들이
판사들과 법률가들이 공유하는 언어 규칙들에 의해 고정되어 있다고 잘못
받아들인다. 이 언어 규칙들은 `law'라는 단어가 특정 체계의 특정 쟁점에
관한 ‘당해 법(the law)’이 무엇인지에 대한 진술들에 나타날 때와
‘법(law)’, 곧 일반적 법이 무엇인지에 관한 진술들에 나타날 때 모두 그
단어의 사용, 따라서 의미를 지배한다고 한다.\footnote{\emph{LE} 31면
  이하.} 이러한 실증주의적 법관에 따르면, 법문제들(questions of law)에
관해 있을 수 있는 유일한 의견불일치들은 그러한 역사적 사실들의 존재 또는
부존재에 관한 것들이며, 무엇이 법의 ‘근거들’을 구성하는지에 관한 이론적
의견불일치나 논쟁은 있을 수 없게 된다.

드워킨은 법실증주의 비판의 많은 해명적 지면을, 무엇이 법의 근거들을
구성하는지에 관한 이론적 의견불일치가 실증주의자의 견해와 반대로 영미법
실천들의 두드러진 특징임을 보이는 데 할애한다. 법률가들과 판사들이
공유하는 언어 규칙들에 의해 그것들이 논쟁의 여지 없이 고정되어 있다는
견해에 맞서, 드워킨은 그것들이 본질적으로 논쟁적이라고 주장한다. 그
가운데에는 역사적 사실들뿐 아니라 매우 자주 논쟁적인 도덕 판단들과 가치
판단들도 있기 때문이다.

드워킨은 나 같은 실증주의자들이 어떻게 이 근본적으로 잘못된 견해를
채택하게 되었는지에 관한 매우 다른 두 설명을 제시한다. 이 설명들 가운데
첫째에 따르면, 실증주의자들은 법의 근거들이 규칙들에 의해 논쟁의 여지
없이 고정되어 있지 않고 이론적 의견불일치를 허용하는 논쟁적 사안이라면,
`law'라는 단어가 서로 다른 사람들에게 서로 다른 것을 \emphksans{뜻하게(mean)}
될 것이고, 그 단어를 사용하면서 그들은 같은 것에 관해 소통하는 것이
아니라 서로 빗나가 말하게 될 것이라고 믿는다. 이렇게 실증주의자에게
귀속된 믿음은 드워킨의 견해로는 전적으로 잘못된 것이며, 그는
\textbf{(p.~246)} 실증주의자가 그것에 기초한다고 상정되는 논쟁적 법의
근거들에 대한 논변을 ‘의미론적 독침(semantic sting)’\footnote{\emph{LE}
  45면.}이라고 부른다. 그것이 'law'라는 단어의 의미에 관한 이론에
기초하기 때문이다. 그래서 그는 \emph{Law's Empire}에서 이 ‘의미론적
독침’을 뽑아내고자 했다.

\emph{Law's Empire}의 첫 장에서 나는 오스틴과 함께 의미론적 이론가로
분류되고, 그래서 'law'라는 단어의 의미에서 명백한-사실 실증주의적
법이론을 도출하며 의미론적 독침을 겪는 것으로 분류된다. 그러나 실제로 내
책이나 내가 쓴 다른 어떤 것에도 내 이론에 관한 그러한 설명을 뒷받침하는
것은 없다. 따라서 발전된 국내법체계들이 법원들이 적용해야 하는 법들의
식별 기준들을 특정하는 승인 규칙을 포함한다는 나의 교설은 틀릴 수
있지만, 나는 어디에서도 이 교설을 모든 법체계들 안에는 그러한 승인
규칙이 있어야 한다는 것이 'law'라는 단어의 의미의 일부라는 잘못된 관념에
기초시키지 않았고, 법의 근거들의 식별 기준들이 논쟁의 여지 없이 고정되어
있지 않다면 'law'가 서로 다른 사람들에게 서로 다른 것을 \emphksans{뜻하게} 될
것이라는 훨씬 더 잘못된 관념에 기초시키지도 않았다.

기실 내게 귀속된 이 마지막 논변은 한 개념의 \emphksans{의미}와 그 개념의
\emphksans{적용} 기준들을 혼동한다. 나는 이것을 받아들이기는커녕
정의(justice)의 개념을 설명하면서 이 책 160쪽에서, 일정한 의미를 가진
개념의 적용 기준들이 달라질 수도 있고 논쟁적일 수도 있다는 사실에
명시적으로 주의를 환기했다. 이것을 명료하게 하기 위해 나는 사실상
드워킨의 후기 작업에서 그토록 두드러지게 나타나는 개념(concept)과 한
개념의 서로 다른 구상들(conceptions) 사이의 바로 그 구별을
그었다.\footnote{이 구분에 대해서는 John Rawls, \emph{A Theory of
  Justice} (1971), 5--6, 10면 참조. {[}Rawls는 정의 개념(concept of
  justice)과 정의의 구상들(conceptions of justice)을 구분하면서 “나는
  H. L. A. Hart의 \emph{The Concept of Law}\ldots{} 155--159면을
  따른다”고 명시함. \emph{A Theory of Justice} 5면 각주 1 참조.{]}}

마지막으로 드워킨은 또한, 실증주의자가 자신의 법이론이 의미론적 이론이
아니라 복합적 사회현상으로서 일반적 법의 독특한 특징들에 관한 기술적
설명이라고 주장하는 것은, 의미론적 이론과의 대비를 제시하지만 그 대비는
공허하고 오도적이라고 고집한다. 그의 논변\footnote{\emph{LE} 418--419면,
  주석 29.}은, 사회현상으로서 법의 독특한 특징들 가운데 하나가
법률가들이 법명제들의 참을 논쟁하고 \textbf{(p.~247)} 이것을 그러한
명제들의 의미를 참조하여 ‘설명한다’는 것이므로, 그러한 기술적 법이론은
결국 의미론적이어야 한다는 것이다.\footnote{\emph{LE} 31--33면 참조.} 이
논변은 내게는 'law'의 의미와 법명제들의 의미를 혼동하는 것으로 보인다.
법에 관한 의미론적 이론은, 'law'라는 단어의 바로 그 의미가 법을 일정한
특정 기준들에 의존하게 만든다는 이론이라고 드워킨에 의해 말해진다.
그러나 법명제들은 전형적으로 'law'가 무엇인지가 아니라 \emphksans{당해 법(the
law)}이 무엇인지, 곧 어떤 체계의 법이 사람들에게 무엇을 하도록
허용하거나, 요구하거나, 권한을 부여하는지에 관한 진술들이다. 그러므로
그러한 법명제들의 의미가 정의들이나 진리조건들에 의해 결정된다 하더라도,
이것은 'law'라는 단어의 바로 그 의미가 법을 일정한 특정 기준들에
의존하게 만든다는 결론으로 이어지지 않는다. 그러한 경우가 되려면 한
체계의 승인 규칙이 제공하는 기준들과 그러한 규칙의 필요성이 'law'라는
단어의 의미에서 도출되어야 할 것이다. 그러나 내 작업에는 그러한 교설의
흔적이 전혀 없다.\footnote{본서 209면 참조. 이곳에서 나는 그러한 교설을
  명시적으로 거부하고 있다.}

드워킨이 나의 법실증주의 형식을 잘못 제시하는 또 하나의 측면이 있다.
그는 나의 승인 규칙 교설을, 승인 규칙이 법의 식별을 위해 제공하는
기준들이 오직 역사적 사실들로만 이루어져야 한다고 요구하는 것으로,
그래서 ‘명백한-사실 실증주의(plain-fact positivism)’의 한 사례로
취급한다.\footnote{{[}이 표현은 하트 자신의 것이며, \emph{LE}에는
  등장하지 않음.{]}} 그러나 승인 규칙이 제공하는 기준들에 관해 내가 든
주된 예들이, 드워킨이 ‘계보(pedigree)’라고 부른 사항들, 곧 법들이
법제도들에 의해 채택되거나 창설되는 방식에만 관련되고 그 내용에는
관련되지 않는 사항들인 것은 사실이지만,\footnote{\emph{TRS} 17면.} 나는
이 책(72쪽)과 앞선 논문 「실증주의와 법과 도덕의 분리(Positivism and the
Separation of Law and Morals)」 양쪽에서, 미국에서처럼 어떤
법체계들에서는 법적 유효성의 궁극적 기준들이 계보뿐 아니라 정의
원칙들이나 실체적 도덕 가치들(substantive moral values)을 명시적으로
편입할 수도 있고, 이것들이 법적 헌법적 제약들의 내용을 형성할 수도
있다고 명시적으로 진술했다.\footnote{\emph{Harvard Law Review} 제71권
  (1958) 598면. \emph{Essays on Jurisprudence and Philosophy}에도
  재수록. 특히 54--55면 참조.} \emph{Law's Empire}에서 드워킨이 내게
‘명백한-사실’ 실증주의를 귀속시키는 것은 내 이론의 이 측면을 무시하는
것이다. 그러므로 \textbf{(p.~248)} 그가 내게 귀속시키는 명백한-사실
실증주의의 ‘의미론적’ 판본은 분명 내 것이 아니며, 내 이론은 어떤 형식의
명백한-사실 실증주의도 아니다.

\subsubsection{\texorpdfstring{(ii) \emphksans{해석적 이론으로서의 실증주의}
(Positivism as an Interpretive
Theory)}{(ii) 해석적 이론으로서의 실증주의 (Positivism as an Interpretive Theory)}}\label{ii-uxd574uxc11duxc801-uxc774uxb860uxc73cuxb85cuxc11cuxc758-uxc2e4uxc99duxc8fcuxc758-positivism-as-an-interpretive-theory}

드워킨의 명백한-사실 실증주의에 대한 둘째 설명은 그것을 의미론적
이론이나 언어적 고려들에 기초한 것으로 다루지 않고, 그가
‘관행주의(conventionalism)’라고 부르는 드워킨식 해석이론의 한 형식으로
재구성하려고 시도한다. 이 이론, 곧 드워킨이 결국 결함 있는 것으로
거부하는 이론에 따르면, 실증주의자는 법을 최선의 빛 아래 보이는 데
전념하는 해석 이론가라는 외양 아래에서, 법의 기준들을 명백한 사실들로
이루어진 것으로 제시한다. 다만 이 사실들은 의미론적 판본에서처럼 법의
어휘에 의해 고정되는 것이 아니라, 판사들과 법률가들이 공유하는 확신에
의해 논쟁의 여지 없이 고정된다. 이것은 법을 수범자들에게 큰 가치가 있는
어떤 것을 확보해 주는 것으로 보임으로써 법에 유리한 빛을 던진다. 곧 법적
강제의 계기들이 모두에게 이용 가능한 명백한 사실들에 의존하도록 만들어,
강제가 사용되기 전에 모두가 공정한 경고를 받을 수 있게 한다는 것이다.
드워킨은 이것을 ‘보호된 기대들의 이상(the ideal of protected
expectations)’이라고 부르지만,\footnote{\emph{LE} 117면.} 그에게 그
이점들은 결국 그 여러 결함들을 능가하지 않는다.

그러나 관행주의로서의 실증주의에 관한 이 해석주의적 설명은 내 법이론의
그럴듯한 판본이나 재구성으로 제시될 수 없다. 두 가지 이유에서 그렇다.
첫째, 내가 이미 말했듯이 내 이론은 명백한-사실 실증주의 이론이 아니다.
법의 기준들 가운데 ‘명백한’ 사실들뿐 아니라 가치들도 인정하기 때문이다.
둘째, 그리고 더 중요하게, 드워킨의 해석적 법이론은 그 모든 형식에서 법과
법적 실천의 요점 또는 목적이 강제를 정당화하는 것이라는 전제에 기초해
있지만,\footnote{{[}\emph{LE} 93면.{]}} 법이 이것을 자신의 요점 또는
목적으로 가진다는 것은 결코 내 견해가 아니며 그런 적도 없다. 다른
실증주의 형식들처럼 내 이론은 법과 법적 실천들 그 자체의 요점 또는
목적을 식별한다고 주장하지 않는다. 그러므로 법의 목적이 강제 사용을
정당화하는 것이라는 드워킨의 견해, 내가 결코 공유하지 않는 그 견해를
뒷받침하는 것은 내 이론 안에 없다. \textbf{(p.~249)} 사실 나는 법 그
자체가 봉사하는 더 특정한 목적을, 인간 행위에 대한 지침들과 그러한
행위에 대한 비판 기준들을 제공하는 것 너머에서 찾는 것은 매우 헛되다고
생각한다. 물론 이것은 같은 일반적 목적들을 가진 다른 규칙들이나
원칙들로부터 법들을 구별하는 데에는 충분하지 않을 것이다. 법의 독특한
특징들은 그 기준들의 식별, 변경, 집행을 위해 이차 규칙들을 통해 마련하는
장치와, 다른 기준들에 대한 우선성(priority)을 내세우는 일반적 주장에
있다. 그러나 내 이론이 법적 강제의 계기들에 대한 사전 고지가 일반적으로
이용 가능하리라는 점을 보증함으로써 기대들을 보호하는 관행주의 형식의
명백한-사실 실증주의에 전적으로 결속되어 있다 하더라도, 이것은 오직 내가
이것을 법이 가진 특정한 도덕적 이점(moral merit)으로 본다는 것만을 보여
줄 뿐, 법 그 자체의 전체 목적이 이것을 제공하는 것임을 보여 주지는 않을
것이다. 법적 강제의 계기들은 주로 그 수범자들의 행위를 인도하는 법의
일차적 기능이 무너진 경우들이므로, 법적 강제는 물론 중요한 사안이지만
부차적 기능이다. 그것의 정당화는 법 그 자체의 요점 또는 목적으로
온당하게 취급될 수 없다.

법적 강제는 그것이 “관행적 이해들(conventional understandings)에 부합할
때”\footnote{\emph{LE} 429면 주석 3.}에만 정당화된다는 주장을 하는
관행주의적 해석이론으로 내 법이론을 재구성하려는 드워킨의 이유들은, 이
책 제V장 제3절의 ‘법의 요소들(Elements of Law)’에 관한 내 설명에
기초한다. 거기서 나는 승인, 변경, 재판의 이차 규칙들을, 오직 의무의 일차
규칙들만으로 이루어진 상상된 단순한 체제의 결함들에 대한
치유책들(remedies)로 제시한다. 이 결함들은 규칙들의 동일성에 관한
\emphksans{불확실성}, 규칙들의 \emphksans{정태적} 성질, 그리고 규칙들이 오직
분산된 사회적 압력에 의해서만 집행될 때의 시간낭비적
\emphksans{비효율성}이다. 그러나 이러한 이차 규칙들을 그러한 결함들에 대한
치유책들로 제시하면서도, 나는 어디에서도 법적 강제가 이러한 규칙들에
부합할 때에만 \emphksans{정당화된다}고 주장하지 않았고, 더더욱 그러한
정당화의 제공이 일반적 법의 요점 또는 목적이라고 주장하지 않았다. 기실
이차 규칙들에 관한 논의에서 내가 강제를 언급하는 유일한 대목은 규칙들의
집행을 \textbf{(p.~250)} 법원들이 운용하는 조직화된 제재들 대신 분산된
사회적 압력에 맡기는 일의 시간낭비적 \emphksans{비효율성}에 관한 것이다.
그러나 분명 비효율성에 대한 치유책은 정당화가 아니다.

물론 의무의 일차 규칙들의 체제에 이차적 승인 규칙이 추가되면, 개인들이
강제의 계기들을 사전에 식별할 수 있게 하는 경우가 많으므로, 강제 사용에
대한 하나의 도덕적 반론을 배제한다는 의미에서 그 사용을 정당화하는 데
도움이 될 것이다. 그러나 승인 규칙이 가져올 법의 요구사항들에 관한
확실성과 사전 지식은 강제가 문제되는 경우에만 중요한 것이 아니다. 그것은
예컨대 유언장을 작성하거나 계약을 체결할 행위능력으로 나타나는 법적
권한들의 지성적 행사와, 일반적으로 사적·공적 삶의 지성적 계획에도 똑같이
핵심적이다. 그러므로 승인 규칙이 기여하는 강제의 정당화는 그 승인 규칙의
일반적 요점 또는 목적으로 제시될 수 없고, 더더구나 법 전체의 일반적 요점
또는 목적으로 제시될 수 없다. 내 이론에는 그것이 그럴 수 있다고 시사하는
것이 없다.

\subsubsection{\texorpdfstring{(iii) \emphksans{연성 실증주의} (Soft
Positivism)}{(iii) 연성 실증주의 (Soft Positivism)}}\label{iii-uxc5f0uxc131-uxc2e4uxc99duxc8fcuxc758-soft-positivism}

드워킨은 내게 ‘명백한-사실 실증주의’ 교설을 귀속시키면서, 내 이론을 승인
규칙의 존재와 권위가 법원들에 의한 그 수용이라는 사실에 의존해야 한다고
요구하는 것으로는, 실제로 그렇듯이, 취급했을 뿐 아니라, 그 규칙이
제공하는 법적 유효성 기준들이 그가 ‘계보(pedigree)’ 사항들이라고 부르는,
법의 창설 또는 채택의 방식과 형식에 관한 특정한 종류의 명백한 사실로만
이루어져야 한다고 요구하는 것으로도, 실제로는 그렇지 않은데, 취급하였다.
이것은 이중으로 잘못되었다. 첫째, 그것은 승인 규칙이 법적 유효성
기준들로서 도덕 원칙들 또는 실체적 가치들(substantive values)에 대한
부합(conformity)을 편입할 수 있다는 나의 명시적 인정을 무시한다. 따라서
내 교설은 이른바 ‘연성 실증주의(soft positivism)’이지, 드워킨의 그
판본에서처럼 ‘명백한-사실’ 실증주의가 아니다. 둘째, 승인 규칙이 제공하는
명백한-사실 기준들이 오직 계보 사항들일 수밖에 없다고 시사하는 것은 내
책에 전혀 없다. 그것들은 대신 종교의 설립이나 투표권 제한에 관한 미국
헌법 제16 수정조항 또는 제19 수정조항처럼 입법 내용에 대한 실체적
제약들(substantive constraints)일 수 있다.

그러나 이 답변은 드워킨의 가장 기본적인 비판들에 답하지 못한다.
\textbf{(p.~251)} 연성 실증주의의 어떤 형식을 채택한 다른 이론가들에게
답하면서,\footnote{E. P. Soper와 J. L. Coleman에 대한 드워킨의 응답을
  \emph{RDCJ} 247면 이하, 252면 이하에서 볼 수 있음.} 그는 그것이
타당하다면 내 이론에도 적용되어 여기서 답을 요구할 중요한 비판들을
제기했기 때문이다.

드워킨의 가장 근본적인 비판은, 법의 식별이 도덕 판단 또는 다른 가치
판단에 대한 부합이라는 논쟁적 사안들에 의존할 수 있도록 허용하는 연성
실증주의와, 논쟁적 도덕 논변들에 의존하지 않고 명백한 사실의 문제들로서
확실하게 식별될 수 있는 신뢰할 만한 공적 행위 기준들을 제공하는 데 법이
본질적으로 관련되어 있다는 일반 실증주의적 ‘그림(picture)’ 사이에 깊은
비일관성이 있다는 것이다.\footnote{\emph{RDCJ} 248면.} 연성 실증주의와
내 이론의 나머지 부분 사이의 그러한 비일관성을 입증하기 위해 드워킨은,
다른 결함들 가운데서도 상상된 전법적(pre-legal) 관습-유형 의무 일차 규칙
체제의 불확실성을 치유하는 것으로 승인 규칙을 설명한 내 설명을 인용할
것이다. 연성 실증주의에 관한 이 비판은 내게는 일관된 실증주의자가 법적
기준들의 한 집합에 귀속해야 한다고 보는 확실성의 정도와, 법적 유효성
기준들이 특정 도덕 원칙들이나 가치들에 대한 부합을 포함할 경우 생길
불확실성 양쪽 모두를 과장하는 것으로 보인다. 승인 규칙의 중요한 기능
하나가 법이 확인될 수 있는 확실성을 증진하는 것임은 물론 참이다. 승인
규칙이 법에 대해 도입한 검사 기준들이 어떤 사례들에서 논쟁적 쟁점들을
제기할 뿐 아니라 모든 사례 또는 대부분의 사례에서 그렇게 한다면, 그것은
이 기능을 수행하지 못할 것이다. 그러나 다른 가치들의 측면에서 어떤
대가를 치르더라도 모든 불확실성을 배제하는 것은 내가 승인 규칙을 위해
생각해 본 목표가 아니다. 이는, 승인 규칙 자체뿐 아니라 그것을 참조하여
식별되는 특정 법규칙들도 다툴 수 있는 불확실성의 ‘반음영(penumbra)’을
가질 수 있다는 내 명시적 진술에 의해서도, 내가 바랐듯이,
분명해진다.\footnote{본서 123면, 147--154면 참조.} 또한 법들이 그 의미에
관해 생길 수 있는 모든 가능한 질문들을 미리 확정할 수 있도록 틀 지어질
수 있다 하더라도, 그러한 법들을 채택하는 것은 법이 소중히 여겨야 할 다른
목적들과 자주 충돌할 것이라는 나의 일반 논변도 있다.\footnote{{[}본서
  128면 참조.{]}} 예견되지 않은 사례의 구성이 알려지고 그 결정에서 걸린
쟁점들이 식별되어 합리적으로 확정될 수 있도록, 많은 법규칙들의 경우
일정한 불확실성의 여지는 \textbf{(p.~252)} 용인되어야 하며 기실
환영되어야 한다. 승인 규칙의 확실성-제공 기능을 최우선적이고 압도적인
것으로 취급할 때에만, 논쟁적일 수 있는 도덕 원칙들이나 가치들에 대한
부합을 법의 기준들 가운데 포함하는 연성 실증주의의 형식이 비일관적인
것으로 간주될 수 있을 것이다. 여기서 바탕에 있는 질문은, 법체계가
일반적으로 신뢰할 수 있고 사전에 식별 가능한 확정적 행위 지침들을
제공함으로써 탈중심화된 관습-유형 규칙 체제로부터 어떤 유의미한 진전을
이루려면, 어느 정도 또는 범위의 불확실성을 용인할 수 있는가에 관한
것이다.

나의 연성 실증주의 판본의 일관성에 대한 드워킨의 둘째 비판은 법의
확정성(determinacy)과 완전성(completeness)에 관한 서로 다르고 더 복잡한
쟁점들을 제기한다. 이 책에서 전개된 내 견해는, 승인 규칙이 제공하는
기준들에 의해 일반적 항들로 식별되는 법규칙들과 법원칙들이 내가 자주
‘개방적 직조(open texture)’라고 부르는 것을 종종 가진다는 것이다. 그래서
주어진 규칙이 특정 사례에 적용되는지가 문제될 때, 법은 어느 쪽의 답도
확정하지 못하고 따라서 부분적으로 불확정적인 것으로 드러난다. 그러한
사례들은 단지 합리적이고 식견 있는 법률가들이 어느 답이 법적으로
올바른지에 관하여 의견불일치할 수 있다는 의미에서 논쟁적인 ‘어려운
사례들’일 뿐 아니라, 그러한 경우 법은 근본적으로 \emphksans{불완전하다}.
그것은 그러한 사례들에서 쟁점이 된 질문들에 \emphksans{아무} 답도 제공하지
않는다. 그것들은 법적으로 규율되지 않으며, 그러한 사례들에서 결정에
이르기 위해 법원들은 내가 ‘재량(discretion)’이라고 부르는 제한된 법창출
기능을 행사해야 한다. 드워킨은 법이 이런 방식으로 불완전할 수 있고
그러한 법창출 재량의 행사로 채워질 틈들을 남길 수 있다는 관념을
거부한다. 그는 이 견해가, 법적 권리 또는 법적 책무의 존재를 단언하는
법명제가 논쟁적일 수 있고 따라서 합리적이고 식견 있는 사람들이 의견을
달리할 수 있으며, 그들이 의견을 달리할 때 그것이 참인지 거짓인지를
확정적으로 논증할 방법이 자주 없다는 사실에서 나온 잘못된 추론이라고
생각한다. 그러한 추론이 잘못인 것은, 법명제가 이렇게 논쟁적일 때에도
그럼에도 그것을 참 또는 거짓으로 만드는 ‘사안의 사실들(facts of the
matter)’이 있을 수 있고, \textbf{(p.~253)} 그 참이나 거짓을 논증할 수
없더라도 그것이 참이라는 논변들이 그것이 거짓이라는 논변들보다 더 나은
것으로 평가될 수 있으며, 그 반대도 가능하기 때문이다. 논쟁적인 법과
불완전하거나 불확정적인 법 사이의 이 구별은 드워킨의 해석이론에 상당히
중요하다. 그 이론에 따르면 법명제는 다른 전제명제들과 결합하여, 법체계의
제도적 역사에 가장 잘 부합하고 그것에 대한 최선의 도덕적 정당화도
제공하는 원칙들에서 따라 나올 때에만 참이기 때문이다. 그러므로
드워킨에게 어떤 법명제의 참은 궁극적으로 무엇이 가장 잘 정당화하는지에
관한 도덕 판단의 참에 의존하며, 그에게 도덕 판단들은 본질적으로
논쟁적이므로 모든 법명제들도 그러하다.

드워킨에게, 적용이 논쟁적인 도덕 판단을 포함하는 법적 유효성 기준이라는
관념은 이론적 어려움을 전혀 제시하지 않는다. 그의 견해에서는 그것이
여전히 이미 존재하는 법(pre-existing law)에 대한 진정한 검사
기준(test)일 수 있다. 그 논쟁적 성격은 그것을 참으로 만드는 사실들, 많은
경우 도덕적 사실들이 존재한다는 점과 완전히 양립 가능하기 때문이다.

그러나 법적 유효성 기준이 부분적으로 도덕적 검사 기준일 수 있음을
허용하는 연성 실증주의는, 드워킨의 주장에 따르면, 위 pp.~251-252에서
이미 논의한 것에 더해 둘째 비일관성에 연루된다. 그것은 법이 확실하게
식별 가능하다는 실증주의적 ‘그림’과 비일관적일 뿐 아니라, 법명제들의
‘객관적 지위(objective standing)’\footnote{\emph{RDCJ} 250면.}를 도덕
판단들의 지위에 관한 어떤 논쟁적 철학 이론에 대한 결속으로부터도
독립시키려는, 그가 실증주의자들에게 귀속시키는 바람과도 비일관적이다.
도덕적 검사 기준이 이미 존재하는 법에 대한 검사 기준일 수 있으려면, 도덕
판단들이 참이 되는 객관적 도덕 사실들이 있어야 한다. 그러나 그러한
객관적 도덕 사실들이 있다는 것은 논쟁적 철학 이론이다. 그러한 사실들이
없다면, 도덕적 검사 기준을 적용하라는 말을 들은 판사는 이것을 오직,
법체계가 이에 부과하는 어떤 제약들에 종속하여, 도덕과 그 요구사항들에
관한 자신의 최선의 이해에 따라 법창출 재량을 행사하라는 요청으로만
취급할 수 있다.

나는 여전히 법이론이 도덕 판단들의 일반적 지위에 관한 \textbf{(p.~254)}
논쟁적 철학 이론들에 결속되는 것을 피해야 하며, 이 책에서 내가
그러하듯(p.~168), 그것들이 드워킨이 ‘객관적 지위’라고 부르는 것을
가지는지에 관한 일반 질문을 열어 두어야 한다고 생각한다. 이 철학적
질문에 대한 답이 무엇이든, 판사의 책무는 동일할 것이기 때문이다. 곧 그가
결정해야 할 수 있는 어떤 도덕적 쟁점들에 대해서도 자신이 할 수 있는
최선의 도덕 판단을 내리는 것이다. 사건들을 이렇게 결정하면서 판사가
도덕에 따라, 법이 부과하는 어떤 제약들에 종속하여, 법을 \emphksans{만들고}
있는 것인지, 아니면 법에 대한 도덕적 검사 기준에 의해 드러나는 이미
\emphksans{존재하는} 법이 무엇인지에 관한 자신의 도덕 판단에 의해 인도되는
것인지는 어떤 실천적 목적에서도 중요하지 않을 것이다. 물론, 내가 그래야
한다고 주장하듯이 법이론이 도덕 판단들의 객관적 지위 문제를 열어 둔다면,
연성 실증주의는 단순히 도덕 원칙들이나 가치들이 법적 유효성의 기준들
가운데 있을 수 있다는 이론으로 특징지어질 수 없다. 도덕 원칙들과
가치들이 객관적 지위를 가지는지가 열린 질문이라면, 그것들에 대한 부합을
현존 법(existing law)에 대한 검사 기준들 가운데 포함하려는 ‘연성
실증주의적’ 규정들이 그러한 효과를 가질 수 있는지, 아니면 오직
법원들에게 도덕에 따라 법을 \emphksans{만들라}는 지시들을 구성할 수 있을
뿐인지도 열린 질문이어야 하기 때문이다.

몇몇 이론가들, 특히 라즈가 다음과 같이 본다는 점은 관찰되어야 한다. 도덕
판단들의 지위가 무엇이든, 법이 법원들에게 법이 무엇인지를 확정하는 데
도덕 기준들을 적용하라고 요구할 때마다, 법은 그로써 법원들에게 재량을
부여하고, 새 법을 만들면서 자기들의 최선의 도덕 판단에 따라 그 재량을
사용하라고 지시한다. 법은 그로써 도덕을 이미 존재하는 법으로 전환하지
않는다.\footnote{J. Raz, “Dworkin: A New Link in the Chain”,
  \emph{California Law Review} 제74권 (1986) 1103면, 1110면 및
  1115--1116면.}

\subsection{3. 규칙의 본성 (THE NATURE OF
RULES)}\label{uxaddcuxce59uxc758-uxbcf8uxc131-the-nature-of-rules}

\subsubsection{\texorpdfstring{(i) \emphksans{규칙의 실천 이론} (The Practice
Theory of
Rules)}{(i) 규칙의 실천 이론 (The Practice Theory of Rules)}}\label{i-uxaddcuxce59uxc758-uxc2e4uxcc9c-uxc774uxb860-the-practice-theory-of-rules}

이 책의 여러 지점에서 나는 법의 내부 진술(internal statements)과 외부
진술(external statements) 사이의 구분, 그리고 법의 내부 측면(internal
aspects)과 외부 측면(external aspects) 사이의 구분에 주의를 기울였다.

이러한 구분들과 그 중요성을 설명하기 위해 나는 (56--7쪽)
\textbf{(p.~255)} 제정된 규칙들과 관습-유형 규칙들 양쪽을 포함하는
고도로 복잡한 법체계의 사례에서 출발하지 않고, 그보다 단순한 사례, 곧
크든 작든 어떤 사회 집단의 관습-유형 규칙들의 사례에서 출발했다.
내부/외부 구분은 이 사례에도 똑같이 적용되며, 나는 이러한 규칙들을
‘사회적 규칙들(social rules)’이라고 부른다. 내가 그것들에 관해 제시한
설명은 ‘규칙의 실천 이론(the practice theory)’으로 알려지게 되었는데,
그것은 한 집단의 사회적 규칙들을 사회적 실천의 한 형식에 의해 구성되는
것으로 취급하기 때문이다. 이 사회적 실천은 집단 구성원 대부분이
반복적으로 일양적인 따름을 보이는 행위 양식들과, 그러한 행위 양식들에
대한 독특한 규범적 태도, 곧 내가 ‘수용(acceptance)’이라고 부른 것을 함께
포함한다. 이 수용은 그러한 행위 양식들을 자신의 미래 행위의 지침들로,
그리고 부합(conformity)에 대한 요구들과 여러 형태의 압력을 정당화할 수
있는 비판 기준들로 삼으려는 개인들의 지속적 성향으로 이루어진다. 사회적
규칙들에 대한 외부 관점은 그 실천의 관찰자의 관점이고, 내부 관점은 그
실천의 참여자, 곧 규칙들을 행위의 지침들과 비판 기준들로 수용하는
참여자의 관점이다.

나의 사회적 규칙들의 실천 이론은 드워킨에 의해 광범위하게 비판되어 왔다.
앞서 말했듯이 그는 이와 유사하지만 실제로는 여러 점에서 매우 다른 구별,
곧 한 공동체의 사회적 규칙들에 대한 사회학자의 외부 기술과, 자신의
행위와 타인의 행위에 대한 평가와 비판을 위해 규칙들에 호소하는 참여자의
내부 관점 사이의 구별을 제시한다.\footnote{{[}\emph{LE} 13--14면
  참조.{]}} 사회적 규칙들에 대한 나의 원래 설명에 관해 드워킨이 제기한
비판 가운데 일부는 분명 타당하며 법의 이해에도 중요하다. 이어서 나는
이제 필요하다고 보는 나의 원래 설명의 상당한 수정들을 지목한다.

\begin{enumerate}
\def\labelenumi{(\roman{enumi})}
\item
  드워킨이 주장했듯이, 나의 설명은 한 집단의 관행적 규칙들에서 드러나는
  \emphksans{관행(convention)}의 합의와, 한 집단의 동시적 실천들에서 드러나는
  독립적 \emphksans{확신(conviction)}들의 합의 사이의 중요한 차이를
  무시한다는 점에서 결함이 있다. 어떤 규칙들에 대한 집단의 일반적
  부합(conformity)이 그 개별 구성원들이 그 규칙들을 수용하는 이유들의
  일부라면, 그 규칙들은 관행적 사회 실천들이다. 이와 대조적으로
  \textbf{(p.~256)} 한 집단의 공유 도덕과 같은 단지 동시적인 실천들은
  관행이 아니라, 집단 구성원들이 특정 방식으로 행위할 동일하지만
  독립적인 이유들을 가지고 대체로 그 이유들에 따라 행위한다는 사실에
  의해 구성된다.
\item
  드워킨이 또한 옳게 주장했듯이, 나의 사회적 규칙 설명은 내가 이제
  설명한 의미에서 관행적인 규칙들에만 적용 가능하다. 이것은 나의 실천
  이론의 범위를 상당히 좁히며, 나는 이제 그것을 개인적 도덕이든 사회적
  도덕이든 도덕에 대한 건전한 설명으로 보지 않는다. 그러나 그 이론은
  일반적 사회 관습들, 곧 법적 효력을 가진 것으로 승인될 수도 있고 그렇지
  않을 수도 있는 관습들뿐 아니라, 승인 규칙을 포함한 일정한 중요한
  법규칙들도 포함하는 관행적 사회적 규칙들에 대한 충실한 설명으로는 남아
  있다. 승인 규칙은 사실상 법원들의 법-식별 및 법-적용 작용에서 수용되고
  실천될 때에만 존재하는 사법적 관습 규칙의 한 형식이다. 이와 대조적으로
  제정된 법규칙들은, 승인 규칙이 제공하는 기준들에 의해 유효한
  법규칙들로 식별될 수 있지만, 그것들을 실천할 계기가 생기기 전에도
  제정의 순간부터 법규칙들로 존재할 수 있으므로, 실천 이론은 그것들에는
  적용되지 않는다.
\end{enumerate}

드워킨의 규칙의 실천 이론에 대한 핵심 비판은, 그것이 사회적 규칙을 그
사회적 실천에 의해 구성되는 것으로 잘못 받아들이며, 그래서 그러한 규칙이
존재한다는 진술을 단지 그 규칙의 존재를 위한 실천-조건들이 충족되었다는
외부 사회학적 사실의 진술로 취급한다는 것이다.\footnote{{[}\emph{TRS}
  48--58면 참조.{]}} 드워킨의 논변에 따르면 그러한 설명은 가장 단순한
관행적 규칙조차 가지는 \emphksans{규범적} 성격을 설명할 수 없다. 그러한
규칙들은 \emphksans{책무들}과 \emphksans{행위의 이유들}을 확립하며, 그러한
규칙들이 행위 비판과 행위 요구의 뒷받침에서 통상 인용될 때에는 바로
이것들에 호소하기 때문이다. 규칙들의 이러한 이유-제공 및 책무-확립
특징이 그 독특한 규범적 성격을 구성하며, 그것은 규칙들의 존재가 실천
이론에 따르면 사회적 규칙의 존재를 구성하는 실천들과 태도들이 그러하듯이
단지 사실적 사태로 이루어질 수 없음을 보여 준다. 드워킨에 따르면, 이러한
독특한 특징들을 가진 규범적 규칙은 ‘일정한 규범적 사태’가 있을 때에만
존재할 수 있다.\footnote{\emph{TRS} 51면.} \textbf{(p.~257)} 나는 인용된
이 말들이 답답할 만큼 모호하다고 본다. ‘교회 출석자들의 규칙’(남성은
교회에서 머리를 드러내야 한다)의 예에 관한 논의에서,\footnote{{[}\emph{TRS}
  50--58면; 본서 124--125면 참조.{]}} 드워킨은 규범적 사태로 규칙이
요구하는 것을 할 좋은 도덕적 근거들 또는 정당화의 존재를 뜻하는 듯하다.
그래서 그는 교회 출석자들이 교회에서 모자를 벗는 단순한 반복적으로
일양적인 실천은 그 규칙을 구성할 수 없지만, 불쾌감을 주는 방식들을
창출하고 교회에서 모자를 벗도록 요구하는 규칙을 위한 좋은 근거들인
기대들을 일으킴으로써 그 규칙을 정당화하는 데 도움을 줄 수 있다고
논변한다. 규범적 규칙의 단언을 보증하기 위해 요구되는 규범적 사태로
드워킨이 뜻하는 것이 이것이라면, 사회적 규칙의 존재 조건에 관한 그의
설명은 내게 지나치게 강해 보인다. 그것은 규칙들에 책무를 확립하거나
행위의 이유들을 제공하는 것으로 호소하는 참여자들이 규칙들에 부합할 좋은
도덕적 근거들 또는 정당화가 있다고 믿어야 할 뿐 아니라, 실제로 그러한
좋은 근거들이 있어야 한다고 요구하는 것처럼 보이기 때문이다. 분명 어떤
사회는 그 구성원들이 수용하지만 도덕적으로 사악한 규칙들, 예컨대 특정
피부색을 가진 사람들이 공원이나 해수욕장 같은 공공시설을 이용하는 것을
금지하는 규칙들을 가질 수 있다. 기실 사회적 규칙의 존재를 위한 일반
조건으로, 참여자들이 그 규칙에 부합할 좋은 도덕적 근거들이 있다고
\emphksans{믿어야} 한다는 더 약한 조건조차도 지나치게 강하다. 어떤 규칙들은
단순히 전통에 대한 존중, 타인들과 동일시하려는 바람, 또는 사회가
개인들에게 무엇이 유리한지를 가장 잘 안다는 믿음에서 수용될 수 있기
때문이다. 이러한 태도들은 그 규칙들이 도덕적으로 문제 있다는 어느 정도
생생한 자각과 공존할 수 있다. 물론 관행적 규칙은 도덕적으로 건전하고
정당화될 수도 있으며 그렇게 믿어질 수도 있다. 그러나 관행적 규칙들을
자신의 행위 지침이나 비판 기준들로 수용한 사람들이 왜 그렇게 했는지가
문제될 때, 나는 제시될 수많은 답변들(이 책 203쪽, 232쪽 참조) 가운데
규칙들의 도덕적 정당화에 대한 믿음을 유일하게 가능하거나 충분한 답변으로
골라낼 이유가 없다고 본다.

마지막으로, 드워킨은 규칙의 실천 이론이 \textbf{(p.~258)} 관행적
규칙들로 제한되더라도 폐기되어야 한다고 논변한다. 그것은 관행적 규칙의
범위가 논쟁적일 수 있고 따라서 의견불일치의 대상이 될 수 있다는 관념을
수용할 수 없기 때문이다.\footnote{{[}\emph{TRS} 58면.{]}} 그는
반복적으로 일양적인 실천과 수용에 의해 구성되는 논쟁의 여지 없는
규칙들이 일부 있다는 점은 부정하지 않지만, 그렇게 구성되는 규칙들은 어떤
게임들의 규칙들처럼 상대적으로 중요하지 않은 사례들만 포함한다고
주장한다. 그러나 이 책에서는 법체계의 기본 승인 규칙처럼 중요하고도 거의
논쟁의 여지가 없는 규칙이 법원들이 그것을 자신들의 법-적용 및 법-집행
작용의 지침으로 수용하는 일양적 실천에 의해 구성되는 규칙으로 취급된다.
이에 맞서 드워킨은 어려운 사례들에서 어떤 주제에 관한 당해 법이
무엇인지에 관해 판사들 사이에 이론적 의견불일치가 자주 있으며, 이것들이
논쟁 없음과 일반적 수용의 외양이 환상임을 보여 준다고 주장한다. 물론
그러한 의견불일치들의 빈도와 중요성은 부정될 수 없다. 그러나 그것들의
존재에 호소하여 승인 규칙에 실천 이론이 적용될 수 없다고 논변하는 것은
그 규칙의 기능에 대한 오해에 기초한다. 그것은 승인 규칙이 특정 사례들의
법적 결과를 완전히 확정하도록 의도되어, 어떤 사례에서 생기는 어떤 법적
쟁점도 그 규칙이 제공하는 기준들이나 검사 기준들에 대한 단순한
호소만으로 해결될 수 있다고 가정한다. 그러나 이것은 오해이다. 그 규칙의
기능은 현대 법체계들에서 올바른 법적 결정들이 충족해야 하는 일반
조건들만을 확정하는 것이다. 그 규칙은 대개 드워킨이 계보 사항들이라고
부르는 유효성 기준들, 곧 법의 내용이 아니라 법들이 창설되거나 채택되는
방식과 형식에 관련되는 기준들을 제공함으로써 이 기능을 수행한다. 그러나
내가 말했듯이(250쪽), 승인 규칙은 그러한 계보 사항들에 더해 법들의
사실적 내용이 아니라 실체적 도덕 가치들(substantive moral values)이나
원칙들에 대한 그 부합과 관련되는 검사 기준들을 제공할 수도 있다. 물론
특정 사례들에서 판사들은 그러한 검사 기준들이 충족되는지 여부에 관해
의견불일치할 수 있고, 승인 규칙 안의 도덕적 검사 기준은 그러한
의견불일치를 해소하지 않을 것이다. 판사들은 특정 사례들에서
\textbf{(p.~259)} 그 검사 기준들이 무엇을 요구하는지에 관해서는
의견불일치하면서도, 그러한 검사 기준들의 관련성이 확립된 사법 실천에
의해 확정된 것이라는 데에는 동의할 수 있다. 이렇게 이해된 승인 규칙에는
규칙의 실천 이론이 온전히 적용 가능하다.

\subsubsection{\texorpdfstring{(ii) \emphksans{규칙들과 원칙들} (Rules and
Principles)}{(ii) 규칙들과 원칙들 (Rules and Principles)}}\label{ii-uxaddcuxce59uxb4e4uxacfc-uxc6d0uxce59uxb4e4-rules-and-principles}

오랫동안 이 책에 대한 드워킨의 비판 가운데 가장 잘 알려진 것은, 이 책이
법을 오직 ‘전부-아니면-무(all-or-nothing)’ 규칙들로만 이루어진 것으로
잘못 제시하고, 법적 추론과 재판에서 중요하고 독특한 역할을 하는 다른
종류의 법적 기준, 곧 법적 원칙들을 무시한다는 것이었다. 내 작업에서 이
결함을 발견한 일부 비판자들은 그것을 다소 고립된 흠으로 보았고, 내가
법체계의 구성 요소들로 법규칙들과 함께 법원칙들을 포함하기만 하면 그것을
고칠 수 있으며, 책의 주요 논제들 가운데 어느 것도 포기하거나 중대하게
수정하지 않고도 그렇게 할 수 있다고 생각했다. 그러나 이 비판 노선을 처음
강하게 제기한 드워킨은, 법원칙들은 나의 법이론 안에 포함될 수 있지만 그
대가는 그 중심 교설들의 포기라고 주장해 왔다. 그에 따르면 내가 법이
부분적으로 원칙들로 이루어져 있음을 인정한다면, 나는 한 체계의 당해 법이
법원들의 실천에서 수용된 승인 규칙이 제공하는 기준들에 의해 식별된다는
주장도, 현존하는 명시적 법이 결정을 지시하지 못하는 경우 법원들이
진정한, 비록 간극적인, 법창출 권한 또는 재량을 행사한다는 주장도, 법과
도덕 사이에 중요한 필연적 또는 개념적 연결이 없다는 주장도 일관되게
유지할 수 없다. 이 교설들은 내 법이론의 중심일 뿐 아니라 현대
법실증주의의 핵심을 구성하는 것으로 자주 받아들여지므로, 그것들의 포기는
상당한 중대성을 가진 문제일 것이다.

이 응답의 이 절에서 나는 내가 법원칙들을 무시했다는 비판의 여러 측면들을
검토하고, 그 비판에서 타당한 것이 무엇이든 내 이론 전체에 중대한 귀결
없이 수용될 수 있음을 보이고자 한다. 그러나 나는 이제, 재판과 법적
추론이라는 주제에 관해, 특히 내 비판자들이 법원칙들이라고 부르는
것들로부터의 논변들에 관해 이 책에서 내가 너무 적게 말했다는 점을 기꺼이
인정한다. 원칙들이 이 책에서 지나가듯 다루어졌을 뿐이라는 점은
결함이라는 데 나는 이제 동의한다.

그러나 정확히 내가 무시했다고 고발되는 것은 무엇인가? \textbf{(p.~260)}
법원칙들은 무엇이며, 그것들은 법규칙들과 어떻게 다른가? 법률가들이
사용하는 ‘원칙들’은 흔히 방대한 이론적·실천적 고려사항들을 포함하며, 그
가운데 일부만이 드워킨이 제기하려 한 쟁점들과 관련된다. ‘원칙’이라는
표현을, 사건들을 결정하는 법원들의 행위를 포함한 행위 기준들로 제한해
받아들인다 하더라도, 규칙들과 그러한 원칙들 사이의 대비를 긋는 방식들은
서로 다를 수 있다. 그러나 내가 원칙들을 무시했다고 비난한 모든
비판자들은 적어도 두 특징이 그것들을 규칙들과 구별한다는 데 동의할
것이라고 생각한다. 첫째는 정도의 문제이다. 원칙들은 규칙들에 비해 넓고
일반적이며 비특정적이다. 흔히 다수의 구별되는 규칙들로 여겨질 것들이
하나의 원칙의 예증들 또는 실례들로 제시될 수 있다는 의미에서 그렇다.
둘째 특징은, 원칙들이 어떤 목적, 목표, 권리자격(entitlement) 또는 가치를
어느 정도 명시적으로 참조하기 때문에, 어떤 관점에서는 유지하거나
고수하는 것이 바람직한 것으로 간주되며, 따라서 그것들을 예증하는
규칙들의 설명 또는 이론적 근거(rationale)를 제공할 뿐 아니라 적어도 그
정당화에도 기여하는 것으로 간주된다는 점이다.

규칙들과 관련하여 원칙들이 수행하는 설명적·정당화적 역할을 설명하는,
넓이와 어떤 관점에서의 바람직성이라는 이 두 비교적 논쟁의 여지가 없는
특징들 외에도, 세 번째 구별 특징이 있다. 나 자신은 이것이 정도의
문제라고 생각하지만, 드워킨은 그것을 핵심적이라고 본다. 그에 따르면
규칙들은 그것들을 적용하는 사람들의 추론에서 ‘전부-아니면-무 방식’으로
기능한다. 즉 어떤 규칙이 유효하고 주어진 사례에 적용 가능하다면, 그것은
법적 결과 또는 귀결을 ‘필연화한다’, 곧 확정적으로 결정한다.\footnote{{[}\emph{TRS}
  24면.{]}} 그가 법규칙들의 예로 든 것들 가운데에는 유료도로에서 시속
60마일의 최고속도를 정하는 규칙이나, 유언장의 작성·증명·효력을 규율하는
제정법들, 예컨대 두 증인이 서명하지 않으면 유언장은 유효하지 않다는
제정법 규칙이 있다. 드워킨에 따르면 법원칙들은 그러한 전부-아니면-무
규칙들과 다르다. 그것들은 적용될 때 결정을 ‘필연화하지’ 않고, 결정 쪽을
가리키거나 \textbf{(p.~261)} 결정에 유리하게 작용하거나, 압도될 수
있지만 법원들이 어느 한 방향으로 기울게 하는 것으로 고려하는 이유를
진술하기 때문이다. 짧게 말해, 나는 원칙들의 이러한 특징을 그
‘비확정적(non-conclusive)’ 성격이라고 부르겠다. 드워킨이 든 그러한
비확정적 원칙들의 예들 가운데에는 ‘법원들은 소비자와 공적 이익들이
공정하게 다루어지는지 보기 위해 {[}자동차{]} 구매 계약들을 면밀히
검토해야 한다’는 것처럼 상대적으로 특정적인 것들도 있고,\footnote{\emph{TRS}
  24면. 다음 판례에서 인용함: \emph{Henningsen v. Bloomfield Motors,
  Inc.}, 32 NJ 358, 161 A.2d 69 (1960), 387면, 85면.} ‘누구도 자신의
그른 행위로부터 이익을 얻어서는 안 된다’는 것처럼 훨씬 넓은 범위를
가지는 것들도 있다.\footnote{\emph{TRS} 25--26면.} 실제로 미국 헌법
제1·5·14 수정조항의 규정들처럼 미국 연방의회의 권한들과 주 입법에 대한
가장 중요한 헌법적 제한들의 상당수는 비확정적 원칙들로
기능한다.\footnote{{[}드워킨은 \emph{TRS} 27면에서 수정헌법 제1조가
  규칙(rule)인지 원칙(principle)인지 논의함.{]}} 드워킨에 따르면
법원칙들은 유효성의 차원이 아니라 \emphksans{무게}의 차원을 가진다는 점에서
규칙들과 다르다.\footnote{{[}\emph{TRS} 26면.{]}} 그래서 더 큰 무게의
다른 원칙과 충돌하면 한 원칙은 압도되어 결정을 확정하지 못할 수 있지만,
그럼에도 온전하게 존속하여 다른 사례들에서 더 작은 무게의 다른 원칙과
경쟁할 때 이길 수도 있다. 반면 규칙들은 유효하거나 유효하지 않을 뿐 이
무게의 차원을 가지지 않는다. 그래서 처음 정식화된 규칙들이 충돌하면,
드워킨에 따르면 그 가운데 오직 하나만 유효할 수 있고, 다른 규칙과의
경쟁에서 진 규칙은 그 경쟁 규칙과 일관되도록 다시 정식화되어야 하며
따라서 주어진 사례에는 적용될 수 없게 된다.\footnote{\emph{TRS}
  24--27면.}

나는 법원칙들과 법규칙들 사이의 이 날카로운 대비도, 유효한 규칙이 주어진
사례에 적용 가능하다면 원칙과 달리 항상 그 사례의 결과를 결정해야 한다는
견해도 받아들일 이유가 없다고 본다. 법체계가, 어떤 유효한 규칙은 그것이
적용 가능한 사례들에서 결과를 결정하지만, 더 중요하다고 판단되는 다른
규칙이 같은 사례에도 적용되는 경우에는 예외라고 승인하지 못할 이유는
없다. 따라서 주어진 사례에서 \textbf{(p.~262)} 더 중요한 규칙과의
경쟁에서 패배한 규칙은 원칙처럼, 다른 경쟁 규칙보다 더 중요하다고
판단되는 다른 사례들에서는 결과를 결정하기 위해 존속할 수
있다.\footnote{이 중요한 지점은 Raz와 Waluchow가 강조했지만, 나는 이에
  충분한 주의를 환기하지 못했다. J. Raz, “Legal Principles and the
  Limits of the Law”, \emph{Yale Law Journal} 제81권 (1972) 823면,
  832--834면; W. J. Waluchow, “Herculean Positivism”, \emph{Oxford
  Journal of Legal Studies} 제5권 (1985) 187면, 189--192면 참조.}

그러므로 드워킨에게 법은 전부-아니면-무 규칙들과 비확정적 원칙들 양쪽을
포함하며, 그는 이 차이를 정도의 문제로 보지 않는다. 그러나 나는 드워킨의
입장이 정합적일 수 있다고 생각하지 않는다. 그의 가장 이른 예들은
규칙들이 원칙들과 충돌할 수 있고, 어떤 원칙은 때로 규칙과의 경쟁에서
이기고 때로 질 것임을 함축한다. 그가 인용한 사례들에는 \emph{Riggs} v.
\emph{Palmer}가 포함된다.\footnote{115 N.Y. 506, 22 N.E. 188 (1889);
  \emph{TRS} 23면; \emph{LE} 15면 이하 참조.} 그 사건에서 사람은 자신의
그른 행위로부터 이익을 얻도록 허용되어서는 안 된다는 원칙은, 유언의
효력을 지배하는 제정법 규칙들의 명확한 문언에도 불구하고, 살인자가
희생자의 유언에 따라 상속하는 것을 배제하는 것으로 판시되었다. 이것은
원칙이 규칙과의 경쟁에서 이기는 예이지만, 그러한 경쟁의 존재 자체가
규칙들이 전부-아니면-무 성격을 가지지 않음을 분명히 보여 준다. 규칙들은
자신보다 더 무거울 수 있는 원칙들과 그러한 충돌 속으로 들어갈 수 있기
때문이다. 우리가 그러한 사례들을 (드워킨이 때때로 시사하듯이) 규칙들과
원칙들 사이의 충돌이 아니라, 문제된 규칙을 설명하고 정당화하는 원칙과
다른 어떤 원칙 사이의 충돌로 서술하더라도, 전부-아니면-무 규칙들과
비확정적 원칙들 사이의 날카로운 대비는 사라진다. 이 견해에서는 어떤
규칙이 그 문언에 따르면 적용 가능한 사례에서 결과를 확정하지 못하는 것은
그 정당화 원칙이 다른 원칙에 의해 무게에서 압도되기 때문이다. (드워킨이
또한 시사하듯이) 우리가 어떤 원칙을 명확히 정식화된 법규칙에 대한 새로운
해석을 위한 이유를 제공하는 것으로 생각하는 경우에도
마찬가지이다.\footnote{{[}드워킨의 논의는 \emph{TRS} 22--28면, \emph{LE}
  15--20면 참조.{]}}

법체계가 전부-아니면-무 규칙들과 비확정적 원칙들 양쪽으로 이루어진다는
주장 안의 이 비정합성은, 그 구별이 정도의 문제임을 인정하면 치유될 수
있다. \textbf{(p.~263)} 확실히, 적용 조건들의 충족이 몇몇 사례들, 곧 그
규정들이 더 중요하다고 판단되는 다른 규칙과 충돌할 수 있는 사례들을
제외하고는 법적 결과를 확정하기에 충분한 거의 확정적인 규칙들과, 단지
결정을 향해 가리키지만 매우 자주 그것을 확정하지 못할 수 있는 일반적으로
비확정적인 원칙들 사이에는 합리적 대비가 이루어질 수 있다.

나는 그러한 비확정적 원칙들로부터의 논변들이 재판과 법적 추론의 중요한
특징이며, 적절한 용어법으로 표시되어야 한다고 확실히 생각한다. 그것들의
중요성과 법적 추론에서의 역할을 보이고 예시한 점에서 드워킨에게는 큰
공이 있다. 그리고 내가 그것들의 비확정적 힘을 강조하지 않은 것은 분명
중대한 실수였다. 그러나 나는 ‘규칙’이라는 말을 사용하면서 법체계들이
오직 ‘전부-아니면-무’ 규칙들 또는 거의 확정적인 규칙들만 포함한다고
주장하려 의도하지 않았다. 나는 이 책 130--3쪽에서 내가 (아마도 썩
적절하지 않게) ‘가변적 법적 기준들’이라고 부른 것들, 곧 고려되어 다른
것들과 저울질되어야 할 요소들을 특정하는 기준들에 주의를 환기했을 뿐
아니라, 왜 어떤 행위 영역들은 ‘상당한 주의’와 같은 가변적 기준들이
아니라 드문 사례들을 제외하고 동일한 특정 행위들을 금지하거나 요구하는
거의 확정적인 규칙들에 의한 규율에 적합한지를 설명하려고도 했다(133--4쪽
참조). 그래서 우리는 살인과 절도에 반대하는 규칙들을 가지고 있는 것이지,
인간 생명과 재산에 대한 적절한 존중을 요구하는 원칙들만 가지고 있는 것이
아니다.

\subsection{4. 원칙들과 승인 규칙 (PRINCIPLES AND THE RULE OF
RECOGNITION)}\label{uxc6d0uxce59uxb4e4uxacfc-uxc2b9uxc778-uxaddcuxce59-principles-and-the-rule-of-recognition}

\subsubsection{\texorpdfstring{\emphksans{계보와 해석} (Pedigree and
Interpretation)}{계보와 해석 (Pedigree and Interpretation)}}\label{uxacc4uxbcf4uxc640-uxd574uxc11d-pedigree-and-interpretation}

드워킨은 법원들의 실천에서 드러나는 승인 규칙이 제공하는 기준들에 의해
법원칙들이 식별될 수 없으며, 원칙들은 법의 필수 요소들이므로 승인 규칙
교설은 폐기되어야 한다고 주장해 왔다. 그에 따르면 법원칙들은 오직 구성적
해석에 의해서만, 곧 한 법체계의 확립된 법 전체의 제도적 역사에 가장 잘
부합하고 그것을 가장 잘 정당화하는 유일한 원칙 집합의 구성원들로서만
식별될 수 있다. 물론 영국이든 미국이든 어떤 법원도 \textbf{(p.~264)}
당해 법을 식별하기 위해 그러한 체계-전체적 총체 기준을 명시적으로 채택한
적은 없으며, 드워킨은 자신의 신화적 이상 판사 ‘헤라클레스’와 구별되는
실제 인간 판사는 누구도 한 번에 자기 나라의 모든 법에 대한 해석을
구성하는 위업을 성취할 수 없다는 점을 인정한다. 그럼에도 그의 견해에서
법원들은 제한된 방식으로 ‘헤라클레스를 모방하려’ 하는 것으로 이해될 때
가장 잘 해명되며, 그는 자신들의 판결을 이렇게 바라보는 일이 ‘숨겨진
구조’를 드러내는 데 기여한다고 생각한다.\footnote{\emph{LE} 265면.}

제한된 형식의 구성적 해석에 의해 원칙들이 식별되는 가장 유명한 예, 영국
법률가들에게 익숙한 예는 \emph{Donoghue} v. \emph{Stevenson} 사건에서
애킨 경이 이전에는 정식화되지 않았던 ‘이웃 원칙(neighbour principle)’을,
여러 상황에서 주의책무를 확립하는 다양한 개별 규칙들의 밑바탕에 놓인
원칙으로 정식화한 것이다.\footnote{{[}1932{]} A.C. 562.}

나는 그러한 제한된 구성적 해석의 수행들에서 판사들이 헤라클레스의 총체적
체계-전체 접근을 모방하려 한다고 이해하는 것이 그럴듯하다고 보지 않는다.
그러나 여기서 나의 현재 비판은, 구성적 해석에 대한 몰두가 드워킨으로
하여금 많은 법원칙들이 그 지위를 확립된 법의 해석으로 기능하는 그 내용
덕분이 아니라, 그가 그 원칙들의 ‘계보’라고 부르는 것, 곧 인정된 권위
있는 원천에 의해 그것들이 창설되거나 채택된 방식 덕분에 얻는다는 사실을
무시하게 만들었다는 것이다. 이 몰두는 실제로 그를 이중의 오류로
이끌었다고 생각한다. 첫째는 법원칙들이 그 계보에 의해 식별될 수 없다는
믿음이고, 둘째는 승인 규칙이 오직 계보 기준들만 제공할 수 있다는
믿음이다. 이 두 믿음은 모두 잘못이다. 첫 번째 믿음이 잘못인 것은
원칙들의 비확정적 성격에도, 그 밖의 특징들에도, 그것들이 계보 기준들에
의해 식별되는 것을 배제하는 것은 전혀 없기 때문이다. 분명 성문헌법이나
헌법 수정조항 또는 제정법의 한 규정은 원칙들에 특징적인 비확정적
방식으로 작동하도록 의도된 것으로, 곧 어떤 다른 규칙이나 원칙이 대안적
결정을 위한 더 강한 이유들을 제시하는 사례들에서는 압도될 수 있는 결정
이유들을 제공하는 것으로 받아들여질 수 있다. 드워킨 자신도 의회가 표현의
자유를 제한해서는 안 된다고 규정하는 미국 헌법 제1수정조항이 바로 그런
방식으로 \textbf{(p.~265)} 해석되어야 한다고 상정했다.\footnote{{[}\emph{TRS}
  27면 참조.{]}} 또한 자기 자신의 그른 행위로부터 이익을 얻어서는 안
된다는 원칙 같은 커먼로의 몇몇 기본 원칙들을 포함한 일부 법원칙들은,
반대 방향을 가리키는 이유들에 의해 어떤 사례들에서는 압도될 수 있음에도
고려되어야 할 결정 이유들을 제공하는 것으로, 여러 다른 사례 범위에서
법원들이 일관되게 원용해 왔다는 점에서 ‘계보’ 검사 기준에 의해 법으로
식별된다. 계보 기준들에 의해 식별되는 법원칙들의 이러한 예들 앞에서,
법의 일부로 원칙들을 포함하는 것이 승인 규칙 교설의 포기를 함축한다는
일반 논변은 성공할 수 없다. 사실 내가 아래에서 보이듯이, 그것들의 포함은
그 교설과 양립 가능할 뿐 아니라 실제로 그 교설의 수용을 요구한다. 승인
규칙이 제공하는 계보 기준들에 의해 적어도 일부 법원칙들이 법으로
‘포착’되거나 식별될 수 있다는 점이 인정된다면, 그리고 분명 인정되어야
한다면, 드워킨의 비판은 더 겸손한 주장으로 축소되어야 한다. 곧 그렇게
포착될 수 없는 법원칙들이 많이 있는데, 그것들은 너무 많거나, 너무
일시적이거나, 변화나 수정에 너무 열려 있거나, 법체계의 제도적 역사와
실천들에 가장 잘 부합하고 그것들을 가장 잘 정당화하는 정합적 원칙 체계에
속한다는 검사 기준 이외의 어떤 검사 기준을 참조해서도 법원칙들로 식별될
수 있게 하는 특징을 가지지 않는다는 주장이다. 첫눈에 이 해석주의적 검사
기준은 승인 규칙이 제공하는 기준의 대안이 아니라, 몇몇 비판자들이
주장했듯이,\footnote{예: E. P. Soper, “Legal Theory and the Obligation
  of a Judge”, \emph{RDCJ} 3면, 16면; J. Coleman, “Negative and
  Positive Positivism”, \emph{RDCJ} 28면; D. Lyons, “Principles,
  Positivism and Legal Theory”, \emph{Yale Law Journal} 제87권 (1977)
  415면.} 원칙들을 그 계보가 아니라 그 내용에 의해 식별하는 그러한
기준의 복합적 ‘연성-실증주의적’ 형식일 뿐인 듯하다. 그러한 해석 기준을
포함하는 승인 규칙이, 위 251쪽 이하에서 논의한 이유들 때문에, 드워킨에
따르면 실증주의자가 원할 법 식별의 확실성 정도를 확보할 수 없다는 것은
사실이다. 그럼에도 해석적 검사 기준이 법-식별의 관행적 양식의 일부였음을
보이는 것은 여전히 \textbf{(p.~266)} 그것의 법적 지위에 대한 좋은 이론적
설명일 것이다. 그러므로 법의 일부로 원칙들을 인정하는 것과 승인 규칙
교설 사이에는 드워킨이 주장하는 그런 양립불가능성이 분명 없다.

앞의 두 단락의 논변은, 드워킨의 논지와 달리 원칙들을 법의 일부로
수용하는 것이 승인 규칙 교설과 양립 가능하다는 점을 보이기에 충분하다.
설령 드워킨의 해석적 검사 기준이 그가 주장하듯 원칙들을 식별하기에
유일하게 적절한 기준이라고 하더라도 그렇다. 그러나 사실 더 강한 결론이
보증된다. 곧 법원칙들이 그러한 기준에 의해 식별되려면 승인 규칙이
필요하다는 결론이다. 왜냐하면 드워킨의 해석적 검사 기준에 의해 드러날
어떤 법원칙을 식별하기 위한 출발점은, 그 원칙이 부합하고 정당화하는 데
도움을 주는 확립된 법의 어떤 특정 영역이기 때문이다. 그러므로 그 기준의
사용은 확립된 법의 식별을 전제하며, 그것이 가능하려면 법의 원천들과 그들
사이에 성립하는 우위와 종속의 관계들을 특정하는 승인 규칙이 필요하다.
\emph{Law's Empire}의 용어법으로 말하면, 밑바탕에 놓인 또는 암시적
법원칙들을 식별하는 해석 과제의 출발점들을 구성하는 법규칙들과 실천들은
‘전해석적 법(preinterpretive law)’을 구성하며, 드워킨이 그것에 관해
말하는 많은 것은, 그것의 식별을 위해 이 책에서 서술된 것처럼 법의 권위
있는 원천들을 식별하는 승인 규칙과 매우 유사한 것이 필요하다는 견해를
지지하는 듯하다. 여기서 내 견해와 드워킨의 견해 사이의 주된 차이는, 내가
법의 원천들을 식별하기 위한 기준들에 관해 판사들 사이에서 발견되는
일반적 합의를 그러한 기준들을 제공하는 \emphksans{규칙들}의 공유된 수용에
귀속시키는 반면, 드워킨은 규칙들이 아니라 같은 해석 공동체의 구성원들이
공유하는 ‘합의’,\footnote{{[}\emph{LE} 65--66면, 91--92면.{]}}
‘패러다임들’,\footnote{{[}\emph{LE} 72--73면.{]}} ‘가정들’\footnote{{[}\emph{LE}
  47, 67면.{]}}에 관해 말하기를 선호한다는 점이다. 물론 드워킨이 분명히
했듯이, 합의의 각 당사자가 일치하는 데 가지는 이유의 일부로 다른 이들의
일치가 들어 있지 않은 독립적 확신들의 합의와, 그것이 그러한 일부인
관행의 합의 사이에는 중요한 구별이 있다. 승인 규칙은 내 책에서 분명
\textbf{(p.~267)} 관행적 형식의 사법적 합의에 기초하는 것으로 취급된다.
그것이 그렇게 기초해 있다는 점은 적어도 영국법과 미국법에서는 아주
분명해 보인다. 영국 판사가 의회의 입법을, 또는 미국 판사가 헌법을, 다른
원천들에 대해 최고성을 갖는 법의 원천으로 취급하는 이유에는 분명 그의
사법 동료들이 그들의 선임자들이 그랬던 것처럼 이에 동의한다는 사실이
포함되기 때문이다. 기실 드워킨 자신도 입법 최고성 교설을 판사의 확신이
수행할 수 있는 역할을 제한하는 법적 역사의 날것의 사실로
말하고,\footnote{{[}\emph{LE} 401면.{]}} ‘해석적 태도는 같은 해석
공동체의 구성원들이 “무엇이 실천의 일부로 간주되는가”에 관해 적어도
대략 동일한 가정들을 공유하지 않으면 존속할 수 없다’고
진술한다.\footnote{\emph{LE} 67면.} 그러므로 나는 드워킨이 말하는
규칙들과 ‘가정들’, ‘합의’, ‘패러다임들’ 사이에 어떤 차이들이 남아 있든,
법의 원천들에 대한 사법적 식별에 관한 그의 설명은 실질적으로 나의 설명과
같다고 결론짓는다.

그러나 나의 견해와 드워킨의 견해 사이에는 큰 이론적 차이들이 남아 있다.
드워킨은 법원칙들을 위한 자신의 해석적 검사 기준을, 어떤 법체계들에서 그
법원들의 수용에 의존하여 존재와 권위를 가지는 관행적 승인 규칙이 취하는
특정 형식으로 다루는 나의 처리를 분명 거부할 것이다. 그의 견해에서는
이것이 법을 최선의 도덕적 빛 아래 보이도록 설계된 ‘구성적’ 해석의 기획,
곧 드워킨의 견해에서 법의 식별에 관여하는 기획을 완전히 잘못 제시하고
격하하는 것이기 때문이다. 그는 이러한 해석 양식을, 특정 법체계들의
판사들과 법률가들이 수용한 단순한 관행적 규칙이 요구하는 법-식별의
방법으로 구상하지 않는다. 대신 그는 그것을 법 이외에도 많은 사회사상과
사회적 실천의 중심 특징으로, 그리고 문학비평에서 이해되는 해석과 심지어
자연과학에서 이해되는 해석까지 포함하여 ‘모든 형식의 해석 사이의 깊은
연결’을 보여 주는 것으로 제시한다.\footnote{\emph{LE} 53면.} 그러나 이
해석 기준이 단지 관행적 규칙이 요구하는 법-식별의 양식이 아니고 다른
분과들에서 이해되는 해석과 친연성 및 연결을 가진다 하더라도, 드워킨의
총체적 해석 기준이 \textbf{(p.~268)} 실제로 법원칙들을 식별하는 데
사용되는 법체계들이 있다면 그 체계들에서 그 기준이 관행적 승인 규칙에
의해 제공될 수 있음은 여전히 사실이다. 그러나 이 완전한 총체 기준이
사용되는 실제 법체계들은 없고, \emph{Donoghue} v. \emph{Stevenson} 같은
사례들에서 잠재적 법원칙들을 식별하기 위해 더 겸손한 구성적 해석
수행들이 이루어지는 영국법과 미국법 같은 체계들만 있을 뿐이다. 그러므로
고려되어야 할 유일한 질문은 그러한 수행들이 관행적 승인 규칙이 제공하는
기준의 적용으로 이해되어야 하는지, 아니면 다른 어떤 방식으로 이해되어야
하는지, 그리고 그렇다면 그것들의 법적 지위가 무엇인지이다.

\subsection{5. 법과 도덕 (LAW AND
MORALITY)}\label{uxbc95uxacfc-uxb3c4uxb355-law-and-morality}

\subsubsection{\texorpdfstring{(i) \emphksans{권리들과 책무들} (Rights and
Duties)}{(i) 권리들과 책무들 (Rights and Duties)}}\label{i-uxad8cuxb9acuxb4e4uxacfc-uxcc45uxbb34uxb4e4-rights-and-duties}

이 책에서 나는, 법과 도덕 사이에는 서로 다른 많은 우연적 연결들이 있지만
법의 내용과 도덕 사이에는 필연적 개념적 연결들이 없으며, 따라서
도덕적으로 사악한 규정들도 법규칙들 또는 법원칙들로서 유효할 수 있다고
논변한다. 이러한 법과 도덕의 분리 형식의 한 측면은, 어떤 도덕적 정당화나
힘도 전혀 가지지 않는 법적 권리들과 책무들이 있을 수 있다는 점이다.
드워킨은 이 관념을 거부하고, 법적 권리들과 책무들의 존재에 관한
단언들에는 적어도 일응의 도덕적 근거들이 있어야 한다는 견해를 선호한다.
그래서 그는 ‘법적 권리들은 도덕적 권리들의 {[}한{]} 종으로 이해되어야
한다’는 관념을 자신의 법이론에서 ‘핵심적(crucial)’\footnote{\emph{RDCJ}
  260면.} 요소로 여기며, 그에 반대되는 실증주의 교설은 도덕적 근거 또는
힘이 전혀 없는 법적 권리들과 책무들이 있을 수 있다는 것을 분석 이전에
그저 아는 것으로 우리에게 주어지는 ‘법적 본질주의의 특이한
세계’\footnote{\emph{RDCJ} 259면.}에 속한다고 말한다. 일반 기술적
법리학이 법의 이해에 할 수 있는 기여의 종류를 이해하기 위해서는, 그의
일반 해석이론의 이점이 무엇이든, 법적 권리들과 책무들이 도덕적 힘이나
정당화를 결여할 수 있다는 교설에 대한 드워킨의 비판이 잘못되었음을 보는
것이 중요하다고 나는 생각한다. 그 이유는 다음과 같다. \textbf{(p.~269)}
법적 권리들과 책무들은 법이 그 강제적 자원들로 각각 개인의 자유를
보호하고 제한하거나, 개인들에게 법의 강제적 장치를 이용할 권한을
부여하거나 부인하는 지점이다. 그러므로 법들이 도덕적으로 좋든 나쁘든,
정의롭든 부정의하든, 권리들과 책무들은 법의 작용에서 인간들에게 최고로
중요한 초점들로서, 법들의 도덕적 이점들과 독립적으로 주목을 요구한다.
따라서 법적 권리들과 책무들의 진술들이 현실 세계에서 의미를 가질 수
있으려면 그 존재를 단언할 어떤 도덕적 근거가 있어야만 한다는 말은 참이
아니다.

\subsubsection{\texorpdfstring{(ii) \emphksans{법의 식별} (The Identification
of the
Law)}{(ii) 법의 식별 (The Identification of the Law)}}\label{ii-uxbc95uxc758-uxc2dduxbcc4-the-identification-of-the-law}

이 책에서 전개된 법이론과 드워킨의 이론 사이에서 법과 도덕의 연결들에
관련된 가장 근본적인 차이는 법의 식별에 관한 것이다. 내 이론에 따르면
법의 존재와 내용은 법의 사회적 원천들, 예컨대 입법, 사법 결정들, 사회
관습들을 참조하여 식별될 수 있으며, 그렇게 식별된 법이 그 자체로 법의
식별을 위한 도덕적 기준들을 편입한 경우를 제외하면 도덕에 대한 참조 없이
그렇게 될 수 있다. 반면 드워킨의 해석이론에서는 어떤 주제에 관한 당해
법이 무엇인지 진술하는 모든 법명제가 필연적으로 도덕 판단을 포함한다.
그의 총체적 해석이론에 따르면 법명제들은 다른 전제명제들과 함께, 법의
사회적 원천들을 참조하여 식별된 한 체계의 확립된 법 전체에 가장 잘
부합하고 그것에 대한 최선의 도덕적 정당화를 제공하는 원칙들의 집합에서
따라 나올 때에만 참이기 때문이다. 그러므로 이 전체적 총체 해석이론은
이중 기능을 가진다. 그것은 법을 식별하는 데에도, 그 법에 대한 도덕적
정당화를 제공하는 데에도 쓰인다.

이것이 \emph{Law's Empire}에서 ‘해석적’ 법과 ‘전해석적’ 법 사이의 구별을
도입하기 전 드워킨의 이론을 간략히 요약한 모습이었다. 법의 존재와 내용이
도덕에 대한 참조 없이 식별될 수 있다는 실증주의자의 이론에 대한 대안으로
고려할 때, 원래 상태의 드워킨 이론은 다음 비판에 취약했다. 사회적
원천들에 대한 참조로 식별된 법이 도덕적으로 사악한 곳에서는,
\textbf{(p.~270)} 그것을 가장 잘 ‘정당화’하는 원칙들은 그 법에 부합하는
원칙들 가운데 가장 덜 사악한 것들일 수밖에 없다. 그러나 그러한 가장 덜
사악한 원칙들은 어떤 정당화하는 힘도 가질 수 없으며, 무엇이 법으로
간주될 수 있는지에 어떤 도덕적 한계나 제약도 구성할 수 없다. 그리고
그것들이 아무리 사악한 법체계에도 부합할 수밖에 없기 때문에, 그것들을
참조하여 법을 식별한다고 주장하는 이론은 도덕에 대한 어떤 참조도 없이
법이 식별될 수 있다는 실증주의 이론과 구별되지 않는다. 드워킨이 ‘배경
도덕(background morality)’이라고 부른 것의 기준상 도덕적으로 건전한
원칙들, 단지 법에 부합하는 원칙들 가운데 도덕적으로 가장 건전한 것이
아닌 원칙들은 기실 무엇이 법으로 간주될 수 있는지에 도덕적 한계나 제약을
제공할 수 있다.\footnote{{[}\emph{TRS} 112, 128면, \emph{TRS} 93면
  참조.{]}} 나는 그 명제에 어떤 식으로도 반대하지 않지만, 그것은 법이
도덕에 대한 참조 없이 식별될 수 있다는 나의 주장과 완전히 양립 가능하다.

드워킨은 후기에 해석적 법과 전해석적 법 사이의 구별을 도입하면서, 우리가
도덕적으로 수용 가능하다고 볼 수 있는 그 법들의 어떤 해석도 불가능할
정도로 사악한 법체계들이 있을 수 있음을 인정한다. 그러한 경우에는 그가
설명하듯이 ‘내부 회의주의’\footnote{\emph{LE} 78--79면.}에 의지하여
그러한 체계들이 법임을 부정할 수 있다. 그러나 그러한 상황들을 서술하기
위한 우리의 자원들은 매우 유연하기 때문에, 우리는 대신 아무리 사악한
법체계들도 전해석적 의미에서는 법이라고 말할 수 있을 때 반드시 그 결론에
이르도록 구속되어 있지 않다.\footnote{{[}\emph{LE} 103면.{]}} 그러므로
우리는 가장 나쁜 나치 법들에 대해서조차 그것들이 법이 아니라고 말하도록
강제되지 않는다. 그 법들은 법창출의 형식들, 재판과 집행의 형식들과 같은
법의 많은 독특한 특징들을 도덕적으로 수용 가능한 체제들의 법들과
공유하면서, 오직 그 사악한 도덕적 내용에서만 그것들과 다를 수 있기
때문이다. 많은 맥락들과 많은 목적들에서는 그 도덕적 차이를 무시하고,
실증주의자와 함께 그러한 사악한 체계들이 법이라고 말할 충분한 이유들이
있을 수 있다. 이에 드워킨은 말하자면 자신의 해석적 관점에 대한 일반적
고수를 드러내는 단서, 곧 그러한 사악한 체계들은 오직 전해석적 의미에서만
법이라는 단서만을 덧붙일 것이다.

나는 우리 언어의 유연성에 대한 이 호소와 \textbf{(p.~271)} 이 지점에서
해석적 법과 전해석적 법 사이의 구별을 도입하는 일이 실증주의자의 논거를
약화하기보다 오히려 인정한다고 본다. 그것은, 그가 기술적 법리학에서는
법이 도덕에 대한 참조 없이 식별될 수 있다고 고집하지만, 확립된 법을
무엇이 가장 잘 정당화하는지에 관한 도덕 판단을 법의 식별이 항상
포함한다고 보는 정당화적 해석 법리학에서는 사정이 다르다는 메시지를
전달하는 것 이상의 일을 거의 하지 않기 때문이다. 물론 이 메시지는
실증주의자가 자신의 기술적 기획을 포기할 어떤 이유도 주지 않으며, 그렇게
하려는 의도도 없다. 그러나 이 메시지조차 한정되어야 한다. 법이 너무
사악하여 ‘내부 회의주의’가 적절한 경우에는 법의 해석이 어떤 도덕 판단도
포함하지 않으며, 드워킨이 이해하는 해석은 포기되어야 하기
때문이다.\footnote{{[}\emph{LE} 105면.{]}}

드워킨이 자신의 해석이론에 가한 또 하나의 수정은 그의 법적 권리 설명에
중요한 관련을 가진다. 원래 제시된 그의 총체 이론에서는 법의 식별과 그
정당화가 모두 한 체계의 확립된 법 모두에 가장 잘 부합하고 그것을 가장 잘
정당화하는 유일한 원칙 집합에서 따라 나오는 것으로 취급되었다. 그러므로
그러한 원칙들은, 내가 말했듯이, 이중 기능을 가진다. 그러나 한 체계의
확립된 법이 너무 사악해서 그 법에 대한 전반적 정당화 해석이 불가능할 수
있으므로, 드워킨은 이 두 기능이 분리될 수 있고, 도덕에 대한 어떤 참조도
없이 식별되는 법원칙들만 남을 수 있다고 관찰했다. 그러나 그러한 법은
드워킨이 모든 법적 권리들이 가진다고 주장하는 일응의 도덕적 힘을 가지는
어떤 권리들도 확립할 수 없다. 그럼에도 드워킨이 나중에 인정했듯이, 그
체계가 너무 사악해서 법 전체에 대한 어떤 도덕적 또는 정당화적 해석도
불가능한 곳에서도, 개인들이 적어도 일응의 도덕적 힘을 가진 권리들을
가진다고 적절히 말해질 수 있는 상황들이 여전히 있을 수 있다.\footnote{{[}\emph{LE}
  105--106면.{]}} 그것은 그 체계가, 예컨대 계약의 형성과 집행에 관련된
법들처럼 그 체계의 일반적 사악함에 영향을 받지 않을 수 있는 법들을
포함하고, 개인들이 자기 삶을 계획하거나 재산 처분들을 할 때 그러한
법들에 의존했을 수 있는 경우일 것이다. 그러한 상황들에 대응하기 위해
드워킨은, \textbf{(p.~272)} 일응의 도덕적 힘을 가진 법적 권리들과
책무들이 법에 대한 일반 해석이론에서 흘러나와야 한다는 자신의 원래
관념을 한정하고, 그러한 상황들이 자신의 일반 이론과 독립하여 개인들에게
어떤 도덕적 힘을 가진 법적 권리들을 귀속시키기 위한 ‘특별한 이유들’을
구성한다고 인정한다.

\subsection[6. 사법적 재량 (JUDICIAL DISCRETION)]{\texorpdfstring{6.
사법적 재량 (JUDICIAL
DISCRETION)\footnote{{[}본 절의 도입 문단의 대체 버전은 미주(endnote)에
  수록됨.{]}}}{6. 사법적 재량 (JUDICIAL DISCRETION)}}\label{uxc0acuxbc95uxc801-uxc7acuxb7c9-judicial-discretion73}

이 책의 법이론과 드워킨의 이론 사이의 가장 날카로운 직접 충돌은, 어떤
법체계에도 언제나 법적으로 규율되지 않은 일정한 사례들이 있을 것이며, 그
사례들에서는 어떤 지점에 관해 어느 쪽의 결정도 법에 의해 지시되지 않고
따라서 법이 부분적으로 불확정적이거나 불완전하다는 나의 논지로부터
생긴다. 그러한 사례들에서 판사가 결정을 내려야 하고, 벤담이 한때
주장했듯이 관할권을 부인하거나 현존 법에 의해 규율되지 않은 쟁점들을
입법부에 회부하여 결정하게 하지 않는다면, 판사는 이미 존재하는 확립된
법을 단지 적용하는 것이 아니라 자기 \emphksans{재량}을 행사하여 그 사례를
위한 법을 \emphksans{만들어야} 한다. 그러므로 그러한 법적으로 마련되지
않았거나 규율되지 않은 사례들에서 판사는 새로운 법을 만들면서, 동시에
자신의 법창출 권한들을 부여하고 제약하는 확립된 법을 적용한다.

법이 부분적으로 불확정적이거나 불완전하고 판사가 제한된 법창출 재량을
행사하여 그 틈들을 메운다는 이 구상은, 드워킨에 의해 법과 사법 추론
양쪽에 대한 오도적 설명으로 거부된다. 그는 사실상 불완전한 것은 법이
아니라 실증주의자의 법 그림이며, 이 점은 사회적 원천들을 참조하여
식별되는 \emphksans{명시적} 확립된 법 외에도 그 명시적 법에 가장 잘
부합하거나 정합하고 또한 그것에 대한 최선의 도덕적 정당화를 제공하는
원칙들인 \emphksans{암시적} 법원칙들을 포함하는 법에 대한 자신의 ‘해석적’
설명에서 드러난다고 주장한다. 이 해석적 견해에서는 법이 결코
불완전하거나 불확정적이지 않으므로, 판사는 결정을 내리기 위해 법
바깥으로 나가 법창출 권한을 행사할 계기를 결코 가지지 않는다. 따라서
\textbf{(p.~273)} 법의 사회적 원천들이 어떤 법적 쟁점에 관해 결정을
확정하지 못하는 ‘어려운 사례들’에서 법원들이 향해야 할 것은, 그 도덕적
차원들을 가진 그러한 암시적 원칙들이다.

법에 의해 부분적으로 규율되지 않은 채 남은 사례들을 규율하기 위해 내가
판사들에게 귀속시키는 법창출 권한들은 입법부의 권한들과 다르다는 점이
중요하다. 판사의 권한들은 입법부가 상당히 자유로울 수 있는 많은 제약들,
곧 \emphksans{그의 선택을 좁히는} 제약들에 종속될 뿐 아니라, 특정한 당해
사례들을 처리하기 위해서만 행사되므로 대규모 개혁이나 새로운 법전들을
도입하는 데 사용될 수 없다. 그래서 그의 권한들은 많은 실체적
제약들(substantive constraints)에 종속될 뿐 아니라 \emphksans{간극적}이다.
그럼에도 현존 법이 어떤 결정도 올바른 것으로 지시하지 못하는 지점들이
있을 것이며, 그런 경우들을 결정하기 위해 판사는 자신의 법창출 권한들을
행사해야 한다. 그러나 그는 이것을 자의적으로 해서는 안 된다. 곧 그는
언제나 자기 결정을 정당화하는 어떤 일반적 이유들을 가져야 하고, 자신의
믿음들과 가치들에 따라 결정함으로써 양심적 입법자가 하듯이 행위해야
한다. 그러나 그가 이 조건들을 충족한다면, 그는 법에 의해 지시되지 않고
유사한 어려운 사례들에 직면한 다른 판사들이 따른 것들과 다를 수 있는
결정 기준들이나 이유들을 따를 자격이 있다.

법에 의해 불완전하게 규율된 채 남은 사례들을 처리하기 위해 법원들이
그러한 제한된 재량적 권한을 행사한다는 나의 설명에 맞서, 드워킨은 세
가지 주요 비판을 제기한다. 첫째는 이 설명이 사법 절차와 ‘어려운
사례들’에서 법원들이 하는 일에 대한 거짓 기술이라는 것이다.\footnote{{[}\emph{TRS}
  81면; \emph{LE} 37--39면 비교.{]}} 이를 보이기 위해 드워킨은 판사들과
법률가들이 판사의 과제를 서술할 때 사용하는 언어와 사법적 의사결정의
현상학에 호소한다. 사건들을 결정하는 판사들과 자신들에게 유리하게
결정하라고 그들을 촉구하는 법률가들은, 새로운 사례들에서조차 판사가 법을
‘만든다’고 말하지 않는다는 것이다. 그러한 사례들 가운데 가장 어려운
것들에서조차 판사는, 실증주의자가 시사하듯이, 결정 과정에 완전히 다른 두
단계가 있다는 자각을 드러내지 않는 경우가 많다. 곧 판사가 먼저 현존 법이
어느 쪽의 결정도 지시하지 못한다고 발견하는 단계와, 그 다음 현존 법에서
돌아서서 무엇이 최선인지에 관한 자신의 관념에 따라 당사자들을 위해
\emphksans{처음부터} 그리고 \emphksans{사후적으로} 법을 만드는 단계 말이다.
\textbf{(p.~274)} 대신 법률가들은 판사가 언제나 현존 법을 발견하고
집행하는 데 관여하는 것처럼 판사에게 말하고, 판사는 법이 모든 사례에
대한 해법이 판사의 발명이 아니라 발견을 기다리는, 빈틈없는 권리자격들의
체계인 것처럼 말한다.

사법 절차의 익숙한 수사가 발전된 법체계에는 법적으로 규율되지 않은
사례들이 없다는 관념을 북돋는다는 데에는 의심의 여지가 없다. 그러나
이것은 얼마나 진지하게 받아들여져야 하는가? 물론 입법자와 판사의 구별을
극화하고, 판사는 현존 법이 명확할 때 그러하듯 언제나 자신이 만들거나
빚어내지 않는 법의 단순한 ‘입’이라고 주장하는 오랜 유럽 전통과 권력분립
교설이 있다. 그러나 법정에서 사건들을 결정할 때 판사들과 법률가들이
사용하는 의례적 언어와, 사법 절차에 관한 그들의 더 성찰적인 일반
진술들을 구별하는 것이 중요하다. 미국의 올리버 웬델 홈스와 카도조,
영국의 맥밀런 경이나 래드클리프 경 또는 리드 경과 같은 격의 판사들,
그리고 학계와 실무계 양쪽의 수많은 다른 법률가들은, 법에 의해 불완전하게
규율된 채 남은 사례들이 있어서 판사에게 피할 수 없는, 비록 ‘간극적’인,
법창출 과제가 있으며, 법에 관한 한 많은 사례들이 어느 쪽으로든 결정될 수
있다고 주장해 왔다.

하나의 주요한 고려사항은 판사들이 때로 법을 만들면서 동시에 적용한다는
주장에 대한 저항을 설명하는 데 도움을 주며, 사법적 법창출을 입법부의
법창출과 구별하는 주요 특징들도 해명한다. 그것은 규율되지 않은 사례들을
결정할 때 법원들이, 자신들이 만드는 새로운 법이 비록 \emphksans{새로운}
법이기는 하지만 이미 현존 법에 발판을 가진 것으로 인정된 원칙들이나
밑받침 이유들과 부합하도록 보장하기 위해 유비에 따라 진행하는 데
전형적으로 부여하는 중요성이다. 특정 제정법들이나 선례들이 불확정적인
것으로 드러나거나 명시적 법이 침묵할 때, 판사들이 단지 법서들을 밀쳐
두고 법으로부터 더 이상의 지침 없이 입법하기 시작하지 않는 것은
사실이다. 그러한 사례들을 결정하면서 그들은 매우 자주 현존 법의 상당한
관련 영역이 예증하거나 증진하는 것으로 이해될 수 있고 당해 어려운 사례에
대한 확정적 답을 향해 가리키는 어떤 일반 원칙이나 일반적 목표 또는
목적을 인용한다. \textbf{(p.~275)} 이것은 기실 드워킨 재판 이론의 매우
두드러진 특징인 ‘구성적 해석’의 바로 그 핵심이다. 그러나 이 절차는
사법적 법창출의 순간을 분명 유예하지만, 제거하지는 않는다. 어떤 어려운
사례에서도 경쟁하는 유비들을 뒷받침하는 서로 다른 원칙들이 나타날 수
있고, 판사는 법에 의해 자신에게 이미 규정된 우선성의 질서가 아니라
양심적 입법자처럼 무엇이 최선인지에 관한 자신의 감각에 의존하여 그들
사이에서 선택해야 하는 경우가 많기 때문이다. 그러한 모든 사례들에 대해,
경쟁하는 저차-질서 원칙들에 상대적 무게나 우선성을 배정하는 고차-질서
원칙들의 어떤 유일한 집합이 현존 법 안에서 언제나 발견될 수 있을 때에만,
사법적 법창출의 순간은 단순히 유예되는 데 그치지 않고 제거될 것이다.

드워킨의 나의 사법 재량 설명에 대한 다른 비판들은 그것을 기술적으로
거짓이라고 비난하는 것이 아니라, 비민주적이고 부정의한 법창출 형식을
지지한다고 비난한다.\footnote{{[}\emph{TRS} 84--85면.{]}} 판사들은 보통
선출되지 않으며, 민주주의에서는 오직 선출된 국민의 대표자들만 법창출
권한들을 가져야 한다고 주장된다. 이 비판에는 많은 답변들이 있다. 법이
규율하지 못한 분쟁들을 처리하기 위해 판사들에게 법창출 권한들이 맡겨져야
한다는 것은, 입법부에 회부하는 것 같은 대안적 규율 방법들의 불편을
피하기 위해 치러야 할 필수적 대가로 여겨질 수 있다. 그리고 판사들이 이
권한들의 행사에서 제약되어 법전들이나 광범위한 개혁들을 형성할 수 없고
특정 사례들이 던지는 특정 쟁점들을 처리하기 위한 규칙들만 만들 수
있다면, 그 대가는 작아 보일 수 있다. 둘째, 제한된 입법 권한들을 행정부에
위임하는 것은 현대 민주주의들의 익숙한 특징이며, 그러한 위임이 사법부에
이루어진다고 해서 민주주의에 더 큰 위협으로 보이지는 않는다. 두 위임
형식 모두에서 선출된 입법부는 보통 잔여적 통제권을 가지며, 받아들일 수
없다고 보는 어떤 종속된 법들도 폐지하거나 수정할 수 있다. 미국에서처럼
입법부의 권한들이 성문헌법에 의해 제한되고 법원들이 광범위한 심사
권한들을 가지는 경우, \emphksans{민주적으로 선출된} 입법부가 어떤 사법적 입법
하나를 뒤집을 수 없음을 알게 될 수 있다는 것은 사실이다. 그러면
\textbf{(p.~276)} 궁극적 민주적 통제는 헌법 수정이라는 번거로운 장치를
통해서만 확보될 수 있다. 그것이 정부에 대한 법적 제약들을 위해 치러야
하는 대가이다.

드워킨은 사법적 법창출이 부정의하며, 물론 보통 부정의한 것으로 여겨지는
소급적 또는 \emphksans{사후적} 법창출의 한 형식으로 그것을 비난한다는 추가
고발을 한다. 그러나 소급적 법창출을 부정의한 것으로 보는 이유는, 행위할
때 자기 행위의 법적 귀결들이 그 행위 당시 확립되어 있던 알려진 법의
상태에 의해 결정될 것이라는 가정에 의존했던 사람들의 정당한 기대들을
그것이 좌절시키기 때문이다. 하지만 이 반론은, 분명히 확립된 법에 대한
법원의 소급적 변경이나 번복에 대해서는 힘을 가질 수 있다 하더라도,
어려운 사례들에서는 전혀 관련 없어 보인다. 그것들은 법이 불완전하게
규율한 채 남겨 둔 사례들이고, 기대들을 정당화할 명확히 확립된 알려진
법의 상태가 없는 경우들이기 때문이다.

\newpage

\subsection{후기 주석}\label{uxd6c4uxae30-uxc8fcuxc11d}


{\footnotesize
\noteentry 272쪽. {[}이 절의 대체 시작문이 여기에 포함되어 있으며, 폐기되지
않음.{]}

\begin{quote}
드워킨(Dworkin)은 오랜 시간에 걸친 재판(adjudication)에 관한 저술
전반에서, 법원들이 현존 법에 의해 불완전하게 규율된 채 남은 사례들을
결정할 법창출 권한(law-creating power)이라는 의미의 재량(discretion)을
가진다는 것을 부인하는 입장을 한결같이 견지해 왔다. 기실 그는, 사소한
예외를 제외하면 그러한 사례들은 존재하지 않으며, 유명하게 말했듯이 어떤
사례에서 생기는 어떤 법적 쟁점에 관해 당해 법이 무엇인지에 대한 의미
있는 질문에는 언제나 단 하나의 ‘정답(right answer)’이 있다고 논변해
왔다. \footnote{{[}그의 글 `No Right Answer?' in P. M. S. Hacker and J.
  Raz (eds.), \emph{Law, Morality and Society} (1977), pp.~58--84. 개정
  후 재수록: `Is There Really No Right Answer in Hard Cases?'
  \emph{AMP}, 제5장 참조.{]}}
\end{quote}

\begin{quote}
그러나 이러한 불변하는 듯한 교설의 외양에도 불구하고, 드워킨이 후기에
해석적 관념들(interpretive ideas)을 자신의 법이론에 도입하고 모든
법명제들(propositions of law)이 그가 이 표현에 부여한 특수한 의미에서
‘해석적’이라고 주장한 것은, 라즈가 처음 분명히 했듯이,\footnote{{[}J.
  Raz, `Dworkin: A New Link in the Chain', \emph{74 California Law
  Review}, 1103 (1986), pp.~1110, 1115--16 참조.{]}} 법원들이 실제로
법창출 재량을 가지며 자주 행사한다는 점을 인정하게 함으로써 이 입장의
실질을 나의 입장에 매우 가까이 가져왔다. 논변해 볼 수 있듯이, 해석적
관념들이 그의 이론에 도입되기 전에는, 드워킨의 강한 의미의 사법 재량
부정과 항상 정답이 있다는 그의 주장이 사건들을 결정하는 판사의 역할은
\emphksans{현존 법을 분별해 내고} \emphksans{집행하는 것}이라는 관념과 결합되어
있었기 때문에, 우리의 각 재판 설명 사이에는 큰 차이가 있는 듯했다.
그러나 물론 법원들이 사건들을 결정하면서 법창출 재량을 자주 행사한다는
나의 주장과 매우 날카롭게 충돌했던 이 초기 구상은 더 이상 이 절에서 전혀
등장하지 않는다.
\end{quote}

{[}섹션 6의 대체 시작문의 텍스트는 여기에서 끝난다.{]}

\subsection{후기 제3판
주석}\label{uxd6c4uxae30-uxc81c3uxd310-uxc8fcuxc11d}

이것은 아마도 어떤 후기(postscript)에 관해 책 한 권 전체가 나온 유일한
사례일 것이다:

Jules Coleman 엮음, \emph{Hart's Postscript}. 수록된 에세이들은
후기(postscript)뿐만 아니라 하트의 법이론 전반을 이해하는 데 매우
유익하다. 또한, Ronald Dworkin, “Hart's Postscript and the Character of
Political Philosophy” (\emph{Oxford Journal of Legal Studies} 제24권,
2004년, 1면)을 참조하라.
}


\clearpage
\backmatter

\end{document}
